\documentclass[14pt,a4paper,oneside]{extreport}
% --- OAK: A4, Times New Roman 14, 1,5 interval, hoshiyalar 30/15/20/20 mm
\usepackage[left=30mm,right=15mm,top=20mm,bottom=20mm,footskip=10mm,includefoot=false]{geometry}
\usepackage{fontspec}
\IfFontExistsTF{Times New Roman}{\setmainfont{Times New Roman}}{\setmainfont{Liberation Serif}}
\usepackage{amsmath,amssymb}
\usepackage{unicode-math}
\IfFontExistsTF{STIX Two Math}{\setmathfont{STIX Two Math}}{\setmathfont{Latin Modern Math}}
\usepackage{setspace}\setstretch{1.42}% Word "1,5 interval" (14 pt x 1,15 x 1,5 = 24,15 pt) ga teng
\setlength{\parindent}{1.27cm}\setlength{\parskip}{0pt}
\usepackage{indentfirst}
\usepackage{graphicx,xcolor,longtable,booktabs,array,calc,etoolbox,fancyhdr,tocloft,needspace}
\usepackage[hidelinks]{hyperref}
\frenchspacing\hyphenpenalty=10000\exhyphenpenalty=10000\emergencystretch=3em\sloppy
\providecommand{\tightlist}{\setlength{\itemsep}{0pt}\setlength{\parskip}{0pt}}
\def\LTcaptype{none}
\newcommand{\real}[1]{#1}
\renewcommand{\arraystretch}{1.15}\setlength{\tabcolsep}{4pt}
\AtBeginEnvironment{longtable}{\renewcommand{\baselinestretch}{1}\selectfont\hyphenpenalty=50\exhyphenpenalty=50\lefthyphenmin=2\righthyphenmin=3}
\setlength{\LTpre}{4pt}\setlength{\LTpost}{8pt}
% --- sahifa raqami: pastki o'ng burchakda; titul varag'i hisoblanadi, lekin raqam qo'yilmaydi
\pagestyle{fancy}\fancyhf{}\fancyfoot[R]{\fontsize{12}{12}\selectfont\thepage}
\renewcommand{\headrulewidth}{0pt}
\fancypagestyle{plain}{\fancyhf{}\fancyfoot[R]{\fontsize{12}{12}\selectfont\thepage}\renewcommand{\headrulewidth}{0pt}}
% --- sarlavhalar
\newcommand{\oakbob}[1]{\clearpage\phantomsection\addcontentsline{toc}{chapter}{#1}%
  \begin{center}\bfseries\MakeUppercase{#1}\end{center}\vspace{6pt}}
\newcommand{\oakpar}[1]{\Needspace{5\baselineskip}\phantomsection\addcontentsline{toc}{section}{#1}%
  \vspace{6pt}\begin{center}\bfseries #1\end{center}\vspace{2pt}}
\newcommand{\oakilova}[1]{\clearpage\begin{flushright}\bfseries #1\end{flushright}}
\newcommand{\oakilovanom}[1]{{\centering\setstretch{1}\bfseries #1\par}\nopagebreak\vspace{2pt}}
\newcommand{\oakjadvalnum}[1]{\par\Needspace{12\baselineskip}\vspace{2pt}\noindent\hfill #1\par\nopagebreak}
\newcommand{\oakjadval}[2]{\par\Needspace{12\baselineskip}\vspace{2pt}\noindent\hfill #1\par\nopagebreak%
  {\centering\setstretch{1}\bfseries #2\par}\nopagebreak\vspace{-2pt}}
\newcommand{\oakrasm}[1]{\vspace{-6pt}\begin{center}\setstretch{1}#1\end{center}\vspace{4pt}}
\newcommand{\oakimg}[1]{\par\Needspace{14\baselineskip}\begin{center}\includegraphics[width=0.94\linewidth,height=0.62\textheight,keepaspectratio]{#1}\end{center}\nopagebreak}
% suzuvchi rasm (asosiy matn): bo'sh joy qoldirmaslik uchun; bob oxirida \clearpage barcha rasmlarni chiqaradi
\newcommand{\oakfig}[2]{\begin{figure}[!tbp]\centering\includegraphics[width=0.94\linewidth,height=0.66\textheight,keepaspectratio]{#1}\par\vspace{2pt}{\setstretch{1}#2\par}\end{figure}}
\renewcommand{\topfraction}{.92}\renewcommand{\bottomfraction}{.6}\renewcommand{\textfraction}{.06}\renewcommand{\floatpagefraction}{.72}
\setcounter{topnumber}{2}\setcounter{bottomnumber}{1}\setcounter{totalnumber}{3}
\setlength{\textfloatsep}{10pt plus 2pt minus 2pt}\setlength{\floatsep}{8pt plus 2pt}
\newcommand{\toldir}[1]{\colorbox{yellow}{\strut #1}}
% --- MUNDARIJA
\renewcommand{\contentsname}{MUNDARIJA}
\renewcommand{\cfttoctitlefont}{\hfill\bfseries}\renewcommand{\cftaftertoctitle}{\hfill}
\setlength{\cftbeforetoctitleskip}{0pt}\setlength{\cftaftertoctitleskip}{12pt}
\renewcommand{\cftchapfont}{\bfseries}\renewcommand{\cftchappagefont}{\bfseries}
\renewcommand{\cftchapleader}{\cftdotfill{\cftdotsep}}
\setlength{\cftbeforechapskip}{4pt}\setlength{\cftsecindent}{0.8cm}\setlength{\cftchapnumwidth}{0pt}
\renewcommand{\cftchapfont}{\normalfont}\renewcommand{\cftchappagefont}{\normalfont}
\makeatletter\renewcommand{\@tocrmarg}{2.55em plus1fil}\makeatother
\usepackage{enumitem}\setlist{itemsep=0pt,parsep=0pt,topsep=2pt,partopsep=0pt}\setlist[enumerate]{leftmargin=*,labelindent=1.27cm}\setlist[itemize]{leftmargin=*,labelindent=1.27cm}

\begin{document}
\begin{titlepage}\setstretch{1.15}\centering
{\bfseries OʻZBEKISTON RESPUBLIKASI QISHLOQ XOʻJALIGI VAZIRLIGI\\[4pt] ANDIJON QISHLOQ XOʻJALIGI VA AGROTEXNOLOGIYALAR INSTITUTI\par}
\vspace{20mm}
\hfill\begin{tabular}{l}Qoʻlyozma huquqida\\ UDK: \toldir{\hspace{25mm}}\end{tabular}
\vspace{25mm}

{\bfseries SIROCHOV ABDUMUTALIB\par}
\vspace{12mm}
{\bfseries\fontsize{16}{20}\selectfont GAT TEXNOLOGIYALARI ASOSIDA MELIORATIV TIZIMLAR SAMARADORLIGI VA ISHONCHLILIGINI TAʼMINLASH\par}
\vspace{4mm}
{\bfseries (Andijon viloyati Marhamat tumani misolida)\par}
\vspace{14mm}
05.09.06 – Gidrotexnika va melioratsiya qurilishi\par
\vspace{10mm}
Texnika fanlari boʻyicha falsafa doktori (PhD)\\ dissertatsiyasi\par
\vfill
\hfill\begin{tabular}{l}Ilmiy rahbar:\\ t.f.n., dotsent A.Xakimov\end{tabular}
\vfill
ANDIJON – 2026
\end{titlepage}
\setcounter{page}{2}
\tableofcontents
\oakbob{KIRISH}

\textbf{Dissertatsiya mavzusining dolzarbligi va zarurati.} Jahonda
oziq-ovqat xavfsizligini taʼminlashda sugʻoriladigan yerlar hal qiluvchi
oʻrin tutadi. Biroq aynan shu yerlar ikkilamchi shoʻrlanish va zax
bosish xavfi ostida turadi. Birlashgan Millatlar Tashkilotining
Oziq-ovqat va qishloq xoʻjaligi tashkiloti (FAO) maʼlumotlariga koʻra,
dunyoda shoʻrlangan tuproqlar maydoni 1 381 mln gektarni, yaʼni
quruqlikning 10,7 foizini tashkil etadi. Sugʻoriladigan ekin
maydonlarining taxminan 10 foizi shoʻrlanishdan zarar koʻrgan {[}53{]}.
Bunga qarshi asosiy chora kollektor-drenaj tizimlaridir. Ular ildiz
qatlamidan ortiqcha suv va tuzni chiqarib, sizot suvlari sathini xavfsiz
chuqurlikda ushlab turadi. Tizimlarning koʻpchiligi oʻtgan asrning
oʻrtalarida qurilgan, ularni ishchi holatda saqlash har yili katta
mablagʻ talab qiladi. Shuning uchun tizimlarning haqiqiy holatini, sizot
suvlariga taʼsirini va taʼmir natijasini ishonchli maʼlumotlar asosida
baholash muhim vazifadir.

Jahonda meliorativ tizimlar holatini baholashda geoaxborot
texnologiyalari (GAT), masofadan zondlash va geostatistika usullari keng
qoʻllanilmoqda. Sizot suvlari sathi va minerallashuvini hududlar va
yillar boʻyicha xaritalash, tuproq shoʻrlanishini sunʼiy yoʻldosh
maʼlumotlari bilan baholash, drenaj ishini hisoblash va qaror qabul
qilishni qoʻllab-quvvatlash shular jumlasidandir. Biroq bu
tadqiqotlarning barchasi kuzatuv maʼlumotlariga tayanadi.
Maʼlumotlarning oʻzi qanchalik ishonchli ekanini baholash, drenaj
tarmogʻining joylashuvini sizot suvlari bilan bogʻlash va taʼmir
samarasini oʻlchash masalalari hali yetarlicha ishlanmagan.

Respublikamizda sugʻoriladigan yerlarning meliorativ holatini
yaxshilash, suv xoʻjaligi obyektlarini modernizatsiya qilish va sohaga
raqamli texnologiyalarni joriy etish boʻyicha keng qamrovli
chora-tadbirlar amalga oshirilmoqda. Oʻzbekiston Respublikasi suv
xoʻjaligini rivojlantirishning 2020--2030-yillarga moʻljallangan
konsepsiyasida yerlarning meliorativ holatini yaxshilash uchun umumiy
uzunligi 142,9 ming km boʻlgan kollektor-drenaj tarmogʻidan
foydalanilayotgani qayd etilgan. Unda sugʻorish tizimlarining foydali
ish koeffitsiyentini 0,63 dan 0,73 ga yetkazish, suv xoʻjaligi
obyektlarida suvni hisobga olish va boshqarish jarayonlarini
raqamlashtirish vazifalari belgilangan {[}2{]}. Bu vazifalarni bajarish
uchun meliorativ tizimlar va sizot suvlari toʻgʻrisidagi maʼlumotlarni
yagona raqamli bazaga oʻtkazish, ularning sifatini nazorat qilish va
sarflangan mablagʻ natijasini oʻlchash kerak.

Andijon viloyati Marhamat tumani misolida bu masalaning amaliy ahamiyati
yaqqol koʻrinadi. Tumandagi davlat hisobidagi 90 ta kollektorning
uzunlik boʻyicha 90 foizi 50 yildan ortiq xizmat qilgan. Xoʻjalik ichki
zovurlarining uchdan ikki qismida taʼmir qayd etilmagan.
2017--2020-yillarda kollektorlarni taʼmirlashga 2026-yil narxlarida 48,2
mlrd soʻm sarflangan. Shu bilan birga monitoring maʼlumotlarining 12
foizi tekshiruvdan oʻtmaydi. Sizot suvi chuqur joylarda kuzatuv
quduqlari haqiqiy sathni oʻlchash uchun yetarlicha chuqur emas. Drenaj
tarmogʻining raqamli xaritasi yoʻq. Bunday sharoitda tizim boʻyicha
qabul qilinadigan qarorlar ishonchli maʼlumotga tayanmaydi.

Oʻzbekiston Respublikasi Prezidentining 2020-yil 10-iyuldagi PF-6024-son
«Oʻzbekiston Respublikasi suv xoʻjaligini rivojlantirishning
2020--2030-yillarga moʻljallangan konsepsiyasini tasdiqlash
toʻgʻrisida»gi Farmoni {[}2{]}, 2019-yil 23-oktabrdagi PF-5853-son
«Oʻzbekiston Respublikasi qishloq xoʻjaligini rivojlantirishning
2020--2030-yillarga moʻljallangan strategiyasini tasdiqlash
toʻgʻrisida»gi Farmoni {[}3{]}, 2017-yil 27-noyabrdagi PQ-3405-son
«2018--2019-yillarda irrigatsiyani rivojlantirish va sugʻoriladigan
yerlarning meliorativ holatini yaxshilash Davlat dasturi toʻgʻrisida»gi
qarori {[}5{]}, 2023-yil 11-sentabrdagi PF-158-son «``Oʻzbekiston ---
2030'' strategiyasi toʻgʻrisida»gi Farmoni {[}1{]} hamda mazkur
faoliyatga tegishli boshqa meʼyoriy-huquqiy hujjatlarda belgilangan
vazifalarni amalga oshirishga ushbu dissertatsiya tadqiqoti muayyan
darajada xizmat qiladi.

\textbf{Tadqiqotning respublika fan va texnologiyalari rivojlanishining
ustuvor yoʻnalishlariga mosligi.} Mazkur tadqiqot respublika fan va
texnologiyalar rivojlanishining V. «Qishloq xoʻjaligi, biotexnologiya,
ekologiya va atrof-muhit muhofazasi» ustuvor yoʻnalishi doirasida
bajarilgan.

\textbf{Muammoning oʻrganilganlik darajasi.} Sugʻoriladigan yerlarni
drenajlash nazariyasi va hisob usullari A.N.Kostyakov, S.F.Averyanov,
I.P.Aydarov, A.I.Golovanov, B.S.Maslov ishlarida ishlab chiqilgan
{[}9--11, 13, 20, 21{]}. Xorijda drenajlar orasidagi masofani hisoblash
va drenaj tizimlarini loyihalash masalalari S.B.Hooghoudt, L.F.Ernst,
D.Kirkham, J.N.Luthin, J. van Schilfgaarde, H.P.Ritzema, L.K.Smedema,
W.F.Vlotman, R.W.Skaggs va R.J.Oosterbaan tomonidan oʻrganilgan {[}52,
69, 78, 85, 99, 108, 120--122, 127, 129{]}. Sizot suvlari harakatining
nazariy asoslari P.Ya.Polubarinova-Kochina, V.M.Shestakov, J.Bear,
R.A.Freeze va J.A.Cherry ishlarida berilgan {[}26, 27, 36, 37, 56{]};
sugʻorishning sizot suvlari rejimiga taʼsiri D.M.Kats va S.F.Averyanov
tomonidan koʻrsatilgan {[}10, 17{]}.

Tuproqlarning shoʻrlanishi, uning tasnifi va shoʻrlangan yerlarni
yaxshilash masalalari V.A.Kovda, V.V.Yegorov, N.G.Minashina,
Ye.I.Pankova, L.A.Richards, I.Szabolcs, J.D.Rhoades, D.L.Corwin,
G.I.Metternicht va J.A.Zinck tadqiqotlarida yoritilgan {[}15, 18, 19,
22, 25, 46, 47, 93, 106, 107, 123{]}. Oʻzbekiston sharoitida zax
qochirish melioratsiyasi X.A.Ahmedov {[}12{]}, meliorativ monitoring va
kadastrni yuritish R.Q.Ikromov, F.A.Barayev, Gʻ.U.Yusupov va
N.M.Karimova {[}16{]}, suv va yer resurslarini boshqarish V.A.Duxovniy,
V.I.Sokolov va R.Kulmatov {[}14, 80{]} ishlarida bayon etilgan.

Geoaxborot tizimlarining nazariy asoslari P.A.Burrough, R.A.McDonnell,
M.F.Goodchild va P.A.Longley {[}41, 42, 58, 83{]}, geostatistik
modellashtirish usullari G.Matheron, D.G.Krige, A.G.Journel, E.H.Isaaks,
N.Cressie, P.Goovaerts, R.Webster va M.A.Oliver {[}48, 60, 71, 73, 79,
91, 97, 131{]}, vaqt qatorlarida tendensiyani aniqlash usullari
H.B.Mann, M.G.Kendall, P.K.Sen va R.M.Hirsch {[}66, 68, 74, 90, 117{]}
tomonidan ishlab chiqilgan. Markaziy Osiyo sharoitida sizot suvlari
rejimi va tuproq shoʻrlanishini GAT yordamida baholash M.Ibrakhimov,
I.Forkutsa, A.Akramxonov, Sh.Kenjabayev va A.Pulatov tadqiqotlarida
amalga oshirilgan {[}28, 55, 70, 75, 76, 103{]}.

Bu tadqiqotlarda meliorativ tizimlarni loyihalash nazariyasi va sizot
suvlarini xaritalash usullari chuqur ishlangan. Biroq quyidagi masalalar
yetarlicha oʻrganilmagan: koʻp yillik monitoring maʼlumotlarining
ishonchliligini baholash; drenaj tarmogʻining raqamli xaritasi boʻlmagan
hududlar uchun GAT bazasini tiklash; sizot suvi sathining drenaj
tarmogʻiga bogʻliqligini har bir quduq boʻyicha aniqlash; kollektor
taʼmirining samarasini monitoring maʼlumotlari bilan oʻlchash.

\textbf{Tadqiqotning dissertatsiya bajarilgan oliy taʼlim muassasasining
ilmiy-tadqiqot ishlari rejalari bilan bogʻliqligi.} Dissertatsiya
tadqiqoti Andijon qishloq xoʻjaligi va agrotexnologiyalar institutining
ilmiy-tadqiqot ishlari rejasiga muvofiq bajarilgan \toldir{{[}mavzu}
\toldir{yoki} \toldir{loyihaning} \toldir{nomi,} \toldir{raqami}
\toldir{va} \toldir{bajarilish} \toldir{muddati} \toldir{kiritiladi{]}}.

\textbf{Tadqiqotning maqsadi} GAT texnologiyalari asosida meliorativ
monitoring maʼlumotlarining ishonchliligini baholash, sizot suvlari
sathining hududlar va yillar boʻyicha oʻzgarishini va uning drenaj
tarmogʻiga bogʻliqligini aniqlash, kollektor-drenaj tizimining
ishonchliligi va taʼmir samarasini baholash hamda shu asosda meliorativ
tizimlar samaradorligi va ishonchliligini taʼminlash boʻyicha tavsiyalar
ishlab chiqishdan iborat.

\textbf{Tadqiqotning vazifalari:}

meliorativ tizimlar ishonchliligi va samaradorligini baholashning mavjud
usullarini tahlil qilish va hal etilmagan masalalarni aniqlash;

mavjud xaritalar va meliorativ kadastr asosida Marhamat tumani
meliorativ tizimining GAT bazasini yaratish va uning aniqligini
baholash;

monitoring maʼlumotlarini tekshirish tartibini ishlab chiqish,
2009--2024-yillar (2025-yil bahori bilan) kuzatuv maʼlumotlarining
ishonchliligini baholash va natijalarni Andijon viloyatining boshqa
tumanlari bilan solishtirish;

tekshirilgan maʼlumotlar asosida sizot suvlari sathining yillar,
mavsumlar va hududlar boʻyicha oʻzgarishini massivlar va quduqlar
kesimida aniqlash, sayoz sath va shoʻrlanish xavfi yuqori hududlarni
ajratish;

sizot suvi sathining drenaj tarmogʻi, sugʻorish kanallari va relyefga
bogʻliqligini har bir quduq boʻyicha baholash;

tizim elementlarining ishonchlilik koʻrsatkichlarini aniqlash, talab
etiladigan drenaj modulini suv balansi boʻyicha hisoblash va kadastrdagi
koʻrsatkichlar bilan solishtirish;

kollektor taʼmirining sizot sathiga taʼsirini, taʼmirning haqiqiy
narxini va xarajatlarning massivlar boʻyicha taqsimotini baholash;

monitoring tizimini takomillashtirish, taʼmir samarasini oʻlchash va
mablagʻlarni birinchi navbatdagi massivlarga yoʻnaltirish boʻyicha
tavsiyalar ishlab chiqish.

\textbf{Tadqiqotning obyekti} Andijon viloyati Marhamat tumani
sugʻoriladigan yerlarining meliorativ tizimi: umumiy uzunligi 265,8 km
boʻlgan 90 ta kollektor, 370,7 km uzunlikdagi 332 ta xoʻjalik ichki
zovuri va sakkizta massivdagi 188 ta kuzatuv qudugʻi. Solishtirish uchun
Andijon viloyatining 14 ta tumanidagi 1 972 ta kuzatuv qudugʻi
maʼlumotlari olingan.

\textbf{Tadqiqotning predmeti} monitoring maʼlumotlarining
ishonchliligi, sizot suvlari sathi va minerallashuvining hududlar va
yillar boʻyicha oʻzgarishi, ularning drenaj tarmogʻiga bogʻliqligi hamda
tizim ishonchliligi va taʼmir samarasini baholash usullari.

\textbf{Tadqiqotning usullari.} Tadqiqotda GAT tahlili (xaritalarni
koordinataga bogʻlash, tarmoqni xaritadan raqamlash, massiv
chegaralarini tiklash), geostatistika (variogramma va kriging, xarita
aniqligini tekshirish, ehtimollik xaritalari), matematik statistika
(Mann--Kendall va Pettitt mezonlari, Sen qiyaligi, korrelyatsiya va
regressiya), taʼmirdan oldingi va keyingi sathni taʼmirlangan va
taʼmirlanmagan kollektorlar yonidagi quduqlarda solishtirish usuli,
FAO-56 usulida suv balansi va Manning formulasi boʻyicha gidravlik hisob
qoʻllanilgan.

\textbf{Tadqiqotning ilmiy yangiligi} quyidagilardan iborat:

koʻp yillik monitoring maʼlumotlarining ishonchliligini baholash usuli
ishlab chiqilgan: yil ichidagi oʻzgarmaslik, oʻlchov birligi va kuzatuv
tarmogʻi tarkibi boʻyicha uchta tekshiruv qoidasi va quduq qurigan
holatlarni aniqlash tartibi asoslangan; maʼlumotlardagi nuqsonlar
tasodifiy emas, sath chuqurligiga bogʻliq ekani Marhamat tumani
massivlari va Andijon viloyatining 14 ta tumani misolida koʻrsatilgan;

sizot suvi sathining koʻp yillik oʻzgarishini faqat doimiy kuzatilgan
quduqlar boʻyicha baholash tartibi asoslangan; tekshiruvsiz hisob
oʻzgarish tezligini 1,4 baravar oshirib koʻrsatishi va tuman oʻrtachasi
massivlardagi qarama-qarshi oʻzgarishlarni yashirishi aniqlangan;

oddiy sath xaritasi (kriging) sayoz sathli maydonni oʻn baravardan ortiq
kam koʻrsatishi aniqlangan; uning oʻrniga sathning kritik chuqurlikdan
sayoz boʻlish ehtimoli xaritasi va xarita aniqligini oshirish uchun
zarur quduqlar sonini aniqlash usuli taklif etilgan;

mavjud xaritalardan drenaj tarmogʻining GAT bazasini tiklash usuli
ishlab chiqilgan: kuzatuv quduqlari tayanch nuqta boʻlib xizmat qilgan,
massiv chegaralari tiklangan, kollektorlar kadastr bilan bogʻlangan,
tarmoq koʻrsatkichlari va taʼmir xarajatlari massivlar boʻyicha
taqsimlangan;

sizot suvi sathining kollektorgacha masofa va tarmoq zichligiga
bogʻliqligi har bir quduq boʻyicha aniqlangan; bir yildagi turli
joylarni solishtirish drenaj samarasini koʻrsatmasligi asoslangan;
kollektordan uzoqda sath tezroq chuqurlashgani aniqlangan;

kollektor taʼmirining sizot sathiga taʼsirini monitoring maʼlumotlari
bilan baholash usuli taklif etilgan: taʼmirdan oldingi va keyingi sath
taʼmirlangan va taʼmirlanmagan kollektorlar yonidagi quduqlarda
solishtiriladi; usulning qoʻllanish chegaralari va taʼsirni aniqlash
uchun zarur quduqlar soni asoslangan.

\textbf{Tadqiqotning amaliy natijalari} quyidagilardan iborat:

Marhamat tumani meliorativ tizimining GAT bazasi yaratilgan: kuzatuv
quduqlari, kollektorlar, kanallar va massiv chegaralari kadastr
maʼlumotlari bilan birga;

2009--2024-yillar monitoring maʼlumotlarining tekshirilgan bazasi va
tekshiruv dasturi yaratilgan, dastur viloyat miqyosida (854 676 ta
oʻlchov) sinovdan oʻtkazilgan;

sizot suvi sathi, xarita xatosi va yozgi sayoz sath ehtimoli xaritalari
tuzilgan; 2019--2024-yillardagi yozgi sath va 2004--2010-yillardagi
minerallashuv boʻyicha shoʻrlanish xavfi ostidagi 56 ta quduq aniqlangan
va massivlar uchta ustuvorlik guruhiga ajratilgan (birinchi navbatda
Z.M.Bobur va Istiqlol massivlari);

kollektor taʼmirining 2026-yil narxlaridagi haqiqiy narxi (1 km uchun
oʻrtacha 348 mln soʻm) va taʼmir turlari orasidagi narx farqi
aniqlangan; taʼmir turini oʻzan holatiga qarab tanlash 2017--2020-yillar
xarajatlarini 22--44 foizga kamaytirishi mumkinligi koʻrsatilgan;

monitoring tizimini takomillashtirish, taʼmir navbatini aniqlash va
taʼmir samarasini oʻlchash boʻyicha 15 ta tavsiya ishlab chiqilgan.

\textbf{Tadqiqot natijalarining ishonchliligi} quyidagilar bilan
izohlanadi: rasmiy monitoring va kadastr maʼlumotlari ishlatilgan va
oldindan tekshirilgan; xaritani bogʻlash aniqligi 181 ta tayanch nuqta
boʻyicha baholangan; xaritalar aniqligi qoʻshimcha tekshiruv bilan
tasdiqlangan; asosiy natijalar chegaralar va shartlar oʻzgartirilganda
ham saqlanishi tekshirilgan; statistik xulosalar ishonch oraliqlari
bilan berilgan; natijalar viloyatning boshqa tumanlarida sinalgan;
barcha hisoblarni dastur orqali takrorlash mumkin. Tadqiqotning
cheklovlari ham koʻrsatilgan: dala oʻlchovlari oʻtkazilmagan,
minerallashuv boʻyicha faqat 2004--2010-yillar maʼlumotlari yaroqli,
kadastrda har bir kollektor uchun faqat oxirgi taʼmir yili qayd etilgan,
kollektorlarning koʻndalang kesimi va oʻlchangan sarflari yoʻq.

\textbf{Tadqiqot natijalarining ilmiy va amaliy ahamiyati.} Ilmiy
ahamiyati shundaki, meliorativ tizimlarni baholashda monitoring
maʼlumotlarini tekshirish alohida bosqich sifatida asoslangan,
maʼlumotlardagi nuqsonlar tasodifiy emasligi koʻrsatilgan, sizot suvi
sathi va drenaj tarmogʻi orasidagi bogʻlanish har bir quduq boʻyicha
aniqlangan, drenaj samarasini turli joylarni emas, bir joyning oʻzini
vaqt boʻyicha solishtirib baholash kerakligi koʻrsatilgan.

Amaliy ahamiyati shundaki, suv xoʻjaligi tashkilotlari monitoring
maʼlumotlarini tekshirish vositasi, drenaj tarmogʻining raqamli bazasini
tiklash usuli, meliorativ xavf xaritalari, birinchi navbatdagi massivlar
roʻyxati, taʼmirning haqiqiy narxi va taʼmir samarasini oʻlchash uchun
kuzatuv tarmogʻiga talablarni oladi. Dissertatsiya natijalaridan
Marhamat tumani sugʻoriladigan yerlarining meliorativ holatini oʻrganish
va rejimini maqbullashtirish boʻyicha tumanga topshirilgan tavsiyanomada
foydalanilgan {[}33{]}.

\textbf{Tadqiqot natijalarining aprobatsiyasi.} Tadqiqot natijalari
\toldir{{[}soni{]}} ta xalqaro va \toldir{{[}soni{]}} ta respublika
ilmiy-amaliy anjumanlarida muhokamadan oʻtkazilgan
\toldir{{[}anjumanlar} \toldir{nomi,} \toldir{joyi} \toldir{va}
\toldir{sanasi} \toldir{kiritiladi{]}}.

\textbf{Tadqiqot natijalarining eʼlon qilinganligi.} Dissertatsiya
mavzusi boʻyicha jami \toldir{{[}soni{]}} ta ilmiy ish chop etilgan,
shulardan Oʻzbekiston Respublikasi Oliy attestatsiya komissiyasining
dissertatsiyalar asosiy ilmiy natijalarini chop etish tavsiya etilgan
ilmiy nashrlarida \toldir{{[}soni{]}} ta maqola, jumladan
\toldir{{[}soni{]}} tasi respublika va \toldir{{[}soni{]}} tasi xorijiy
jurnallarda nashr etilgan.

\textbf{Dissertatsiyaning tuzilishi va hajmi.} Dissertatsiya kirish,
toʻrtta bob, umumiy xulosalar, 139 nomdagi foydalanilgan adabiyotlar
roʻyxati va ilovalardan iborat. Dissertatsiyaning asosiy matni 120 betni
tashkil etadi, unda 9 ta jadval, 16 ta rasm (shundan 14 tasi GAT asosida
tuzilgan xarita va sxema) va 21 ta formula keltirilgan; yordamchi jadval
va rasmlar 74 ta ilovada berilgan.

Birinchi bobda meliorativ tizimlar ishonchliligi va samaradorligini
baholashning mavjud usullari, meliorativ monitoringda GAT va
geostatistik usullarning qoʻllanilishi tahlil qilingan, hal etilmagan
masalalar aniqlangan hamda tadqiqotning maqsad va vazifalari
belgilangan.

Ikkinchi bobda tadqiqot obyektining tabiiy-meliorativ sharoiti,
meliorativ tizimning tarkibi, GAT bazasini yaratish tartibi, monitoring
maʼlumotlarini tekshirish qoidalari va qoʻllanilgan tahlil usullari
bayon etilgan.

Uchinchi bobda monitoring maʼlumotlarining ishonchliligi baholangan.
Sizot suvlari sathining yillar va mavsumlar boʻyicha oʻzgarishi, hududiy
taqsimoti, sayoz sath ehtimoli va drenaj tarmogʻi, kanallar va relyefga
bogʻliqligi aniqlangan. Minerallashuv va shoʻrlanish xavfi tahlil
qilingan, natijalar Andijon viloyatining 14 ta tumani bilan
solishtirilgan.

Toʻrtinchi bobda tizim elementlarining ishonchlilik koʻrsatkichlari,
talab etiladigan va haqiqiy drenaj moduli, kollektorlarning oʻtkazish
qobiliyati, taʼmirning sizot sathiga taʼsiri, narxi va massivlar
boʻyicha taqsimoti baholangan, ustuvor massivlar aniqlangan, monitoring
tizimini takomillashtirish boʻyicha tavsiyalar hamda ularning iqtisodiy
asosi keltirilgan.

\oakbob{I BOB. MELIORATIV TIZIMLAR SAMARADORLIGI VA ISHONCHLILIGINI BAHOLASHNING
HOZIRGI HOLATI}

\oakpar{1.1-§. Sugʻoriladigan yerlarda kollektor-drenaj tizimlarining holati va
muammolari}

\textbf{Shoʻrlanish va zax bosish --- sugʻorma dehqonchilikning doimiy
xavfi.} Qurgʻoqchil mintaqalarda sugʻorma dehqonchilik ikki oʻzaro
bogʻliq jarayon --- ikkilamchi shoʻrlanish va zax bosish bilan birga
kechadi. Sugʻorish suvi bilan dalaga har yili maʼlum miqdorda tuz
keladi. Suv tuproq va oʻsimlik orqali bugʻlanib ketadi, tuz esa tuproqda
qoladi. Sugʻorish suvining bir qismi va kanallardan sizib chiqqan suv
sizot suvlariga qoʻshilib, ularning sathini koʻtaradi. Sath maʼlum
chuqurlikdan sayozlashganda suv kapillyarlar orqali tuproq yuzasigacha
koʻtariladi va bugʻlanib, tuzni yuzada qoldiradi. Bu jarayonning nazariy
asoslari V.A.Kovda, V.V.Yegorov va N.G.Minashina ishlarida ishlab
chiqilgan {[}15, 18, 19, 22{]}. Tuproq shoʻrlanishini aniqlash va
tasniflashning xalqaro amaliyoti L.A.Richards tahriridagi qoʻllanmaga
hamda I.Szabolcs ishlariga tayanadi {[}107, 123{]}. Sugʻorish suvining
sifati va ekinlarning tuzga chidamliligi boʻyicha meʼyoriy
koʻrsatkichlar R.S.Ayers va D.W.Westcot, E.V.Maas va G.J.Hoffman
tomonidan umumlashtirilgan {[}35, 86{]}.

Muammoning koʻlami global xarakterga ega. FAOning 2024-yilgi baholashiga
koʻra, dunyoda shoʻrlangan tuproqlar 1 381 mln gektarni, yaʼni
quruqlikning 10,7 foizini egallaydi; sugʻoriladigan ekin maydonlarining
taxminan 10 foizi shoʻrlanishdan zarar koʻrgan {[}53{]}. Shoʻrlanishning
sabablari, tarqalishi va unga qarshi kurash strategiyalari F.Ghassemi va
hammualliflari, P.Rengasamy, M.Qadir va hammualliflari, S.A.Shahid va
hammualliflari ishlarida tahlil qilingan {[}57, 104, 105, 118{]}.
Ularning umumiy xulosasi: sugʻoriladigan yerlarda tuzni boshqarishning
ishonchli yoʻli ortiqcha suv va tuzni drenaj orqali chiqarishdir,
drenajsiz sugʻorish esa uzoq muddatda yerni ishdan chiqaradi.

\textbf{Drenaj --- asosiy muhandislik chorasi.} Drenaj tizimi sizot
suvlari sathini ekin uchun xavfsiz chuqurlikda ushlab turadi va shoʻr
yuvish suvlarini hududdan chiqaradi. Ochiq va yopiq gorizontal drenaj,
vertikal drenaj quduqlari va kollektorlar yagona tizim boʻlib ishlaydi:
dala zovuri yigʻgan suv xoʻjaliklararo kollektorga, undan magistral
kollektorga oʻtadi. Tizimning biror boʻgʻinida suv oʻtkazish qobiliyati
pasaysa, undan yuqoridagi barcha boʻgʻinlarda suv sathi koʻtariladi va
drenajning ishlash chuqurligi kamayadi. Drenajni loyihalash nazariyasi
A.N.Kostyakov, S.F.Averyanov, I.P.Aydarov ishlarida, xorijda
S.B.Hooghoudt, L.F.Ernst va D.Kirkham tadqiqotlarida yaratilgan {[}9,
11, 20, 52, 69, 78{]}. Zamonaviy loyihalash, qurish va foydalanish
amaliyoti H.P.Ritzema tahriridagi qoʻllanmada hamda L.K.Smedema,
W.F.Vlotman va D.W.Rycroft kitobida jamlangan {[}108, 122, 129{]}.
Drenajning sizot sathiga taʼsiri dala tajribalarida yaxshi koʻrinadi:
Tarim havzasida yopiq drenaj drenajsiz dalaga nisbatan sizot sathini
chuqurlashtirgan va uning tebranishini kamaytirgan, drenlar orasi
qisqargan sari taʼsir kuchaygan {[}133{]}. Biroq bunday tajribalar
natijasini ishlab turgan katta tarmoqqa toʻgʻridan-toʻgʻri koʻchirib
boʻlmaydi.

Ochiq kollektor-drenaj tarmogʻining oʻziga xos xususiyati uning vaqt
oʻtishi bilan muqarrar yomonlashishidir. Oʻzanga loyqa choʻkadi,
qirgʻoqlar oʻpirilib tushadi, tub va yonbagʻirlarni qamish va boshqa
oʻsimliklar bosadi. Natijada koʻndalang kesim torayadi, gʻadir-budurlik
ortadi va kollektor loyihaviy sarfni oʻtkaza olmay qoladi. Shu sababli
ochiq tarmoqni muntazam tozalash va taʼmirlash kerak. Bu ishlar
toʻxtasa, tizim rasman mavjud boʻlgani holda oʻz vazifasini bajarmay
qoʻyadi {[}108, 122{]}.

\textbf{Markaziy Osiyodagi holat.} Orol dengizi havzasi dunyodagi eng
yirik sugʻorma dehqonchilik mintaqalaridan biri boʻlib, bu yerda drenaj
masalasi alohida keskinlikka ega. V.A.Duxovniy va hammualliflarining
maʼlumotiga koʻra, havzada 5,35 mln gektar yerda turli drenaj tizimlari
qurilgan. 1990-yillargacha drenajga katta mablagʻ ajratilgan, keyin
moliyalashtirish keskin kamaygan. Natijada tizimlarga yetarli xizmat
koʻrsatilmay qolgan, shoʻrlangan va zax bosgan maydonlar kengaya
boshlagan {[}51{]}. Mualliflar buning sababini faqat mablagʻ
yetishmasligida emas, drenaj tarmogʻiga masʼul tashkilotlarning
zaiflashganida ham koʻradi. Mintaqada suv va yer resurslarini boshqarish
muammolari V.A.Duxovniy va V.I.Sokolov, R.Kulmatov ishlarida {[}14,
80{]}, tuproqlarning tabiiy va antropogen shoʻrlanishi Ye.I.Pankova va
hammualliflari monografiyasida {[}25{]} umumlashtirilgan. Masofadan
zondlash maʼlumotlariga koʻra, 1990--2018-yillarda Markaziy Osiyoda
tuproqning elektr oʻtkazuvchanligi oʻrtacha 6,86 foizga oshgan {[}92{]}.
Amudaryo deltasida sizot sathi 1999--2002-yillarda 1,89 m dan 3,50 m
gacha pasaygan, 2017-yilgacha esa 2,11 m gacha qaytgan; sugʻoriladigan
hududda uning oʻzgarishi sugʻorish bilan chambarchas bogʻliq {[}100{]}.

Oʻzbekiston drenaj tizimining hozirgi holatini Y.Okuda va hammualliflari
maxsus oʻrgangan {[}96{]}. Ularning kuzatishicha, ayrim kollektorlarda
loyqa choʻkkani sababli suv sathi koʻtarilgan, dala drenajlarining ogʻzi
kollektordagi suvdan pastda qolgan yoki tuproq bilan koʻmilgan. Bunday
holda sizot suvi kollektorga oqib chiqa olmaydi va drenaj sizot sathini
pasaytira olmaydi. Bu kuzatuv muhim xulosaga olib keladi: kollektorning
holati faqat uning oʻziga emas, unga ulangan butun quyi tarmoqning
ishiga taʼsir qiladi.

\textbf{Oʻzbekistondagi holat va davlat siyosati.} Suv xoʻjaligini
rivojlantirishning 2020--2030-yillarga moʻljallangan konsepsiyasiga
koʻra, respublikada sugʻoriladigan yerlarning meliorativ holatini
yaxshilash uchun umumiy uzunligi 142,9 ming km boʻlgan kollektor-drenaj
tarmogʻidan foydalaniladi, shundan 106,2 ming km ochiq va 36,7 ming km
yopiq-yotiq tarmoq. Konsepsiyada 14,5 ming km kollektor-drenaj
tarmogʻini, 93 ta meliorativ nasos stansiyasini va 1 530 ta vertikal
drenaj qudugʻini rekonstruksiya qilish va yangidan qurish zarurligi
koʻrsatilgan {[}2{]}. Bu mavjud tarmoqning taxminan oʻndan bir qismiga
toʻgʻri keladi.

Meliorativ holatni yaxshilash ishlari davlat dasturlari asosida olib
boriladi. 2018--2019-yillarga moʻljallangan Davlat dasturi {[}5{]}
kollektor-drenaj tarmoqlarini tozalash va taʼmirlashni nazarda tutgan.
Marhamat tumani kadastri maʼlumotlariga koʻra, tumandagi 90 ta
kollektorning 55 tasi 2016--2018-yillarda, yana 16 tasi
2019--2020-yillarda taʼmirlangan (4.1-§). 2025-yilda qabul qilingan suv
resurslarini boshqarish va irrigatsiya sektorini rivojlantirishning
2025--2028-yillarga moʻljallangan dasturida meliorativ tadbirlar
hisobiga shoʻrlangan maydonlarni 1 898 ming gektardan 1 732 ming
gektarga, shu jumladan oʻrtacha va kuchli shoʻrlangan yerlarni 578 ming
gektardan 460 ming gektarga kamaytirish maqsad qilingan {[}4{]}. Bunday
aniq maqsadlarga erishilganini bilish uchun meliorativ holat
toʻgʻrisidagi maʼlumotlar ishonchli boʻlishi kerak.

\textbf{Andijon viloyati va Marhamat tumani.} Viloyat meliorativ
ekspeditsiyasining 2021-yil 1-yanvar holatidagi kadastriga koʻra,
Andijon viloyatida 7 616,3 km ochiq va 786,6 km yopiq-yotiq
kollektor-drenaj tarmogʻi mavjud (1-ilova). Ochiq tarmoqning 4 357,4 km
qismi, yaʼni 57 foizi, suv isteʼmolchilari uyushmalari hisobidagi
xoʻjalik ichki zovurlaridir. Tarmoqda 6 076 ta gidrotexnik inshoot qayd
etilgan.

Marhamat tumani ochiq tarmoq uzunligi boʻyicha viloyatda beshinchi,
magistral kollektorlar uzunligi boʻyicha birinchi, yopiq-yotiq drenaj
boʻyicha uchinchi oʻrinda turadi. Tumandagi tarmoq viloyat ochiq
tarmogʻining 8,2 foizini, yopiq drenajining 11,7 foizini tashkil etadi.
Demak, tuman viloyat uchun xos va tarmogʻi murakkab obyekt: unda
tarmoqning barcha turlari --- magistral, tumanlararo va xoʻjaliklararo
kollektorlar, xoʻjalik ichki zovurlari va yopiq drenaj bor.

\textbf{Muammolarning tizimlashtirilishi.} Koʻrib chiqilgan manbalar va
kadastr maʼlumotlari asosida kollektor-drenaj tizimlarining hozirgi
muammolarini toʻrt guruhga ajratish mumkin.

Birinchi guruh --- jismoniy eskirish. Tarmoqning asosiy qismi oʻtgan
asrning oʻrtalarida qurilgan; Marhamat tumanida kollektorlarning yoshi
medianasi 61 yil (4.1-§). Eskirgan tarmoq tezroq loyqalanadi va tez-tez
taʼmir talab qiladi.

Ikkinchi guruh --- masʼuliyatning boʻlinishi. Tarmoqning yarmidan koʻpi
suv isteʼmolchilari uyushmalari hisobida boʻlib, davlat dasturlari
asosan davlat hisobidagi kollektorlarni qamraydi. Natijada tizimning
quyi boʻgʻini --- dala bilan bevosita bogʻlangan ichki zovurlar ---
xizmat koʻrsatishdan chetda qoladi. V.A.Duxovniy va hammualliflari ham
institutsional zaiflikni tizim yomonlashuvining asosiy sabablaridan biri
sifatida koʻrsatgan {[}51{]}.

Uchinchi guruh --- axborot taʼminoti. Tarmoqning holati va sizot suvlari
rejimi toʻgʻrisidagi maʼlumotlar turli jadval va qogʻoz xaritalarda
yuritiladi, yagona raqamli asosga ega emas va rasmiy tekshiruvdan
oʻtkazilmaydi.

Toʻrtinchi guruh --- natijani oʻlchash. Taʼmir ishlarining natijasi
bajarilgan ish hajmi bilan oʻlchanadi; ularning sizot rejimiga taʼsiri
baholanmaydi.

Dastlabki ikki guruh moliyaviy va tashkiliy masalalar, ular ushbu
tadqiqot doirasidan tashqarida. Keyingi ikki guruh esa ilmiy masalalar:
ular hal etilmasa, ajratilgan mablagʻ qanchalik samarali
sarflanayotganini bilib boʻlmaydi. Mazkur dissertatsiya aynan shu ikki
masalaga --- maʼlumotlarning ishonchliligi va natijani oʻlchashga
qaratilgan. Keyingi paragraflarda ular boʻyicha mavjud usullar koʻrib
chiqiladi.

\oakpar{1.2-§. Meliorativ tizimlar ishonchliligi va samaradorligini baholash
usullari}

\textbf{Asosiy tushunchalar.} Texnikada ishonchlilik --- obyektning oʻz
vazifasini belgilangan sharoitda va belgilangan muddat davomida bajara
olishi. U buzilmaslik, chidamlilik, taʼmirlanuvchanlik va
saqlanuvchanlikni oʻz ichiga oladi {[}8{]}. Kollektor-drenaj tarmogʻi
uchun ishonchlilik --- uning foydalanish davomida sizot suvlarini
loyihaviy chuqurlikda ushlab turishi va hisobiy drenaj oqimini oʻtkaza
olishidir. Samaradorlik esa olingan natijaning sarflangan mablagʻga
nisbati. Meliorativ tizim uchun natija --- sizot suvi va tuproq
shoʻrlanishidagi yaxshilanish, sarf --- qurilish, taʼmir va foydalanish
xarajatlari.

Ikki tushuncha oʻzaro bogʻliq, lekin bir xil emas. Ishonchliligi yuqori
tizim ham samarasiz boʻlishi mumkin --- masalan, drenajga ehtiyoj
boʻlmagan joyda qurilgan boʻlsa. Ishonchliligi pasaygan tizim esa
loyihaviy samarani bera olmaydi. Shuning uchun har ikkalasi alohida
baholanishi kerak.

Adabiyotda meliorativ tizimni baholashning uchta yondashuvi shakllangan.
Ular nimani baholashi, qanday maʼlumot talab qilishi va cheklovlari
bilan farq qiladi (1.1-jadval).

\oakjadvalnum{1.1-jadval}

\oakilovanom{\textbf{Meliorativ tizimni baholash yondashuvlarining qiyosi}}

{\fontsize{12}{14}\selectfont\setlength{\tabcolsep}{4pt}

\noindent\begin{longtable}[]{|
  >{\raggedright\arraybackslash}p{(\columnwidth - 8\tabcolsep - 5\arrayrulewidth) * \real{0.3131}}|>{\raggedright\arraybackslash}p{(\columnwidth - 8\tabcolsep - 5\arrayrulewidth) * \real{0.2342}}|>{\raggedright\arraybackslash}p{(\columnwidth - 8\tabcolsep - 5\arrayrulewidth) * \real{0.2184}}|>{\raggedright\arraybackslash}p{(\columnwidth - 8\tabcolsep - 5\arrayrulewidth) * \real{0.2342}}|}
\hline
\begin{minipage}[b]{\linewidth}\raggedright
\textbf{Yondashuv}
\end{minipage} & \begin{minipage}[b]{\linewidth}\raggedright
\textbf{Nimani baholaydi}
\end{minipage} & \begin{minipage}[b]{\linewidth}\raggedright
\textbf{Zarur maʼlumot}
\end{minipage} & \begin{minipage}[b]{\linewidth}\raggedright
\textbf{Asosiy cheklovi}
\end{minipage} \\
\hline
\endhead
\hline
\endlastfoot
Loyihaviy-gidravlik & Tizim nimaga moʻljallangan: drenaj moduli,
drenajlar oraligʻi va chuqurligi, oʻtkazish qobiliyati & Loyiha
hujjatlari, tuproq va gidrogeologik parametrlar & Haqiqiy holat
loyihadan qancha uzoqlashganini koʻrsatmaydi \\ \hline
Rejim koʻrsatkichlari boʻyicha & Tizim ishining natijasi: sizot sathi,
minerallashuv, tuproq shoʻrlanishi & Monitoring quduqlari va shoʻr
sʼyomkasi maʼlumotlari & Rejimga drenajdan tashqari omillar ham taʼsir
qiladi; maʼlumot sifatiga bogʻliq \\ \hline
Ishonchlilik nazariyasi & Elementlarning nosozliksiz ishlashi va
tiklanishi & Har bir element boʻyicha nosozlik va taʼmirlar hisobi &
Bunday hisob amalda yuritilmaydi \\
\end{longtable}

}

\textbf{Loyihaviy-gidravlik yondashuv.} Bu yondashuvda tizim loyihaviy
koʻrsatkichlar bilan baholanadi. Asosiy koʻrsatkich --- drenaj moduli,
yaʼni bir gektardan bir soniyada chiqariladigan suv miqdori. U suv-tuz
balansidan aniqlanadi {[}9, 11, 20{]}. Drenajlar oraligʻi S.B.Hooghoudt,
L.F.Ernst va boshqa tenglamalar boʻyicha hisoblanadi {[}52, 69, 108{]}.
Kollektorlarning oʻtkazish qobiliyati ochiq oʻzanlar gidravlikasi
formulalari bilan tekshiriladi {[}13, 21{]}. Oʻzbekistonda loyihalash
meʼyorlari tegishli shaharsozlik normalarida belgilangan {[}6{]}. Drenaj
tizimlari ishini modellashtirish uchun R.W.Skaggs va hammualliflarining
DRAINMOD, R.J.Oosterbaanning SALTMOD modellari keng qoʻllaniladi {[}99,
120, 121{]}.

Yondashuvning kuchli tomoni --- u fizik qonunlarga asoslangan. Cheklovi
--- u tizimning loyihadagi holatini koʻrsatadi. Loyqa, oʻsimlik va
qirgʻoqlar oʻpirilishi tufayli kollektorning haqiqiy kesimi oʻzgaradi.
Buni hisobga olish uchun har bir kollektorni muntazam oʻlchash kerak,
amalda esa bu qilinmaydi.

\textbf{Rejim koʻrsatkichlari boʻyicha baholash.} Amaliyotda eng keng
qoʻllaniladigan yondashuv shu. Oʻzbekistonda sugʻoriladigan yerlarning
meliorativ holati uchta koʻrsatkich --- sizot suvlarining chuqurligi,
minerallashuvi va tuproqlarning shoʻrlanishi boʻyicha baholanadi. Yerlar
shu koʻrsatkichlarga qarab yaxshi, qoniqarli va qoniqarsiz meliorativ
holat toifalariga ajratiladi {[}16{]}. Davlat dasturlari ham shu
koʻrsatkichlarga tayanadi: masalan, 2018--2019-yillar dasturida
respublikada 1 957 ming gektar sugʻoriladigan yer turli darajada
shoʻrlangani, shundan 542 ming gektari oʻrtacha va 99 ming gektari
kuchli shoʻrlangani koʻrsatilgan {[}5{]}. Monitoring va kadastrni
yuritish tartibi R.Q.Ikromov va hammualliflari qoʻllanmasida batafsil
bayon etilgan {[}16{]}. Soʻnggi yillarda bu koʻrsatkichlar GAT muhitida
tuproq, relyef va tarmoq qatlamlari bilan birga ballar yigʻindisi
koʻrinishida baholanmoqda: Xorazmdagi Qulovot kanali tizimida yerlar shu
yoʻl bilan uch holat toifasiga ajratilgan {[}75{]}, Jizzax viloyatida
tumanlar 14 ta koʻrsatkich boʻyicha baholangan {[}32{]}. Marhamat tumani
massivlari uchun shu usul tavsiyanomada qoʻllanilgan {[}33{]}.

Bu yondashuvning afzalligi --- u meliorativ natijaning oʻzini oʻlchaydi.
Biroq uning ikki jiddiy cheklovi bor. Birinchisi --- sabab masalasi.
Sizot suvi sathi faqat drenajga emas, sugʻorish hajmi va usuliga,
kanallardan sizib chiqqan suvga, relyef va ob-havoga ham bogʻliq {[}10,
17{]}. Shuning uchun holatning oʻzgarishini drenajga bogʻlash uchun
boshqa omillar taʼsirini ajratib olish kerak. Ikkinchisi --- maʼlumot
masalasi: baho toʻliq monitoring maʼlumotlarining toʻgʻriligiga tayanadi
(1.4-§).

\textbf{Ishonchlilik nazariyasi.} Gidromeliorativ inshootlarga
ishonchlilik nazariyasini qoʻllashning asoslari Ts.Ye.Mirtsxulava
ishlarida ishlab chiqilgan {[}23, 24{]}. Bu yondashuvda tizim
elementlari ularning qanchalik tez-tez buzilishi bilan tavsiflanadi. Eng
sodda holda element t vaqt davomida nosozliksiz ishlash ehtimoli

\begin{equation}\tag{1.1}
P(t) = e^{-\lambda t}
\end{equation}

koʻrinishida, taʼmirlanadigan elementning ishga tayyorlik koeffitsiyenti
esa

\begin{equation}\tag{1.2}
K_t = \frac{T_0}{T_0 + T_v}
\end{equation}

koʻrinishida ifodalanadi, bu yerda λ --- nosozliklar jadalligi; T₀ ---
nosozliklar orasidagi oʻrtacha ish muddati; Tᵥ --- tiklashning oʻrtacha
muddati. Ketma-ket ulangan elementlardan iborat tizimning ishonchliligi
elementlar ishonchliligining koʻpaytmasiga teng. Drenaj tarmogʻida dala
zovuri, xoʻjaliklararo va magistral kollektor aynan shunday ketma-ket
zanjir hosil qiladi. Mahalliy adabiyotda bu yondashuv meliorativ
inshootlarni taʼmirlashni xarajat va samara nisbati boʻyicha
rejalashtirish masalasi bilan birga {[}16{]} da keltirilgan.

Nazariya qatʼiy, ammo uni qoʻllash uchun har bir element boʻyicha
nosozliklar va taʼmirlarning koʻp yillik hisobi zarur. Amalda bunday
hisob yuritilmaydi: Marhamat tumani kadastrida har bir kollektor uchun
faqat oxirgi taʼmir yili qayd etilgan, nosozliklar esa umuman hisobga
olinmagan (4.1-§). Bundan tashqari, ochiq kollektor mashina detali kabi
birdan buzilmaydi: oʻzan asta-sekin loyqalanadi va suv oʻtkazishi
sekin-asta pasayadi. Shuning uchun kollektor ishonchliligini buzilishlar
soni bilan emas, vazifasini qanchalik bajarayotgani bilan oʻlchash
toʻgʻri.

\textbf{Xalqaro amaliyot: tizim ishini baholash koʻrsatkichlari.}
Sugʻorish va drenaj tizimlari ishini baholashning umumiy uslubiy asosi
M.G.Bos, M.A.Burton va D.J.Molden tomonidan Xalqaro irrigatsiya va
drenaj komissiyasi ishchi guruhi doirasida ishlab chiqilgan {[}40{]}.
Unda har bir koʻrsatkich haqiqiy qiymatning maqsadli qiymatga nisbati
sifatida aniqlanadi. Suv tanqisligi sharoitida sugʻorishni boshqarish
koʻrsatkichlari L.S.Pereira va hammualliflari tomonidan koʻrib chiqilgan
{[}101{]}.

Drenaj tizimlari uchun maxsus koʻrsatkichlar B.Vincent, W.F.Vlotman,
D.Zimmer va J.W.Hornbuckle hisobotida taklif etilgan {[}128{]}.
Mualliflarning qayd etishicha, drenaj tizimlari ishini baholash keng
tarqalmagan, chunki hozirgacha yaxlit usul yoʻq. Demak, drenaj
tizimlarini baholash usullari sugʻorish tizimlarinikiga qaraganda ancha
kam ishlangan; bu mazkur tadqiqot yoʻnalishini asoslaydi.

\textbf{Taʼmir samarasini oʻlchash masalasi.} Kollektorlarni tozalash va
taʼmirlash ishlari odatda bajarilgan hajm --- tozalangan kilometr,
chiqarilgan grunt hajmi, sarflangan mablagʻ --- bilan hisobot qilinadi.
Bu koʻrsatkichlar xarajatni koʻrsatadi, natijani emas. Taʼmir natijasini
bilish uchun uning sizot suvi sathiga taʼsirini oʻlchash kerak. Buning
uchun taʼmirdan oldingi va keyingi holatni solishtirish yetarli emas,
chunki shu yillar orasida ob-havo, suv taʼminoti va sugʻorish ham
oʻzgaradi. Tadbirning kutilgan samarasi xoʻjaliklarning harakati tufayli
ham oʻzgarishi mumkin: Oʻzbekistonda suv tejovchi texnologiyalar
boʻyicha oʻtkazilgan ekspert baholashi ularning tejamini qisman yoki
toʻliq yoʻqqa chiqaruvchi «qaytish samarasi» ehtimolini koʻrsatgan
{[}63{]}.

Bunday hollarda tadbir oʻtkazilgan va oʻtkazilmagan obyektlardagi
oʻzgarishlar solishtiriladi (iqtisodiyotda bu usul «farqlar farqi» deb
ataladi). Shunda barcha obyektlarga bir xil taʼsir qilgan omillar
(ob-havo, suv taʼminoti) chiqarib tashlanadi {[}34{]}. Usul irrigatsiya
infratuzilmasiga ham qoʻllanilgan: K.Tilekeyev va D.Kirbasheva
Qirgʻizistonning janubida sugʻorish kanallarini rekonstruksiya
qilishning hosildorlikka taʼsirini fermerlarning 2016 va 2019-yillardagi
soʻrovlari asosida shu usulda baholagan {[}125{]}. Biroq kollektorlarni
taʼmirlashning sizot suvlari sathiga taʼsirini kuzatuv quduqlari
maʼlumotlari asosida shu usulda baholagan tadqiqotlar adabiyotda
uchramadi. Usulni drenajga qoʻllashda qiyinchilik ham bor: kollektorlar
yagona tarmoq boʻlib ishlaydi, shuning uchun bir kollektorning taʼmiri
qoʻshni hududlarga ham taʼsir qilishi mumkin va solishtiriladigan
quduqlar taʼsirdan xoli boʻlmaydi.

\textbf{Yondashuvlarning umumiy cheklovi.} Koʻrib chiqilgan barcha
yondashuvlar dastlabki maʼlumotlarni toʻgʻri deb oladi. Loyihaviy hisob
loyiha koʻrsatkichlariga, holat boʻyicha baho monitoring maʼlumotlariga,
ishonchlilik hisobi kadastr yozuvlariga tayanadi. Maʼlumotning oʻzi
tekshirilmasa, eng yaxshi usul ham notoʻgʻri xulosa beradi. Shuning
uchun maʼlumotlar ishonchliligini tizim ishonchliligidan oldin, alohida
bosqich sifatida baholash kerak (1.4-§).

\oakpar{1.3-§. Meliorativ monitoringda GAT va geostatistik usullarning
qoʻllanilishi}

\textbf{Geoaxborot tizimlarining nazariy asoslari.} Geografik axborot
tizimi --- joyga bogʻlangan maʼlumotlarni yigʻish, saqlash, tahlil
qilish va xaritada koʻrsatish uchun moʻljallangan dasturiy majmua. Uning
nazariy asoslari P.A.Burrough va R.A.McDonnell, M.F.Goodchild,
P.A.Longley va hammualliflari, C.D.Tomlin ishlarida ishlab chiqilgan
{[}41, 42, 58, 83, 126{]}. GATning melioratsiya uchun asosiy qimmati
shundaki, u turli maʼlumotlarni --- quduq oʻlchovlari, drenaj tarmogʻi
sxemasi, tuproq xaritasi, relyef, kadastr yozuvlari --- yagona
koordinatalar tizimida birlashtiradi. Shunda ular orasidagi masofa,
qoʻshnilik va ustma-ust tushishni hisoblash mumkin boʻladi.

GATda maʼlumotlar ikki xil saqlanadi. Birinchisida (vektor) obyektlar
nuqta (kuzatuv qudugʻi), chiziq (kollektor, kanal) va maydon (massiv,
kontur) koʻrinishida beriladi, har biriga maʼlumotlar jadvali
biriktiriladi. Ikkinchisida (rastr) hudud bir xil kataklarga boʻlinadi
va har bir katakka qiymat (sizot sathi, balandlik, shoʻrlanish)
beriladi. Melioratsiyada ikkalasi birga ishlatiladi: tarmoq va quduqlar
birinchi, uzluksiz yuzalar ikkinchi shaklda saqlanadi {[}41, 126{]}.

Xarita maʼlumotlarining sifati alohida masala. U toʻliqlik, izchillik,
joylashuv aniqligi, vaqt boʻyicha sifat va qiymatlar aniqligi bilan
tavsiflanadi. Bu tamoyillar S.C.Guptill va J.L.Morrison tahriridagi
toʻplamda ishlab chiqilgan va xalqaro standartda mustahkamlangan {[}7,
61{]}. J.Zhang va M.F.Goodchild xarita maʼlumotlaridagi xatolar
manbasini va ularning natijaga taʼsirini batafsil koʻrib chiqqan
{[}134{]}. Qogʻoz xaritani raqamlashda joylashuv aniqligi xaritani
koordinataga qanchalik toʻgʻri bogʻlashga bogʻliq. Bogʻlash koordinatasi
maʼlum tayanch nuqtalar bilan bajariladi, aniqligi esa shu nuqtalardagi
xato bilan baholanadi {[}41, 83{]}.

\textbf{Geostatistik modellashtirish.} Monitoring alohida nuqtalarda
oʻlchov beradi, qarorlar esa butun hudud uchun kerak. Nuqtalar orasidagi
qiymatlarni topish geostatistika usullari bilan hal etiladi. Bu
usullarning asoslari D.G.Krige va G.Matheron ishlarida yaratilgan {[}79,
91{]} va keyinchalik A.G.Journel va C.J.Huijbregts, E.H.Isaaks va
R.M.Srivastava, N.Cressie, P.Goovaerts, J.-P.Chilès va P.Delfiner,
R.Webster va M.A.Oliver tomonidan rivojlantirilgan {[}45, 48, 60, 71,
73, 131{]}. Soʻnggi yillarda kriging bilan birga mashinali oʻrganish
usullari ham qoʻllanmoqda. Ularning afzalligi --- baho bilan birga uning
xatosini ham koʻrsatib, xavfni hisobga olib qaror qabul qilishga imkon
beradi {[}49{]}.

Geostatistikaning markaziy tushunchasi --- variogramma. U bir-biridan h
masofadagi nuqtalar qiymatlari qanchalik farq qilishini koʻrsatadi:

\begin{equation}\tag{1.3}
\gamma(h) = \frac{1}{2N(h)} \sum_{i=1}^{N(h)} \left[ z(x_i) - z(x_i + h) \right]^2
\end{equation}

bu yerda z(x) --- x nuqtadagi qiymat; N(h) --- bir-biridan h masofadagi
nuqtalar juftlari soni. Oʻlchovlar boʻyicha hisoblangan variogrammaga
nazariy egri chiziq (sferik, eksponensial, Gauss va boshqalar)
moslashtiriladi. Uning uchta koʻrsatkichi bor: nagget --- juda yaqin
nuqtalar orasidagi tasodifiy farq; plato --- qiymatlarning umumiy
tarqoqligi; taʼsir radiusi --- qiymatlar bir-biriga bogʻliq boʻlib
qoladigan masofa {[}97, 131{]}. Nagget va plato nisbati qoʻshni nuqtalar
qanchalik oʻxshashligini koʻrsatadi: C.A.Cambardella va hammualliflari
bu nisbat 25 foizdan kichik boʻlsa bogʻlanishni kuchli, 75 foizdan katta
boʻlsa kuchsiz deb hisoblashni taklif etgan {[}44{]}.

Oddiy kriging usulida oʻlchanmagan x₀ nuqtadagi qiymat qoʻshni
oʻlchovlarning vaznli oʻrtachasi sifatida topiladi:

\begin{equation}\tag{1.4}
\hat{z}(x_0) = \sum_{i=1}^{n} \lambda_i \, z(x_i), \qquad \sum_{i=1}^{n} \lambda_i = 1
\end{equation}

Vaznlar λᵢ variogramma boʻyicha, baho xatosi eng kichik boʻladigan qilib
tanlanadi. Krigingning boshqa usullardan, masalan teskari masofa
usulidan {[}119{]}, asosiy afzalligi --- u har bir nuqtada baho bilan
birga uning xatosini ham beradi {[}60, 98{]}. J.Li va A.D.Heap sharhiga
koʻra, koʻp hollarda kriging boshqa usullardan aniqroq, lekin bu
quduqlar zichligiga bogʻliq {[}82{]}. Xarita aniqligi quyidagicha
tekshiriladi: har bir oʻlchov navbat bilan chiqarib tashlanadi,
qolganlari boʻyicha baholanadi va xato hisoblanadi (kross-validatsiya)
{[}60, 67{]}.

Geostatistik tahlilning amaliy cheklovi kuzatuv tarmogʻining zichligi
bilan bogʻliq. Agar quduqlar bir-biridan juda uzoq boʻlsa, kriging ular
orasidagi qiymatni ishonchli tiklay olmaydi va baho hudud oʻrtachasiga
yaqinlashadi. Shuning uchun xaritani berishdan oldin uning xatosini
maʼlumotlarning oʻz tarqoqligi bilan solishtirish kerak.

\textbf{Kuzatuv tarmogʻini loyihalash.} Kriging xatosidan kuzatuv
tarmogʻini loyihalashda ham foydalaniladi. H.A.Loáiciga va
hammualliflari tarmoqni loyihalash usullarini gidrogeologik va statistik
guruhlarga ajratgan {[}84{]}. Statistik yondashuvda yangi quduq xato eng
katta joyga qoʻyiladi; olib tashlanganda xarita deyarli oʻzgarmaydigan
quduq esa ortiqcha hisoblanadi. Bu yerda cheklov bor: kriging xatosi
asosan quduqlar joylashuviga bogʻliq, oʻlchangan qiymatlarga kam
bogʻliq, shuning uchun u xatoning faqat taxminiy oʻlchovidir. S.Teimoori
va hammualliflari qayd etganidek, sizot suvi sifati kuzatuv tarmoqlariga
qaraganda sath kuzatuv tarmoqlarini loyihalash kam oʻrganilgan va koʻp
hududlarda mavjud tarmoqlar qaror qabul qilish uchun yetarli maʼlumot
bermaydi {[}124{]}.

\textbf{Koʻp yillik oʻzgarishni aniqlash.} Koʻp yillik kuzatuv
maʼlumotlarida oʻzgarish yoʻnalishini aniqlash uchun odatda
Mann--Kendall mezoni qoʻllaniladi {[}74, 90{]}. Mezon statistikasi

\begin{equation}\tag{1.5}
S = \sum_{k=1}^{n-1} \sum_{j=k+1}^{n} \operatorname{sgn}(x_j - x_k)
\end{equation}

koʻrinishida hisoblanadi; S orqali oʻzgarish tasodifiy emasligi
tekshiriladi. Oʻzgarish tezligi P.K.Sen qiyaligi bilan baholanadi: u
barcha juft yillar orasidagi qiyaliklarning medianasiga teng {[}117{]}.
Usulning gidrologiyada qoʻllanilishi D.R.Helsel va R.M.Hirsch,
R.M.Hirsch va hammualliflari ishlarida, qoʻshni yillar qiymatlari
bogʻliq boʻlgan hol uchun tuzatish K.H.Hamed va A.R.Rao ishida berilgan
{[}62, 66, 68{]}. Mann--Kendall mezonining afzalligi --- u chetki
qiymatlarga chidamli. Biroq u barcha yillarda bir xil usul va bir xil
quduqlar bilan oʻlchangan deb oladi. Agar kuzatuv tarmogʻi tarkibi
oʻzgargan boʻlsa, mezon tarmoq oʻzgarishini sathning oʻzgarishi deb
koʻrsatadi.

\textbf{Masofadan zondlash.} Tuproq shoʻrlanishini va ekinlar holatini
katta hududlarda kuzatish uchun sunʼiy yoʻldosh tasvirlari qoʻllaniladi.
Oʻsimlik indekslari (masalan, J.W.Rouse va hammualliflari taklif etgan
NDVI) va shoʻrlanish indekslari tasvirning turli kanallari nisbatidan
hisoblanadi {[}72, 110{]}. Shoʻrlanishni masofadan baholashning
imkoniyatlari va cheklovlari G.I.Metternicht va J.A.Zinck, J.Farifteh va
hammualliflari, A.Allbed va L.Kumar, D.L.Corwin va E.Scudiero
sharhlarida tahlil qilingan {[}30, 47, 54, 93{]}. Ularning umumiy
xulosasi: indekslarning oʻzi yetarli emas, sunʼiy yoʻldosh
maʼlumotlarini yerdagi oʻlchovlar bilan birga ishlatish kerak.

\textbf{Koʻp mezonli baholash.} GATda bir nechta koʻrsatkich boʻyicha
baholash usullari, ayniqsa T.L.Saati usuli, keng qoʻllaniladi; ularning
GAT bilan birga ishlatilishi J.Malczewski ishlarida tizimlashtirilgan
{[}87--89, 112--115{]}. Bu usullarda bir nechta koʻrsatkich
mutaxassislar bergan vaznlar bilan bitta ballga birlashtiriladi.
Usulning kuchli tomoni --- tushunarliligi va mutaxassislar fikrini
hisobga olishi. Cheklovi --- natija vaznlarga bogʻliq va sabab-oqibatni
koʻrsatmaydi, shuning uchun vaznlar oʻzgarganda natija qanday
oʻzgarishini tekshirish kerak {[}95, 116{]}.

\textbf{Markaziy Osiyodagi tadqiqotlar.} Mintaqada GAT va geostatistika
asosan sizot suvlari rejimi va tuproq shoʻrlanishini oʻrganish uchun
qoʻllanilgan. Asosiy ishlar 2-ilovada umumlashtirilgan. Soʻnggi ishlarda
ham xaritalar asosan oddiy usullar bilan tuziladi: Qoraqalpogʻistonning
Xoʻjayli tumanida 142 ta kuzatuv qudugʻi maʼlumotlari boʻyicha
2015--2023-yillardagi sath va minerallashuv xaritalari teskari masofa
usulida tuzilgan {[}81{]}. Sirdaryo boʻyidagi sugʻoriladigan hududda
(Qozogʻiston) tuproq shoʻrlanishining fazoviy taqsimotini eng koʻp
tushuntiruvchi omillar sugʻorish kanallari va drenaj tarmogʻigacha
masofa boʻlgan {[}50{]}. Demak, tarmoqqa nisbatan joylashuv meliorativ
holatga taʼsir qiladi.

Jadvaldan koʻrinadiki, tadqiqotlarning aksariyati Quyi Amudaryo
hududiga, ayniqsa Xorazm viloyatiga toʻgʻri keladi. Fargʻona vodiysi,
xususan Andijon viloyati boʻyicha kuzatuv quduqlari maʼlumotlarini
hududlar va yillar boʻyicha tahlil qilgan ishlar kam.

Xorijiy tadqiqotlarda ham sizot suvlari sathining sugʻorish va drenaj
tarmogʻi bilan bogʻliqligi GAT yordamida oʻrganilgan. Masalan, J.Wang va
hammualliflari Avstraliyaning sugʻoriladigan hududida kuzatuv quduqlari
maʼlumotlarini GAT maʼlumotlar bazasiga yigʻib, toshqin va yogʻinning
sizot suvlarini toʻyintirishga taʼsirini va uning fazoviy koʻlamini
miqdoriy baholagan {[}130{]}.

\textbf{Mavjud ishlarning umumiy xususiyatlari.} Koʻrib chiqilgan
tadqiqotlarda deyarli doim sizot sathi, minerallashuv va tuproq
shoʻrlanishi baholanadi. Drenaj tarmogʻining joylashuvi, zichligi va
holati tahlilga kam kiritiladi: tarmoq odatda xaritada fon sifatida
koʻrsatiladi, lekin sizot suvlari bilan son jihatdan bogʻlanmaydi.

Monitoring maʼlumotlari odatda tayyor holda olinadi: ular yillar
davomida bir xil tartibda oʻlchanganmi, quduqlar tarkibi oʻzgarmaganmi,
qiymatlar fizik jihatdan mumkinmi --- bu tekshirilgani haqida maʼlumot
deyarli uchramaydi.

Tadqiqotlarning aksariyati bir yildagi hududiy taqsimotni yoki hudud
boʻyicha oʻrtacha oʻzgarishni tavsiflaydi. Tarmoqdagi oʻzgarish (yangi
qurilish yoki taʼmir) sizot suvi sathida qanday aks etganini alohida
quduqlarda yillar boʻyicha kuzatgan ishlar uchramadi.

Bundan tashqari, koʻp tumanlarda drenaj tarmogʻining raqamli qatlamlari
mavjud emas: tarmoq qogʻoz xarita va jadvallarda yuritiladi. Bunday
sharoitda mavjud xaritalardan GAT bazasini tiklash va uning aniqligini
baholash usuli amaliy ahamiyatga ega. Ushbu masalalar mazkur
dissertatsiyaning II va III boblarida hal etiladi.

\oakpar{1.4-§. Monitoring maʼlumotlarining sifati masalasi va tadqiqot
vazifalari}

\textbf{Maʼlumot sifatining boshqaruvdagi oʻrni.} Meliorativ monitoring
maʼlumotlari boshqaruv qarorlarining asosi hisoblanadi: yerlarning
meliorativ holati toifalari, taʼmir va rekonstruksiya rejalari, davlat
dasturlarining natijalari shu maʼlumotlar boʻyicha baholanadi {[}4, 5,
16{]}. 1.2-§ da koʻrsatilganidek, holat boʻyicha baholash ham,
ishonchlilik hisobi ham, xaritalash ham dastlabki maʼlumotlarni toʻgʻri
deb oladi. Agar maʼlumotda doimiy nuqson boʻlsa, u barcha hisoblarga
oʻtadi: oʻzgarish tezligi notoʻgʻri chiqadi, xarita haqiqiy holatni
emas, oʻlchash tartibini koʻrsatadi, dastur maqsadlariga erishilgani
notoʻgʻri baholanadi. Shuning uchun maʼlumot sifati alohida ilmiy
masala.

\textbf{Gidrologik kuzatuvlarda xato turlari.} Gidrologik va
gidrogeologik oʻlchovlardagi xatolar odatda uch guruhga ajratiladi
{[}66, 132{]}.

Birinchi guruh --- tasodifiy xatolar. Ular asbob aniqligi, oʻlchovchi
mahorati va sharoitga bogʻliq, qiymatni goh oshirib, goh kamaytirib
yuboradi va koʻp oʻlchovlar oʻrtachasida kamayadi.

Ikkinchi guruh --- doimiy (tizimli) xatolar. Ular koʻp oʻlchovlarni bir
tomonga siljitadi: notoʻgʻri sozlangan asbob, quduq ogʻzi balandligining
oʻzgarishi, oʻlchash usulining almashishi. Bunday xato oʻrtachada
kamaymaydi va sathning oʻzgarishi boʻlib koʻrinishi mumkin.

Uchinchi guruh --- yozish va hisobot xatolari. Ular oʻlchovdan keyin,
maʼlumotni yozish va jamlashda yuzaga keladi: birlikning almashtirilishi
(metr oʻrniga santimetr), raqamlarning oʻrin almashishi, oʻlchanmagan
holatda shartli qiymat yozilishi. Bu xatolar koʻpincha katta va oson
aniqlanadi, lekin faqat maʼlumot maxsus tekshirilganda.

Jahon meteorologiya tashkilotining gidrologik amaliyot boʻyicha
qoʻllanmasida kuzatuv maʼlumotlarini dastlabki tekshirish majburiy
bosqich sifatida belgilangan: qiymat fizik jihatdan mumkinmi, oldingi va
keyingi oʻlchovlarga mosmi, qoʻshni nuqtalarga mosmi --- tekshiriladi
{[}132{]}.

\textbf{Maʼlumotlarning bir jinsliligi.} Koʻp yillik maʼlumotlar bir
jinsli hisoblanadi, agar ulardagi oʻzgarishlar oʻlchash sharoitiga emas,
faqat oʻrganilayotgan jarayonga bogʻliq boʻlsa. Kuzatuv nuqtasining
koʻchirilishi, asbob yoki kuzatuvchining almashishi maʼlumotlarda sunʼiy
sakrash hosil qiladi. Bunday sakrashlarni topish uchun A.N.Pettitt
{[}102{]} va H.Alexandersson {[}29{]} testlari qoʻllaniladi. Iqlim va
gidrologiyada bu testlar odatiy, meliorativ monitoringda esa kam
ishlatiladi.

Meliorativ monitoringda bu masala ayniqsa muhim, chunki kuzatuv
quduqlari vaqt oʻtishi bilan oʻzgaradi: eskilari ishdan chiqadi,
yangilari quriladi, baʼzilari chuqurlashtiriladi. Bunda barcha quduqlar
boʻyicha oʻrtacha sath faqat sizot suvlari emas, quduqlar tarkibi
oʻzgargani uchun ham oʻzgaradi. Mann--Kendall mezoni (1.3-§) esa buni
sathning oʻzgarishi deb koʻrsatadi.

\textbf{Kesilgan maʼlumotlar muammosi.} Kuzatuv qudugʻi maʼlum
chuqurlikka ega. Sizot suvi quduq tubidan pastga tushsa, quduq quriydi
va haqiqiy sathni oʻlchab boʻlmaydi: maʼlum boʻladigan narsa faqat
sathning quduq chuqurligidan pastda ekanidir. Statistikada bunday
qiymatlar kesilgan maʼlumotlar deyiladi; ularni oddiy son sifatida
oʻrtachaga qoʻshish natijani bir tomonga ogʻdiradi {[}65{]}. Agar quruq
quduq uchun quduq chuqurligi yoki shartli qiymat yozilsa, oʻrtacha sath
haqiqatdagidan sayozroq chiqadi; agar bunday quduqlar tahlildan butunlay
chiqarilsa, chuqur sathli hududlar eʼtibordan chetda qoladi. Ikkala
holda ham natija notoʻgʻri boʻladi. Quduqlar sayoz boʻlsa, bu muammo
sath chuqurlashayotgan davrda ayniqsa kuchayadi.

\textbf{Mahalliy amaliyot.} Oʻzbekistonda meliorativ monitoringni
oʻtkazish, kuzatuv quduqlarini jihozlash va hisobotlarni tuzish tartibi
R.Q.Ikromov va hammualliflari qoʻllanmasida bayon etilgan {[}16{]};
monitoring natijalari meliorativ kadastrga kiritiladi. Biroq koʻp yillik
maʼlumotlarni bazaga kiritishdan oldin tekshirish qoidalari va
ishonchlilikning son koʻrsatkichi belgilanmagan. Kuzatuv tarmogʻidagi
oʻzgarishlarni --- yangi quduq qurish, chuqurlashtirish, koʻchirish ---
sanasi bilan qayd etish ham monitoring qatorlari bilan birga
yuritilmaydi.

\textbf{Aniqlangan boʻshliqlar.} 1.1--1.4-§ dagi tahlil quyidagi toʻrtta
hal etilmagan masalani koʻrsatadi (1.2-jadval).

\oakjadvalnum{1.2-jadval}

\oakilovanom{\textbf{Adabiyotlar tahlili asosida aniqlangan hal etilmagan masalalar}}

{\fontsize{12}{14}\selectfont\setlength{\tabcolsep}{4pt}

\noindent\begin{longtable}[]{|
  >{\raggedright\arraybackslash}p{(\columnwidth - 8\tabcolsep - 5\arrayrulewidth) * \real{0.0423}}|>{\raggedright\arraybackslash}p{(\columnwidth - 8\tabcolsep - 5\arrayrulewidth) * \real{0.4138}}|>{\raggedright\arraybackslash}p{(\columnwidth - 8\tabcolsep - 5\arrayrulewidth) * \real{0.2394}}|>{\raggedright\arraybackslash}p{(\columnwidth - 8\tabcolsep - 5\arrayrulewidth) * \real{0.3046}}|}
\hline
\begin{minipage}[b]{\linewidth}\raggedright
\textbf{№}
\end{minipage} & \begin{minipage}[b]{\linewidth}\raggedright
\textbf{Hal etilmagan masala}
\end{minipage} & \begin{minipage}[b]{\linewidth}\raggedright
\textbf{Qayerda koʻrsatilgan}
\end{minipage} & \begin{minipage}[b]{\linewidth}\raggedright
\textbf{Dissertatsiyada hal etiladi}
\end{minipage} \\
\hline
\endhead
\hline
\endlastfoot
1 & Koʻp yillik monitoring maʼlumotlarining ishonchliligini baholash
qoidalari va son koʻrsatkichi ishlab chiqilmagan & 1.2, 1.4-§ & 2.4,
3.1-§ \\ \hline
2 & Drenaj tarmogʻining raqamli xaritasi boʻlmagan hududlar uchun GAT
bazasini mavjud xaritalardan tiklash usuli va uning aniqligi
baholanmagan & 1.3-§ & 2.3-§ \\ \hline
3 & Sizot suvi sathining drenaj tarmogʻiga bogʻliqligi har bir quduq
boʻyicha, tarmoq joylashuvi bilan birga oʻrganilmagan & 1.3-§ & 3.3,
3.4-§ \\ \hline
4 & Kollektor taʼmirining samarasi monitoring maʼlumotlari bilan
oʻlchanmagan & 1.2-§ & 4.3, 4.5-§ \\
\end{longtable}

}

\textbf{Tadqiqotning maqsadi} GAT texnologiyalari asosida meliorativ
monitoring maʼlumotlarining ishonchliligini baholash, sizot suvlari
sathining hududlar va yillar boʻyicha oʻzgarishini va uning drenaj
tarmogʻiga bogʻliqligini aniqlash, kollektor-drenaj tizimining
ishonchliligi va taʼmir samarasini baholash hamda shu asosda meliorativ
tizimlar samaradorligi va ishonchliligini taʼminlash boʻyicha tavsiyalar
ishlab chiqishdan iborat.

\textbf{Tadqiqotning vazifalari:}

meliorativ tizimlar ishonchliligi va samaradorligini baholashning mavjud
usullarini tahlil qilish va hal etilmagan masalalarni aniqlash;

mavjud xaritalar va meliorativ kadastr asosida Marhamat tumani
meliorativ tizimining GAT bazasini yaratish va uning aniqligini
baholash;

monitoring maʼlumotlarini tekshirish tartibini ishlab chiqish,
2009--2024-yillar (2025-yil bahori bilan) kuzatuv maʼlumotlarining
ishonchliligini baholash va natijalarni Andijon viloyatining boshqa
tumanlari bilan solishtirish;

tekshirilgan maʼlumotlar asosida sizot suvlari sathining yillar,
mavsumlar va hududlar boʻyicha oʻzgarishini massivlar va quduqlar
kesimida aniqlash, sayoz sath va shoʻrlanish xavfi yuqori hududlarni
ajratish;

sizot suvi sathining drenaj tarmogʻi, sugʻorish kanallari va relyefga
bogʻliqligini har bir quduq boʻyicha baholash;

tizim elementlarining ishonchlilik koʻrsatkichlarini aniqlash, talab
etiladigan drenaj modulini suv balansi boʻyicha hisoblash va kadastrdagi
koʻrsatkichlar bilan solishtirish;

kollektor taʼmirining sizot sathiga taʼsirini, taʼmirning haqiqiy
narxini va xarajatlarning massivlar boʻyicha taqsimotini baholash;

monitoring tizimini takomillashtirish, taʼmir samarasini oʻlchash va
mablagʻlarni birinchi navbatdagi massivlarga yoʻnaltirish boʻyicha
tavsiyalar ishlab chiqish.

\textbf{Tadqiqot bosqichlari.} Vazifalar toʻrt bosqichda bajariladi va
har bir bosqich oldingisining natijasiga tayanadi.

\begin{enumerate}
\def\labelenumi{\arabic{enumi}.}
\tightlist
\item
  Maʼlumotlarni yigʻish va GAT bazasini yaratish: monitoring jadvallari,
  kadastr, xaritalar yagona koordinatalar tizimida birlashtiriladi (II
  bob).
\item
  Maʼlumotlarni tekshirish: monitoring maʼlumotlari belgilangan qoidalar
  boʻyicha tekshiriladi, ishonchsiz qiymatlar ajratiladi, ishonchlilik
  koʻrsatkichi hisoblanadi (II--III boblar).
\item
  Sizot suvlarini tahlil qilish: tekshirilgan maʼlumotlar asosida
  sathning yillar, mavsumlar va hududlar boʻyicha oʻzgarishi va drenaj
  tarmogʻiga bogʻliqligi aniqlanadi (III bob).
\item
  Tizim ishonchliligi va taʼmir samaradorligini baholash: kadastr
  boʻyicha ishonchlilik koʻrsatkichlari, taʼmirning sizot sathiga
  taʼsiri va narxi baholanadi, tavsiyalar ishlab chiqiladi (IV bob).
\end{enumerate}

Bu tartib tadqiqotning asosiy qoidasini aks ettiradi: tizim
ishonchliligini baholashdan oldin maʼlumotlarning ishonchliligini
baholash kerak.

\oakpar{I bob boʻyicha xulosalar}

\begin{enumerate}
\def\labelenumi{\arabic{enumi}.}
\tightlist
\item
  Sugʻoriladigan yerlarda ikkilamchi shoʻrlanish va zax bosishga qarshi
  asosiy muhandislik chorasi kollektor-drenaj tizimlaridir.
  Oʻzbekistonda 142,9 ming km kollektor-drenaj tarmogʻi ishlatiladi;
  Andijon viloyatida ochiq tarmoq 7 616,3 km, shundan Marhamat tumanida
  621,1 km. Tarmoqning asosiy qismi oʻtgan asrning oʻrtalarida qurilgan
  va muntazam taʼmir talab qiladi.
\item
  Kollektor-drenaj tizimlarining muammolari toʻrt guruhga ajratildi:
  jismoniy eskirish, masʼuliyatning boʻlinishi, axborot taʼminotining
  yetishmasligi va natijani oʻlchashning yoʻqligi. Oxirgi ikki guruh
  ilmiy masala boʻlib, mazkur tadqiqot shularga bagʻishlangan.
\item
  Meliorativ tizimni baholashning uch yondashuvi ---
  loyihaviy-gidravlik, rejim koʻrsatkichlari boʻyicha va ishonchlilik
  nazariyasi --- tahlil qilindi. Loyihaviy yondashuv haqiqiy holatni,
  holat boʻyicha baholash sababni koʻrsatmaydi; ishonchlilik nazariyasi
  uchun zarur buzilishlar hisobi esa amalda yuritilmaydi.
\item
  Drenaj tizimlari ishini baholash usullari sugʻorish tizimlarinikiga
  qaraganda ancha kam ishlangan; xalqaro adabiyotda ham bu sohada yaxlit
  usul yoʻqligi qayd etilgan.
\item
  Taʼmir samarasi amalda bajarilgan ish hajmi bilan oʻlchanadi. Uni
  sizot rejimidagi oʻzgarish orqali oʻlchash uchun tadbir oʻtkazilgan va
  oʻtkazilmagan obyektlarni solishtirish usuli mavjud, ammo u drenaj
  tizimlariga kuzatuv quduqlari maʼlumotlari asosida qoʻllanilmagan.
\item
  Meliorativ monitoringda GAT va geostatistika asosan sizot suvlari va
  shoʻrlanishni xaritalash uchun qoʻllanilgan. Drenaj tarmogʻining oʻzi
  tahlilga kam kiritiladi; Markaziy Osiyodagi tadqiqotlar asosan Xorazm
  viloyatiga tegishli, Fargʻona vodiysi boʻyicha bunday ishlar kam.
\item
  Koʻrib chiqilgan barcha yondashuvlar dastlabki maʼlumotlarni toʻgʻri
  deb oladi. Monitoring maʼlumotlarining bir jinsliligi, quduqlar
  tarkibining oʻzgarishi va quruq quduqlar muammosi meliorativ
  adabiyotda deyarli yoritilmagan.
\item
  Toʻrtta hal etilmagan masala aniqlanib, tadqiqotning maqsadi, sakkizta
  vazifasi va toʻrt bosqichli tartibi belgilandi. Tadqiqotning asosiy
  qoidasi: tizim ishonchliligini baholashdan oldin maʼlumotlarning
  ishonchliligini baholash kerak.
\end{enumerate}

\oakbob{II BOB. TADQIQOT OBYEKTI, MAʼLUMOTLAR VA USULLAR}

Ushbu bobda tadqiqot obyekti, maʼlumot manbalari, GAT bazasini yaratish
tartibi hamda keyingi boblarda qoʻllanilgan tekshiruv va tahlil usullari
bayon etiladi.

\oakpar{2.1-§. Marhamat tumanining tabiiy-meliorativ sharoiti}

\textbf{Geografik oʻrni va maʼmuriy tuzilishi.} Marhamat tumani Andijon
viloyatining janubiy qismida, Fargʻona vodiysining sharqiy chekkasida
joylashgan. Tumanning maydoni 321 km², aholisi 2025-yil holatiga 197,8
ming kishi. Tuman tarkibida Marhamat shahri, oʻnta shaharcha va beshta
qishloq fuqarolar yigʻini bor {[}135{]}. Andijon viloyati sharq va
janubdan Qirgʻiziston Respublikasining Jalolobod va Oʻsh viloyatlari
bilan chegaradosh {[}135{]}. Tadqiqotda qamrab olingan kuzatuv quduqlari
72°12′ dan 72°24′ gacha sharqiy uzunlik va 40°27′ dan 40°37′ gacha
shimoliy kenglik oraligʻida joylashgan. Tumanning viloyatdagi oʻrni,
relyefi va meliorativ tizimi 2.1-rasmda koʻrsatilgan. Xarita tiklangan
GAT bazasi (2.3-§), Copernicus GLO-30 raqamli relyef modeli {[}138{]} va
geoBoundaries maʼmuriy chegaralari {[}111{]} asosida tuzildi.

Sugʻoriladigan yerlar maʼmuriy-xoʻjalik jihatdan sakkizta massivga
boʻlingan: Z.M.Bobur, Mashʼal, Istiqlol, Paxtakor, M.Toʻjiboyev,
T.Mirzayev, U.Qoraboyev va Y.Oxunboboyev. Meliorativ monitoring, kadastr
va ekin joylashtirish maʼlumotlari shu massivlar kesimida yuritiladi.
Massivlarning asosiy tavsifi 3-ilovada, ularning xarita boʻyicha maydoni
va chegaralari 2.3-§ da keltirilgan.

\oakfig{figures/map1_location.png}{2.1-rasm. Tadqiqot hududining joylashuvi: a --- Markaziy Osiyo; b ---
Andijon viloyati tumanlari; c --- Marhamat tumani relyefi, massivlari,
kollektor-drenaj tarmogʻi, sugʻorish kanallari va kuzatuv quduqlari}

\textbf{Iqlim.} Tuman iqlimi keskin kontinental, qurgʻoqchil: qishi
nisbatan sovuq, yozi issiq va quruq. Tumanda meteostansiya yoʻq, shuning
uchun iqlim eng yaqin Andijon meteostansiyasining 1961--1990-yillardagi
oʻrtacha koʻrsatkichlari boʻyicha tavsiflandi {[}136{]}. Marhamat
Andijondan taxminan 30 km janubda va balandroqda joylashgani uchun
harorat biroz farq qilishi mumkin.

Meteostansiya maʼlumotlarida bugʻlanuvchanlik yoʻq. Shuning uchun etalon
bugʻlanuvchanlik (ET₀), yaʼni yaxshi sugʻorilgan oʻtloqdan bugʻlanadigan
suv miqdori, FAO tavsiya etgan Hargreaves--Samani usulida harorat
boʻyicha hisoblandi {[}31{]}:

\begin{equation}\tag{2.1}
ET_0 = 0{,}0023 \, (T_{\mathrm{oʻr}} + 17{,}8) \, \sqrt{T_{\mathrm{max}} - T_{\mathrm{min}}} \; R_a
\end{equation}

bu yerda \(T_{\mathrm{oʻr}}\), \(T_{\mathrm{max}}\),
\(T_{\mathrm{min}}\) --- oylik oʻrtacha, oʻrtacha maksimal va oʻrtacha
minimal harorat, °C; \(R_{a}\) --- atmosfera chegarasidagi quyosh
radiatsiyasi, mm/sut hisobida; u joyning kengligi (40,5°) va yil kuni
boʻyicha hisoblanadi {[}31{]}. Natijalar 4-ilovada keltirilgan.

Bu maʼlumotlardan tumanning suv rejimiga oid quyidagi xususiyatlar kelib
chiqadi.

Yillik yogʻin (323 mm) bugʻlanuvchanlikdan (1 185 mm) 3,7 baravar kam;
bu qurgʻoqchil iqlimga xos. Yillik namlik tanqisligi 862 mm ni tashkil
etadi.

Yogʻin va bugʻlanuvchanlik mavsumlar boʻyicha qarama-qarshi
taqsimlangan. Yogʻinning 65 foizi noyabr--aprelda yogʻadi;
aprel--sentabrdagi vegetatsiya davrida esa yogʻin atigi 126 mm,
bugʻlanuvchanlik 949 mm. Iyun--avgustda yogʻin oyiga 6--17 mm dan
oshmaydi, bugʻlanuvchanlik esa oyiga 170--194 mm. Vegetatsiya davrida
ekinlarning suvga talabi deyarli butunlay sugʻorish hisobiga qoplanadi.

Faqat dekabr--fevralda yogʻin bugʻlanuvchanlikdan biroz koʻp. Bu davrda
tuproq tabiiy ravishda yuviladi, lekin bu juda kam (yiliga taxminan 24
mm). Yil davomida tuproqda tuz toʻplanishga moyil; uni faqat shoʻr
yuvish va drenaj orqali chiqarish mumkin {[}19, 123{]}.

Iqlim koʻrsatkichlarining yillik yoʻli 5-ilovada koʻrsatilgan.

Turli manbalarda Andijon iqlimi koʻrsatkichlari hisoblash davri va
usuliga qarab farq qiladi: yillik yogʻin 232 mm dan 384 mm gacha,
oʻrtacha yillik harorat 13,6--14,9 °C oraligʻida keltiriladi {[}136{]}.
Andijon meteostansiyasining 2015--2023-yillardagi oʻrtacha yillik
yogʻini 196 mm {[}33{]}. Yuqorida bitta manbaning bir davr uchun
maʼlumotlari ishlatildi. Hargreaves--Samani usuli Penman--Monteith
usuliga qaraganda taxminiy; aniq hisob uchun namlik, shamol va quyosh
radiatsiyasi maʼlumotlari kerak {[}31{]}.

\textbf{Relyef.} Tuman hududi janubdan shimolga qiyalangan (2.1-rasm,
c). Copernicus GLO-30 relyef modeli (30 m) boʻyicha tuman hududining
balandligi 504 m dan 1 276 m gacha, oʻrtacha (mediana) 607 m; maydonning
43,5 foizi 600 m dan past, 22 foizi 700 m dan baland joylashgan
{[}138{]}. Eng baland qism janubdagi Y.Oxunboboyev massivi (median
balandlik 810 m) boʻlib, u togʻ oldi tekisligiga tegishli; shimoldagi
Z.M.Bobur massivining median balandligi 559 m. Paxtakor, U.Qoraboyev,
Mashʼal va Z.M.Bobur massivlari tekis: median balandlik 559--598 m,
nishablik medianasi 1,3--1,6 foiz. Istiqlol massivining gʻarbiy qismi
qiya (nishablik medianasi 5,4 foiz), Y.Oxunboboyev massivida esa
nishablik medianasi 15,6 foizga yetadi. Kuzatuv quduqlarining
pasportdagi balandligi 528 m dan 839 m gacha, medianasi 597 m (3-ilova);
ular relyef modeli bilan yaxshi mos keladi: 172 ta quduq boʻyicha farq
medianasi −1,1 m.

Relyef sizot suvlarining oqim yoʻnalishini va kollektorlar qiyaligini
belgilaydi: kollektorlar janubdan shimolga, past joylar tomon yoʻnalgan.
Biroq massivlar orasidagi balandlik farqi sizot suvi sathidagi farqni
toʻliq tushuntirmaydi (3.4-§).

Massivlar median balandlik oʻsishi tartibida joylashtirilgan. Bonitet
balli 2021-yilgi ekin joylashtirish jadvalidagi 1 607 ta kontur
boʻyicha, maydon hisobga olingan oʻrtacha sifatida hisoblandi; ichki
zovurlar uzunligi kadastr boʻyicha keltirilgan. Kadastrda ichki
zovurlarning bir qismi massivga emas, «Markaziy zovur tizimi»ga (52,8
km) biriktirilgan, shuning uchun M.Toʻjiboyev massivi boʻyicha ichki
zovurlar uzunligi alohida koʻrsatilmagan.

\textbf{Gidrografiya va sugʻorish manbalari.} Tuman hududidan Aravonsoy
daryosi hamda Shahrixonsoy, Karkidon va Janubiy Fargʻona kanallari
oʻtadi {[}135{]}; tumanning sharqiy qismida Qoradaryo joylashgan.
Meliorativ ekspeditsiya xaritasidan raqamlangan sugʻorish kanallari
uzunligi 170,9 km (2.3-§). Sugʻorish asosan egatlab oʻtkaziladi; bunda
suvning bir qismi pastga sizib, sizot suvlariga qoʻshiladi {[}10, 17{]}.

\textbf{Gidrogeologik sharoit.} Tuman Fargʻona togʻlararo botigʻining
sharqiy qismida joylashgan. Sizot suvlari daryo va sel olib kelgan
toʻrtlamchi davr yotqiziqlarida toʻplangan. Ular asosan sugʻorish suvi
va kanallardan sizib chiqqan suv hisobiga toʻladi, bugʻlanish va drenaj
orqali sarflanadi.

Sizot suvlarining sugʻorishga bogʻliqligi monitoring maʼlumotlarida
yaqqol koʻrinadi. 2009--2024-yillardagi barcha kuzatuvlar boʻyicha sath
fevralda eng chuqur (2,48 m), iyul--avgustda eng sayoz (2,05--2,06 m)
boʻladi. Sath yogʻin eng koʻp boʻlgan mart--aprelda emas, aksincha
yogʻin deyarli boʻlmagan, ammo sugʻorish eng jadal kechadigan yoz
oylarida koʻtariladi. Bu sizot suvlarining asosiy toʻyinish manbai
yogʻin emas, sugʻorish ekanini koʻrsatadi. Mavsumiy oʻzgarish 3.2-§ da
batafsil tahlil qilinadi.

Meliorativ ekspeditsiyaning 2024-yil 1-aprel holatiga tuzilgan xaritasi
boʻyicha sizot suvlari sathining kuzatuv quduqlari boʻyicha taqsimoti
6-ilovada keltirilgan.

Aprel boshida, yaʼni sugʻorish mavsumi oldidan, quduqlarning 19,2
foizida sath 1,5 m dan sayoz boʻlgan; medianasi 2,34 m. Bu koʻrsatkich
yillik eng sayoz holat emas: koʻp yillik maʼlumotlarga koʻra
iyul--avgustda sath aprel boshiga nisbatan oʻrtacha yana 0,2 m
koʻtariladi (3.2-§).

Sizot suvlarining minerallashuvi 2004--2010-yillarda tuman boʻyicha
oʻrtacha (mediana) 1,02--1,45 g/l boʻlgan va massivlar orasida sezilarli
farq qiladi: quduqlar boʻyicha medianasi Mashʼal massivida 2,47 g/l,
U.Qoraboyev va T.Mirzayev massivlarida 0,80--0,84 g/l (3.5-§).

\textbf{Tuproqlar.} Andijon viloyatining sugʻoriladigan tuproqlari
asosan och tusli boʻz tuproqlar va oʻtloqi-boʻz tuproqlardan iborat.
Tuproq sifatining umumlashgan koʻrsatkichi --- bonitet balli tuman
boʻyicha oʻrtacha 59,7 ni tashkil etadi; konturlarning 10 foizida ball
40 dan past, 16 foizida 70 dan yuqori. Eng past ball Y.Oxunboboyev
massivida (41,3), eng yuqori ball Z.M.Bobur massivida (67,9) qayd
etilgan (3-ilova).

\textbf{Qishloq xoʻjaligidan foydalanish.} 2021-yilgi ekin joylashtirish
jadvaliga koʻra massivlar konturlari boʻyicha jami 9 701 ga maydonning 3
580 ga qismiga gʻoʻza, 4 073 ga qismiga boshoqli don ekilgan; qolgan
maydonlar boshqa ekinlar, bogʻ va tomorqalarga toʻgʻri keladi. Gʻoʻza va
gʻalla almashlab ekishi sugʻorish rejimining yillik yoʻlini belgilaydi:
kuzgi gʻalla bahorda, gʻoʻza esa yoz oylarida eng koʻp suv talab qiladi.

\textbf{Xulosa.} Tumanning tabiiy sharoiti drenajni doim zarur qiladi:
qurgʻoqchil iqlimda ekinlar suvni sugʻorishdan oladi, sugʻorish esa
sizot suvlarini toʻldirib, sathni yozda koʻtaradi. Sizot suvlari tabiiy
yoʻl bilan kam chiqib ketadi, shuning uchun tuzni boshqarishning asosiy
vositasi kollektor-drenaj tarmogʻidir. Tarmoqning tarkibi va holati
2.2-§ da tavsiflanadi.

\oakpar{2.2-§. Meliorativ tizimning tarkibi: kollektorlar, zovurlar, kuzatuv
quduqlari}

Meliorativ tizimning tarkibi Andijon viloyati meliorativ
ekspeditsiyasining 2021-yil 1-yanvar holatidagi kadastri boʻyicha
tavsiflanadi. Kadastr yigʻma jadvallar (tarmoq turlari, inshootlar) va
obyektlar roʻyxatlaridan (har bir kollektor, ichki zovur, yopiq drenaj
va kuzatuv qudugʻi) iborat.

\textbf{Kollektor-drenaj tarmogʻining umumiy tarkibi.} Kadastrning
yigʻma jadvaliga koʻra tumanda 621,1 km ochiq va 91,7 km yopiq-yotiq
kollektor-drenaj tarmogʻi mavjud (2.1-jadval). Tarmoq toʻrt bosqichdan
iborat: magistral, tumanlararo, xoʻjaliklararo kollektorlar va xoʻjalik
ichki zovurlari. Ochiq tarmoqning 57 foizi va yopiq drenajning 61 foizi
suv isteʼmolchilari uyushmalari hisobida, qolgani davlat hisobida.
Viloyat GAT bazasida tumanning kollektor-drenaj tarmogʻi 119 ta zovur,
240,3 km deb berilgan {[}33{]}; manbalar tarmoqning turli qismini
hisobga oladi.

\oakjadvalnum{2.1-jadval}

\oakilovanom{\textbf{Marhamat tumani kollektor-drenaj tarmogʻining tarkibi (kadastr,
2021-yil 1-yanvar)}}

{\fontsize{12}{14}\selectfont\setlength{\tabcolsep}{4pt}

\noindent\begin{longtable}[]{|
  >{\raggedright\arraybackslash}p{(\columnwidth - 10\tabcolsep - 6\arrayrulewidth) * \real{0.2850}}|>{\raggedright\arraybackslash}p{(\columnwidth - 10\tabcolsep - 6\arrayrulewidth) * \real{0.0695}}|>{\raggedright\arraybackslash}p{(\columnwidth - 10\tabcolsep - 6\arrayrulewidth) * \real{0.1547}}|>{\raggedright\arraybackslash}p{(\columnwidth - 10\tabcolsep - 6\arrayrulewidth) * \real{0.1854}}|>{\raggedright\arraybackslash}p{(\columnwidth - 10\tabcolsep - 6\arrayrulewidth) * \real{0.3054}}|}
\hline
\begin{minipage}[b]{\linewidth}\raggedright
\textbf{Tarmoq turi}
\end{minipage} & \begin{minipage}[b]{\linewidth}\raggedright
\textbf{Soni}
\end{minipage} & \begin{minipage}[b]{\linewidth}\raggedright
\textbf{Uzunligi, km}
\end{minipage} & \begin{minipage}[b]{\linewidth}\raggedright
\textbf{Ochiq tarmoqdagi ulushi, \%}
\end{minipage} & \begin{minipage}[b]{\linewidth}\raggedright
\textbf{Hisobida}
\end{minipage} \\
\hline
\endhead
\hline
\endlastfoot
Magistral kollektorlar, jami & 15 & 41,8 & 6,7 & Davlat \\ \hline
shu jumladan davlatlararo & 14 & 36,3 & & Davlat \\ \hline
shu jumladan viloyatlararo & 1 & 5,5 & & Davlat \\ \hline
Tumanlararo kollektorlar & 3 & 21,5 & 3,5 & Davlat \\ \hline
Xoʻjaliklararo kollektorlar & 72 & 202,5 & 32,6 & Davlat \\ \hline
Xoʻjalik ichki zovurlari & 330 & 355,3 & 57,2 & Suv isteʼmolchilari
uyushmalari \\ \hline
Ochiq tarmoq, jami & 420 & 621,1 & 100,0 & \\ \hline
Yopiq-yotiq drenaj & 7 & 36,0 & & Davlat \\ \hline
Yopiq-yotiq drenaj & 7 & 55,7 & & Suv isteʼmolchilari uyushmalari \\
\end{longtable}

}

Tumanning oʻziga xos xususiyati --- davlatlararo kollektorlarning
koʻpligi. Kadastrga koʻra 14 ta magistral kollektor Oʻzbekiston va
Qirgʻiziston hududidan oʻtadi; ularning umumiy uzunligi 36,3 km, shundan
Oʻzbekiston hududida atigi 6,7 km. Masalan, Aravon kollektorining 8,4 km
uzunligidan 1,1 km, Yuqori Rovot kollektorining 3,2 km dan 0,9 km,
Beshtentak kollektorining 5,2 km dan 4,7 km qismi Oʻzbekiston hududida
joylashgan. Bu holat tarmoqni boshqarishda muhim: kollektorning yuqori
qismidagi holat va unga tushadigan suv miqdori qoʻshni davlat hududida
shakllanadi. Polvontosh kollektori (5,5 km) viloyatlararo kollektor
sifatida qayd etilgan.

\textbf{Davlat hisobidagi kollektorlar.} Kadastrda davlat hisobidagi 90
ta kollektorning har biri uchun nomi, uzunligi, oxiridagi suv sarfi,
xizmat koʻrsatish maydoni, foydalanishga topshirilgan yili, oxirgi
taʼmir yili va taʼmirlangan uchastka (piketlar) koʻrsatilgan; 53 tasida
taʼmir xarajati ham bor. Ularning umumiy uzunligi 265,8 km, xizmat
koʻrsatish maydoni 11 079 ga; barchasi III kapitallik sinfiga kiradi va
meliorativ ekspeditsiya tomonidan ishlatiladi.

Kollektorlar oʻlchami boʻyicha keskin farq qiladi (7-ilova). Ularning
yarmi 2,1 km dan qisqa; 10 km dan uzunlari faqat bitta --- Gʻarbiy
Qoraqoʻrgʻon (12,5 km). Oxiridagi suv sarfi boʻyicha ham farq katta: 31
ta kollektorda sarf 0,01 m³/s gacha, faqat 5 tasida 1 m³/s dan ortiq.
Barcha kollektorlar oxiridagi sarflar yigʻindisi 18,3 m³/s.

Eng uzun oʻnta kollektor tarmoqning asosiy boʻgʻinini tashkil etadi
(8-ilova). Ular asosan 1940--1950-yillarda qurilgan va
2012--2020-yillarda taʼmirlangan.

Kollektorlarning yoshi va taʼmirlari 4.1-§ da, suv sarfi va xizmat
maydoni 4.2-§ da tahlil qilinadi.

\textbf{Xoʻjalik ichki zovurlari.} Kadastrning ichki zovurlar roʻyxatida
tuman boʻyicha 332 ta yozuv bor, ularning umumiy uzunligi 370,7 km;
yigʻma jadvalda esa 330 ta zovur va 355,3 km koʻrsatilgan. 15,4 km lik
farq roʻyxat va yigʻma jadval turli sanalarda yangilangani bilan bogʻliq
boʻlishi mumkin. Roʻyxatda zovurlar massivlar va «Markaziy zovur tizimi»
boʻyicha guruhlangan; har bir zovur uchun uzunligi, oxiridagi sarf,
xizmat maydoni, foydalanishga topshirilgan yili va taʼmir yili berilgan.
Ichki zovurlar sizot suvlarini daladan toʻgʻridan-toʻgʻri oladigan quyi
boʻgʻin; ularning holati 4.1-§ da tahlil qilinadi.

\textbf{Yopiq-yotiq drenaj.} Kadastrga koʻra tumanda 14 ta yopiq-yotiq
drenaj tizimi bor, umumiy uzunligi 91,7 km. Roʻyxatda M.Toʻjiboyev
massividagi 6,82 km uzunlikdagi yopiq drenaj alohida qayd etilgan; u
2016-yilda foydalanishga topshirilgan va 60 ga maydonga xizmat
koʻrsatadi. Yopiq drenaj yer ostida boʻlgani uchun xaritada
koʻrsatilmagan va xarita boʻyicha tahlilga kiritilmagan.

\textbf{Gidrotexnik inshootlar.} Kollektor-drenaj tarmogʻida jami 455 ta
inshoot qayd etilgan (9-ilova). Ularning yarmidan koʻpi (238 ta) davlat
hisobidagi nazorat tarmogʻi inshootlaridir. Inshootlar tarmoq
ishonchliligi uchun muhim: koʻprik va quvurli oʻtish joylari ostida
loyqa tez toʻplanadi va ular kollektorda «tor joy» boʻlib qolishi mumkin
{[}108, 122{]}.

Suv oʻlchash inshootlari atigi oltita, yaʼni kollektorlardagi oqim
muntazam oʻlchanmaydi. Kadastrdagi sarflar loyiha yoki bir martalik
oʻlchov qiymatlari boʻlishi mumkin. Shuning uchun tarmoq ishini amalda
faqat sizot suvi sathi orqali baholash mumkin.

\textbf{Kuzatuv quduqlari.} Kadastrda tuman boʻyicha 188 ta kuzatuv
qudugʻining pasporti bor: massivi, qurilgan yili, chuqurligi, diametri,
materiali, biriktirilgan maydoni va mutlaq balandligi. Barcha
quduqlarning diametri 100 mm. Bitta quduqqa oʻrtacha 85 ga maydon
biriktirilgan (48 dan 133 ga gacha); jami biriktirilgan maydon 15 985
ga. 2019-yildan boshlab bu jadvallarda har bir massiv bloki oxirida «27
та» kabi jami qatori (blokdagi quduqlar soni va ularning oʻrtacha
qiymatlari), tuman bloki oxirida esa tuman boʻyicha jami qatori («188»)
turadi; ular quduq sifatida hisobga olinmadi. Keyingi boblarda har bir
hisob qaysi manbaga tayanishi koʻrsatiladi.

Quduqlar tarkibining ikki xususiyati keyingi tahlil uchun muhim.

Kuzatuv quduqlari asosan yangi. 188 ta quduqdan 120 tasi (64 foiz)
2013--2016-yillarda polietilen quvurdan qurilgan; 65 tasi
1971--1990-yillarda, 3 tasi 1991--2012-yillarda metall quvurdan qurilgan
quduqlardir (10-ilova). Bu 3.1-§ da aniqlangan quduqlar tarkibi
oʻzgarishini tushuntiradi.

Quduqlar sayoz. Chuqurlik medianasi 3,2 m; 142 ta quduq (76 foiz)
2,5--3,5 m chuqurlikda, faqat 26 tasi (14 foiz) 4 m dan chuqur. Sizot
suvi quduq tubidan pastga tushganda quduq quriydi va haqiqiy sathni
oʻlchab boʻlmaydi (1.4-§). 2024-yil 1-aprelda quduqlarning 28 foizida
sath 3 m dan chuqur boʻlgan (6-ilova), yaʼni ularning bir qismi quriy
boshlagan. Demak, sath chuqur joylarda monitoring sathni haqiqiydan
sayozroq koʻrsatishi mumkin (3.1-§).

\textbf{Manbalar orasidagi farqlar.} Quduqlar soni manbalar boʻyicha
farq qiladi: pasportda 188 ta, GAT eksportidagi koordinatalar jadvalida
205 ta, 2024-yilgi xaritada 202 ta, viloyat GAT bazasida 213 ta
{[}33{]}, 2009--2024-yillar kuzatuv jadvallarida esa 188 ta quduq bor
(massivlar va tuman boʻyicha jami qatorlari chiqarilgandan keyin).
Farqlarning bir qismi nomlanishdagi tafovutlar bilan (masalan, «104А» va
«104а»), bir qismi pasporti rasmiylashtirilmagan qoʻshimcha quduqlar
bilan izohlanadi. Keyingi boblarda har bir hisob qaysi manbaga tayanishi
koʻrsatiladi.

Kadastr meliorativ tizim haqida batafsil maʼlumot beradi, lekin jadval
shaklida, joylashuvsiz. Kollektorlarni sizot suvlari bilan bogʻlash
uchun ularni GAT bazasiga oʻtkazish kerak (2.3-§).

\oakpar{2.3-§. GAT bazasini yaratish: xaritalarni koordinataga bogʻlash va
tarmoqni raqamlash}

\textbf{Vazifa.} Sizot suvlarini drenaj tarmogʻi bilan bogʻlash uchun
har bir quduq va kollektorning joyi maʼlum boʻlishi kerak. Kadastr
yozuvlarida joylashuv yoʻq (2.2-§), tumanning drenaj tarmogʻi raqamli
xaritasi ham yoʻq edi. Shuning uchun birinchi vazifa mavjud manbalardan
GAT bazasini tiklash boʻldi.

\textbf{Boshlangʻich manbalar.} Bazani yaratish uchun toʻrt turdagi
manba ishlatildi (11-ilova).

Asosiy qiyinchilik shundaki, faqat bitta manbada --- quduqlar
koordinatalari jadvalida --- haqiqiy koordinatalar bor. Xaritalar esa
koordinatasiz rasm. Usulning gʻoyasi: quduqlar ham jadvalda
(koordinatasi bilan), ham xaritada (raqami bilan) bor; demak ularni
tayanch nuqta qilib, xaritani koordinataga bogʻlash mumkin.

\textbf{Bosqichlar.} Baza oltita bosqichda yaratildi (12-ilova). Barcha
bosqichlar dastur orqali bajarildi va takrorlanadi.

\textbf{Tasvir va yozuvlarni ajratish.} Xarita PDF faylida ikki
qatlamdan iborat: rangli tasvir (fon, chiziqlar, belgilar) va matn
(yozuvlar). Matn qatlami alohida oʻqildi: har bir yozuvning matni,
rangi, oʻlchami va joyi saqlandi. Tasvir esa yozuvlarsiz rasmga
aylantirildi. Natijada quduq raqamlari (qora yozuvlar) va oʻlchangan
qiymatlar (qizil yozuvlar) alohida olindi, chiziqlar esa yozuvlar bilan
toʻsilmadi.

\textbf{Quduq belgilarini aniqlash.} Xaritada kuzatuv quduqlari bir xil
havo rangli doira bilan tasvirlangan. Shu rangdagi piksellar ajratilib,
yonma-yon turganlari guruhlarga birlashtirildi; oʻlchami 150--2 500
piksel boʻlgan guruhlar quduq belgisi deb olindi va har birining markazi
topildi. Har bir belgiga 90 piksel (taxminan 380 m) radiusdagi eng yaqin
qora yozuv quduq raqami sifatida, eng yaqin qizil yozuv esa sath yoki
minerallashuv qiymati sifatida biriktirildi. Ikkala xaritada 202 tadan
quduq belgisi aniqlandi; sath xaritasida 184 ta (tekshiruvdan keyin 182
ta), minerallashuv xaritasida 128 ta qiymat juftlandi.

Juftlashda bir nechta xato aniqlandi va tuzatildi. Ikki quduqqa (137 va
141) yonma-yon turgan ikki qiymat birlashib, fizik jihatdan imkonsiz 3
402,04 qiymati biriktirilgan; bu qiymatlar tahlildan chiqarildi. Toʻrtta
belgiga quduq raqami oʻrniga yaqin turgan kanal yoki kollektor nomi
biriktirilgan; ular raqamsiz quduq sifatida saqlandi. Bunday xatolar
avtomatik juftlashda muqarrar; ularni qiymat chegaralari va nom shakli
boʻyicha tekshirib topish mumkin.

\textbf{Xaritani koordinataga bogʻlash.} Raqami koordinatalar
jadvalidagi raqamga mos kelgan 190 ta quduq tayanch nuqta qilib olindi.
Xaritadagi piksel koordinatalari (x, y) va joydagi koordinatalar (E, N)
orasidagi bogʻlanish chiziqli (affin) almashtirish bilan ifodalandi:

\begin{equation}\tag{2.2}
E = a_1 x + a_2 y + a_3, \qquad N = b_1 x + b_2 y + b_3
\end{equation}

Oltita koeffitsiyent eng kichik kvadratlar usulida topildi.
Koordinatalar tizimi sifatida WGS 84 / UTM 43N olindi: tuman UTM ning
43-zonasiga tushadi va bu tizimda masofalar metrda oʻlchanadi. Birinchi
hisobdan keyin xatosi katta boʻlgan 9 ta nuqta chiqarildi va
koeffitsiyentlar qayta topildi. Chiqarilgan nuqtalarning barchasi nom
boʻyicha juftlash xatosi bilan bogʻliq: masalan, xaritada ikki marta
uchraydigan raqamlar yoki «ПЕЗ» qoʻshimchasi bilan farqlanadigan yaqin
quduqlar. Bogʻlash aniqligining koʻrsatkichlari 13-ilovada keltirilgan.

Xatolar taqsimoti 14-ilovada koʻrsatilgan. Nuqtalarning 91 foizida xato
5 m dan kichik, yaʼni bir-ikki piksel. Oʻrtacha kvadratik xato
medianadan ancha katta, chunki uni bir nechta katta xatoli nuqta
oshiradi. Piksel oʻlchami nazariy qiymatga deyarli toʻliq mos keldi:
demak, xarita masshtabi buzilmagan va tanlangan almashtirish yetarli.

\textbf{Tarmoqni raqamlash.} Xaritada kollektorlar yashil, kanallar koʻk
chiziq bilan tasvirlangan. Har bir rang alohida ajratildi, quduq
belgilari ostidagi uzilishlar toʻldirildi, mayda dogʻlar (300 pikseldan
kichik) olib tashlandi. Soʻng chiziqlar bir piksel qalinlikkacha
ingichkalashtirildi va tugunlar orasidagi boʻlaklar siniq chiziqqa
aylantirildi. 40 m dan qisqa boʻlaklar ortiqcha shoxchalar sifatida
chiqarildi. Natijada 332 ta kollektor boʻlagi (214,6 km) va 426 ta kanal
boʻlagi (170,9 km) olindi.

Xaritadan olingan kollektorlar uzunligi kadastrdagi davlat hisobidagi
kollektorlar uzunligidan (265,8 km) 19 foizga kam. Farqning sabablari:
xaritada barcha kollektorlar koʻrsatilmagan; kanal bilan yonma-yon yoki
ustma-ust tushgan uchastkalar bitta chiziq boʻlib koʻrinadi; 14 ta
davlatlararo kollektorning Qirgʻiziston hududidagi qismlari xaritada
yoʻq (2.2-§). Xoʻjalik ichki zovurlari xaritada umuman tasvirlanmagan.

\textbf{Kollektor nomlarini biriktirish.} Xaritada kollektor nomlari
chiziq yoniga yozilgan. Nomni kadastr bilan bogʻlashda xarita va kadastr
turli qisqartmalardan foydalanishi aniqlandi: xaritada oʻzbekcha,
kadastrda ruscha qisqartmalar ishlatilgan. Masalan, xaritadagi «АЗ»
kadastrdagi «АД» ga, «ҒЗ» --- «ЗД» ga, «МЗ» --- «ЦД» ga, «ШЗ» --- «ШД»
ga, «ҒКЗ» --- «ЗКД» ga mos keladi. Shu moslik jadvali asosida xaritadagi
67 ta nom yozuvi kadastrdagi 65 ta kollektor bilan bogʻlandi. Nom
yozuvdan 220 m gacha masofadagi eng yaqin chiziq boʻlagiga biriktirildi,
soʻng tugunlarda toʻgʻri davom etgan (yoʻnalishi 40° dan kam oʻzgargan)
qoʻshni boʻlaklarga oʻtkazildi.

Natijada 60 ta kollektorga chiziq geometriyasi biriktirildi; ularning
xaritadagi uzunligi 114,8 km. Xaritadagi va kadastrdagi uzunliklar
nisbatining medianasi 0,85; 27 ta kollektorda farq 25 foizdan oshmaydi,
19 tasida xaritadagi uzunlik kadastrdagidan ikki baravardan koʻproq
qisqa. Ikkinchi holat asosan uzun kollektorlarda uchraydi: nom
tarmoqlanish tugunida toʻxtaydi va kollektorning faqat bir qismi
qamraladi. Ikki yozuv («Mashʼal» va «Qirgʻiz») notoʻgʻri biriktirilgani
ehtimoli tufayli keyingi tahlildan chiqarildi.

\textbf{Massiv chegaralarini tiklash.} Xaritada massiv chegaralari qizil
chiziq bilan berilgan, lekin ular koʻp joyda kollektor, kanal va aholi
punktlari belgilari ostida qolib, uzuq-yuluq koʻrinadi. Shuning uchun
chegaralar ikki manbadan tiklandi. Qizil chiziq boʻlaklari toʻsiq qilib
olindi: ularni kesib oʻtish oddiy joyga qaraganda 400 baravar «qimmat»
hisoblandi. Har bir massiv quduqlaridan toʻsiqni aylanib oʻtuvchi eng
qisqa masofa hisoblandi va hududning har bir katagi eng yaqin massivga
biriktirildi. Qizil chiziq saqlangan joyda chegara aynan shu chiziqdan
oʻtadi; chiziq koʻrinmaydigan joyda u qoʻshni massivlar quduqlari
oʻrtasidan oʻtadi va taxminiydir.

Natija mustaqil belgi bilan tekshirildi: xaritadagi sakkizta massiv
nomining har biri oʻz chegarasi ichiga tushdi. Pasportli 174 ta
quduqning barchasi oʻz massivi chegarasida joylashgan; pasporti yoʻq 28
ta quduq chegara boʻyicha massivga biriktirildi. Massivlar maydoni
15-ilovada keltirilgan.

Sugʻoriladigan maydon xaritada sizot sathi darajalari bilan boʻyalgan
hudud boʻyicha hisoblandi; aholi punktlari va sugʻorilmaydigan yerlar
unga kirmaydi. Bu maydon (20 662 ga) quduqlarga biriktirilgan maydon
yigʻindisidan (15 985 ga) katta, chunki quduqlarga biriktirilgan maydon
faqat kuzatuv quduqlari xizmat koʻrsatadigan hududni qamraydi.
Meliorativ ekspeditsiya hisobotlarida tumanning sugʻoriladigan maydoni
17 965 ga {[}33{]}: hisobot va xarita maydon chegarasini turlicha oladi.
Kollektor zichligi faqat xaritada tasvirlangan kollektorlar boʻyicha
berilgan, xoʻjalik ichki zovurlari hisobga olinmagan. Y.Oxunboboyev
massivi hududida xaritada kollektor koʻrsatilmagan.

\textbf{GAT bazasining tuzilishi.} Natijada beshta qatlamdan iborat GAT
bazasi yaratildi (2.2-rasm). Qatlamlar ochiq GeoJSON formatida, WGS 84
koordinatalarida saqlangan; hisoblar UTM 43N tizimida bajariladi.

\begin{itemize}
\tightlist
\item
  Kuzatuv quduqlari (202 nuqta): quduq raqami, 2024-yil 1-aprel holatiga
  sath va minerallashuv, pasport maʼlumotlari, eng yaqin kollektor va
  kanalgacha masofa, 1 km radiusdagi kollektorlar uzunligi,
  2009--2024-yillardagi tekshirilgan oʻrtacha sath.
\item
  Kollektorlar (332 boʻlak, 214,6 km): boʻlak uzunligi, kadastrdagi
  nomi, foydalanishga topshirilgan va oxirgi taʼmir yili.
\item
  Sugʻorish kanallari (426 boʻlak, 170,9 km): boʻlak uzunligi.
\item
  Massivlar (8 ta maydon): nomi, umumiy va sugʻoriladigan maydon.
\item
  Tuman chegarasi.
\end{itemize}

Ularga relyef qatlami sifatida Copernicus GLO-30 modeli (30 m) qoʻshildi
{[}138{]}. Qatlamlar va ularning maʼlumotlari 2.2-rasmda koʻrsatilgan.

\oakfig{figures/fig23_gis_layers.png}{2.2-rasm. Tiklangan GAT bazasining tuzilishi: qatlamlar, obyekt turi,
obyektlar soni va asosiy maʼlumotlar (qatlamlar qiya koʻrinishda,
masshtabsiz)}

Massivlar va kuzatuv quduqlari qatlamlari GIS Cloud geoaxborot
platformasiga yuklanib, interaktiv xarita shakllantirildi {[}137{]};
xaritani suv xoʻjaligi tashkilotlari mutaxassislari bilan ulashish va
yangi oʻlchovlar bilan toʻldirish mumkin. Monitoring va kadastr
jadvallari quduq raqami va kollektor nomi orqali qatlamlarga bogʻlanadi.

\textbf{Hosil qilingan quduq koʻrsatkichlari.} Baza har bir quduq uchun
tarmoq koʻrsatkichlarini hisoblash imkonini berdi: eng yaqin
kollektorgacha masofa (medianasi 291 m, quduqlarning yarmi 130--600 m
oraligʻida), eng yaqin kanalgacha masofa va quduq atrofidagi 1 km
radiusda joylashgan kollektorlar uzunligi. Bu koʻrsatkichlar 3.4-§ da
sizot suvi sathi bilan bogʻliqlikni tahlil qilishda ishlatiladi.

\textbf{Cheklovlar.} Tiklangan bazaning ikki cheklovi bor. Uning
aniqligi boshlangʻich xarita aniqligidan yuqori boʻla olmaydi: 1:50 000
masshtabli xaritada chiziq qalinligi joyda 20--30 m ga teng, shuning
uchun kollektorgacha masofa oʻnlab metr aniqlikda baholanadi. Bundan
tashqari, baza xaritada tasvirlangan obyektlar bilan cheklangan:
xoʻjalik ichki zovurlari va yopiq drenaj unga kirmagan.

Shunga qaramay, usulni drenaj tarmogʻining raqamli xaritasi boʻlmagan
har qanday tumanda qogʻoz xarita va quduqlar koordinatalari boʻyicha bir
necha soatda takrorlash mumkin.

\oakpar{2.4-§. Monitoring maʼlumotlarining ishonchliligini baholash usuli}

\textbf{Usulning maqsadi.} 1.4-§ da koʻrsatilganidek, monitoring
maʼlumotlarida tasodifiy xatolardan tashqari doimiy nuqsonlar ham
uchraydi: oʻlchanmaganda shartli qiymat yozilishi, birlikning
almashishi, quduqlar tarkibining oʻzgarishi. Usulning maqsadi --- koʻp
yillik maʼlumotlarni tahlil qilishdan oldin bunday nuqsonlarni aniq
qoidalar boʻyicha topish va ishonchlilikni son bilan baholash.

Qoidalar maʼlumotning fizik maʼnosidan kelib chiqadi va qoʻshimcha
oʻlchov talab qilmaydi. Masalan, sugʻoriladigan hududda sizot sathi yil
boʻyi oʻzgarmay qolishi mumkin emas; bu statistik emas, fizik dalil.

Nuqsonli qiymatlar tuzatilmaydi va taxminiy qiymat bilan
almashtirilmaydi, balki tahlildan chiqariladi. Istisno faqat oʻlchov
birligi xatosi: uning tuzatilishi bir maʼnoli.

Har bir qoida natijasi alohida qayd etiladi, shuning uchun tekshiruvni
boshqa tadqiqotchi takrorlashi va natijasini tekshirishi mumkin.

\textbf{Tekshiruv qoidalari.} Sath oʻlchovlari uchun uchta qoida
qoʻllaniladi (2.2-jadval).

\oakjadvalnum{2.2-jadval}

\oakilovanom{\textbf{Monitoring maʼlumotlarini tekshirish qoidalari}}

{\fontsize{12}{14}\selectfont\setlength{\tabcolsep}{4pt}

\noindent\begin{longtable}[]{|
  >{\raggedright\arraybackslash}p{(\columnwidth - 8\tabcolsep - 5\arrayrulewidth) * \real{0.2394}}|>{\raggedright\arraybackslash}p{(\columnwidth - 8\tabcolsep - 5\arrayrulewidth) * \real{0.2272}}|>{\raggedright\arraybackslash}p{(\columnwidth - 8\tabcolsep - 5\arrayrulewidth) * \real{0.3061}}|>{\raggedright\arraybackslash}p{(\columnwidth - 8\tabcolsep - 5\arrayrulewidth) * \real{0.2272}}|}
\hline
\begin{minipage}[b]{\linewidth}\raggedright
\textbf{Qoida}
\end{minipage} & \begin{minipage}[b]{\linewidth}\raggedright
\textbf{Mezon}
\end{minipage} & \begin{minipage}[b]{\linewidth}\raggedright
\textbf{Asos}
\end{minipage} & \begin{minipage}[b]{\linewidth}\raggedright
\textbf{Chora}
\end{minipage} \\
\hline
\endhead
\hline
\endlastfoot
1. Yil ichidagi oʻzgarmaslik & Quduqning bir yildagi barcha oʻlchovlari
(6 tadan koʻp) bir xil & Sath mavsumiy tebranadi; oʻzgarmas qator quduq
quriganini yoki oʻlchov oʻtkazilmaganini bildiradi & Quduq-yil
chiqariladi \\ \hline
2. Oʻlchov birligi & Qiymat 20 m dan katta & Quduqlar chuqurligi 6 m dan
oshmaydi & Qiymat 100 ga boʻlinadi \\ \hline
3. Tarmoq tarkibi & Yillar orasida quduqlar soni yoki qiymatlar oraligʻi
keskin oʻzgargan & Tarkib oʻzgarsa, yillik oʻrtachalarni solishtirib
boʻlmaydi & Koʻp yillik tahlil faqat doimiy kuzatilgan quduqlar
boʻyicha \\
\end{longtable}

}

Birinchi qoida formula koʻrinishida quyidagicha yoziladi: i-quduqning
t-yildagi oʻlchovlari toʻplami \{$h_{itk}$\}, k = 1, \ldots, $n_{it}$ boʻlsa,
quduq-yil yaroqsiz deb topiladi, agar

\begin{equation}\tag{2.3}
n_{it} > 6 \quad \text{va} \quad \max_k h_{itk} - \min_k h_{itk} = 0
\end{equation}

«6 dan koʻp» chegarasi bir nechta oʻlchov tasodifan bir xil chiqqan
hollarni chiqarib yubormaslik uchun qoʻyilgan. Natija bu chegaraga
deyarli bogʻliq emas: chegara 3 dan 12 gacha oʻzgartirilganda
chiqariladigan quduq-yillar soni 267--272 oraligʻida qoladi, doimiy
kuzatilgan quduqlar (84 ta) va sathning oʻzgarish tezligi (yiliga 0,033
m) oʻzgarmaydi.

Uchinchi qoida boʻyicha ikki koʻrsatkich kuzatiladi: har yili kuzatilgan
quduqlar soni va chuqur qiymatlar ulushi (masalan, 5 m dan chuqur
oʻlchovlar). Ulardan birining keskin oʻzgarishi yangi quduqlar
qoʻshilganini yoki oʻlchash sharoiti oʻzgarganini bildiradi. Bunday
holda barcha quduqlar boʻyicha yillik oʻrtachalarni solishtirish
notoʻgʻri.

\textbf{Quduq qurigan holatlar (kesilgan qiymatlar).} Quruq quduq
haqiqiy sathni emas, faqat «sath quduq tubidan chuqur» degan maʼlumotni
beradi (1.4-§). Birinchi qoida bunday holatlarning bir qismini --- yil
boʻyi bir xil qiymat yozilganlarini --- topadi. Qolganlarini topish
uchun oʻlchangan sath quduq chuqurligi bilan solishtiriladi: sath quduq
tubiga 0,1 m dan yaqin boʻlsa, oʻlchov «kesilgan» deb belgilanadi. Bu
chegara 0,05 m dan 0,3 m gacha oʻzgartirilganda kesilgan oʻlchovlar
ulushi 7,0--11,6 foiz, sathning oʻzgarish tezligi esa yiliga
0,029--0,026 m boʻladi, yaʼni xulosalar oʻzgarmaydi. Bunday oʻlchovlar
koʻp yillik oʻrtachada alohida koʻrib chiqiladi.

\textbf{Ishonchlilik koʻrsatkichi.} Maʼlumotlar ishonchliligining umumiy
bahosi sifatida tekshiruvdan oʻtgan oʻlchovlar ulushi olinadi:

\begin{equation}\tag{2.4}
K_i = \frac{N_y}{N} \cdot 100\,\%
\end{equation}

bu yerda N --- tahlilga kiritilgan davrdagi oʻlchovlar soni; \(N_{y}\)
--- uch qoida boʻyicha yaroqli deb topilgan oʻlchovlar soni. Koʻrsatkich
yillar, massivlar va quduqlar boʻyicha ham hisoblanadi; bu nuqsonlar
qayerda toʻplanganini koʻrsatadi va monitoringni qayerda yaxshilash
kerakligini bildiradi.

\textbf{Doimiy kuzatilgan quduqlar (barqaror panel).} Koʻp yillik tahlil
uchun barcha yillarda kuzatilgan va hech bir yilda oʻzgarmas qiymat
bermagan quduqlar tanlanadi. Shunda quduqlar tarkibi oʻzgarishi va
qurigan quduqlar natijaga taʼsir qilmaydi. Bu quduqlar boʻyicha yillik
oʻrtacha sath

\begin{equation}\tag{2.5}
\bar{h}_t = \frac{1}{|P|} \sum_{i \in P} \bar{h}_{it}
\end{equation}

formula bilan hisoblanadi, bu yerda P --- doimiy kuzatilgan quduqlar
toʻplami; h̄ᵢₜ --- i-quduqning t-yildagi oʻrtacha sathi.

Bu usulning ham kamchiligi bor: bu guruhga sathi sayozroq quduqlar
koʻproq tushadi, chunki chuqur sathli quduqlar tez-tez quriydi va
guruhdan chiqib qoladi. Shuning uchun ularning oʻrtacha sathi tuman
oʻrtachasidan sayozroq. Bu guruh sathning oʻzini emas, uning yillar
boʻyicha oʻzgarishini baholash uchun ishlatiladi.

\textbf{Tekshiruv tartibi.} Maʼlumotlar quyidagi ketma-ketlikda qayta
ishlanadi.

\begin{enumerate}
\def\labelenumi{\arabic{enumi}.}
\tightlist
\item
  Monitoring jadvallaridan tuman quduqlari ajratiladi va bitta jadvalga
  keltiriladi: yil, quduq, oʻn kunlik, sath.
\item
  Quduq raqamlari yagona yozuvga keltiriladi (masalan, «104А» va «104а»
  bir quduq deb olinadi).
\item
  Ikkinchi qoida boʻyicha birlik xatolari tuzatiladi.
\item
  Birinchi qoida boʻyicha oʻzgarmas quduq-yillar belgilanadi va
  chiqariladi.
\item
  Uchinchi qoida boʻyicha quduqlar tarkibi tahlil qilinadi va doimiy
  kuzatilgan quduqlar tanlanadi.
\item
  Ishonchlilik koʻrsatkichi hisoblanadi.
\end{enumerate}

Tekshiruv dastur orqali bajariladi va bir tuman maʼlumotini bir necha
daqiqada qayta ishlaydi. Natijalari 3.1-§ da keltiriladi.

\oakpar{2.5-§. Statistik va geostatistik tahlil usullari}

Tekshirilgan maʼlumotlar besh guruh usul bilan tahlil qilindi. Har bir
usul qaysi masala uchun va qaysi paragrafda qoʻllanilgani 16-ilovada
koʻrsatilgan.

\textbf{Koʻp yillik oʻzgarish.} Sath yillar davomida bir tomonga
oʻzgarayotganini Mann--Kendall mezoni bilan tekshirildi {[}74, 90{]}.
Mezon koʻrsatkichi S (1.3-§) oʻzgarish boʻlmaganda nolga yaqin boʻladi;
uning tarqoqligi

\begin{equation}\tag{2.6}
\operatorname{Var}(S) = \frac{n(n-1)(2n+5)}{18}
\end{equation}

va normallashtirilgan qiymati

\begin{equation}\tag{2.7}
Z = \frac{S - \operatorname{sgn}(S)}{\sqrt{\operatorname{Var}(S)}}
\end{equation}

orqali oʻzgarish tasodifiy emasligi baholanadi; bu yerda n --- yillar
soni. Bundan tashqari Kendall koeffitsiyenti τ = S / {[}n(n − 1)/2{]}
hisoblandi. Oʻzgarish tezligi Sen qiyaligi bilan baholandi:

\begin{equation}\tag{2.8}
\beta = \operatorname{mediana} \left( \frac{x_j - x_k}{t_j - t_k} \right), \quad j > k
\end{equation}

Sen qiyaligi chetki qiymatlarga chidamli {[}117{]}. Solishtirish uchun
oddiy toʻgʻri chiziq (chiziqli regressiya) qiyaligi ham hisoblandi.

Oʻzgarish pogʻonali emasligini tekshirish uchun yillar ikki davrga
ajratildi va har bir davr ichidagi oʻzgarish alohida baholandi. Agar
davrlar ichida oʻzgarish sezilmasa-yu, davrlar orasidagi farq katta
boʻlsa, oʻzgarish bir tekis emas, pogʻonali hisoblanadi.

\textbf{Mavsumiy oʻzgarish.} Doimiy kuzatilgan quduqlar boʻyicha har bir
oʻn kunlik uchun koʻp yillik oʻrtacha sath hisoblandi. Yil ichidagi
tebranish (amplituda) eng chuqur va eng sayoz oʻn kunliklar farqi
sifatida aniqlandi.

\textbf{Sath xaritasini tuzish.} Sath xaritasi oddiy kriging usulida
tuzildi (1.3-§). Oʻlchovlar boʻyicha hisoblangan variogrammaga uchta
nazariy egri chiziq --- sferik, eksponensial va Gauss ---
moslashtirildi. Xarita aniqligi quyidagicha tekshirildi
(kross-validatsiya): har bir quduq navbat bilan chiqarib tashlanadi,
qolganlari boʻyicha uning sathi baholanadi va xato $e_{i}$ = ẑᵢ − $z_{i}$
hisoblanadi. Uchta koʻrsatkich qoʻllanildi:

\begin{equation}\tag{2.9}
ME = \frac{1}{n}\sum e_i, \qquad RMSE = \sqrt{\frac{1}{n}\sum e_i^2}, \qquad RMSSE = \sqrt{\frac{1}{n}\sum \left(\frac{e_i}{\sigma_i}\right)^2}
\end{equation}

bu yerda σᵢ --- krigingning i-nuqtadagi standart xatosi. Oʻrtacha xato
(ME) xarita sathni muntazam oshirib yoki kamaytirib koʻrsatishini,
oʻrtacha kvadratik xato (RMSE) aniqligini koʻrsatadi; RMSSE birga yaqin
boʻlsa, xarita oʻz xatosini toʻgʻri baholaydi {[}60, 67{]}. Xarita
foydali ekanini bilish uchun RMSE quduqlardagi sath tarqoqligi bilan
solishtirildi. Xaritalar faqat eng yaqin quduqqacha masofa 1,5 km dan
oshmaydigan hudud uchun tuzildi.

\textbf{Bogʻliqlik tahlili.} Sathning tarmoq koʻrsatkichlari, kanalgacha
masofa va balandlikka bogʻliqligi avval Spearman korrelyatsiyasi bilan
baholandi. Bu usul bogʻlanish toʻgʻri chiziqli boʻlishini talab qilmaydi
va chetki qiymatlarga chidamli. Soʻng omillarning har birining alohida
taʼsirini ajratish uchun bir nechta omilli regressiya tuzildi:

\begin{equation}\tag{2.10}
h_i = b_0 + b_1 \lg(d_i + 10) + b_2 \lg(c_i + 10) + b_3 z_i + \varepsilon_i
\end{equation}

bu yerda $h_{i}$ --- i-quduqning koʻp yillik oʻrtacha sathi, m; $d_{i}$ --- eng
yaqin kollektorgacha masofa, m; $c_{i}$ --- eng yaqin kanalgacha masofa, m;
$z_{i}$ --- mutlaq balandlik, m. Masofalar logarifmlandi, chunki ularning
taʼsiri yaqin masofada kuchli va uzoqlashgan sari susayadi. Ikkinchi
variantda masofa oʻrniga 1 km radiusdagi kollektorlar uzunligi olindi.
Xatolar tarqoqlik notekis boʻlgan holga mos usulda hisoblandi.

\textbf{Taʼmir taʼsirini baholash.} Taʼmirlangan kollektor yonidagi
quduqlarda (300 m gacha) taʼmirdan keyingi oʻzgarish kollektorlardan
uzoq quduqlardagi (600 m dan ortiq) oʻzgarish bilan solishtirildi
{[}34{]}. Hisobda har bir quduqning oʻzgarmas xususiyatlari (relyef,
tuproq, chuqurlik) va barcha quduqlarga bir xil taʼsir qilgan yillik
omillar (ob-havo, suv taʼminoti) chiqarib tashlanadi. Bir quduqning
turli yillardagi oʻlchovlari bir-biriga bogʻliq, shuning uchun xatolar
quduqlar boʻyicha guruhlab hisoblandi. Formula 4.3-§ da keltirilgan.

Usulning asosiy sharti: taʼmir boʻlmaganda kollektor yonidagi va
uzoqdagi quduqlarda sath bir xil oʻzgargan boʻlar edi. Bu shart
taʼmirdan oldingi yillardagi oʻzgarishlarni solishtirib tekshiriladi.

\textbf{Zarur quduqlar soni.} Δ taʼsirni ishonchli aniqlash uchun (5
foizlik ahamiyat darajasi, 80 foiz ehtimol bilan) bahoning standart
xatosi Δ / 2,8 dan oshmasligi kerak. Xato kollektor yonidagi quduqlar
soni ortgan sari kamayadi, shuning uchun zarur quduqlar soni

\begin{equation}\tag{2.11}
n_{zarur} = n_0 \left( \frac{2{,}8 \, SE_0}{\Delta} \right)^2
\end{equation}

formula bilan taxminan topiladi, bu yerda n₀ va SE₀ --- hozirgi quduqlar
soni va standart xato. Formula soddalashtirilgan, shuning uchun natija
haqiqiy quduq maʼlumotlaridan 1 000 marta tasodifiy tanlab oʻtkazilgan
sinov hisobi (Monte Karlo usuli) bilan aniqlashtirildi.

\textbf{Dasturiy taʼminot.} Barcha hisoblar Python tilida bajarildi:
jadvallar bilan ishlash uchun pandas, statistik modellar uchun
statsmodels, tendensiya testlari uchun pymannkendall, kriging uchun
pykrige, fazoviy amallar uchun shapely va pyproj, tasvirlarni qayta
ishlash uchun scikit-image kutubxonalari ishlatildi. Interaktiv xarita
GIS Cloud platformasida shakllantirildi {[}137{]}. Hisob dasturlari va
tekshirilgan maʼlumotlar saqlangan; ular boʻyicha natijalarni mustaqil
takrorlash mumkin.

\oakpar{II bob boʻyicha xulosalar}

\begin{enumerate}
\def\labelenumi{\arabic{enumi}.}
\tightlist
\item
  Tadqiqot obyekti --- Marhamat tumanining sakkizta massivdan iborat
  sugʻoriladigan yerlari. Iqlim qurgʻoqchil: yillik yogʻin 323 mm,
  bugʻlanuvchanlik 1 185 mm.
\item
  Sizot suvlarining asosiy toʻyinish manbai sugʻorish: sath yogʻin eng
  koʻp boʻlgan bahorda emas, sugʻorish eng jadal boʻlgan iyul--avgustda
  eng sayoz holatga keladi. 2024-yil 1-aprel holatiga quduqlarning 19
  foizida sath 1,5 m dan sayoz.
\item
  Kadastrga koʻra tumanda 621,1 km ochiq va 91,7 km yopiq drenaj
  tarmogʻi, 455 ta gidrotexnik inshoot bor. 14 ta magistral kollektor
  davlatlararo boʻlib, ularning asosiy qismi Qirgʻiziston hududidan
  oʻtadi. Tarmoqda atigi 6 ta suv oʻlchash inshooti bor, shuning uchun
  uning ishi asosan sizot suvi sathi orqali baholanadi.
\item
  Kuzatuv quduqlari 2013--2016-yillarda deyarli yangilangan: ularning 64
  foizi shu yillarda qurilgan. Quduqlar chuqurligi oʻrtacha 3,2 m,
  shuning uchun sath chuqur joylarda ular quriydi.
\item
  Mavjud xaritalar va quduqlar koordinatalari asosida GAT bazasini
  tiklash usuli ishlab chiqildi. Xarita 181 ta quduq orqali koordinataga
  bogʻlandi (oʻrtacha xato 1,9 m); 214,6 km kollektor va 170,9 km kanal
  raqamlandi, 60 ta kollektorga kadastr maʼlumotlari biriktirildi,
  sakkizta massiv chegarasi tiklandi.
\item
  Monitoring maʼlumotlarini tekshirishning uch qoidasi, quduq qurigan
  holatlarni aniqlash tartibi va ishonchlilik koʻrsatkichi taklif
  etildi; ular qoʻshimcha oʻlchov talab qilmaydi.
\end{enumerate}

\oakbob{III BOB. MONITORING MAʼLUMOTLARINING ISHONCHLILIGI VA SIZOT SUVI
SATHINING HUDUDLAR VA YILLAR BOʻYICHA OʻZGARISHI}

Ushbu bobda Marhamat tumanining 2009--2024-yillardagi monitoring
maʼlumotlari tekshiriladi va tekshirilgan maʼlumotlar asosida sizot suvi
sathining yillar va hududlar boʻyicha oʻzgarishi aniqlanadi. Avval
maʼlumotga qanchalik ishonish mumkinligi tekshiriladi (3.1-§). Soʻng
faqat tekshiruvdan oʻtgan maʼlumotlar boʻyicha sathning koʻp yillik va
mavsumiy oʻzgarishi (3.2-§), xaritasi (3.3-§) va drenaj tarmogʻi,
kanallar va relyefga bogʻliqligi (3.4-§) oʻrganiladi. Minerallashuv
alohida koʻrib chiqiladi (3.5-§). Oxirida natijalar Andijon viloyatining
boshqa tumanlari bilan solishtiriladi (3.6-§).

\oakpar{3.1-§. Monitoring maʼlumotlari tekshiruvi natijalari}

\textbf{Tahlil qilingan maʼlumotlar.} Tahlil uchun meliorativ
ekspeditsiyaning 2009--2024-yillardagi kuzatuv jadvallari olindi.
Jadvallarda har bir kuzatuv qudugʻi uchun sizot suvi sathi oyiga uch
marta, yaʼni yiliga 36 ta oʻn kunlik oʻlchov koʻrinishida qayd etilgan.
2010, 2013 va 2017-yillar boʻyicha mustaqil kuzatuv maʼlumotlari mavjud
emas, shuning uchun tahlilga 13 yil kiritildi. Shu yillar uchun Marhamat
tumani boʻyicha 188 ta quduq va 73 868 ta oʻlchov mavjud (massivlar va
tuman boʻyicha jami qatorlari hisobga olinmagan, 2.2-§).

Maʼlumotlar 2.4-§ da bayon etilgan uchta qoida boʻyicha tekshirildi: yil
ichidagi oʻzgarmaslik, oʻlchov birligi va kuzatuv tarmogʻi tarkibi.
Bundan tashqari har bir oʻlchov quduq chuqurligi bilan solishtirilib,
quduq qurigan holatlar (kesilgan qiymatlar) aniqlandi. Yillar boʻyicha
natijalar 17-ilovada keltirilgan.

\textbf{Yil ichidagi oʻzgarmaslik.} 270 ta quduq-yilda (9 107 ta
oʻlchov) yil davomidagi barcha qiymatlar bir xil. Sugʻoriladigan hududda
sath mavsum boʻyicha oʻzgaradi (3.2-§), shuning uchun bunday qiymatlar
haqiqiy oʻlchov emas. Oʻzgarmas qatorlar 56 ta quduqda uchraydi, shundan
36 tasida kamida uch yil davomida.

Yozilgan qiymatlarning oʻzi ularning kelib chiqishini koʻrsatadi
(18-ilova). Eng koʻp 3,0 m uchraydi (65 ta quduq-yil), keyin 2,8; 3,2 va
2,6 m --- bular kuzatuv quduqlarining odatiy chuqurligi (2.2-§).
Pasporti bor quduqlarda oʻzgarmas qiymatlarning 53 foizi quduq
chuqurligidan 0,3 m dan kam farq qiladi. Demak, quduq quriganda uning
chuqurligi yoki unga yaqin shartli qiymat yozilgan. Boshqa mumkin
boʻlgan sabab --- kuzatuvchining hisobot odati. Ularni ajratish uchun
oʻzgarmas qiymatlar quduq tubigacha zaxira (quduq chuqurligi va oʻrtacha
sath farqi) bilan solishtirildi. Zaxira 0,5 m dan kam boʻlgan
quduq-yillarning 44 foizida qiymat oʻzgarmas, zaxira 1,5--2,0 m
boʻlganda atigi 2 foizida (1 924 ta quduq-yil; p \textless{} 0,001).
Zaxirasi 2 m dan katta quduqlarda ham 9 foiz oʻzgarmas qiymat bor:
ularning bir qismini hisobot odati tushuntiradi.

\textbf{Quduq qurigan holatlar (kesilgan qiymatlar).} Qiymati oʻzgarib
turgan quduqlarda ham quduq qurigan holatlar uchraydi: bir necha oʻn
kunlik davomida sath quduq tubiga yetib, keyin yana koʻtariladi.
Pasportda chuqurligi koʻrsatilgan quduqlarning yaroqli oʻlchovlari
orasida 5 326 tasi (8,2 foiz) quduq chuqurligidan 0,1 m dan kamroq farq
qiladi. Bu oʻlchovlar haqiqiy sathni emas, «sath quduq tubidan chuqur»
degan holatni aks ettiradi. Kesilgan oʻlchovlar 2018-yildan keyin keskin
koʻpaygan: 2009--2016-yillarda yiliga 29--155 ta, 2018--2024-yillarda
500--1 023 ta.

Kesilgan oʻlchovlar sathni sayozroq koʻrsatadi: haqiqiy sath yozilgan
qiymatdan chuqurroq. Shuning uchun ular chiqarilmadi, balki alohida
belgilandi; ularning taʼsiri 3.2 va 3.3-§ larda hisobga olinadi.

\textbf{Oʻlchov birligi xatolari.} 13 ta qiymat 20 dan katta (masalan,
400, 532, 197). Ular metr oʻrniga santimetrda yozilgan, shuning uchun
100 ga boʻlindi.

\textbf{Kuzatuv tarmogʻining oʻzgarishi.} Quduqlar tarkibi kuzatilgan
davrda oʻzgarib turgan (19-ilova; u yerda har bir tahlilga kiritilgan
quduqlar soni ham bor). 2009--2016-yillarda har yili 128--138 ta quduq
kuzatilgan, quduqlar tarkibi yildan-yilga biroz oʻzgargan: yiliga 1--12
ta quduq qoʻshilgan va 1--8 tasi chiqib ketgan. 2018-yilda esa tarmoqqa
birdaniga 60 ta yangi quduq qoʻshilgan va kuzatilgan quduqlar soni 188
taga yetgan; 2020-yilda yana 7 ta quduq qoʻshilgan. Bu
2013--2016-yillarda 120 ta yangi quduq qurilgani bilan mos keladi
(2.2-§): yangi quduqlar kuzatuv jadvallariga asosan 2018-yildan
kiritilgan.

Bu oʻzgarish oʻlchovlarda aniq iz qoldirgan. 2018-yilgacha 5 m dan
chuqur oʻlchovlar deyarli uchramaydi (0--0,1 foiz), 2018-yildan keyin
esa ular har yili 3,2--4,3 foizni tashkil etadi. Kesilgan oʻlchovlar
soni ham 2018-yildan koʻpaygan. Yangi quduqlarning bir qismi sizot suvi
chuqurroq joylashgan hududlarga qoʻyilgan va ularning bir qismi
chuqurroq (4--6 m) qurilgan. Shuning uchun yillarni barcha quduqlar
oʻrtachasi boʻyicha solishtirish notoʻgʻri: 2018-yildagi oʻrtacha sath
oʻzgarishi asosan quduqlar tarkibi oʻzgarganini koʻrsatadi. Koʻp yillik
tahlil faqat doimiy kuzatilgan quduqlar boʻyicha bajarilishi kerak
(3.2-§).

\textbf{Massivlar boʻyicha ishonchlilik.} Maʼlumotlar ishonchliligi
massivlar orasida keskin farq qiladi (20-ilova).

Ishonchlilik sizot suvi sathi bilan bevosita bogʻliq. Sath sayoz boʻlgan
Mashʼal va Istiqlol massivlarida yaroqli oʻlchovlar 97--98 foiz,
kesilganlar 3--6 foiz. Sath chuqur boʻlgan T.Mirzayev va Y.Oxunboboyev
massivlarida esa oʻlchovlarning faqat 50--60 foizi yaroqli,
yaroqlilarining ham 21--29 foizi kesilgan. Sababi aniq: sath chuqur
joylarda sayoz quduqlar tez-tez quriydi.

Demak, monitoring xatosi tasodifiy emas: u sizot suvi chuqur, yaʼni
meliorativ jihatdan qulay hududlarda toʻplanadi. Natijada hudud boʻyicha
oʻrtacha sath haqiqatdagidan sayozroq, meliorativ holat esa ogʻirroq
koʻrinadi. T.Mirzayev va Y.Oxunboboyev massivlari uchun mavjud kuzatuv
tarmogʻi sathni ishonchli oʻlchashga yaroqsiz; bu massivlarda quduqlarni
chuqurlashtirish birinchi navbatdagi vazifa hisoblanadi (4.5-§).

\textbf{Ishonchsiz maʼlumotlarning joylashuvi.} Koordinatasi aniqlangan
188 ta quduq boʻyicha yaroqli oʻlchovlar ulushi xaritaga tushirildi
(3.1-rasm, a). Ishonchsiz quduqlar tuman boʻylab tasodifiy tarqalmagan.
Yaroqli oʻlchovlari 75 foizdan kam boʻlgan 33 ta quduqdan 27 tasi va
kamida uch yil oʻzgarmas qatorga ega 36 ta quduqdan 29 tasi tumanning
janubi-sharqidagi T.Mirzayev, Y.Oxunboboyev, M.Toʻjiboyev va U.Qoraboyev
massivlarida joylashgan.

Bu taqsimotning sababi 3.1-rasm, b da koʻrinadi. Unda quduq tubigacha
zaxira, yaʼni pasportdagi chuqurlik bilan 2019--2024-yillardagi oʻrtacha
sath orasidagi farq koʻrsatilgan. 182 ta quduq boʻyicha zaxira medianasi
1,14 m, 80 ta quduqda u 1 m dan kam, 5 ta quduqda esa oʻrtacha sath
pasportdagi chuqurlikdan chuqurroq qayd etilgan, yaʼni pasport maʼlumoti
yoki oʻlchovning oʻzi notoʻgʻri. Zaxira qancha kichik boʻlsa, yaroqli
oʻlchovlar shuncha kam: zaxirasi 1 m dan kam quduqlarda yaroqli
oʻlchovlar oʻrtacha 81,5 foiz, qolganlarida 91,7 foiz. Sath chuqur
joylarda quduqlar quriy boshlaydi, shuning uchun maʼlumot sifati quduq
qayerda joylashganiga bogʻliq.

\oakfig{figures/map2_data_quality.png}{3.1-rasm. Monitoring maʼlumotlari ishonchliligining quduqlar boʻyicha
taqsimoti: a --- quduqlar boʻyicha yaroqli oʻlchovlar ulushi,
2009--2024-yillar (massiv yozuvlarida --- massiv boʻyicha ulush); b ---
quduq tubigacha zaxira (massiv yozuvlarida --- mediana)}

\textbf{Umumiy baho.} Tekshiruvdan soʻng 64 761 ta oʻlchov (umumiy
hajmning 87,7 foizi) tahlil uchun yaroqli deb topildi; ularning 8,2
foizi kesilgan qiymat sifatida belgilandi. Haqiqiy sathni aks ettiruvchi
toʻliq oʻlchovlar umumiy hajmning taxminan 80 foizini tashkil etadi. Bu
Marhamat tumani monitoring maʼlumotlari ishonchliligining son bilan
ifodalangan bahosi.

\textbf{2025-yil bahori.} Tekshiruv 2025-yil mart va aprel oylaridagi
188 ta quduq oʻlchovlariga ham qoʻllanildi (56-ilova). Shu ikki oyning
oltita oʻn kunligida qiymati bir xil boʻlgan quduqlar ulushi 32,4 foiz;
2021--2024-yillarning shu oylarida bu ulush 12,8--14,9 foiz edi. Demak,
tekshiruv yangi maʼlumotlar uchun ham zarur. Yaroqli quduqlarda bahorgi
oʻrtacha sath 2,03 m boʻlib, 2023--2024-yillar darajasida (2,04--2,09 m)
qoldi.

Monitoringning asosiy kamchiliklari: quduqlar sizot suvi sathiga
nisbatan sayoz, quduq qurigan holatlar va quduqlar tarkibidagi
oʻzgarishlar hujjatlarda qayd etilmaydi. Ularni bartaraf etish boʻyicha
tavsiyalar 4.5-§ da berilgan.

\oakpar{3.2-§. Sizot suvi sathining koʻp yillik va mavsumiy oʻzgarishi}

Ushbu paragrafda sizot suvi sathining 2009--2024-yillardagi oʻzgarishi
uch darajada koʻrib chiqiladi: tuman, massivlar va alohida quduqlar. Har
birida ikki savolga javob izlanadi: sath qaysi tomonga va qanday
tezlikda oʻzgargan, bu natija 3.1-§ dagi maʼlumot nuqsonlariga qanchalik
bogʻliq. Oʻzgarish Mann--Kendall mezoni {[}74, 90{]} va Sen qiyaligi
{[}117{]} bilan baholandi; bu usullar chetki qiymatlarga chidamli va
gidrologiyada keng qoʻllanadi {[}66, 68{]}. Sen qiyaligi haqiqiy yillar
boʻyicha hisoblandi, shuning uchun 2010, 2013 va 2017-yillar oʻrnidagi
boʻshliqlar natijani buzmaydi.

\textbf{Doimiy kuzatilgan quduqlar.} Tuman boʻyicha koʻp yillik
oʻzgarish 13 yilning barchasida kuzatilgan va hech bir yilda oʻzgarmas
qiymat bermagan quduqlar boʻyicha baholandi; ular 84 ta. Shunda har yili
aynan bir xil quduqlar solishtiriladi. 2018-yilda qoʻshilgan 60 ta yangi
quduqni (3.1-§) oʻrtachaga qoʻshish yillar farqini quduqlar tarkibi
farqi bilan aralashtirib yuboradi. Natijalar 21-ilova va 3.2-rasmda
keltirilgan. Yillik oʻrtachadan tashqari iyul--avgust (19--24-dekadalar)
oʻrtacha sathi ham hisoblandi, chunki drenaj aynan shu davrda eng katta
yuk bilan ishlaydi.

Bu quduqlarda yillik oʻrtacha sath 2009-yildagi 1,60 m dan 2024-yilda
1,99 m gacha chuqurlashgan. Oʻzgarish tasodifiy emas (Mann--Kendall
mezoni, p = 0,0005). Sen qiyaligi yiliga 0,033 m (95 foizlik ishonch
oraligʻi 0,022--0,046 m), oddiy toʻgʻri chiziq boʻyicha yiliga 0,032 m;
ikki baho mos keladi. Iyul--avgustdagi sath ham xuddi shunday tezlikda
chuqurlashgan: 1,44 m dan 1,80 m gacha, yiliga 0,031 m. Yozgi sath
yillik oʻrtachadan har yili 0,11--0,28 m sayozroq, lekin ikkalasi bir
xil oʻzgaradi. Demak, chuqurlashish sugʻorish mavsumiga ham tegishli.

\oakfig{figures/ch3_dynamics.png}{3.2-rasm. Sizot suvi sathi: a --- koʻp yillik oʻzgarish (tekshirilmagan
maʼlumot va doimiy kuzatilgan quduqlar); b --- yil ichidagi oʻrtacha
yoʻli oʻn kunliklar boʻyicha}

\textbf{Pogʻonali oʻzgarish.} Sath bir tekis emas, pogʻona bilan
oʻzgargan. 2009--2016-yillarda oʻrtacha sath 1,63 m, 2018--2024-yillarda
1,92 m; farq 0,30 m. Har bir davr ichida esa sezilarli oʻzgarish yoʻq
(Sen qiyaligi yiliga 0,021 va 0,022 m; p = 0,26 va 0,23). Pettitt mezoni
{[}102{]} boʻyicha sakrash 2016-yildan keyin yuz bergan (p = 0,023),
yozgi sathda ham shunday. Oʻzgarishning asosiy qismi 2016 va 2018-yillar
orasida yuz bergan; 2017-yil boʻyicha mustaqil kuzatuv maʼlumotlari
mavjud emasligi sababli uzilishni bir yil aniqligida belgilab boʻlmaydi.

Bu pogʻonaning ikki mumkin boʻlgan izohi bor, mavjud maʼlumot ularni
toʻliq ajratishga imkon bermaydi. Birinchisi haqiqiy meliorativ
oʻzgarish: kadastrga koʻra tumandagi 90 ta kollektorning 55 tasi aynan
2016--2018-yillarda taʼmirlangan. Ikkinchisi oʻlchash tartibining
oʻzgarishi: 2018-yilda kuzatuv tarmogʻi qayta tuzilgan (3.1-§). Doimiy
kuzatilgan quduqlarni olish quduqlar tarkibi taʼsirini yoʻqotadi, lekin
oʻsha quduqlarda oʻlchash tartibi oʻzgargan boʻlsa, buni yoʻqota
olmaydi. Taʼmirning taʼsiri alohida usul bilan, taʼmirlangan va
taʼmirlanmagan kollektorlar yaqinidagi quduqlarni solishtirish orqali
4.3-§ da baholanadi.

\textbf{Quduq qurigan holatlarning taʼsiri.} Quduq quriganda sath
oʻrniga quduq chuqurligi yoki unga yaqin qiymat yoziladi va bu sathni
sayozroq koʻrsatadi {[}65{]}. Doimiy kuzatilgan quduqlarda bunday
oʻlchovlar 3,5 foiz, lekin ular yillar boʻyicha notekis. Shuning uchun
oʻzgarish tezligi ularsiz ham hisoblandi: sath 1,60 m dan 1,93 m gacha,
yiliga 0,028 m (p = 0,0008), yaʼni 15 foiz kam. Davrlar orasidagi farq
0,30 m oʻrniga 0,26 m.

\textbf{Qoʻshni yillar bogʻliqligi.} Ketma-ket yillar sathi bir-biriga
kuchli bogʻliq, asosan pogʻona tufayli. Buni hisobga oluvchi Hamed--Rao
tuzatishi bilan ham oʻzgarish ishonchli qoladi (p = 0,0005) {[}62{]};
boshqa tuzatish usulida (pre-whitening) esa p = 0,11, chunki bu usul 13
yillik qisqa qatorda oʻzgarishni sezmay qoladi. Shuning uchun natija
ikki koʻrsatkich bilan beriladi: yillik tezlik va ikki davr orasidagi
pogʻona (0,30 m; p \textless{} 0,001).

Kesilgan qiymatlarni olib tashlash eng chuqur holatlarni chiqarib
yuboradi va tezlikning pastki chegarasini beradi; ularni saqlash esa
yuqori chegarasiga yaqin. Haqiqiy tezlik shu ikki baho oraligʻida ---
yiliga 0,028--0,033 m. Chuqurlashish xulosasi oʻzgarmaydi, tezlik esa
±15 foiz aniqlikda.

\textbf{Tekshiruvsiz hisobning xatosi.} Agar barcha quduqlarning
maʼlumoti tekshiruvsiz olinsa, oʻrtacha sath 1,89 m dan (2009) 2,47 m
gacha (2024) oʻzgargan boʻlib chiqadi, tezlik yiliga 0,045 m. Bu doimiy
kuzatilgan quduqlar bahosidan (yiliga 0,032 m) 1,4 baravar katta.
Tekshiruvsiz oʻrtacha sathning oʻzi ham har yili 0,29--0,53 m chuqur.
Sababi --- 2018-yilda chuqurroq sathli yangi quduqlar qoʻshilgani va
oʻzgarmas shartli qiymatlar.

\textbf{Massivlar boʻyicha oʻzgarish.} Tuman oʻrtachasi massivlar
orasidagi farqlarni yashiradi. Shuning uchun doimiy kuzatilgan quduqlar
massivlar boʻyicha guruhlandi (3.1-jadval). Y.Oxunboboyev massivida
bunday quduq yoʻq: bu massivda oʻlchovlarning yarmi yaroqsiz (3.1-§),
shuning uchun u jadvalga kiritilmadi.

\oakjadvalnum{3.1-jadval}

\oakilovanom{\textbf{Doimiy kuzatilgan quduqlar boʻyicha sizot suvi sathining
massivlardagi oʻzgarishi}}

{\fontsize{10}{12}\selectfont\setlength{\tabcolsep}{2.5pt}

\noindent\begin{longtable}[]{|
  >{\raggedright\arraybackslash}p{(\columnwidth - 18\tabcolsep - 10\arrayrulewidth) * \real{0.1663}}|>{\raggedright\arraybackslash}p{(\columnwidth - 18\tabcolsep - 10\arrayrulewidth) * \real{0.1084}}|>{\raggedright\arraybackslash}p{(\columnwidth - 18\tabcolsep - 10\arrayrulewidth) * \real{0.0563}}|>{\raggedright\arraybackslash}p{(\columnwidth - 18\tabcolsep - 10\arrayrulewidth) * \real{0.0563}}|>{\raggedright\arraybackslash}p{(\columnwidth - 18\tabcolsep - 10\arrayrulewidth) * \real{0.1227}}|>{\raggedright\arraybackslash}p{(\columnwidth - 18\tabcolsep - 10\arrayrulewidth) * \real{0.1227}}|>{\raggedright\arraybackslash}p{(\columnwidth - 18\tabcolsep - 10\arrayrulewidth) * \real{0.1619}}|>{\raggedright\arraybackslash}p{(\columnwidth - 18\tabcolsep - 10\arrayrulewidth) * \real{0.0832}}|>{\raggedright\arraybackslash}p{(\columnwidth - 18\tabcolsep - 10\arrayrulewidth) * \real{0.1224}}|}
\hline
\begin{minipage}[b]{\linewidth}\raggedright
\textbf{Massiv}
\end{minipage} & \begin{minipage}[b]{\linewidth}\raggedright
\textbf{Quduqlar soni}
\end{minipage} & \begin{minipage}[b]{\linewidth}\raggedright
\textbf{Sath 2009 y., m}
\end{minipage} & \begin{minipage}[b]{\linewidth}\raggedright
\textbf{Sath 2024 y., m}
\end{minipage} & \begin{minipage}[b]{\linewidth}\raggedright
\textbf{2009--2016 oʻrtacha, m}
\end{minipage} & \begin{minipage}[b]{\linewidth}\raggedright
\textbf{2018--2024 oʻrtacha, m}
\end{minipage} & \begin{minipage}[b]{\linewidth}\raggedright
\textbf{Mann--Kendall τ}
\end{minipage} & \begin{minipage}[b]{\linewidth}\raggedright
\textbf{p}
\end{minipage} & \begin{minipage}[b]{\linewidth}\raggedright
\textbf{Sen qiyaligi, m/yil}
\end{minipage} \\
\hline
\endhead
\hline
\endlastfoot
Istiqlol & 15 & 1,51 & 2,22 & 1,76 & 2,22 & 0,641 & 0,003 & +0,047 \\ \hline
M.Toʻjiboyev & 10 & 1,79 & 1,98 & 1,62 & 1,88 & 0,667 & 0,002 &
+0,040 \\ \hline
Mashʼal & 17 & 1,46 & 1,82 & 1,46 & 1,80 & 0,641 & 0,003 & +0,031 \\ \hline
Paxtakor & 9 & 1,85 & 2,43 & 1,92 & 2,32 & 0,667 & 0,002 & +0,041 \\ \hline
T.Mirzayev & 6 & 1,46 & 2,11 & 1,48 & 1,85 & 0,795 & 0,0002 & +0,051 \\ \hline
U.Qoraboyev & 9 & 1,64 & 2,17 & 1,57 & 2,04 & 0,718 & 0,0008 & +0,053 \\ \hline
Z.M.Bobur & 9 & 1,89 & 1,51 & 1,84 & 1,55 & −0,487 & 0,024 & −0,031 \\ \hline
Tuman & 84 & 1,60 & 1,99 & 1,63 & 1,92 & 0,744 & 0,0005 & +0,033 \\
\end{longtable}

}

Yetti massivdan oltitasida sath ishonchli chuqurlashgan, tezlik yiliga
0,031 m dan 0,053 m gacha. Eng tez chuqurlashish tumanning sharqiy va
janubiy qismlaridagi U.Qoraboyev va T.Mirzayev massivlarida kuzatiladi.
Z.M.Bobur massivida jarayon teskari: sath 1,89 m dan 1,51 m gacha
koʻtarilgan, yiliga 0,031 m tezlikda sayozlashgan (p = 0,024). Bu massiv
tumanning shimoliy qismida joylashgan va 2018--2024-yillarda oʻrtacha
sathi 1,55 m bilan tumanda eng sayoz massivga aylangan. Massivlar
orasidagi tafovut ham kuchaygan: 2009-yilda massivlarning oʻrtacha sathi
1,46--1,89 m oraligʻida boʻlgan boʻlsa, 2024-yilda 1,51--2,43 m
oraligʻida. Tuman oʻrtachasi esa bu ikki qarama-qarshi jarayonni bitta
songa aralashtiradi.

Z.M.Bobur massividagi sayozlashish meliorativ jihatdan xavfli: sathning
1,5 m atrofigacha koʻtarilishi Oʻzbekiston sharoitida qabul qilingan
kritik chuqurlikka yaqinlashish demakdir: bu chuqurlikda kapillyar
koʻtarilish orqali tuz toʻplanishi tezlashadi {[}55, 70{]}. Massivdagi
oʻzgarish tumanning umumiy pogʻonasidan farqli vaqtda yuz bergan:
2009--2018-yillarda sath 1,65--1,95 m oraligʻida boʻlgan, 2019-yildan
boshlab esa 1,27--1,71 m oraligʻiga koʻtarilgan va 2022-yilda 1,27 m ga
yetgan. Shu sababli bu massiv drenaj tarmogʻini birinchi navbatda
tekshirish talab etiladigan hudud sifatida ajratiladi (4.5-§).

\textbf{Quduqlar boʻyicha.} Massiv ichidagi farqni koʻrish uchun
oʻzgarish har bir quduq boʻyicha hisoblandi. Bunda kamida 8 yil yaroqli
maʼlumoti bor barcha quduqlar olindi; ular 122 ta. Natijalar 22-ilova va
3.3-rasmda keltirilgan.

122 ta quduqdan 86 tasida (70 foiz) sath chuqurlashgan, 36 tasida
sayozlashgan. Ishonchli oʻzgarish 54 ta quduqda chuqurlashish, 9 tasida
sayozlashish koʻrinishida. Koʻp quduqni bir vaqtda tekshirganda
tasodifiy natijalarni chiqaruvchi tuzatishdan keyin ham
(Benjamini--Hochberg {[}38{]}) 37 ta quduqda chuqurlashish va 6 tasida
sayozlashish ishonchli qoladi. Oʻzgarish tezligining medianasi yiliga
0,023 m. Doimiy kuzatilgan 84 ta quduqda ham shunday: 65 tasida
chuqurlashish, 19 tasida sayozlashish. Demak, tuman boʻyicha sath
chuqurlashayotgan boʻlsa ham, quduqlarning taxminan toʻrtdan birida sath
koʻtarilmoqda.

Sayozlashayotgan quduqlar hudud boʻyicha tasodifiy tarqalmagan
(3.3-rasm). Ularning katta qismi Z.M.Bobur massivida, tumanning
shimolida toʻplangan: bu massivdagi 23 ta quduqdan faqat 5 tasida sath
chuqurlashgan. Qolgan massivlarda chuqurlashayotgan quduqlar ulushi
74--100 foiz. Rasmda toʻldirilgan belgilar ishonchli oʻzgarishli
quduqlarni bildiradi; ular sharqiy va janubi-gʻarbiy qismlarda zich.
3.3-rasm, b da 3.1-jadvaldagi massivlar boʻyicha oʻzgarish koʻrsatilgan:
Z.M.Bobur massivi sath koʻtarilayotgan yagona massiv (yiliga −3,1 sm),
qolgan massivlarda sath yiliga 3,1--5,3 sm chuqurlashmoqda.

\oakfig{figures/fig34_trend_map.png}{3.3-rasm. Sizot suvi sathining oʻzgarish tezligi (Sen qiyaligi, sm/yil):
a --- quduqlar boʻyicha (▲ --- koʻtarilish, ▼ --- chuqurlashish;
toʻldirilgan belgilar --- ishonchli oʻzgarish, p \textless{} 0,05); b
--- massivlar boʻyicha doimiy kuzatilgan quduqlar (3.1-jadval)}

Quduqdagi oʻzgarish qaysi omillarga bogʻliqligi Spearman korrelyatsiyasi
bilan tekshirildi. Eng kuchli bogʻlanish quduqning hozirgi sathi bilan:
2019--2024-yillarda sathi chuqur quduqlarda chuqurlashish tezroq (ρ =
+0,55). Bu bogʻlanishning bir qismi tabiiy: tez chuqurlashgan quduq
oxirgi yillarda chuqur sath beradi. Shu bilan birga sathi allaqachon
sayoz boʻlgan joylarda chuqurlashish deyarli kuzatilmaydi: tuman
boʻyicha oʻrtacha yaxshilanish eng muammoli quduqlarga yetib bormagan.
Kollektorgacha masofa bilan bogʻlanish kuchsizroq, lekin ishonchli (ρ =
+0,25): kollektordan uzoqdagi quduqlarda sath tezroq chuqurlashgan. Bu
bogʻlanish 3.4-§ da batafsil koʻrib chiqiladi.

\textbf{Mavsumiy oʻzgarish.} Tekshirilgan maʼlumotlarda sath yil ichida
aniq mavsumiy oʻzgaradi (3.2-rasm, b). Sath yanvar--fevralda eng chuqur
(2,03--2,04 m), mart oxiridan koʻtarila boshlaydi va iyul oxirida eng
sayoz holatga yetadi (1,56 m). May oyining ikkinchi yarmidan avgust
oxirigacha sath 1,56--1,62 m oraligʻida turadi, sentabrdan boshlab yana
chuqurlashadi. Yil ichidagi tebranish 0,48 m. Bu yoʻl sugʻorish mavsumi
bilan mos keladi: 2.1-§ da koʻrsatilganidek, yogʻinlarning 65 foizi
noyabr--aprelda tushadi, lekin sath shu davrda eng chuqur, koʻtarilish
esa sugʻorish mavsumiga toʻgʻri keladi. Demak, sathni asosan sugʻorish
va kanallardan sizib chiqqan suv belgilaydi {[}130{]}; lekin quduqlar
atrofida berilgan suv hajmi nomaʼlum, shuning uchun bu son bilan
tekshirilmadi.

Massivlar boʻyicha mavsumiy oʻzgarish sezilarli farq qiladi (23 va
24-ilovalar): tebranish 0,32 m dan (T.Mirzayev) 0,85 m gacha
(U.Qoraboyev).

Massivlarni mavsumiy oʻzgarish boʻyicha ikki guruhga ajratish mumkin.
Birinchi guruhda (Istiqlol, M.Toʻjiboyev, U.Qoraboyev) sath sugʻorish
mavsumining avjida, iyul oxirida eng sayoz holatga yetadi. Bu
massivlarda tebranish katta (0,43--0,85 m), yozgi koʻtarilish asosan
sugʻorish suvining sizib oʻtishi bilan bogʻliq. Ikkinchi guruhda
(Mashʼal, T.Mirzayev, Z.M.Bobur) sath may oyida eng sayoz boʻladi va
yozda koʻtarilish toʻxtaydi. Bunday yoʻl vegetatsiya oldidan
oʻtkaziladigan shoʻr yuvish va bahorgi sugʻorishlarning taʼsiri
kuchliroq ekanini, yozda esa drenaj tarmogʻi ortiqcha suvni olib keta
olishini koʻrsatadi. Paxtakor massivi oraliq holatda: eng sayoz sath
iyun oʻrtasida.

U.Qoraboyev massivining yozgi eng sayoz sathi (1,45 m) tumandagi eng
past koʻrsatkich, qishki eng chuqur sathi esa (2,29 m) eng chuqurlaridan
biri. Massivning yillik oʻrtacha sathi (3.1-jadval) qoniqarli koʻrinsa
ham, yozda sath kritik chuqurlikka yetadi. Bu massivda yillik oʻrtacha
sath eng tez chuqurlashgan (yiliga 0,053 m), shuning uchun yillik
oʻrtacha koʻrsatkich undagi holatni haqiqatdagidan yaxshiroq koʻrsatadi.
Shu sababli drenaj tizimi samaradorligini baholashda yillik oʻrtacha
emas, yozgi eng sayoz holat asosiy mezon boʻlishi kerak.

\oakpar{3.3-§. Sizot suvi sathi xaritasi va uning aniqligi}

Quduqlar sathni faqat alohida nuqtalarda beradi, qarorlar esa maydon
boʻyicha qabul qilinadi: qaysi massivning qancha qismida sath kritik
chuqurlikdan sayoz, qayerda drenajni tekshirish kerak. Quduqlar
orasidagi sathni topish uchun kriging usuli qoʻllanildi {[}60, 91,
131{]}. Uning afzalligi --- har bir nuqtada baho bilan birga uning
xatosini ham beradi {[}48, 71{]}. Shuning uchun sath xaritasidan
tashqari xato xaritasi va sathning kritik chuqurlikdan sayoz boʻlish
ehtimoli xaritasi tuziladi. Hisoblar 2.4-§ da bayon etilgan tartibda
bajarildi.

\textbf{Boshlangʻich maʼlumotlar.} Xarita ikki koʻrsatkich boʻyicha
tuzildi. Birinchisi har bir quduqning 2019--2024-yillardagi tozalangan
oʻrtacha sathi; bunday quduqlar 168 ta. Ikkinchisi shu yillardagi
iyul--avgust oylari oʻrtacha sathi (167 ta quduq). Yozgi sath drenajni
baholashning asosiy mezoni (3.2-§), shuning uchun u alohida xaritalandi.
Koʻp yillik oʻrtacha bitta yilning tasodifiy xatolarini kamaytiradi.
Qoʻshni quduqlar orasidagi masofaning medianasi 606 m.

Ikkala koʻrsatkichda ham qiymatlar notekis taqsimlangan (25-ilova).
Yillik sathning oʻrtachasi 2,39 m, medianasi 2,25 m: sathi 4--6 m
boʻlgan kam sonli quduqlar oʻrtachani oshiradi. Yozgi sath medianasi
1,93 m. Kriging uchun bu muhim emas, lekin ehtimollik xaritasida hisobga
olinadi.

\textbf{Variogramma.} Variogramma 500 m qadam bilan 10 km gacha
hisoblandi (26-ilova, a). Eng qisqa masofada (0,5--1,0 km) u 0,52--0,57
m², yaʼni qiymatlar umumiy tarqoqligining (0,97 m²) qariyb yarmi. Demak,
qoʻshni quduqlar sathi ham bir-biridan ancha farq qiladi: 500 m
oraliqdagi ikki quduq orasidagi farq hududning istalgan ikki qudugʻi
orasidagi farqning yarmiga teng. Masofa ortgan sari variogramma toʻgʻri
chiziq boʻylab oʻsadi va toʻxtamaydi. Bu sathda tuman boʻylab umumiy
yoʻnalish borligini koʻrsatadi {[}97, 131{]}.

Bu umumiy yoʻnalish koordinatalar boʻyicha regressiya bilan tasdiqlandi:
u sathdagi farqlarning 18 foizini tushuntiradi (p \textless{} 0,001).
Sath sharq tomonga har kilometrda 0,10 m, janub tomonga har kilometrda
0,07 m chuqurlashadi. Yozgi sath uchun ham manzara oʻxshash. Bu
yoʻnalish tumanning umumiy relyefiga va 3.2-§ da aniqlangan massivlar
farqiga mos keladi: sayoz sathli Z.M.Bobur, Istiqlol va Mashʼal
massivlari tumanning shimoliy va gʻarbiy qismida, chuqur sathli
Y.Oxunboboyev va T.Mirzayev massivlari janubda joylashgan.

Turli yoʻnalishlar boʻyicha variogrammalar (26-ilova, b) deyarli farq
qilmaydi: gʻarb--sharq yoʻnalishida sath biroz silliqroq oʻzgaradi,
lekin farq kichik va barqaror emas, shuning uchun hisobga olinmadi.

\textbf{Egri chiziqni tanlash.} Variogrammaga uchta nazariy egri chiziq
(sferik, eksponensial va Gauss) moslashtirildi {[}97{]}. Bunda 10 km
gacha masofa olindi, chunki kriging bahosini asosan eng yaqin quduqlar
belgilaydi. Solishtirish uchun dasturning avtomatik tanlovi ham
keltirilgan. Har bir variantning aniqligi quduqlarni navbat bilan
chiqarib tekshirildi (27-ilova).

Eng kichik xatoni ikkala eksponensial variant berdi (0,813--0,819 m).
Biroq ular oʻz xatosini turlicha baholaydi: avtomatik variant xatoni
taxminan 20 foiz kam koʻrsatadi, qoʻlda moslashtirilgan variant (nagget
0,50 m²) esa haqiqatga yaqin. Xato va ehtimollik xaritalari uchun aynan
xatoni toʻgʻri baholash muhim, shuning uchun barcha xaritalar shu
variant boʻyicha tuzildi.

\textbf{Boshqa usullar bilan solishtirish.} Oddiy kriging yana ikki usul
bilan solishtirildi. Universal krigingda sathning tuman boʻylab umumiy
yoʻnalishi hisobga olinadi {[}48, 131{]}. Regressiya-krigingda sath
avval drenaj tarmogʻi koʻrsatkichlari (kollektorgacha masofa, 1 km
radiusdagi kollektorlar uzunligi) boʻyicha baholanadi, qolgan farq esa
kriging bilan toʻldiriladi {[}67{]}. Bundan tashqari ikki oddiy usul ---
eng yaqin quduq qiymati va barcha quduqlar oʻrtachasi --- ham hisoblandi
(28-ilova).

Umumiy yoʻnalishni yoki tarmoq koʻrsatkichlarini qoʻshish xaritani
yaxshilamaydi: farq 0,01--0,02 m va oddiy kriging foydasiga. Tarmoq
koʻrsatkichlari sath bilan bogʻlangan (3.4-§), lekin bashorat uchun juda
kuchsiz. Eng yaqin quduq qiymati esa hudud oʻrtachasidan ham yomon baho
beradi (1,052 m ga qarshi 0,991 m). Demak, alohida quduq oʻz atrofini
toʻliq ifodalamaydi. Shu sababli bitta quduq oʻlchovi boʻyicha unga
qoʻshni dalaning meliorativ holati haqida xulosa chiqarish xatoga olib
keladi. Kriging hudud oʻrtachasiga nisbatan xatoni atigi 17 foizga
kamaytiradi: yillik sath uchun 0,99 m dan 0,82 m gacha, yozgi sath uchun
0,93 m dan 0,77 m gacha.

\textbf{Oʻrtacha yillik sath xaritasi.} Natijada ikki xarita olindi
(3.4-rasm): sath xaritasi va uning xatosi xaritasi. Xaritalar tuman
chegarasi ichida, faqat eng yaqin quduqqacha masofa 1,5 km dan
oshmaydigan hudud uchun tuzildi; bu hudud tuman maydonining 82 foizini
qamraydi. Unda baholangan sath 1,57--4,52 m oraligʻida, hududning 32
foizida sath 2,0 m dan sayoz. Sayoz sathli zona tumanning
shimoli-sharqidan janubi-gʻarbiga choʻzilgan va kollektor tarmogʻi zich
joylashgan hududga toʻgʻri keladi; chuqur sathli zonalar janubda va
sharqda joylashgan. Xato 0,79--1,11 m, quduqlar siyrak joylarda va hudud
chekkalarida eng katta. Xato xaritasi kuzatuv tarmogʻini toʻldirish
uchun asos boʻladi: yangi quduq xato eng katta joyga qoʻyilganda eng
koʻp maʼlumot beradi (4.5-§).

\oakfig{figures/fig35_kriging.png}{3.4-rasm. Oddiy kriging natijalari: a --- sizot suvi sathining
2019--2024-yillardagi oʻrtacha xaritasi (chiziqlar --- 0,5 m oraliqdagi
teng chuqurlik chiziqlari); b --- xarita xatosi. Nuqtalar --- kuzatuv
quduqlari, yashil chiziqlar --- kollektorlar}

\textbf{Kriging xaritasi sayoz joylarni yashiradi.} Kriging bahosi
qoʻshni quduqlarning oʻrtachasi, shuning uchun xarita haqiqatdan
silliqroq va eng sayoz hamda eng chuqur qiymatlar unda yoʻqoladi {[}60,
71{]}. Qoʻshni quduqlar farqi katta boʻlgani uchun bu hududda silliqlash
kuchli. Yozda 167 ta quduqdan 39 tasida (23 foiz) sath 1,5 m dan sayoz,
kriging xaritasida esa bunday maydon hududning 1,1 foizi: xarita sayoz
sathli maydonni oʻn baravardan ortiq kam koʻrsatadi. Bu farq quduqlar
sayoz joylarda zichroq boʻlgani uchun emas: hudud kataklarga boʻlinib
tenglashtirilganda ham sayoz quduqlar 20--24 foiz. Har bir quduq
qolganlari boʻyicha baholanganda kriging 167 tadan faqat bittasida
sathni 1,5 m dan sayoz deb topadi. Demak, kam koʻrsatish usulning
oʻzidan kelib chiqadi.

\textbf{Sayoz sath ehtimoli xaritasi.} Bu muammoni hal etish uchun
xaritaga sath emas, ehtimol chiqarildi: har bir katakda yozgi sathning
kritik chuqurlikdan (1,5 m) sayoz boʻlish ehtimoli kriging bahosi va
uning xatosi orqali hisoblandi:

\begin{equation}\tag{3.1}
P(z < z_{kr}) = \Phi\left(\frac{z_{kr} - \hat z}{\sigma_K}\right)
\end{equation}

bu yerda ẑ --- kriging bahosi, \(\sigma_{K}\) --- bahoning standart
xatosi, \(z_{\mathrm{kr}}\) = 1,5 m, Φ --- standart normal taqsimot
funksiyasi. Qiymatlar notekis taqsimlangani uchun bu ehtimol biroz
oshirib baholanadi, yaʼni natija xavfni kam koʻrsatmaydi. Natijalar
3.5-rasm va 29-ilovada keltirilgan.

\oakfig{figures/fig36_summer_prob.png}{3.5-rasm. Iyul--avgustdagi sizot suvi sathi, 2019--2024-yillar
oʻrtachasi: a --- kriging bahosi; b --- sathning 1,5 m dan sayoz boʻlish
ehtimoli (qora chiziq --- 0,3; massiv yozuvlarida --- ehtimollik 0,3 dan
katta maydon ulushi)}

Ehtimollik xaritasi boʻyicha Mashʼal, Z.M.Bobur va Istiqlol massivlari
ajralib turadi: bu massivlar maydonining 55--71 foizida yozgi sathning
kritik chuqurlikdan sayoz boʻlish ehtimoli 0,3 dan katta, 1,5 m dan
sayoz quduqlar ulushi 33--46 foiz. Tumandagi 1,5 m dan sayoz 39 ta
quduqning 31 tasi aynan shu uch massivga toʻgʻri keladi. Paxtakor,
U.Qoraboyev, T.Mirzayev va M.Toʻjiboyev massivlarida xavfli maydon
17--28 foiz, Y.Oxunboboyev massivida esa yozda ham sath 2,0 m dan
chuqur. Biroq bu massiv maydonining atigi 57 foizi xarita bilan
qoplangan va 3.1-§ ga koʻra uning maʼlumotlari eng kam ishonchli;
shuning uchun bu massiv boʻyicha xulosa shartli.

Massivlarni tartiblashda ehtimollik sath xaritasidan ishonchliroq.
Masalan, U.Qoraboyev massivida maydonning 45 foizida baholangan sath 2,0
m dan sayoz, lekin sathning 1,5 m dan sayoz boʻlish ehtimoli 0,3 dan
katta boʻlgan maydon atigi 19 foiz. Istiqlol massivida esa bu ikki
koʻrsatkich bir-biriga yaqin (61 va 55 foiz): U.Qoraboyevda sath
koʻpincha 1,5--2,0 m oraligʻida turadi, Istiqlolda esa kritik chegaradan
pastga tushadi. 3.2-§ dagi mavsumiy tahlil (U.Qoraboyevda yozgi minimum
1,45 m) bilan birga olinganda, U.Qoraboyev uchun xavf qisqa muddatli,
Istiqlol, Mashʼal va Z.M.Bobur uchun esa barqaror ekani koʻrinadi.

\textbf{Kritik chuqurlik va usul tanlovi.} 1,5 m chegarasi Oʻzbekiston
sharoiti uchun adabiyotdan olingan {[}55, 70{]}, biroq tuproq va
minerallashuvga qarab u 2,0--2,5 m gacha olinadi. Shu sababli hisob 2,0
va 2,5 m uchun ham bajarildi (29-ilova). Chegara 2,0 m boʻlganda yozda
sathi undan sayoz quduqlar 55 foiz, ehtimoli 0,3 dan katta maydon
tumanning 71 foizi; 2,5 m boʻlganda mos ravishda 78 va 86 foiz. Foizlar
chegaraga kuchli bogʻliq, massivlar tartibi esa yoʻq: krigingning boshqa
turlarida ham Mashʼal, Z.M.Bobur va Istiqlol birinchi uchta oʻrinda.
Ulardan biri (indikator kriging) xavfli maydonni ikki baravar kichik
koʻrsatadi (19 foizga qarshi 40 foiz), shuning uchun foizlar taxminiy,
massivlar tartibi esa ishonchli.

\textbf{Quduqlar sonining xarita aniqligiga taʼsiri.} Kuzatuv tarmogʻini
rejalashtirish (4.5-§) uchun asosiy savol: quduqlar koʻpaytirilsa,
xarita qanchalik aniqlashadi? Bunga quduqlarni sunʼiy kamaytirish
tajribasi bilan javob berildi. 168 ta quduqdan tasodifiy 34 tasi nazorat
uchun ajratildi, qolgan 134 tasidan esa 25, 40, 55, 70, 85 va 100 foizi
tanlanib, ular boʻyicha nazorat quduqlaridagi sath baholandi. Tajriba
har bir zichlik uchun 40 marta takrorlandi. Natijalar 30-ilovada
keltirilgan.

34 ta quduq (qoʻshni quduqlar orasi 1,2 km) boʻyicha xato 0,98 m boʻlib,
hududiy oʻrtacha bilan baholashdan farq qilmaydi. Quduqlar soni 54 taga
(980 m) yetganda xato 0,90 m gacha, 94 taga (800 m) yetganda 0,86 m
gacha kamayadi. Shundan keyin xato deyarli oʻzgarmaydi: 114 va 134 ta
quduqda 0,84--0,86 m. Demak, quduqlar orasini 800 m dan qisqartirish
xaritani aniqlashtirmaydi. Sababi --- qoʻshni quduqlar orasidagi katta
farq: 500 m dan kichik masofadagi oʻzgarishni ushlash uchun quduqlar
orasini bir necha baravar qisqartirish kerak boʻlardi.

Kuzatuv tarmogʻini xarita aniqligi uchun bir tekis zichlashtirish
samarasiz: mavjud quduqlar soni bilan hudud boʻyicha umumiy manzara
olinadi, undan keyingi har bir yangi quduq xaritani deyarli
yaxshilamaydi. Yangi quduqlarni aniq savolga javob beradigan joylarga
qoʻyish kerak: taʼmirlanadigan kollektorlar yaqiniga (taʼmir samarasini
oʻlchash uchun, 4.3-§), sayoz sath ehtimoli yuqori massivlarga va xato
eng katta chekka hududlarga (4.5-§).

\oakpar{3.4-§. Sizot suvi sathining drenaj tarmogʻi, kanallar va relyefga
bogʻliqligi}

Drenaj nazariyasiga koʻra kollektor va zovurlar sizot suvi sathini oʻz
tubi atrofidagi chuqurlikda ushlab turadi, ularning taʼsiri esa masofa
ortishi bilan susayadi {[}69, 78, 108{]}. Shu sababli toʻgʻri
ishlayotgan tarmoqda sath kollektorga yaqin joyda chuqurroq, oraliqlarda
esa sayozroq boʻlishi kutiladi {[}122{]}. Ushbu paragrafda bu kuzatuv
maʼlumotlari bilan tekshiriladi. 2.3-§ dagi GAT bazasi har bir quduq
uchun quyidagilarni hisoblash imkonini berdi: eng yaqin kollektorgacha
masofa, eng yaqin kanalgacha masofa, quduq atrofidagi 1 km radiusda
joylashgan kollektorlar uzunligi (tarmoq zichligi) va quduqning mutlaq
balandligi. Bundan tashqari, kadastr bilan bogʻlangan kollektorlar uchun
eng yaqin kollektorning ishga tushirilgan yili, drenaj moduli (kollektor
oxiridagi sarfning xizmat maydoniga nisbati) va solishtirma uzunligi
olindi. Tahlil 2019--2024-yillardagi oʻrtacha sathi maʼlum boʻlgan 168
ta quduq boʻyicha bajarildi.

\textbf{Kollektorgacha masofa.} Sath kollektorga yaqinlashgan sari
sayozlashadi (3.2-jadval). Kollektordan 150 m gacha boʻlgan quduqlarda
oʻrtacha sath 1,98 m, 800 m dan uzoqdagilarda 3,10 m. Spearman
korrelyatsiyasi ρ = +0,40 (95 foizlik ishonch oraligʻi 0,26--0,53).
Bogʻlanish qishda ham, yozda ham saqlanadi (ρ = +0,38 va +0,37). Yozda
sathi 1,5 m dan sayoz quduqlar ulushi kollektorga eng yaqin guruhda 38
foiz, eng uzoq guruhda esa nolga teng.

\oakjadvalnum{3.2-jadval}

\oakilovanom{\textbf{Sizot suvi sathi, uning mavsumiy va koʻp yillik oʻzgarishining
kollektorgacha masofaga bogʻliqligi}}

{\fontsize{9}{11}\selectfont\setlength{\tabcolsep}{2.5pt}

\noindent\begin{longtable}[]{|
  >{\raggedright\arraybackslash}p{(\columnwidth - 16\tabcolsep - 9\arrayrulewidth) * \real{0.1647}}|>{\raggedright\arraybackslash}p{(\columnwidth - 16\tabcolsep - 9\arrayrulewidth) * \real{0.0941}}|>{\raggedright\arraybackslash}p{(\columnwidth - 16\tabcolsep - 9\arrayrulewidth) * \real{0.0941}}|>{\raggedright\arraybackslash}p{(\columnwidth - 16\tabcolsep - 9\arrayrulewidth) * \real{0.1529}}|>{\raggedright\arraybackslash}p{(\columnwidth - 16\tabcolsep - 9\arrayrulewidth) * \real{0.1294}}|>{\raggedright\arraybackslash}p{(\columnwidth - 16\tabcolsep - 9\arrayrulewidth) * \real{0.1412}}|>{\raggedright\arraybackslash}p{(\columnwidth - 16\tabcolsep - 9\arrayrulewidth) * \real{0.1059}}|>{\raggedright\arraybackslash}p{(\columnwidth - 16\tabcolsep - 9\arrayrulewidth) * \real{0.1176}}|}
\hline
\begin{minipage}[b]{\linewidth}\raggedright
\textbf{Kollektorgacha masofa}
\end{minipage} & \begin{minipage}[b]{\linewidth}\raggedright
\textbf{Quduqlar soni}
\end{minipage} & \begin{minipage}[b]{\linewidth}\raggedright
\textbf{Yillik oʻrtacha sath, m}
\end{minipage} & \begin{minipage}[b]{\linewidth}\raggedright
\textbf{Yanvar--fevral sathi, m}
\end{minipage} & \begin{minipage}[b]{\linewidth}\raggedright
\textbf{Iyul--avgust sathi, m}
\end{minipage} & \begin{minipage}[b]{\linewidth}\raggedright
\textbf{Mavsumiy koʻtarilish, m}
\end{minipage} & \begin{minipage}[b]{\linewidth}\raggedright
\textbf{Yozda 1,5 m dan sayoz quduqlar, \%}
\end{minipage} & \begin{minipage}[b]{\linewidth}\raggedright
\textbf{Sen qiyaligi medianasi, sm/yil}
\end{minipage} \\
\hline
\endhead
\hline
\endlastfoot
0--150 m & 47 & 1,98 & 2,21 & 1,81 & 0,39 & 38 & +1,7 \\ \hline
150--400 m & 57 & 2,30 & 2,54 & 2,08 & 0,46 & 28 & +2,2 \\ \hline
400--800 m & 33 & 2,48 & 2,79 & 2,22 & 0,57 & 15 & +2,2 \\ \hline
800 m dan uzoq & 31 & 3,10 & 3,36 & 2,71 & 0,65 & 0 & +6,3 \\
\end{longtable}

}

Izoh: mavsumiy koʻrsatkichlar 2019--2024-yillar boʻyicha (oxirgi guruhda
30 ta quduq); Sen qiyaligi kamida 8 yil maʼlumotli 112 ta quduq
boʻyicha.

\textbf{Masofa zonalari xaritasi.} Masofa zonalari tumandagi barcha
kollektor va zovurlar atrofida chizildi (3.6-rasm). Tuman maydonining 18
foizi kollektordan 150 m gacha, 19 foizi 150--400 m, 15 foizi 400--800 m
oraligʻida, 47 foizi esa 800 m dan uzoqda joylashgan. Eng katta uzoq
zona janubdagi Y.Oxunboboyev massivi va tumanning togʻ oldi chekkalariga
toʻgʻri keladi. Kuzatuv quduqlari esa asosan tarmoq yaqinida: 168 ta
quduqdan 104 tasi (62 foiz) kollektordan 400 m gacha, 800 m dan uzoqda
atigi 31 tasi (18 foiz). Shu sababli monitoring tarmoq taʼsiridan
tashqaridagi hududni, yaʼni tuman maydonining qariyb yarmini, kam
qamraydi.

Xaritada 3.2-jadvaldagi qonuniyat yaqqol koʻrinadi. Yozda 1,5 m dan
sayoz 39 ta quduqning 31 tasi Z.M.Bobur, Mashʼal va Istiqlol
massivlarida, asosan 0--400 m zonalarda joylashgan (3.6-rasm, a). Sath
eng tez chuqurlashayotgan massivlar esa tarmoq siyrak boʻlgan janub va
janubi-sharqda: Sen qiyaligining medianasi Y.Oxunboboyevda yiliga 6,8
sm, M.Toʻjiboyevda 4,1 sm, T.Mirzayevda 3,7 sm (3.6-rasm, b).

\oakfig{figures/map3_collector_zones.png}{3.6-rasm. Kollektorlardan masofa zonalari: a --- yozgi (iyul--avgust)
sizot suvi sathi, 2019--2024-yillar; b --- sathning oʻzgarish tezligi
(Sen qiyaligi; toʻldirilgan belgilar --- ishonchli oʻzgarish, p
\textless{} 0,05)}

\textbf{Tarmoq zichligi.} Zichlik bilan bogʻlanish yanada kuchli: ρ =
−0,47. Regressiya boʻyicha quduq atrofidagi 1 km radiusda kollektor
uzunligi har 1 km ga ortganda sath oʻrtacha 0,18 m ga sayozroq boʻladi.
Zichlik yozgi sath bilan ham xuddi shunday bogʻlangan (ρ = −0,45).

\textbf{Bogʻlanish shakli.} Sath va kollektorgacha masofa orasidagi
bogʻlanish chiziqli emas (31-ilova, a). Ikki qismli toʻgʻri chiziq bilan
sinish nuqtasi 850 m da topildi: shu masofagacha sath har 100 m da
oʻrtacha 0,12 m chuqurlashadi, undan keyin deyarli oʻzgarmaydi (har 100
m da 0,02 m). Bu oddiy toʻgʻri chiziqdan yaxshiroq (p = 0,018), lekin
sinish nuqtasining ishonch oraligʻi keng (100--1 075 m), chunki 800 m
dan uzoqda quduqlar kam (31 ta). Shuning uchun 850 m kollektor taʼsir
qiladigan hududning taxminiy kengligi sifatida qaraladi.

\textbf{Kanallar va relyef.} Kanalgacha masofa bilan bogʻlanish kuchsiz
(ρ = −0,18) va boshqa omillar hisobga olinganda yoʻqoladi. Balandlik
bilan bogʻlanish ham kuchsiz (ρ = +0,23): balandlik sathdagi farqlarning
atigi 2 foizini tushuntiradi. Demak, tekis vodiyda relyef sizot
oqimining yoʻnalishini belgilaydi, lekin sath chuqurligini belgilamaydi.

\textbf{Uch omil birga.} Uch omil birga olinganda (32-ilova) faqat
tarmoq koʻrsatkichi ishonchli boʻlib qoladi. Ular birgalikda sathdagi
farqlarning 15--20 foizini tushuntiradi.

\textbf{Qoʻshni quduqlar taʼsirini tekshirish.} Yaqin quduqlar sathi
oʻxshash (3.3-§), shuning uchun ular toʻliq mustaqil kuzatuv emas va
oddiy hisob natijani ishonchliroq koʻrsatishi mumkin. Buni tekshirish
uchun hudud 2×2 km kataklarga boʻlinib, hisob kataklar boʻyicha
takrorlandi (33-ilova; Moran indeksi {[}94{]}). Sathning masofa va
zichlik bilan bogʻlanishi bu holda ham saqlanadi: ρ ning ishonch
oraligʻi +0,26 \ldots{} +0,52 va −0,58 \ldots{} −0,33. Massivlar ichida
U.Qoraboyev va M.Toʻjiboyevda bogʻlanish ishonchli qoladi, Istiqlol va
Z.M.Bobur massivlarida esa chegarada.

\textbf{Massivlar ichidagi bogʻlanish.} Topilgan bogʻlanish massivlar
orasidagi farqning natijasi boʻlishi mumkin edi: masalan, Y.Oxunboboyev
massivida quduqlar kollektordan oʻrtacha 3,2 km uzoqda va sath ham eng
chuqur. Buni tekshirish uchun bogʻlanish faqat bir massiv ichidagi
quduqlarni solishtirib baholandi {[}34{]} (33-ilova).

Bunda bogʻlanish atigi 9--10 foizga kuchsizlanadi va ishonchli qoladi,
yaʼni uni massivlar orasidagi emas, massiv ichidagi farq belgilaydi.
Kollektorgacha masofa bilan bogʻlanish toʻrtta massivda ishonchli:
U.Qoraboyev (ρ = +0,60), M.Toʻjiboyev (ρ = +0,50), Istiqlol (ρ = +0,40)
va Z.M.Bobur (ρ = +0,39). T.Mirzayev massivida yoʻnalish shu, lekin
quduqlar kam (ρ = +0,42; p = 0,13). Mashʼal va Paxtakor massivlarida esa
bogʻlanish yoʻq (ρ = −0,22 va −0,08). Mashʼal massivida sath 3.2-§ va
3.3-§ ga koʻra tumandagi eng sayozlar qatorida va quduqlar kollektorga
yaqin joylashgan (oʻrtacha 284 m), Paxtakor esa eng kichik massiv (19 ta
quduq). Bu massivlarda sathni kollektorlar joylashuvidan koʻra boshqa
mahalliy omillar belgilaydi.

\textbf{Mavsumiy koʻtarilish.} Agar kollektor sathning yozgi
koʻtarilishini cheklasa, kollektorga yaqin quduqlarda mavsumiy
koʻtarilish kichikroq boʻlishi kerak. Maʼlumotlarda shunday manzara
kuzatiladi: kollektorgacha 150 m boʻlgan quduqlarda sath
yanvar--fevraldan iyul--avgustgacha oʻrtacha 0,39 m ga, 800 m dan
uzoqdagilarda 0,65 m ga koʻtariladi (3.2-jadval va 31-ilova, b). Biroq
bu farqni kollektor taʼsiri deb talqin qilib boʻlmaydi. Mavsumiy
koʻtarilish birinchi navbatda qishki sathga bogʻliq: sathi chuqur
quduqda koʻtarilish uchun joy koʻp. Qishki sath va masofa birga
olinganda qishki sath har 1 m chuqurlashganda koʻtarilish 0,23 m ga
ortadi, masofaning taʼsiri esa yoʻqoladi. Bundan tashqari, kollektordan
uzoqdagi quduqlarda kesilgan oʻlchovlar ulushi ancha yuqori (9,9 foizga
qarshi 29,7 foiz), shuning uchun u yerda qishki sath haqiqatda yanada
chuqurroq va mavsumiy koʻtarilish kam baholangan. Mavjud maʼlumotlar
kollektorlarning sath koʻtarilishini cheklash taʼsirini alohida ajratib
koʻrsatishga imkon bermaydi.

\textbf{Kollektorlarning kadastr koʻrsatkichlari.} 87 ta quduq kadastrda
nomi bilan aniqlangan 41 ta kollektordan 500 m gacha masofada
joylashgan. Bu kollektorlar 1948--1985-yillarda ishga tushirilgan,
oʻrtacha yoshi 68 yil. Loyiha drenaj moduli asosan 0,40--1,54 l/(s·ga),
medianasi 0,83 l/(s·ga). Bu koʻrsatkichlarning hech biri quduqdagi sath
bilan bogʻlanmagan (kollektor yoshi, drenaj moduli va solishtirma
uzunlik bilan ρ = +0,03\ldots+0,11). Guruhlar boʻyicha oʻrtachalar ham
bir tomonlama oʻzgarmaydi (34-ilova).

Ular birga olinganda ham faqat kollektorgacha masofa ishonchli, uchala
koʻrsatkich esa sathdagi farqlarning atigi 5 foizini tushuntiradi.
Kadastr koʻrsatkichlari kollektorning hozirgi ish holatini aks
ettirmaydi: drenaj moduli loyihaviy qiymat boʻlib, kollektorning loyqa
bosishi, oʻsimlik bilan qoplanishi va chiqish qismidagi suv sathi
hisobga olinmagan. Kollektorning haqiqiy holatini baholash uchun
gidrometrik oʻlchovlar (oqim sarfi, suv sathi) va texnik koʻrik
maʼlumotlari zarur {[}40, 128{]}. Bu koʻrsatkichlarni monitoring
tizimiga kiritish taklifi 4.5-§ da beriladi.

\textbf{Koʻp yillik oʻzgarish va tarmoq.} Sathning koʻp yillik
oʻzgarishi ham kollektorgacha masofaga bogʻliq (3.2-jadval va 31-ilova,
c). Kollektordan 150 m gacha boʻlgan quduqlarda Sen qiyaligining
medianasi yiliga 1,7 sm, 800 m dan uzoqdagilarda esa yiliga 6,3 sm:
uzoqdagi quduqlarda sath qariyb toʻrt baravar tez chuqurlashgan. Farq
yiliga 4,7 sm, lekin uning ishonch oraligʻi keng (0,4--7,8 sm; p =
0,006), shuning uchun «toʻrt baravar» taxminiy. Bogʻlanish ishonchli
(masofa bilan ρ = +0,25, zichlik bilan ρ = −0,34) va massiv ichida
solishtirganda ham saqlanadi.

Bu bogʻliqlik 3.2-§ da aniqlangan chuqurlashishning sababi haqida
maʼlumot beradi. Agar chuqurlashish asosan 2016--2018-yillardagi
kollektorlar taʼmirining natijasi boʻlganida, u kollektorlarga yaqin
joylarda eng kuchli boʻlishi kerak edi. Kuzatuvda esa aksincha:
chuqurlashish kollektordan uzoqda eng kuchli. Bunga ikki izoh bor.
Birinchisi: chuqurlashishni tuman miqyosidagi omil, masalan sugʻorish
suvi hajmi yoki sugʻorish rejimining oʻzgarishi belgilaydi va bu omil
drenaj taʼsiridagi joylarda kamroq seziladi. Ikkinchisi: kollektor
yaqinida sath kollektor tubining chuqurligi bilan cheklangan va u yerda
sath kam oʻzgaradi, kollektordan uzoqda esa sath erkinroq oʻzgaradi.
Ikkala izoh ham bir xulosaga olib keladi: tumandagi sath
chuqurlashishini toʻgʻridan-toʻgʻri drenaj tarmogʻi taʼmiriga bogʻlab
boʻlmaydi. Taʼmir samarasi 4.3-§ da alohida, taʼmirlangan va
taʼmirlanmagan kollektorlar yonidagi quduqlarni solishtirib baholanadi.

\textbf{Natijaning talqini.} Sath va tarmoq orasidagi bogʻlanish
drenajdan kutilganiga teskari: drenaj sathni pasaytirishi kerak,
kuzatuvda esa tarmoq zich joyda sath sayoz. Buni «kollektor sathni
koʻtaradi» deb talqin qilish notoʻgʻri. Eng ehtimoliy izoh teskari
bogʻliqlikdir: kollektorlar sizot suvi tabiatan sayoz boʻlgan, zax
bosish xavfi yuqori joylarga qurilgan {[}51, 122{]}. Massiv ichidagi
solishtiruv natijasi bu izohga zid emas, chunki kollektorlar odatda
massiv ichida ham eng past va zax joylar boʻylab loyihalanadi.
Qurilishdan oldingi sath boʻyicha maʼlumot yoʻqligi sababli bu izohni
bevosita tekshirib boʻlmaydi; shunga qaramay, tarmoq joylashuvini
sathning sababi emas, koʻproq oqibati deb qarash asosliroq.

Bir yildagi turli joylarni solishtirib drenaj samarasini baholab
boʻlmaydi: tarmogʻi zich hududda sath sayoz boʻlishi tarmoq yomon
ishlayotganini bildirmaydi. Massivlar oʻrtachasini solishtirishga
tayangan baholar shu sababli notoʻgʻri xulosaga olib kelishi mumkin.
Kadastrdagi loyihaviy koʻrsatkichlar (yosh, drenaj moduli) kollektorning
hozirgi samaradorligini baholash uchun yaroqsiz. Drenaj taʼsirini
aniqlash uchun vaqt boʻyicha taqqoslash zarur: bir xil quduqda tarmoq
holati oʻzgarishidan oldingi va keyingi sath. Bunday tahlil 4.3-§ da
bajariladi.

Koʻrib chiqilgan omillar sathdagi farqlarning 67--85 foizini
tushuntirmaydi. Bu 3.3-§ dagi xulosa bilan mos: sath qisqa masofada
mahalliy omillar (sugʻorish rejimi, tuproqning suv oʻtkazuvchanligi,
ichki zovurlar holati) taʼsirida oʻzgaradi va bu omillar boʻyicha quduq
darajasidagi maʼlumot mavjud emas.

\oakpar{3.5-§. Sizot suvlari minerallashuvi va shoʻrlanish xavfi}

Sizot suvining xavfliligi faqat uning sathi bilan emas, minerallashuvi
bilan ham belgilanadi: bir xil chuqurlikda minerallashgan suv chuchuk
suvga qaraganda tuproqqa koʻp tuz olib chiqadi {[}57, 70, 107{]}.
Shuning uchun meliorativ holatni baholashda sath va minerallashuv
birgalikda qaraladi. Ushbu paragrafda minerallashuv maʼlumotlari 3.1-§
dagi tartibda tekshiriladi, ishonchli davr boʻyicha uning hududlar va
yillar boʻyicha oʻzgarishi tahlil qilinadi va sath bilan birga
shoʻrlanish xavfi baholanadi.

\textbf{Maʼlumotlar va ularning tekshiruvi.} Minerallashuv boʻyicha ikki
manba mavjud: 1999--2015-yillar uchun yiliga uch marta (1-aprel, 1-iyul,
1-oktabr holatiga) oʻlchangan qiymatlar jadvali (123 ta quduq, 4 961 ta
qiymat) va 2009--2024-yillar uchun yillik qiymatlar. Har ikkisi 3.1-§
dagi qoidalar boʻyicha tekshirildi. Natijalar 35-ilovada keltirilgan.

Ikki manba bir-biri bilan ham solishtirildi. 2009-yilda ikki manba
yaxshi mos keladi (ρ = 0,92; 105 ta quduq), 2010-yilda esa kuchsiz (ρ =
0,37; 98 ta quduq), yaʼni yillik manba kamroq ishonchli. Natijada 26
yillik maʼlumotdan faqat 2004--2010-yillar (7 yil, 122 ta quduq, 2 136
ta qiymat) ishonchli deb topildi va keyingi tahlil shu davr boʻyicha
bajarildi.

\textbf{Minerallashuv darajasi.} 2004--2010-yillarda oʻlchangan
qiymatlarning 30,7 foizi 1 g/l dan kam (chuchuk suv), 60,4 foizi 1--3
g/l (kam minerallashgan), 8,8 foizi 3--5 g/l (oʻrtacha minerallashgan),
atigi 0,2 foizi 5 g/l dan yuqori. Yillar boʻyicha taqsimot 36-ilovada
keltirilgan. Sugʻorish suvi sifati mezonlari boʻyicha {[}35{]} 1--3 g/l
minerallashuv sizot suvi kapillyar koʻtarilish orqali tuproqqa yetib
kelganda oʻrtacha shoʻrlanish xavfini, 3 g/l dan yuqori minerallashuv
esa jiddiy xavfni bildiradi.

\textbf{Vaqt boʻyicha oʻzgarish.} Tuman boʻyicha yillik mediana
1,02--1,45 g/l oraligʻida oʻzgargan. 2004--2009-yillarda pasaygan (1,38
dan 1,02 g/l gacha), 2010-yilda esa 1,35 g/l ga qaytgan; tuman boʻyicha
ishonchli oʻzgarish yoʻq (p = 0,13). Quduqlar boʻyicha esa pasayish
ustun: kamida 12 ta qiymatli 102 ta quduqdan 76 tasida minerallashuv
kamaygan, 24 tasida oshgan; ishonchli pasayish 9 ta, ishonchli oʻsish 1
ta quduqda. Oʻzgarish tezligining medianasi yiliga −0,05 g/l. Bitta
quduqda ham minerallashuv yildan-yilga ancha farq qiladi (oʻrtacha 39
foiz). Bunday oʻzgaruvchanlik sinamalarni olish va tahlil qilish
xatoligi bilan ham bogʻliq boʻlishi mumkin, shuning uchun alohida yil
boʻyicha xulosa chiqarishda ehtiyot boʻlish kerak.

\textbf{Mavsumiy farq.} Tuman boʻyicha medianalar aprelda 1,27, iyulda
1,25, oktabrda 1,32 g/l. Bir xil quduq va yil boʻyicha solishtirilganda
(675 juft) iyuldagi qiymat aprelnikidan atigi 0,04 g/l, oktabrdagisi
iyulnikidan 0,05 g/l yuqori (p = 0,29 va 0,28). Demak, minerallashuvda
sezilarli mavsumiy oʻzgarish yoʻq. Sathning mavsumiy yoʻli (3.2-§) bilan
solishtirilganda farq aniq koʻrinadi: sath yozda 0,5 m atrofida
koʻtarilsa ham, sugʻorish suvi kelishi minerallashuvni sezilarli
suyultirmaydi. Bundan sizot suvining tuz zaxirasi katta va sugʻorish
mavsumida u yangilanmasligi kelib chiqadi.

\textbf{Massivlar boʻyicha farq.} Massivlar orasidagi farq katta
(37-ilova va 3.7-rasm). Quduqlar boʻyicha minerallashuv medianasi
Mashʼal massivida 2,47 g/l, U.Qoraboyev va T.Mirzayev massivlarida
0,80--0,84 g/l: eng ogʻir massiv bilan eng yengili orasida uch baravar
farq bor. Massivlar orasidagi farq ishonchli va katta (Kruskal--Wallis
mezoni, p \textless{} 0,001). Chuchuk sizot suvi tumanning sharqiy va
janubiy qismlarida (U.Qoraboyev, M.Toʻjiboyev, T.Mirzayev),
minerallashgan suv esa gʻarbiy va shimoliy qismlarida (Mashʼal,
Z.M.Bobur) tarqalgan. 3 g/l dan yuqori minerallashuvli toʻrtta quduqning
ikkitasi Mashʼal, ikkitasi Z.M.Bobur massivida joylashgan. Y.Oxunboboyev
massivi boʻyicha 2004--2010-yillarda minerallashuv oʻlchovlari mavjud
emas.

\oakfig{figures/fig310_mineralization.png}{3.7-rasm. Sizot suvlari minerallashuvi, 2004--2010-yillar: a ---
quduqlar boʻyicha medianasi (113 ta quduq); b --- massivlar boʻyicha
medianasi (37-ilova)}

\textbf{Minerallashuvning sath va tarmoq bilan bogʻliqligi.} Quduqlar
boʻyicha minerallashuv sath bilan kuchsiz bogʻlangan (ρ =
−0,12\ldots−0,16; ishonchli emas): sath sayoz joyda minerallashuv biroz
yuqori. Kollektorgacha masofa, tarmoq zichligi va kanalgacha masofa
bilan bogʻlanish umuman yoʻq. Minerallashuvning hududiy taqsimotini
massivlar miqyosidagi gidrogeologik sharoit (sizot oqimining yoʻnalishi,
suvli qatlam tarkibi, sizot suvining hududda turib qolish vaqti)
belgilaydi {[}57, 70{]}. Tumanning gʻarbiy qismida minerallashuv
yuqoriligi 3.3-§ dagi sathning umumiy yoʻnalishi bilan mos: sizot suvi
shimoli-gʻarbiy qismda sayozroq va u yerda bugʻlanish orqali tuz
toʻplanishi kuchliroq.

\textbf{Shoʻrlanish xavfi.} Sath va minerallashuv birgalikda baholash
uchun quduqlar ikki belgi boʻyicha guruhlandi: yozgi sath
(2019--2024-yillar, iyul--avgust) va minerallashuv medianasi
(2004--2010-yillar). Natijalar 38-ilovada keltirilgan.

Yozgi sath 2,0 m dan sayoz va minerallashuv 1 g/l dan yuqori boʻlgan
quduqlar 56 ta (50 foiz); shundan 29 tasida (26 foiz) yozgi sath 1,5 m
dan ham sayoz. Bu quduqlar shoʻrlanish xavfi eng yuqori nuqtalar
hisoblanadi. Ular massivlar boʻyicha notekis taqsimlangan (37-ilova):
Mashʼal va Z.M.Bobur massivlarida 17 tadan, Istiqlol massivida 12 ta,
qolgan toʻrtta massivda jami 10 ta; xavfli quduqlarning 82 foizi uchta
massivga toʻgʻri keladi. Bu 3.3-§ dagi ehtimollik xaritasiga mos: oʻsha
uchta massivda yozda sath kritik chuqurlikdan sayoz boʻlish ehtimoli eng
yuqori. Sath va minerallashuv boʻyicha alohida baholar bir xil
massivlarni ajratadi.

Xavfli quduqlarning joylashuvi va massivlar boʻyicha ulushi 3.8-rasmda
koʻrsatilgan. Xaritada sath va minerallashuv bitta rang shkalasida birga
berilgan, xavf guruhi belgi shakli bilan ajratilgan. Xavfli quduqlar
ulushi Mashʼal massivida 81 foiz (21 tadan 17 tasi), Z.M.Bobur massivida
74 foiz, Istiqlolda 63 foiz, Paxtakor va M.Toʻjiboyev massivlarida
27--29 foiz, U.Qoraboyev va T.Mirzayev massivlarida esa 8--11 foiz.
Y.Oxunboboyev massivi boʻyicha 2004--2010-yillarda minerallashuv
oʻlchanmagan.

\oakfig{figures/map4_salinity_risk.png}{3.8-rasm. Shoʻrlanish xavfi: a --- quduqlar boʻyicha yozgi sath
(2019--2024) va minerallashuv (2004--2010) birgalikda; b --- massivlar
boʻyicha xavfli quduqlar ulushi (yozgi sath \textless{} 2,0 m va
minerallashuv ≥ 1 g/l)}

Bu bahoning cheklovi shundaki, sath va minerallashuv turli davrlarga
tegishli. 2010-yildan keyingi minerallashuv oʻzgarishi nomaʼlum va bu
davrda minerallashuv oʻzgargan boʻlishi mumkin. Shu sababli xavf bahosi
minerallashuvning nisbiy darajasiga (massivlar orasidagi farqqa)
tayanadi; massivlar orasidagi farq gidrogeologik sharoit bilan
belgilangani uchun vaqt oʻtishi bilan saqlanib qolgan boʻlishi ehtimoli
yuqori.

\textbf{Monitoring tizimi uchun xulosa.} 2010-yildan keyingi
minerallashuv oʻzgarishi haqida ishonchli xulosa chiqarib boʻlmaydi. Bu
tadqiqotning cheklovi sifatida qayd etiladi; ayni paytda u monitoring
tizimining jiddiy kamchiligini koʻrsatadi: soʻnggi 14 yil davomida tuman
boʻyicha sizot suvlari minerallashuvining haqiqiy holati
hujjatlashtirilmagan. Minerallashuvni oʻlchashni tiklash uchun barcha
quduqlarda oʻlchash shart emas: bir quduq ichidagi mavsumiy farq
ahamiyatsiz boʻlgani uchun yiliga bir marta oʻlchash yetarli, xavfli deb
ajratilgan 56 ta quduq esa birinchi navbatda qamrab olinishi kerak.
Elektr oʻtkazuvchanlikni dalada oʻlchaydigan asboblar bu ishni
arzonlashtiradi {[}106{]}. Bu takliflar 4.5-§ da monitoring tizimini
takomillashtirish tavsiyalariga kiritiladi.

\oakpar{3.6-§. Marhamat tumani natijalarini Andijon viloyati tumanlari bilan
solishtirish}

Oldingi natijalar bitta tumanga tegishli. Ularni umumlashtirish uchun
ikki savolga javob kerak: Marhamatdagi maʼlumot sifati viloyat uchun
odatiymi yoki istisnomi, va «sath chuqur joylarda maʼlumot ishonchsiz»
degan qonuniyat (3.1-§) boshqa tumanlarda ham bormi. Buning uchun
meliorativ ekspeditsiyaning Andijon viloyatidagi 14 ta tuman boʻyicha
2009--2023-yillardagi kuzatuv jadvallari tahlil qilindi. Jadvallar
Marhamat tumani jadvallari bilan bir xil shaklda: har bir quduq uchun
yiliga 36 ta oʻn kunlik oʻlchov. Jami 1 972 ta quduq va 854 676 ta
oʻlchov mavjud (massivlar va tumanlar boʻyicha jami qatorlari
chiqarilgan).

Barcha tumanlarga 3.1-§ dagi tekshiruv qoidalari bir xil tartibda
qoʻllanildi. Har bir tumanda 1--3 yil boʻyicha mustaqil kuzatuv
maʼlumotlari mavjud emas, bu yillar tahlildan chiqarildi. Qolgan yillar
boʻyicha yil ichidagi oʻzgarmas qatorlar (7 tadan ortiq oʻlchovi boʻlgan
va barcha oʻlchovlari bir xil qiymatga ega quduq-yillar) va oʻlchov
birligi xatolari aniqlandi. Qoʻshimcha koʻrsatkich sifatida bir xonali
yaxlit qiymatlar (masalan, 2,0; 2,5; 3,0 m) ulushi hisoblandi: santimetr
aniqligida oʻlchanadigan koʻrsatkichda bunday qiymatlarning koʻpligi
oʻlchov taxminiy yozilganini bildirishi mumkin. Natijalar 39-ilovada
keltirilgan. Marhamat tumani boʻyicha bu yerda olingan koʻrsatkichlar
3.1-§ dagilar bilan mos keladi (yaroqli oʻlchovlar ulushi 88,2 va 87,7
foiz), farq faqat davr chegarasi (2024-yil viloyat jadvalida yoʻq) bilan
bogʻliq.

\textbf{Maʼlumotlar sifatining tumanlar orasidagi farqi.} Tumanlar
maʼlumotlar sifati boʻyicha katta farq qiladi. Oltinkoʻl, Boʻston,
Shahrixon, Ulugʻnor, Buloqboshi va Baliqchi tumanlarida oʻzgarmas
qatorlar 5 foizdan kam, yaroqli oʻlchovlar 96--100 foiz. Asaka,
Paxtaobod, Qoʻrgʻontepa va Jalaquduq tumanlarida esa quduq-yillarning
22--29 foizida yil davomida bitta qiymat yozilgan. Marhamat tumani 14 ta
tuman ichida 9-oʻrinda, viloyat medianasidan (10,6 foiz oʻzgarmas qator)
biroz yomonroq. Marhamat tumanidagi maʼlumot nuqsonlari istisno emas:
viloyatning toʻrtdan bir qismidan koʻprogʻida ular yanada jiddiyroq.
Yaxlit qiymatlar ulushi oʻzgarmas qiymatlar ulushi bilan bogʻlangan (ρ =
+0,65): sifati past tumanlarda oʻlchovlar ham taxminiyroq yozilgan.

\textbf{Ishonchlilik va sath chuqurligi.} 3.1-§ da Marhamat tumani
massivlari boʻyicha aniqlangan qonuniyat viloyat miqyosida ham toʻliq
tasdiqlanadi. Oʻzgarmas qiymatlar ulushi tumandagi sath chuqurligi bilan
kuchli bogʻlangan (ρ = +0,85; 95 foizlik oraliq 0,48--0,99), yaroqli
oʻlchovlar ulushi esa teskari bogʻlangan (ρ = −0,82). Sathi sayoz
tumanlarda (Boʻston, Oltinkoʻl, Baliqchi; 2019--2023-yillardagi mediana
1,6--1,8 m) maʼlumotlar deyarli nuqsonsiz, sathi chuqur tumanlarda
(Asaka, Paxtaobod; 2,7--2,9 m) esa eng koʻp nuqsonli. Bu 3.1-§ dagi
izohni tasdiqlaydi: oʻzgarmas qatorlar asosan sath quduq tubidan pastga
tushganda quduqning qurib qolishi va kuzatuvchining haqiqiy sath oʻrniga
shartli qiymat yozishi. Kuzatuv quduqlarining chuqurligi sath chuqur
hududlar uchun yetarli emas va bu muammo bitta tuman emas, butun viloyat
monitoring tarmogʻiga xos.

Maʼlumotlar nuqsoni tasodifiy emas: u aynan sath chuqur joylarda va
yillarda toʻplanadi. Shuning uchun nuqsonli qiymatlarni oddiygina
chiqarib tashlash sathni sayozroq koʻrsatadi, ularni tekshiruvsiz
qoldirish esa shartli qiymatlarni haqiqiy oʻlchov sifatida qabul
qilishga olib keladi. Ikkala holda ham xato bir tomonga yoʻnalgan. Bu
muammoni faqat quduqlarni chuqurlashtirish yoki kesilgan qiymatlarni
maxsus usullar bilan hisobga olish orqali hal etish mumkin {[}65{]}.

\textbf{Sizot suvi koʻrsatkichlari.} Har bir tuman uchun 3.2-§ dagi
tartibda doimiy kuzatilgan quduqlar tanlandi va sathning koʻp yillik
oʻzgarishi, 2019--2023-yillardagi yozgi sath va mavsumiy tebranish
hisoblandi (40-ilova va 3.9-rasm).

\oakfig{figures/fig312_region.png}{3.9-rasm. Andijon viloyati tumanlari: a --- yaroqli oʻlchovlar ulushi; b
--- doimiy kuzatilgan quduqlar boʻyicha sathning oʻzgarish tezligi va
yozgi sath medianasi (qalin chiziq --- Marhamat tumani)}

Sath chuqurlashishi butun viloyatga xos: 14 ta tumanning 12 tasida
ishonchli chuqurlashish aniqlandi, qolgan ikkitasida (Oltinkoʻl,
Boʻston) ham yoʻnalish shu, lekin ishonchli emas. Hech bir tumanda sath
koʻtarilmagan. Tumanlar boʻyicha chuqurlashish tezligining medianasi
yiliga 3,3 sm; Marhamat tumani (yiliga 3,3 sm) aynan viloyat medianasida
turadi. Bu 3.2-§ va 3.4-§ dagi xulosani kuchaytiradi: Marhamat
tumanidagi chuqurlashish tumanning oʻz drenaj tarmogʻidagi
oʻzgarishlardan koʻra viloyat miqyosidagi omillar, xususan suv taʼminoti
va sugʻorish rejimi bilan bogʻliq boʻlishi ehtimoli yuqori. Biroq bu
qiyos taʼmir taʼsirini toʻliq istisno etmaydi, chunki boshqa tumanlarda
ham shu yillarda kollektorlar taʼmiri olib borilgan boʻlishi mumkin va
bu haqda viloyat boʻyicha maʼlumot tahlil qilinmadi; taʼmir samarasi
4.3-§ da alohida baholanadi.

Chuqurlashish tezligi tumandagi yozgi sath bilan bogʻlangan (ρ = +0,70;
95 foizlik oraliq 0,21--0,91): sathi chuqur Paxtaobod (yiliga 10,3 sm)
va Asaka (6,4 sm) tumanlarida sath eng tez chuqurlashmoqda, sathi sayoz
Boʻston, Oltinkoʻl va Ulugʻnor tumanlarida deyarli oʻzgarmaydi; quduqlar
darajasida ham (3.2-§) chuqur quduqlarda chuqurlashish tezroq. Biroq
tezlik maʼlumot sifati bilan ham bogʻlangan (ρ = −0,74): eng tez
chuqurlashayotgan tumanlarda hisob eng kam va eng ishonchsiz quduqlarga
(25--57 ta) tayanadi. Maʼlumot sifati hisobga olinganda tezlik va sath
chuqurligi bogʻlanishi yoʻqoladi (r = 0,19; p = 0,52). Shuning uchun
tumanlar boʻyicha tezliklarni ehtiyotkorlik bilan talqin qilish kerak.

Marhamat tumani yozgi sathning medianasi boʻyicha viloyatda oʻrtacha
oʻrinda (1,92 m; 7-oʻrin), lekin sathi 1,5 m dan sayoz quduqlar ulushi
boʻyicha Boʻston tumanidan keyin ikkinchi (22,9 foiz; viloyat jadvallari
boʻyicha 2019--2023-yillar). Mavsumiy tebranish boʻyicha ham Marhamat
Xoʻjaobod tumanidan keyin ikkinchi (0,48 m). Marhamatning oʻziga xosligi
--- sath hududlar va mavsumlar boʻyicha kuchli oʻzgaradi: tumanda bir
vaqtning oʻzida sathi chuqur va sayoz hududlar mavjud, yozda esa sath
boshqa tumanlarga qaraganda koʻproq koʻtariladi. Bu 3.2--3.3-§ larda
massivlar boʻyicha aniqlangan farqlar bilan mos va tadqiqotda massivlar
darajasidagi yondashuv tanlanganini asoslaydi.

\oakpar{III bob boʻyicha xulosalar}

\begin{enumerate}
\def\labelenumi{\arabic{enumi}.}
\tightlist
\item
  Marhamat tumani boʻyicha 2009--2024-yillardagi sizot suvi sathi
  kuzatuvlari (13 yil, 188 ta quduq, 73 868 ta oʻlchov) ishlab chiqilgan
  tartib boʻyicha tekshirildi. 270 ta quduq-yilda yil davomida bitta
  qiymat yozilgani, 2018-yilda kuzatuv tarmogʻi qayta tuzilgani (60 ta
  yangi quduq) va yaroqli oʻlchovlarning 8,2 foizi quduq qurigan holatni
  aks ettiruvchi kesilgan qiymat ekani aniqlandi. Oʻlchovlarning 87,7
  foizi tahlil uchun yaroqli, haqiqiy sathni toʻliq aks ettiruvchi
  oʻlchovlar esa umumiy hajmning taxminan 80 foizi.
\item
  Maʼlumotlar nuqsoni tasodifiy emas: ishonchlilik sath chuqurligiga
  bogʻliq. Sathi sayoz Mashʼal va Istiqlol massivlarida yaroqli
  oʻlchovlar 97--98 foiz, sathi chuqur T.Mirzayev va Y.Oxunboboyev
  massivlarida 50--60 foiz. Bu qonuniyat Andijon viloyatining 14 ta
  tumani miqyosida ham tasdiqlandi: oʻzgarmas qiymatlar ulushi tumandagi
  sath chuqurligi bilan kuchli bogʻlangan (ρ = +0,85). Uning sababi
  kuzatuv quduqlarining chuqurligi sath chuqur hududlar uchun yetarli
  emasligi.
\item
  Doimiy kuzatilgan 84 ta quduq boʻyicha sizot suvi sathi 2009-yildagi
  1,60 m dan 2024-yilda 1,99 m gacha chuqurlashgan, tezlik yiliga
  0,028--0,033 m. Bu natija quduq qurigan holatlarga va hisoblash
  usuliga bogʻliq emas. Oʻzgarish asosan 2016-yildan keyin pogʻona
  koʻrinishida yuz bergan: 2009--2016-yillarda oʻrtacha sath 1,63 m,
  2018--2024-yillarda 1,92 m.
\item
  Tekshiruvsiz maʼlumot boʻyicha oʻzgarish tezligi (yiliga 0,045 m)
  toʻgʻri bahodan (yiliga 0,032 m) 1,4 baravar katta, oʻrtacha sathning
  oʻzi esa 0,29--0,53 m chuqur. Demak, maʼlumotlarni tekshirmasdan
  oʻzgarishni toʻgʻri baholab boʻlmaydi.
\item
  Tuman oʻrtachasi massivlardagi qarama-qarshi oʻzgarishlarni yashiradi.
  Oltita massivda sath yiliga 0,031--0,053 m chuqurlashgan, Z.M.Bobur
  massivida esa yiliga 0,031 m koʻtarilgan va 2019-yildan keyin
  1,27--1,71 m oraligʻiga yetgan. Kamida 8 yil maʼlumotli 122 ta
  quduqdan 36 tasida sath sayozlashgan, ularning asosiy qismi Z.M.Bobur
  massivida joylashgan.
\item
  Sathning yil ichidagi tebranishi tuman boʻyicha 0,48 m, massivlar
  boʻyicha 0,32 m dan (T.Mirzayev) 0,85 m gacha (U.Qoraboyev). Eng
  chuqur holat yanvar--fevralda, eng sayoz holat may--iyulda kuzatiladi.
  Yogʻinning asosiy qismi noyabr--aprelda boʻlsa-da, sath yozda
  koʻtariladi: demak, uni asosan sugʻorish va kanallardan sizib chiqqan
  suv belgilaydi. Yillik oʻrtacha sath ayrim massivlardagi yozgi xavfli
  holatni yashiradi, shuning uchun drenaj tizimini baholashda yozgi eng
  sayoz holat asosiy mezon boʻlishi kerak.
\item
  Sath qoʻshni quduqlar orasida ham keskin farq qiladi: eng yaqin quduq
  qiymati sathni hudud oʻrtachasidan ham yomonroq baholaydi. Kriging
  xaritasining xatosi 0,82 m, hudud oʻrtachasidan atigi 17 foiz kichik;
  umumiy yoʻnalish yoki drenaj tarmogʻi koʻrsatkichlarini qoʻshish
  xaritani yaxshilamaydi.
\item
  Kriging xaritasi sayoz sathli maydonni oʻn baravardan ortiq kam
  koʻrsatadi: yozda quduqlarning 23 foizida sath 1,5 m dan sayoz,
  xaritada esa bunday maydon 1,1 foiz. Shu sababli meliorativ holatni
  baholash uchun sathning kritik chuqurlikdan sayoz boʻlish ehtimoli
  xaritasi taklif etildi. Unga koʻra Mashʼal, Z.M.Bobur va Istiqlol
  massivlari maydonining 55--71 foizida yozgi sath 1,5 m dan sayoz
  boʻlish ehtimoli 0,3 dan yuqori.
\item
  Quduqlarni sunʼiy kamaytirish tajribasiga koʻra, quduqlar soni 94
  tadan (quduqlar orasi taxminan 800 m) oshirilganda xarita aniqligi
  oshmaydi. Kuzatuv tarmogʻini bir tekis zichlashtirish samarasiz; yangi
  quduqlar muayyan meliorativ savolga javob beradigan joylarga
  qoʻyilishi kerak.
\item
  Sizot suvi kollektorlarga yaqin va tarmoq zich joylarda sayozroq (ρ =
  +0,40 va −0,47); bogʻlanish taxminan 850 m gacha kuzatiladi va
  massivlar ichida ham saqlanadi. Buning sababi, ehtimol, tarmoq
  tabiatan sayoz sathli joylarga qurilgani; shuning uchun bir yildagi
  turli joylarni solishtirib drenaj samarasini baholab boʻlmaydi.
  Kollektorlarning kadastrdagi koʻrsatkichlari (yoshi, drenaj moduli,
  solishtirma uzunligi) sath bilan bogʻlanmagan; relyef va kanalgacha
  masofa sath oʻzgaruvchanligini deyarli tushuntirmaydi.
\item
  Sath kollektordan 800 m dan uzoqdagi quduqlarda kollektor yaqinidagiga
  qaraganda tezroq chuqurlashgan (medianalar yiliga 6,3 va 1,7 sm;
  farqning 95 foizlik ishonch oraligʻi 0,4--7,8 sm). Viloyatda ham sath
  14 ta tumanning 12 tasida ishonchli chuqurlashgan va Marhamat
  tumanidagi tezlik (yiliga 3,3 sm) viloyat medianasiga teng. Bu
  natijalar tumandagi chuqurlashishni drenaj tarmogʻi taʼmirining
  natijasi deb hisoblashga asos bermaydi; taʼmir samarasi alohida usul
  bilan baholanishi kerak.
\item
  Minerallashuv boʻyicha 26 yillik maʼlumotdan faqat 2004--2010-yillar
  ishonchli. Bu davrda qiymatlarning 60 foizi 1--3 g/l, ishonchli koʻp
  yillik va mavsumiy oʻzgarish aniqlanmagan; massivlar orasidagi farq
  uch baravarga yetadi (Mashʼalda 2,47 g/l, T.Mirzayev va U.Qoraboyevda
  0,80--0,84 g/l). 2010-yildan keyingi minerallashuv dinamikasi
  hujjatlashtirilmagan.
\item
  2019--2024-yillardagi yozgi sath va 2004--2010-yillardagi
  minerallashuv birgalikda baholanganda quduqlarning yarmi (113 tadan 56
  tasi) shoʻrlanish xavfi ostida (yozgi sath 2,0 m dan sayoz,
  minerallashuv 1 g/l dan yuqori), ularning 82 foizi Mashʼal, Z.M.Bobur
  va Istiqlol massivlarida joylashgan. Ehtimollik xaritasi va
  shoʻrlanish xavfi bir xil ustuvor massivlarni ajratadi; bu massivlar
  IV bobda drenaj tarmogʻini birinchi navbatda tekshirish va kuzatuv
  tarmogʻini takomillashtirish hududi sifatida qaraladi.
\end{enumerate}

\oakbob{IV BOB. MELIORATIV TIZIM ELEMENTLARINING ISHONCHLILIGI VA TAʼMIR
SAMARADORLIGI}

Ushbu bobda kadastr boʻyicha tizim elementlarining ishonchliligi
baholanadi, kollektorlarning suv oʻtkazishi tahlil qilinadi, taʼmirning
sizot sathiga taʼsiri va narxi oʻrganiladi hamda monitoringni yaxshilash
boʻyicha tavsiyalar beriladi.

\oakpar{4.1-§. Kollektor va zovurlarning yoshi, taʼmir tarixi va ishonchlilik
koʻrsatkichlari}

Texnik obyektning ishonchliligi --- uning oʻz vazifasini belgilangan
sharoitda va muddat davomida bajara olishi {[}8{]}. Gidromeliorativ
inshootlar uchun u toʻrt qismdan iborat: buzilmaslik, uzoq muddatlilik,
taʼmirboplik va saqlanuvchanlik {[}23, 24{]}. Har biri oʻz
koʻrsatkichlari bilan baholanadi va ular uchun maʼlum maʼlumotlar kerak.
Shuning uchun avval mavjud maʼlumotlar qaysi koʻrsatkichlarni
hisoblashga yetishi aniqlandi (41-ilova).

Kadastr maʼlumotlari ishonchlilikning odatdagi koʻrsatkichlarini
(buzilmasdan ishlash ehtimoli, buzilishlar soni) hisoblashga yetmaydi.
Shuning uchun ikki oʻrinbosar koʻrsatkich olindi: elementning yoshi va
oxirgi taʼmirdan beri oʻtgan muddat. Tizim oʻz vazifasini qanchalik
bajarayotgani III bobda sizot suvi sathi boʻyicha baholangan. Barcha
hisoblar 2021-yil 1-yanvar holatidagi meliorativ kadastr maʼlumotlari
asosida bajarildi.

\textbf{Kollektorlar.} Tumanda davlat hisobidagi 90 ta kollektor bor,
umumiy uzunligi 265,8 km, xizmat koʻrsatadigan maydoni 11 079 ga.
Ularning barchasi Meliorativ ekspeditsiya balansida. Kollektorlarning
yoshi 36 yildan 103 yilgacha, medianasi 61 yil. Kollektorlarning 82
foizi (uzunlik boʻyicha 90 foizi) 50 yildan eski (42-ilova). Eng koʻp
qurilish 1950--1960-yillarga toʻgʻri keladi: shu ikki oʻn yillikda 61 ta
kollektor (159,3 km) foydalanishga topshirilgan. 1980-yillardan keyin
tumanda yangi kollektor qurilmagan.

Shu bilan birga kollektorlarning barchasi 2010--2020-yillar oraligʻida
kamida bir marta taʼmirlangan. Oxirgi taʼmirdan beri oʻtgan muddatning
medianasi 4 yil; uzunlikning 79,5 foizi soʻnggi besh yilda taʼmirlangan.
Taʼmirlar vaqt boʻyicha notekis: 55 ta kollektor (124,7 km)
2016--2018-yillarda, yana 13 tasi (65,0 km) 2019-yilda taʼmirlangan;
2016--2019-yillarga umumiy uzunlikning 71 foizi toʻgʻri keladi
(43-ilova). Kollektorlar eski, lekin ularga yaqinda xizmat koʻrsatilgan.

\textbf{Taʼmir tartibi.} Ishonchlilik nazariyasiga koʻra taʼmir
elementning holati yoki yoshiga qarab rejalashtiriladi: eski va eskirgan
elementlar birinchi navbatda taʼmirlanadi {[}23, 128{]}. Kadastrda
bunday tartib koʻrinmaydi. Oxirgi taʼmir yili kollektorning yoshi,
uzunligi va sarfi bilan bogʻlanmagan (ρ = −0,06\ldots+0,17; ishonchli
emas). 2016--2019-yillarda taʼmirlangan 68 ta kollektorning oʻrtacha
yoshi 60,5 yil, qolgan 22 tasiniki 61 yil --- deyarli bir xil. Yosh
guruhlari boʻyicha ham manzara bir xil: 55 yoshgacha, 55--65 yoshli va
65 yoshdan katta kollektorlarda oxirgi taʼmirdan beri oʻtgan muddatning
medianasi 4 yil, soʻnggi besh yilda taʼmirlanganlar ulushi 77--80 foiz.

Taʼmirlar kollektorning holatiga qarab emas, taʼmir dasturlari boʻyicha
oʻtkazilgan: maʼlum yillarda tarmoqning katta qismi bir vaqtda
taʼmirlangan. Bunday tartibda taʼmir resurslari eng koʻp ehtiyoj boʻlgan
joyga yoʻnaltirilganini tasdiqlab boʻlmaydi. Bundan tashqari, taʼmirlar
bir vaqtga toʻplangani sababli keyingi yillarda tarmoqning katta qismi
bir vaqtda eskiradi va taʼmirga muhtoj boʻladi. Kadastrdagi maʼlumotlar
boʻyicha 2010--2020-yillarda yiliga oʻrtacha 24,2 km kollektor
taʼmirlangan, bu butun tarmoqni taxminan 11 yilda bir marta taʼmirlashga
mos keladi. Haqiqatda taʼmir tez-tezroq boʻlgan boʻlishi mumkin, chunki
kadastrda faqat oxirgi taʼmir yili bor.

\textbf{Ichki zovurlar.} Xoʻjalik ichki tarmogʻida vaziyat boshqacha.
332 ta zovur (370,7 km) 1929--2012-yillarda qurilgan, yoshining
medianasi kollektorlarniki bilan bir xil (61 yil); ularning 112 tasi
(127,3 km) 1950-yillarda qurilgan. Zovurlarning 217 tasida (239,7 km,
uzunlikning 64,7 foizi) taʼmir umuman qayd etilmagan. Qolgan 115 tasi
faqat ikki yilda taʼmirlangan: 2017-yilda 87 ta (110,2 km) va 2018-yilda
28 ta (20,8 km). Taʼmirlangan va taʼmirlanmagan zovurlarning yoshi bir
xil (61 yil), uzunligi ham deyarli farq qilmaydi (1,0 va 0,9 km). Ichki
zovurlarni taʼmirga tanlashda ham yosh yoki holat mezoni qoʻllanilmagan.

Taʼmirsiz qolgan zovurlar ulushi tizimlar boʻyicha katta farq qiladi
(44-ilova): Mashʼal massivida 21 foiz, Markaziy zovur tizimi va
Y.Oxunboboyev massivida 100 foiz.

Tizimning ishonchliligi boʻyicha eng zaif boʻgʻin kollektorlar emas,
xoʻjalik ichki zovurlaridir: ular kadastrdagi ochiq tarmoq uzunligining
58 foizini tashkil etadi, uchdan ikki qismida esa taʼmir qayd etilmagan.

\textbf{Kollektorlar yoshi va taʼmirining xaritasi.} Kadastrda nomi
aniqlangan 60 ta kollektor (xaritada 114,8 km) foydalanishga
topshirilgan yili va oxirgi taʼmir yili boʻyicha xaritaga tushirildi
(4.1-rasm). 1950-yilgacha qurilgan eng eski kollektorlar (38,6 km)
asosan Z.M.Bobur, Mashʼal va Istiqlol massivlarida joylashgan,
T.Mirzayev massivida esa 1950-yillarda qurilganlari ustun (15,1 km).

Taʼmir yillari massivlar boʻyicha bloklarga ajraladi. Mashʼal massivida
kollektorlarning asosiy qismi 2017--2018-yillarda (18,9 km), T.Mirzayev
massivida 2015--2016-yillarda (18,1 km) taʼmirlangan.
2010--2014-yillarda taʼmirlangan, yaʼni eng uzoq vaqt xizmat
koʻrsatilmagan kollektorlar Istiqlol, M.Toʻjiboyev va Z.M.Bobur
massivlarida. Bu taqsimot taʼmirlar alohida kollektorning holatiga qarab
emas, massivlar boʻyicha dasturlar asosida oʻtkazilganini tasdiqlaydi:
xaritadagi 60 ta kollektor boʻyicha ham oxirgi taʼmir yili qurilish
yiliga bogʻliq emas. Taʼmirsiz ichki zovurlar ulushi esa kollektorlari
yaqinda taʼmirlangan massivlarda ham yuqori: Z.M.Bobur massivida 82
foiz, Istiqlol va Paxtakorda 74 foiz (4.1-rasm, b).

\oakfig{figures/map5_collector_age_repair.png}{4.1-rasm. Kollektorlar: a --- foydalanishga topshirilgan yil boʻyicha; b
--- oxirgi taʼmir yili boʻyicha (massiv yozuvlarida --- taʼmir qayd
etilmagan ichki zovurlar ulushi). Kadastrda nomi aniqlangan kollektorlar}

\textbf{Massivlar boʻyicha ishonchlilik koʻrsatkichlari.} Kadastr
kollektorlarni massivlarga bogʻlamaydi, shuning uchun bu GAT bazasi
orqali tiklandi: xaritadan raqamlashtirilgan kollektor boʻlaklari massiv
chegaralari bilan kesishtirildi va kadastrda nomi aniqlangan boʻlaklarga
kollektorning yoshi va taʼmir yili biriktirildi (2.3-§). Massivlarning
sugʻoriladigan maydoni bilan birga bu maʼlumotlar har bir massiv uchun
drenaj tarmogʻining solishtirma uzunligini va holatini baholash imkonini
berdi (45-ilova va 46-ilova).

Massivlar drenaj tarmogʻining taʼminlanganligi va holati boʻyicha katta
farq qiladi. Ochiq tarmoqning umumiy solishtirma uzunligi (kollektorlar
va ichki zovurlar) Istiqlol massivida 48,6 m/ga, Mashʼal massivida 38,7
m/ga, M.Toʻjiboyev massivida esa atigi 6,8 m/ga. Ichki zovurlar holati
boʻyicha Mashʼal massivi boshqalardan ajralib turadi: u yerda
zovurlarning 79 foizi taʼmirlangan, qolgan massivlarda esa 0--53 foizi.
Toʻrtta koʻrsatkich (tarmoq zichligi, taʼmirsiz ichki zovurlar ulushi,
kollektorlar yoshi va oxirgi taʼmirdan beri oʻtgan muddat) boʻyicha
massivlarning oʻrtacha oʻrni hisoblanganda eng zaif tarmoq
Y.Oxunboboyev, M.Toʻjiboyev va Z.M.Bobur massivlarida, eng yaxshisi
Mashʼal massivida ekani aniqlandi.

Bu bahoni III bob natijalari bilan solishtirish mumkin. Z.M.Bobur
massivi tarmogʻi zaif massivlar qatorida (umumiy solishtirma uzunlik
20,6 m/ga, ichki zovurlarning 82 foizi taʼmirsiz) va aynan shu massivda
sizot suvi sathi koʻtarilmoqda (3.2-§), shoʻrlanish xavfi ostidagi
quduqlar soni esa eng koʻp (3.5-§). Istiqlol massivida tarmoq zich,
lekin ichki zovurlarning 74 foizi taʼmirsiz va kollektorlar tumandagi
eng eskilari (oʻrtacha 75 yil). Mashʼal massivida tarmoq eng yaxshi
holatda boʻlsa ham, sath sayoz va minerallashuv yuqori: bu yerda holatni
tarmoq emas, tabiiy sharoit belgilaydi. Y.Oxunboboyev va M.Toʻjiboyev
massivlarida tarmoq siyrak, lekin sath chuqur, shuning uchun bu yerda
drenajga ehtiyoj ham kam. Tarmoq holati va meliorativ holat birgalikda
qaralganda Z.M.Bobur va Istiqlol massivlari birinchi navbatdagi hudud
sifatida ajraladi; bu xulosa 4.5-§ dagi tavsiyalarga asos boʻladi.

\textbf{Cheklov.} Kadastrda faqat oxirgi taʼmir yili yozilgan; undan
oldingi taʼmirlar va nosozliklar hisobi yoʻq. Shuning uchun
ishonchlilikning odatdagi koʻrsatkichlarini (buzilishlar soni,
buzilishlar orasidagi oʻrtacha muddat) hisoblab boʻlmaydi. «Taʼmir qayd
etilmagan» degan yozuv taʼmir boʻlmaganini ham, hisobga olinmaganini ham
bildirishi mumkin. Kollektorlarni massivlarga bogʻlash xaritadagi
nomlangan boʻlaklarga tayanadi, shuning uchun massivlar boʻyicha yosh va
taʼmir koʻrsatkichlari taxminiy. Bu kamchiliklarni bartaraf etish uchun
kadastrga nosozliklar va barcha taʼmirlar jurnalini kiritish taklif
etiladi (4.5-§).

\oakpar{4.2-§. Drenaj moduli va kollektorlarning suv oʻtkazish qobiliyati}

Drenaj tizimi oʻz vazifasini bajarishi uchun daladan ortiqcha suvni
kerakli miqdorda olib chiqa olishi kerak. Buni ikki koʻrsatkich
tavsiflaydi: drenaj moduli (bir gektardan chiqariladigan suv) va
kollektorning suv oʻtkazish qobiliyati {[}108, 122{]}. Ushbu paragrafda
kadastrdagi koʻrsatkichlar suv balansi boʻyicha hisoblangan talab bilan
solishtiriladi, natija sizot suvining yozgi koʻtarilishi bilan
tekshiriladi va kollektor holatining suv oʻtkazishga taʼsiri baholanadi.

\textbf{Kadastr boʻyicha haqiqiy koʻrsatkichlar.} Kadastrda har bir
kollektorning oxiridagi suv sarfi va xizmat maydoni berilgan. Ularning
nisbati bir gektarga toʻgʻri keladigan sarfni (solishtirma sarf) beradi:

\begin{equation}\tag{4.1}
q = \frac{1000 \, Q}{F}
\end{equation}

bu yerda q --- solishtirma sarf, l/(s·ga); Q --- kollektor oxiridagi
sarf, m³/s; F --- xizmat koʻrsatish maydoni, ga.

90 ta kollektor boʻyicha solishtirma sarfning medianasi 0,53 l/(s·ga),
kollektorlarning yarmida 0,19--1,35 l/(s·ga). Chetki qiymatlar orasidagi
farq juda katta: 0,03 dan 50 l/(s·ga) gacha. 13 ta kollektorda
koʻrsatkich 0,1 l/(s·ga) dan kichik, 5 tasida 5 l/(s·ga) dan katta.
Barcha kollektorlar oxiridagi sarflar yigʻindisi 18,27 m³/s, bu xizmat
maydoniga (11 079 ga) nisbatan 1,65 l/(s·ga) ni tashkil etadi.

Farqning bir sababi --- yirik kollektorlar oʻz maydonidan tashqari
yuqoridan keladigan suvni ham oʻtkazadi. Eng katta sarfli beshta
kollektorda (Gʻarbiy Qoraqoʻrgʻon 4,6 m³/s, Oʻqchi 2,2 m³/s, Gʻarbiy 2,0
m³/s, Markaziy 1,5 m³/s, Nayman 1,2 m³/s) solishtirma sarf 2,5--10,2
l/(s·ga) ga yetadi; ulardagi sarfni xizmat maydoniga boʻlish drenaj
modulini bermaydi. Bu beshta kollektor chiqarilganda qolgan 85 tasi
boʻyicha mediana 0,50 l/(s·ga). Kadastr yozuvlarining oʻzida ham
nomuvofiqliklar bor: masalan, Joʻga kollektorining uzunligi 5,0 km,
xizmat maydoni esa 4 ga deb koʻrsatilgan. Bundan tashqari kadastrdagi
sarf loyihaviy yoki bir martalik oʻlchov qiymati boʻlib, uning qachon va
qanday sharoitda aniqlangani koʻrsatilmagan.

\textbf{Tarmoq zichligi va hudud nishabligi.} Kollektorlarning
solishtirma uzunligi medianasi 33 m/ga. Kollektorlar va ichki zovurlar
birgalikda (636,5 km) kadastrdagi xizmat maydoniga (11 079 ga) nisbatan
57,5 m/ga ni tashkil etadi. Hudud nishabligi 173 ta kuzatuv qudugʻining
balandligi (528--686 m) boʻyicha baholandi: yer yuzasi shimolga qarab
oʻrtacha 2,4 m/km (i = 0,0024) nishablikda pasayadi. Bu kollektorlar
tubi uchun 0,0005--0,002 oraligʻidagi nishablikni taʼminlashga imkon
beradi va hudud tabiiy oqim sharoiti boʻyicha drenaj uchun qulay.

\textbf{Talab etiladigan drenaj moduli.} Talab etiladigan modul
vegetatsiya davrining suv balansidan aniqlandi. Drenaj sugʻorish
suvining dalada va ichki kanallarda yoʻqolib, sizot suviga qoʻshilgan
qismini olib chiqishi kerak {[}108, 122{]}:

\begin{equation}\tag{4.2}
D = \beta \left[ \frac{N}{E_a} - N + \frac{N}{E_a E_c} - \frac{N}{E_a} \right], \qquad q_t = \frac{D}{8{,}64 \, T}
\end{equation}

bu yerda N --- sof sugʻorish talabi, mm; Eₐ --- dalada sugʻorish
samaradorligi; \(E_{c}\) --- xoʻjalik ichki tarmogʻining foydali ish
koeffitsiyenti; β --- yoʻqotishlarning sizot suviga yetib boradigan
ulushi (0,7 qabul qilindi); D --- drenaj yuklamasi, mm; T --- hisobiy
davr, sutka; \(q_{t}\) --- talab etiladigan drenaj moduli, l/(s·ga).

Sof sugʻorish talabi FAO-56 usulida hisoblandi {[}31{]}: oylik
bugʻlanuvchanlik (2.1-§) ekin koeffitsiyentiga koʻpaytirildi va undan
foydali yogʻin ayirildi. Ekinlar tarkibi meliorativ kadastrning kontur
jadvalidan olindi: 25 405 ga maydonning 36,6 foizi gʻoʻza, 42,5 foizi
gʻalla, 20,9 foizi boshqa ekinlar (bogʻ, sabzavot, yem-xashak). Ekin
koeffitsiyentlari FAO-56 boʻyicha oylar kesimida olindi {[}31{]},
foydali yogʻin USDA SCS usulida aniqlandi. Hisob natijalari 47-ilova va
48-ilovada keltirilgan.

Asosiy hisobda gʻalladan keyingi maydonlar iyul--sentabrda boʻsh deb
qabul qilingan, shuning uchun iyul--avgustda oʻrtacha ekin
koeffitsiyenti 0,62 ga tushadi. Amalda gʻalladan keyin takroriy ekinlar
ekiladi; ular hisobga olinganda yillik ekinlar bugʻlanishi 891 mm ga,
sof sugʻorish talabi 696 mm ga, iyuldagi talab esa 178 mm ga oshadi.
Shuning uchun hisob ikki variantda (takroriy ekinlarsiz va bilan) va
sugʻorish samaradorligining uch darajasi boʻyicha bajarildi
(4.1-jadval).

\oakjadvalnum{4.1-jadval}

\oakilovanom{\textbf{Sugʻorish samaradorligiga qarab drenaj yuklamasi va talab
etiladigan drenaj moduli}}

{\fontsize{12}{14}\selectfont\setlength{\tabcolsep}{4pt}

\noindent\begin{longtable}[]{|
  >{\raggedright\arraybackslash}p{(\columnwidth - 12\tabcolsep - 7\arrayrulewidth) * \real{0.2042}}|>{\raggedright\arraybackslash}p{(\columnwidth - 12\tabcolsep - 7\arrayrulewidth) * \real{0.1972}}|>{\raggedright\arraybackslash}p{(\columnwidth - 12\tabcolsep - 7\arrayrulewidth) * \real{0.1007}}|>{\raggedright\arraybackslash}p{(\columnwidth - 12\tabcolsep - 7\arrayrulewidth) * \real{0.1658}}|>{\raggedright\arraybackslash}p{(\columnwidth - 12\tabcolsep - 7\arrayrulewidth) * \real{0.1660}}|>{\raggedright\arraybackslash}p{(\columnwidth - 12\tabcolsep - 7\arrayrulewidth) * \real{0.1662}}|}
\hline
\begin{minipage}[b]{\linewidth}\raggedright
\textbf{Samaradorlik darajasi}
\end{minipage} & \begin{minipage}[b]{\linewidth}\raggedright
\textbf{Dalada samaradorlik Eₐ}
\end{minipage} & \begin{minipage}[b]{\linewidth}\raggedright
\textbf{Ichki tarmoq FIK \(E_{c}\)}
\end{minipage} & \begin{minipage}[b]{\linewidth}\raggedright
\textbf{Yalpi sugʻorish, mm}
\end{minipage} & \begin{minipage}[b]{\linewidth}\raggedright
\textbf{Yillik drenaj yuklamasi, mm}
\end{minipage} & \begin{minipage}[b]{\linewidth}\raggedright
\textbf{Iyuldagi talab etiladigan modul, l/(s·ga)}
\end{minipage} \\
\hline
\endhead
\hline
\endlastfoot
Yuqori samaradorlik & 0,75 & 0,85 & 822 / 1 091 & 209 / 277 & 0,16 /
0,27 \\ \hline
Oʻrta samaradorlik & 0,65 & 0,80 & 1 008 / 1 338 & 339 / 450 & 0,26 /
0,43 \\ \hline
Past samaradorlik & 0,55 & 0,75 & 1 270 / 1 687 & 522 / 694 & 0,40 /
0,66 \\
\end{longtable}

}

Izoh: birinchi qiymat --- takroriy ekinlarsiz, ikkinchisi --- takroriy
ekinlar bilan.

Yalpi sugʻorish hajmi 8,2--16,9 ming m³/ga; oʻrta daraja (10,1--13,4
ming m³/ga) Fargʻona vodiysidagi odatiy suv olishga mos keladi {[}51{]}.
Iyuldagi talab etiladigan drenaj moduli 0,16--0,66 l/(s·ga), oʻrta
darajada 0,26--0,43 l/(s·ga). Bu qiymatlar adabiyotda sugʻoriladigan
quruq hududlar uchun tavsiya etilgan loyihaviy drenaj sarfi oraligʻiga
(1--3 mm/sut yoki 0,12--0,35 l/(s·ga)) toʻgʻri keladi yoki undan biroz
yuqori {[}108, 122{]}. Drenajga tushadigan yuk sugʻorish samaradorligiga
kuchli bogʻliq: samaradorlik pastdan yuqoriga oshirilsa, talab 2,5
baravar kamayadi. Demak, drenaj muammosini sugʻorishni yaxshilash orqali
ham hal etish mumkin.

\textbf{Haqiqiy va talab etiladigan modul qiyosi.} Kadastrdagi
solishtirma sarf talab etiladigan modul bilan solishtirildi (48-ilova,
b). Oʻrta daraja boʻyicha (0,26 l/(s·ga)) 90 ta kollektordan 27 tasining
solishtirma sarfi talabdan past; ular 60,2 km uzunlikka va 2 863 ga
xizmat maydoniga ega. Takroriy ekinlar hisobga olinganda (0,43 l/(s·ga))
ular 38 ta (4 501 ga), past samaradorlikda esa 48 ta (5 845 ga, xizmat
maydonining 53 foizi). Kadastr koʻrsatkichlari boʻyicha kollektorlarning
uchdan biridan yarmigacha qismi sugʻorish mavsumining eng ogʻir oyida
talab etilgan sarfni oʻtkazmaydi. Biroq bu xulosa kadastrdagi sarf
qiymatining ishonchliligiga bogʻliq; uni tasdiqlash uchun kollektorlarda
sarfni oʻlchash zarur (4.5-§).

\textbf{Talabni qoplamaydigan kollektorlarning joylashuvi.} Kadastr nomi
xaritadagi boʻlaklarga biriktirilgan 60 ta kollektor solishtirma sarfi
va talab etiladigan modul boʻyicha xaritaga tushirildi (4.2-rasm).
Ulardan 10 tasining sarfi oʻrta daraja talabidan ham past, 9 tasi
takroriy ekinlar hisobga olinganda, 8 tasi past samaradorlikda talabni
qoplamaydi. 28 tasi barcha variantlarda talabni qoplaydi, 5 tasi esa
boshqa hududlar suvini ham oʻtkazuvchi kollektorlar.

Talabdan past kollektorlar bitta massivda toʻplanmagan, tuman boʻylab
tarqoq joylashgan. Ularning nomi aniqlangan kollektorlar uzunligidagi
ulushi Paxtakorda 0 foiz, T.Mirzayevda 15 foiz, U.Qoraboyevda esa 47
foiz. Massivlar boʻyicha ulush faqat kadastr nomi aniqlangan boʻlaklarga
tayanadi: Istiqlol, Paxtakor va M.Toʻjiboyev massivlarida u atigi 5--6
km kollektorni qamraydi, shuning uchun bu massivlar uchun baho taxminiy.

\oakfig{figures/map6_drainage_capacity.png}{4.2-rasm. Kollektorlarning kadastr boʻyicha solishtirma sarfi va
iyuldagi talab etiladigan drenaj moduli (4.1-jadval): a --- kollektorlar
boʻyicha; b --- massivlar boʻyicha talabdan (0,43 l/(s·ga)) past
kollektorlar uzunligining ulushi}

\textbf{Sizot suvi sathi boʻyicha tekshiruv.} Suv balansi hisobi
taxminiy maʼlumotlarga tayanadi, shuning uchun u kuzatilgan sizot suvi
sathi bilan tekshirildi. 3.2-§ ga koʻra sath mart oxiridan iyul
oxirigacha, taxminan 120 sutka davomida, oʻrtacha 0,48 m ga koʻtariladi.
Sathning koʻtarilishi sizot suviga tushayotgan suv drenaj orqali
chiqayotgan suvdan koʻpligini bildiradi. Bu ortiqcha suv tuproqning suv
beruvchanlik koeffitsiyenti μ orqali baholanadi: qumoq tuproqlar uchun μ
= 0,05--0,15 qabul qilinganda u 24--72 mm ni yoki 0,02--0,07 l/(s·ga) ni
tashkil etadi. Bu qiymat talab etiladigan drenaj modulining 5--45
foiziga teng. Mavjud tarmoq sugʻorish mavsumida sizot suviga tushayotgan
suvning asosiy qismini olib chiqadi, lekin toʻliq emas: drenaj
quvvatidagi taqchillik tuman oʻrtachasida talabning 5--45 foizini
tashkil etadi. Bu baho yuqoridagi kadastr qiyosidan koʻra yumshoqroq,
chunki u butun tuman boʻyicha oʻrtacha holatni aks ettiradi; massivlar
boʻyicha esa mavsumiy koʻtarilish 0,32--0,85 m oraligʻida (3.2-§),
shuning uchun taqchillik U.Qoraboyev massivida oʻrtachadan qariyb ikki
baravar, Paxtakor massivida 1,4 baravar katta.

\textbf{Suv oʻtkazish qobiliyati.} Kollektorning suv oʻtkazish
qobiliyati Manning formulasi boʻyicha aniqlanadi {[}85, 122{]}:

\begin{equation}\tag{4.3}
Q = \frac{1}{n} \, \omega \, R^{2/3} \, i^{1/2}
\end{equation}

bu yerda ω --- jonli kesim yuzasi, m²; R --- gidravlik radius, m; i ---
tub nishabligi; n --- gʻadir-budurlik koeffitsiyenti. Kadastrda
kollektorlarning koʻndalang kesim oʻlchamlari va tub nishabligi
berilmagan, shuning uchun aniq kollektorlar boʻyicha hisob bajarib
boʻlmaydi. Buning oʻrniga namunaviy trapetsiya kesim (tub kengligi 1,0
m, qiyalik koeffitsiyenti 1,5, chuqurligi 1,2 m) uchun oʻzan holatining
oʻtkazish qobiliyatiga taʼsiri baholandi (49-ilova).

Oʻzanni oʻsimlik qoplashi oʻtkazish qobiliyatini 29 foizga, oʻsimlik
bilan birga 0,3 m loyqa qatlami 42 foizga, kuchli oʻsimlik va 0,5 m
loyqa esa 61 foizga kamaytiradi: texnik xizmat koʻrsatilmagan kollektor
toza holatdagi sarfning faqat 40--70 foizini oʻtkazadi. 4.1-§ ga koʻra
taʼmir kollektor holatiga qarab rejalashtirilmaydi, shuning uchun
taʼmirlar orasidagi 10 yil atrofidagi davrda kollektorlarning bir qismi
loyihaviy oʻtkazish qobiliyatining ancha qismini yoʻqotishi mumkin.
Shuning uchun kollektorlarning texnik holatini (oʻsimlik bilan
qoplanish, loyqa qatlami qalinligi) muntazam qayd etish va uni taʼmir
navbatini aniqlash mezoni sifatida qoʻllash taklif etiladi (4.5-§).

\textbf{Cheklovlar.} Hisob uchta taxminga tayanadi: ekinlar tarkibi
kadastr kontur jadvali boʻyicha olingan; sugʻorish samaradorligi va
yoʻqotishlarning sizot suviga yetib borish ulushi stsenariy sifatida
berilgan; suv beruvchanlik koeffitsiyenti adabiyotdagi oraliqdan
olingan. Bu taxminlar natijani 2--3 baravar oʻzgartiradi, lekin asosiy
xulosalarni (kadastrdagi sarflar juda turlicha, tarmoqning bir qismi
talabni qoplamaydi, oʻzan holati suv oʻtkazishga kuchli taʼsir qiladi)
oʻzgartirmaydi. Hisobni aniqlashtirish uchun kanallardan suv olish
hajmlari, kollektorlardagi oʻlchangan sarflar va kollektorlar kesimining
oʻlchamlari kerak.

\oakpar{4.3-§. Kollektor taʼmirining sizot sathiga taʼsirini baholash}

3.4-§ da koʻrsatilganidek, drenaj taʼsirini bir yildagi turli joylarni
solishtirib aniqlab boʻlmaydi: kollektorlar sath tabiatan sayoz joylarga
qurilgan. Shuning uchun taʼmir samarasi vaqt boʻyicha baholandi: bir xil
quduqda taʼmirdan oldingi va keyingi sath solishtirildi, barcha
quduqlarga bir xil taʼsir qilgan yillik oʻzgarishlar esa kollektordan
uzoqdagi quduqlar yordamida chiqarib tashlandi. Bu usul «farqlar farqi»
deb ataladi va tadbirlar samarasini baholashda keng qoʻllanadi {[}34{]}.

\textbf{Usul.} Xaritadagi kollektor chiziqlariga kadastrdagi nom va
taʼmir yili biriktirildi; bu 60 ta kollektor uchun bajarildi (2.3-§).
Quduqlar ikki guruhga ajratildi. Birinchisi --- nomi va taʼmir yili
maʼlum kollektordan 300 m gacha masofadagi quduqlar (36 ta, 28 ta
kollektor yonida). Ikkinchisi (solishtirish guruhi) --- har qanday
kollektordan 600 m dan uzoqdagi quduqlar (49 ta). 300--600 m oraliqdagi
quduqlar guruhlar aralashmasligi uchun olinmadi. Jami 819 ta quduq-yil;
oʻzgarmas qiymatli quduq-yillar va mustaqil maʼlumoti boʻlmagan yillar
3.1-§ dagi tartibda chiqarildi. Hisob formulasi:

\begin{equation}\tag{4.4}
h_{it} = \alpha_i + \gamma_t + \beta \, D_{it} + \varepsilon_{it}
\end{equation}

bu yerda h --- i-quduqning t-yildagi oʻrtacha sathi, m; α --- quduqning
doimiy xususiyati; γ --- yilning umumiy taʼsiri; D --- taʼmirdan keyingi
yillar uchun 1 ga, boshqa hollarda 0 ga teng koʻrsatkich; β --- taʼmir
taʼsiri. Musbat β sathning chuqurlashganini bildiradi. Bir quduqning
turli yillardagi oʻlchovlari bogʻliq boʻlgani uchun xatolar quduqlar
boʻyicha guruhlab hisoblandi.

Usulning asosiy sharti: taʼmir boʻlmaganda ikki guruhda sath bir xil
oʻzgargan boʻlar edi {[}34, 109{]}. Buni toʻgʻridan-toʻgʻri tekshirib
boʻlmaydi, lekin taʼmirdan oldingi yillarda guruhlar farqi
oʻzgarmaganini koʻrsatish orqali tasdiqlash mumkin. Kollektorlar turli
yillarda (2010--2020) taʼmirlangani uchun oldin taʼmirlanganlar keyin
taʼmirlanganlar uchun notoʻgʻri solishtirish boʻlib qolishi mumkin
{[}59{]}. Shuning uchun yillar boʻyicha bahoda faqat hech qachon taʼmir
taʼsiriga tushmagan uzoq quduqlar bilan solishtirildi.

\textbf{Asosiy natija va qoʻshimcha tekshiruvlar.} Asosiy variant
tahlildan oldin belgilandi: yillik oʻrtacha sath boʻyicha yuqoridagi
formula. Qolgan variantlar natija shartlar oʻzgarganda saqlanishini
tekshiradi; xulosa ular orasidan tanlab olinmaydi (4.2-jadval).

\oakjadvalnum{4.2-jadval}

\oakilovanom{\textbf{Kollektor taʼmirining sizot sathiga taʼsiri bahosi}}

{\fontsize{12}{14}\selectfont\setlength{\tabcolsep}{4pt}

\noindent\begin{longtable}[]{|
  >{\raggedright\arraybackslash}p{(\columnwidth - 10\tabcolsep - 6\arrayrulewidth) * \real{0.4414}}|>{\raggedright\arraybackslash}p{(\columnwidth - 10\tabcolsep - 6\arrayrulewidth) * \real{0.0977}}|>{\raggedright\arraybackslash}p{(\columnwidth - 10\tabcolsep - 6\arrayrulewidth) * \real{0.1727}}|>{\raggedright\arraybackslash}p{(\columnwidth - 10\tabcolsep - 6\arrayrulewidth) * \real{0.0782}}|>{\raggedright\arraybackslash}p{(\columnwidth - 10\tabcolsep - 6\arrayrulewidth) * \real{0.2100}}|}
\hline
\begin{minipage}[b]{\linewidth}\raggedright
\textbf{Variant}
\end{minipage} & \begin{minipage}[b]{\linewidth}\raggedright
\textbf{β, m}
\end{minipage} & \begin{minipage}[b]{\linewidth}\raggedright
\textbf{95 foizlik ishonch oraligʻi}
\end{minipage} & \begin{minipage}[b]{\linewidth}\raggedright
\textbf{p}
\end{minipage} & \begin{minipage}[b]{\linewidth}\raggedright
\textbf{Quduq-yillar}
\end{minipage} \\
\hline
\endhead
\hline
\endlastfoot
Asosiy model, yillik oʻrtacha sath & −0,12 & −0,33 \ldots{} +0,09 & 0,26
& 819 \\ \hline
Yozgi sath (iyul--avgust) & −0,07 & −0,27 \ldots{} +0,13 & 0,50 & 803 \\ \hline
Mavsumiy koʻtarilish (qishki sath − yozgi sath) & −0,05 & −0,20 \ldots{}
+0,11 & 0,56 & 803 \\ \hline
Faqat hech qachon oʻzgarmas qator bermagan quduqlar & −0,20 & −0,39
\ldots{} −0,02 & 0,04 & 618 \\ \hline
Faqat 2009--2016-yillar & −0,03 & −0,14 \ldots{} +0,08 & 0,61 & 289 \\ \hline
Faqat 2018--2024-yillar & −0,16 & −0,44 \ldots{} +0,12 & 0,27 & 530 \\ \hline
Faqat taʼmirlangan kollektor yonidagi quduqlar & +0,08 & −0,08 \ldots{}
+0,25 & 0,33 & 424 \\ \hline
Kollektordan 150 m gacha (17 ta quduq) & −0,09 & −0,33 \ldots{} +0,14 &
0,44 & 819 \\ \hline
Kollektordan 150--300 m (19 ta quduq) & −0,15 & −0,45 \ldots{} +0,15 &
0,33 & 819 \\ \hline
Soxta taʼmir yili (3 yil oldin) & +0,04 & −0,08 \ldots{} +0,16 & 0,51 &
543 \\ \hline
Taʼmir yili boʻyicha guruhlab (Callaway--SantʼAnna), keyingi 0--3 yil &
−0,11 & −0,29 \ldots{} +0,10 & 0,26 & 819 \\ \hline
Taʼmir yili boʻyicha guruhlab, barcha keyingi yillar & −0,27 & −0,49
\ldots{} −0,05 & 0,02 & 819 \\ \hline
Taʼmirsiz kutiladigan sath bilan solishtirish & −0,18 & −0,39 \ldots{}
+0,02 & 0,08 & 819 \\
\end{longtable}

}

Izoh: oxirgi uch qatorda ishonch oraligʻi quduqlarni takroriy tanlash
usuli bilan hisoblangan.

Hech bir variantda taʼmirdan keyin sath chuqurlashgani aniqlanmadi.
Asosiy baho −0,12 m, ishonch oraligʻi nolni oʻz ichiga oladi. Oraliqning
yuqori chegarasi +0,09 m, yaʼni taʼmirlangan kollektor yonida sath
uzoqdagi quduqlarga nisbatan 0,09 m dan ortiq chuqurlashgan boʻlishi
ehtimoldan uzoq. Yozgi sath va mavsumiy koʻtarilish boʻyicha ham taʼsir
yoʻq: agar taʼmir kollektorning oʻtkazish qobiliyatini tiklagan boʻlsa
(4.2-§), u eng avvalo yozgi koʻtarilishni kamaytirishi kerak edi, lekin
mavsumiy koʻtarilish oʻzgarmagan (−0,05 m; p = 0,56). Kollektorga
yaqinroq quduqlarda (150 m gacha) taʼsir uzoqroqdagilardan katta emas,
garchi shunday boʻlishi kutilgan boʻlsa ham.

Variantlar orasida ishora oʻzgaradi, yaʼni barqaror taʼsir yoʻq. Bitta
variantdagi ishonchli manfiy baho (−0,20 m) uzoqdagi quduqlarda
2018-yildan keyin sath kuchliroq chuqurlashgani bilan bogʻliq va taʼmir
taʼsiri deb qaralmaydi; bu 3.4-§ dagi xulosaga mos. Soxta taʼmir yili
bilan tekshiruvda (taʼmir yili 3 yil oldinga surilib, faqat haqiqiy
taʼmirdan oldingi yillar olinganda) taʼsir nolga yaqin (+0,04 m): usul
yoʻq taʼsirni «topib» qoʻymaydi.

\textbf{Kollektorlar turli yillarda taʼmirlangani hisobga olingan
tekshiruv.} Kollektorlar 2010--2020-yillarda turli yillarda
taʼmirlangan. Bunday holda oddiy hisob notoʻgʻri natija berishi mumkin:
oldin taʼmirlangan quduqlar keyin taʼmirlanganlar uchun solishtirish
boʻlib qoladi {[}59{]}. Shuning uchun taʼsir yana ikki usulda
hisoblandi. Birinchisida (Callaway va SantʼAnna {[}43{]}) bir yilda
taʼmirlangan kollektorlar yonidagi quduqlar faqat hech qachon
taʼmirlanmagan kollektorlardan uzoqdagi quduqlar bilan solishtirildi.
Ikkinchisida {[}39{]} taʼmir boʻlmaganda sath qanday boʻlishi taʼmirdan
oldingi yillar va uzoqdagi quduqlar boʻyicha hisoblab topildi va haqiqiy
sath bilan solishtirildi. Taʼmirdan keyingi 0--3 yilda natijalar −0,11 m
(−0,29 \ldots{} +0,10 m) va −0,18 m (−0,39 \ldots{} +0,02 m), asosiy
hisobga mos. Barcha keyingi yillar olinganda baho −0,27 m gacha ortadi;
uning manbai quyidagi yillar boʻyicha bahodagi uzoq muddatli qiymat
bilan bir xil. Hech bir usul taʼmirdan keyin sath chuqurlashganini
koʻrsatmadi.

\textbf{Taʼmirdan oldingi va keyingi yillar boʻyicha baho.} Taʼmir
taʼsiri vaqt oʻtib namoyon boʻlishi mumkin, shuning uchun taʼmir yiliga
nisbatan har bir yil alohida baholandi (4.3-rasm). Taʼmirdan oldingi
yillarda ikki guruh farqi nolga teng (p = 0,18): bu usulning asosiy
shartini tasdiqlaydi. Taʼmir yili va undan keyingi uch yilda farq −0,07
\ldots{} −0,10 m va ishonchli emas. Faqat taʼmirdan 4 va undan koʻp yil
oʻtganda farq −0,25 m (p = 0,04): sath biroz sayozlashgan. Bu baho
2010--2015-yillarda taʼmirlangan oz sonli kollektorlar yonidagi
quduqlarga tayanadi, sakkizta bahodan bittasi tasodifan ishonchli
chiqishi ham mumkin; shuning uchun u taʼmir taʼsiri deb qaralmaydi.

\oakfig{figures/ch4_event_study.png}{4.3-rasm. Kollektor taʼmirining sizot sathiga taʼsiri yillar boʻyicha:
taʼmir yiliga nisbatan har bir yilda uzoqdagi quduqlarga nisbatan sath
farqi va 95 foizlik ishonch oraligʻi (asos --- taʼmirdan bir yil oldin)}

Taʼmirlangan kollektorlar taʼmir yili boʻyicha ham ajratildi.
2017-yilgacha taʼmirlangan kollektorlar yonidagi 28 ta quduq boʻyicha
taʼsir −0,06 m (p = 0,59), 2018--2020-yillarda taʼmirlanganlar yonidagi
8 ta quduq boʻyicha −0,26 m (p = 0,11). Hech bir guruhda chuqurlashish
yoʻq.

\textbf{Talqin va cheklovlar.} Natija «taʼmir samarasiz» degan xulosani
bermaydi. U faqat mavjud maʼlumotda mahalliy taʼsir koʻrinmasligini
bildiradi. Buning toʻrt sababi boʻlishi mumkin.

\begin{itemize}
\tightlist
\item
  Quduqlar kam. Taʼmirlangan kollektorlar yonida 36 ta quduq bor; bunday
  namuna bilan 0,29 m dan kichik taʼsirni ishonchli aniqlab boʻlmaydi
  (4.5-§).
\item
  Taʼmir taʼsiri mahalliy emas, hududiy boʻlishi mumkin. Kollektorlar
  yagona tarmoq boʻlib ishlaydi; 2016--2018-yillarda 55 ta kollektor
  birdaniga taʼmirlangan. Bunday holda uzoqdagi quduqlar ham taʼsir
  ostida qoladi va farq yoʻqoladi. Biroq bu izohga ikki natija qarshi
  turadi: sath chuqurlashishi kollektordan uzoqdagi quduqlarda kuchliroq
  (3.4-§) va u viloyatning 14 ta tumanidan 12 tasida ham kuzatiladi
  (3.6-§).
\item
  2018-yilda kuzatuv tarmogʻi oʻzgargan, bu taʼmirlar bilan bir vaqtga
  toʻgʻri kelib, ikki taʼsirni ajratishni qiyinlashtiradi.
\item
  Kadastrda faqat oxirgi taʼmir yili bor; ilgarigi taʼmirlar hisobga
  olinmagan, taʼmir hajmi va turi esa koʻpchilik kollektorlar uchun
  koʻrsatilmagan. Taʼmir kollektorning qaysi qismida bajarilgani ham
  nomaʼlum, shuning uchun baʼzi quduqlar taʼmirlanmagan qism yonida
  boʻlishi mumkin. Bu bahoni nolga tomon kamaytiradi, yaʼni haqiqiy
  taʼsir baholangandan kattaroq boʻlishi mumkin.
\end{itemize}

Tumanda 2017--2020-yillarda taʼmirga 2026-yil narxlarida 48,2 mlrd soʻm
sarflangan (4.4-§), ammo mavjud monitoring tizimi bu ishlarning
meliorativ natijasini oʻlchashga imkon bermaydi. Mavjud maʼlumotlarga
koʻra taʼmirdan keyin kollektor yonidagi sath uzoqdagi quduqlarga
nisbatan chuqurlashmagan: barcha baholarda ishonch oraligʻining yuqori
chegarasi +0,10 m dan oshmaydi. Biroq oraliqlar keng va 0,2 m atrofidagi
taʼsirni istisno qilmaydi; 2016--2018-yillardagi tuman boʻyicha
pogʻonali chuqurlashish esa taʼmirdan koʻra hududiy omillar bilan
yaxshiroq mos keladi. Taʼmir samarasini ishonchli oʻlchash uchun
monitoring tarmogʻini taʼmir rejalari bilan muvofiqlashtirish zarur;
buning talablari 4.5-§ da keltirilgan.

\oakpar{4.4-§. Taʼmir xarajatlari tahlili}

Taʼmir xarajatlari meliorativ tizim samaradorligining iqtisodiy tomonini
tavsiflaydi: qancha mablagʻ sarflangani, u tarmoqning qaysi qismiga
yoʻnaltirilgani va sarflangan mablagʻ qanday natija bergani {[}40,
128{]}. Ushbu paragrafda kadastrdagi xarajatlar uch jihatdan tahlil
qilinadi: 1 km taʼmir narxi va unga taʼsir qiluvchi omillar, xarajatning
kollektor qiymati va maydonga nisbati, xarajatlarning massivlar boʻyicha
taqsimoti va meliorativ ehtiyojga mosligi.

\textbf{Umumiy koʻrsatkichlar.} Kadastrda 53 ta kollektor boʻyicha
2017--2020-yillardagi taʼmir xarajati va taʼmirlangan uchastka
(piketlar) koʻrsatilgan. Jami 163,3 km uchastkaga kadastrda oʻsha
yillarning narxlarida 21 314 mln soʻm sarflangan. Turli yillardagi
xarajatlarni solishtirish va bugungi rejalashtirishda qoʻllash uchun
ular 2026-yil oktabr narxlariga keltirildi. Buning uchun isteʼmol
narxlari indeksidan foydalanildi (yillik inflyatsiya 2017-yilda 14,4,
2018-yilda 14,3, 2019-yilda 15,2, 2020-yilda 11,1, 2021--2025-yillarda
7,3--12,3 foiz, 2026-yil yanvar--oktabrida 4,8 foiz) {[}139{]}. Har bir
taʼmir yilining xarajati shu yilning oʻrtacha narxlarida deb olindi;
qayta hisoblash koeffitsiyenti 2017-yil uchun 2,60, 2018-yil uchun 2,27,
2019-yil uchun 1,98, 2020-yil uchun 1,75. 2026-yil narxlarida xarajat 48
176 mln soʻmni tashkil etadi. Jami xarajatning jami uzunlikka nisbati
295 mln soʻm/km, kollektorlar boʻyicha mediana 348 mln soʻm/km. Bundan
keyin barcha xarajatlar 2026-yil narxlarida keltiriladi; yillar boʻyicha
taqsimot ikkala narxda 50-ilovada berilgan.

\textbf{Narx tarkibi.} Narxlar ikki guruhga aniq ajraladi. 2019-yilda
taʼmirlangan sakkizta kollektorda (Qushchi, Nayman, Gʻarbiy, ND-2-1,
Shoʻr qishloq, Ambarquloq, Marhamat, ND-7) narx 29--59 mln soʻm/km; bu
yengil tozalash ishlariga mos keladi. Qolgan 45 ta kollektorda narx
asosan 270--680 mln soʻm/km oraligʻida; bu mexanizmlar bilan
chuqurlashtirish va tiklash ishlariga mos keladi. Ikki guruhning
oʻrtacha narxi taxminan oʻn baravar farq qiladi (34 va 368 mln soʻm/km).
Eng yuqori narx ZKD-6-1-1 kollektorida qayd etilgan: 1 335 mln soʻm/km.

Narx kollektorning suv sarfi va yoshiga bogʻliq emas. Biroq taʼmirlangan
uchastka uzunligi bilan aniq bogʻliqlik bor: 1 km dan qisqa
uchastkalarda narx medianasi 586 mln soʻm/km, 1--3 km uchastkalarda 348
mln soʻm/km, 3 km dan uzunlarida 289 mln soʻm/km. Uchastka uzunligi ikki
baravar oshganda 1 km narxi taxminan 28 foizga kamayadi (p = 0,001).
Sababi --- ish hajmiga bogʻliq boʻlmagan doimiy xarajatlar (texnikani
tashish, tayyorgarlik ishlari). Shu sababli qoʻshni kichik
kollektorlarni bitta shartnoma doirasida birgalikda taʼmirlash xarajatni
kamaytiradi.

\textbf{Balans qiymati va xizmat maydoniga nisbat.} Taʼmirlangan
kollektorlarda taʼmir xarajati (oʻsha yillarning narxlarida) kadastrdagi
balans qiymatining oʻrtacha 33 foizini tashkil etadi, yaʼni bitta taʼmir
kollektor qiymatining uchdan biriga teng; bunday nisbat taʼmirning
koʻpchilik hollarda yengil tozalash emas, kapital tiklash xususiyatida
boʻlganini tasdiqlaydi. Balans qiymati qayta baholanmagan boʻlishi
mumkinligi sababli bu nisbat taxminiy. Xizmat maydoniga nisbatan taʼmir
xarajati 53 ta kollektor boʻyicha (7 909 ga) oʻrtacha 6,09 mln soʻm/ga,
kollektorlar boʻyicha mediana 9,28 mln soʻm/ga.

\textbf{Xarajatlarning massivlar boʻyicha taqsimoti.} Kadastr
xarajatlarni massivlar boʻyicha koʻrsatmaydi, shuning uchun taqsimot GAT
bazasi orqali tiklandi: har bir kollektorning xarajati uning xaritadagi
boʻlaklari massivlar boʻyicha uzunligiga mutanosib taqsimlandi (4.1-§).
Bu usul 36 ta kollektor va umumiy xarajatning 80 foizi (38 582 mln soʻm)
uchun bajarildi. Natijalar 51-ilova va 4.4-rasmda III bobda aniqlangan
meliorativ xavf koʻrsatkichlari bilan birga keltirilgan.

\oakfig{figures/map7_repair_cost.png}{4.4-rasm. Kollektorlar taʼmir xarajatlari (2017--2020-yillar): a --- 1
ga sugʻoriladigan maydonga xarajat va taʼmirlangan kollektorlar taʼmir
turi boʻyicha; b --- meliorativ xavf (yozda sath 1,5 m dan sayoz boʻlish
ehtimoli 0,3 dan katta maydon ulushi, 3.3-§) va xarajat (doira yuzasi)}

Taʼmir xarajatlari massivlar orasida 14 baravar farq qiladi: 1 ga
sugʻoriladigan maydonga hisoblaganda Istiqlol massivida 3 648 ming soʻm,
M.Toʻjiboyev massivida 261 ming soʻm; Y.Oxunboboyev massivida kollektor
yoʻq. Xarajat va meliorativ xavf oʻrtasida ishonchli bogʻlanish yoʻq
(yettita massiv; p = 0,29). Bu nomuvofiqlik xaritada yaqqol koʻrinadi
(4.4-rasm, b): eng katta doiralar xavf oʻrtacha boʻlgan Istiqlol va xavf
past (19 foiz) boʻlgan U.Qoraboyev massivlarida, xavf eng yuqori Mashʼal
va Z.M.Bobur massivlarida esa xarajat oʻrtacha darajada. Xaritadagi 36
ta taʼmirlangan kollektordan 34 tasida narx kapital taʼmirga mos keladi
(4.4-rasm, a). Z.M.Bobur massivi meliorativ xavf boʻyicha birinchi
oʻrinlarda (sayoz sath ehtimoli yuqori maydon 58 foiz, shoʻrlanish
xavfli quduqlar 17 ta, sath koʻtarilmoqda), lekin unga toʻgʻri kelgan
xarajat tuman oʻrtachasidan (1 867 ming soʻm/ga) past: 1 793 ming
soʻm/ga. U.Qoraboyev massivida esa xavf past (19 foiz, 1 ta quduq),
xarajat esa ikkinchi oʻrinda (2 715 ming soʻm/ga). 4.1-§ dagi kabi bu
yerda ham taʼmir mablagʻlari kollektorlarning holatiga ham, hududning
meliorativ ehtiyojiga ham bogʻliq holda taqsimlanmagan.

\textbf{Amaliy natijalar.} 1 km taʼmirning 2026-yil narxlaridagi haqiqiy
narxi (mediana 348 mln soʻm, uchastka uzunligiga qarab 289--586 mln
soʻm) rejalashtirish hisoblari uchun asos boʻladi. Yengil tozalash va
kapital taʼmir narxi taxminan oʻn baravar farq qiladi; 4.2-§ ga koʻra
oʻzandagi oʻsimlik va loyqa oʻtkazish qobiliyatini 29--61 foizga
kamaytiradi, shuning uchun koʻp hollarda arzon tozalash ishlari
oʻtkazish qobiliyatining katta qismini tiklashi mumkin. Taʼmir
mablagʻlarini taqsimlashda meliorativ xavf koʻrsatkichlarini (III bob)
hisobga olish zarur: 2017--2020-yillardagi xarajatlar xavf yuqori
massivlarga ustuvor yoʻnaltirilmagan. Biroq qaysi tur va qaysi joydagi
taʼmir qancha meliorativ natija berishi oʻlchanmagan ekan (4.3-§), bu
tanlovni miqdoriy asoslab boʻlmaydi.

\textbf{Cheklov.} Kadastrda taʼmir turi faqat uchta kollektor uchun
yozilgan; yengil tozalash va kapital taʼmirga ajratish narx boʻyicha
qilingan taxmindir. Xarajatlar 2026-yil narxlariga isteʼmol narxlari
indeksi boʻyicha keltirilgan; qurilish ishlari narxi isteʼmol
narxlaridan farqli surʼatda oʻzgargan boʻlishi mumkin, shuning uchun
2026-yil narxlaridagi qiymatlar taxminiy. Xarajatlarni massivlarga
taqsimlash xaritadagi nomlangan boʻlaklarga tayanadi va xarajatning 20
foizi taqsimlanmagan.

\oakpar{4.5-§. Monitoring tizimini takomillashtirish va amaliy tavsiyalar}

III bob va 4.1--4.4-§ natijalariga koʻra meliorativ tizim
samaradorligini taʼminlash uchun avvalo uni oʻlchash tizimini yoʻlga
qoʻyish kerak. Hozirgi monitoring sizot suvi sathini qayd etadi, lekin
maʼlumot sifati nazorat qilinmaydi, tarmoq holati va taʼmirlar esa sath
bilan bogʻlanmagan. Natijada 2026-yil narxlarida 48,2 mlrd soʻmlik
taʼmir ishlarining meliorativ natijasini baholab boʻlmaydi (4.3-§),
mablagʻlar esa meliorativ ehtiyojga mos taqsimlanmaydi (4.4-§). Shu
uzilishni bartaraf etish uchun 15 ta tavsiya ishlab chiqildi; ularning
toʻliq roʻyxati va har birining tadqiqotdagi asosi 52-ilovada
keltirilgan. Tavsiyalar uch guruhga boʻlingan.

Maʼlumotlar sifati boʻyicha (T-1 -- T-4): monitoring jadvallarini qabul
qilishda tekshiruv tartibini majburiy qoʻllash, quruq quduq holatini
qiymat oʻrniga maxsus belgi bilan qayd etish, kuzatuv tarmogʻidagi
oʻzgarishlarni sanasi bilan hujjatlashtirish va hisobotlarda doimiy
kuzatilgan quduqlar boʻyicha oʻrtachani va yozgi eng sayoz sathni
alohida koʻrsatish.

Kuzatuv tarmogʻi boʻyicha (T-5 -- T-9): sath chuqur hududlarda
quduqlarni 5--6 m gacha chuqurlashtirish, yangi quduqlarni bir tekis
emas, maqsadli joylashtirish, meliorativ holatni sayoz sath ehtimoli
xaritasi bilan baholash, minerallashuv oʻlchovlarini tiklash va asosiy
kollektorlar oxirida sarfni oʻlchash.

Tizimni boshqarish boʻyicha (T-10 -- T-15): taʼmir va nosozliklar
daftarini yuritish, taʼmir navbatini kollektor holati va meliorativ xavf
boʻyicha aniqlash, qoʻshni kollektorlarni birgalikda taʼmirlash va
taʼmir turini oʻzan holatiga qarab tanlash, taʼmir rejasini monitoring
bilan muvofiqlashtirish, ichki zovurlarni inventarizatsiya qilish va
sugʻorish samaradorligini oshirish.

\textbf{Ustuvor massivlar.} III bob va 4.1--4.4-§ koʻrsatkichlari
massivlar boʻyicha birlashtirildi (4.3-jadval): meliorativ xavf (yozgi
sayoz sath ehtimoli, shoʻrlanish xavfi, sathning oʻzgarishi), tarmoq
holati va 2017--2020-yillardagi xarajatlar (2026-yil narxlarida). Xavf
yuqori deb olindi, agar ehtimoli 0,3 dan katta maydon 50 foizdan ortiq
yoki shoʻrlanish xavfli quduqlar 10 tadan koʻp boʻlsa. I guruh: xavf
yuqori, tarmoq zaif yoki oʻrta; II guruh: xavf yuqori, tarmoq yaxshi;
III guruh: xavf past; maʼlumotlarining yarmidan koʻpi yaroqsiz massiv
alohida (monitoring) guruhda. Chegaralar oʻzgartirilganda (ehtimollik
29--54 foiz, quduqlar soni 5--11) va krigingning boshqa turlarida
(29-ilova) ham guruhlar tarkibi oʻzgarmaydi. Kritik chuqurlik 2,0 m
olinganda Paxtakor massivi ehtimollik boʻyicha yuqori xavf chegarasidan
oʻtadi, lekin shoʻrlanish xavfli quduqlar soni (4 ta) boʻyicha past
qoladi.

\oakjadvalnum{4.3-jadval}

\oakilovanom{\textbf{Massivlar boʻyicha meliorativ xavf, drenaj tarmogʻi holati va
ustuvorlik}}

{\fontsize{9}{11}\selectfont\setlength{\tabcolsep}{2.5pt}

\noindent\begin{longtable}[]{|
  >{\raggedright\arraybackslash}p{(\columnwidth - 14\tabcolsep - 8\arrayrulewidth) * \real{0.1639}}|>{\raggedright\arraybackslash}p{(\columnwidth - 14\tabcolsep - 8\arrayrulewidth) * \real{0.1029}}|>{\raggedright\arraybackslash}p{(\columnwidth - 14\tabcolsep - 8\arrayrulewidth) * \real{0.1377}}|>{\raggedright\arraybackslash}p{(\columnwidth - 14\tabcolsep - 8\arrayrulewidth) * \real{0.1685}}|>{\raggedright\arraybackslash}p{(\columnwidth - 14\tabcolsep - 8\arrayrulewidth) * \real{0.1630}}|>{\raggedright\arraybackslash}p{(\columnwidth - 14\tabcolsep - 8\arrayrulewidth) * \real{0.1380}}|>{\raggedright\arraybackslash}p{(\columnwidth - 14\tabcolsep - 8\arrayrulewidth) * \real{0.1261}}|}
\hline
\begin{minipage}[b]{\linewidth}\raggedright
\textbf{Massiv}
\end{minipage} & \begin{minipage}[b]{\linewidth}\raggedright
\textbf{Yozgi sayoz sath ehtimoli 0,3 dan katta maydon, \%}
\end{minipage} & \begin{minipage}[b]{\linewidth}\raggedright
\textbf{Shoʻrlanish xavfli quduqlar}
\end{minipage} & \begin{minipage}[b]{\linewidth}\raggedright
\textbf{Sathning oʻzgarishi, m/yil}
\end{minipage} & \begin{minipage}[b]{\linewidth}\raggedright
\textbf{Tarmoq holati (4.1-§)}
\end{minipage} & \begin{minipage}[b]{\linewidth}\raggedright
\textbf{Taʼmir xarajati, ming soʻm/ga (2026-yil narxlarida)}
\end{minipage} & \begin{minipage}[b]{\linewidth}\raggedright
\textbf{Ustuvorlik}
\end{minipage} \\
\hline
\endhead
\hline
\endlastfoot
Z.M.Bobur & 58 & 17 & −0,031 (koʻtarilish) & zaif & 1 793 & I \\ \hline
Istiqlol & 55 & 12 & +0,047 & oʻrta: zich, lekin eski, ichki
zovurlarning 74 \% taʼmirsiz & 3 648 & I \\ \hline
Mashʼal & 71 & 17 & +0,031 & yaxshi & 2 424 & II \\ \hline
U.Qoraboyev & 19 & 1 & +0,053 & oʻrta & 2 715 & III \\ \hline
T.Mirzayev & 28 & 1 & +0,051 & oʻrta & 1 639 & III \\ \hline
Paxtakor & 21 & 4 & +0,041 & oʻrta & 1 461 & III \\ \hline
M.Toʻjiboyev & 17 & 4 & +0,040 & zaif (siyrak) & 261 & III \\ \hline
Y.Oxunboboyev & 0 & --- & --- & zaif (kollektor yoʻq) & 0 &
monitoring \\
\end{longtable}

}

Birinchi ustuvorlik guruhiga Z.M.Bobur va Istiqlol massivlari kiradi.
Ularda meliorativ xavf yuqori, drenaj tarmogʻi holati esa zaif yoki
oʻrta. Z.M.Bobur massivi alohida: bu tumandagi sath koʻtarilayotgan
yagona massiv va unga ajratilgan taʼmir mablagʻlari tuman oʻrtachasidan
past. Bu massivlarda birinchi navbatda kollektorlar va ichki
zovurlarning texnik holatini tekshirish (T-10), sarfni oʻlchash (T-9) va
zarur boʻlsa taʼmir oʻtkazish tavsiya etiladi. Ikkinchi guruhga Mashʼal
massivi kiradi: xavf yuqori, lekin tarmoq tumandagi eng yaxshi holatda.
Bu yerda tarmoqni kengaytirish emas, sugʻorish samaradorligini oshirish
(T-15) va minerallashuv monitoringini tiklash (T-8) ustuvor. Uchinchi
guruhdagi massivlarda sath chuqurlashmoqda va xavf past; ularda mavjud
tizimni saqlash yetarli, U.Qoraboyev massivida esa yozgi koʻtarilish
katta (0,85 m) boʻlgani uchun yozgi sathni alohida kuzatish kerak.
Y.Oxunboboyev massivi boʻyicha maʼlumotlarning yarmi yaroqsiz, shuning
uchun bu yerda avval kuzatuv quduqlarini chuqurlashtirish (T-5) kerak.

\textbf{Tuman uchun tayyorlangan tavsiyanoma bilan muvofiqlik.}
Dissertatsiya natijalaridan Marhamat tumani uchun tayyorlangan
tavsiyanomada foydalanildi {[}33{]}. Tavsiyanomada massivlar meliorativ
ekspeditsiya hisobotlari, tuproq va relyef koʻrsatkichlari boʻyicha ball
bilan baholangan. Ikkala baho ham Mashʼal, Z.M.Bobur va Istiqlol
massivlarini xavfli deb koʻrsatadi. Farq Paxtakor massivida:
tavsiyanomada u 2024-yilda minerallashgan sizot suvli yerlar koʻpaygani
uchun birinchi navbatga kiritilgan, quduqlardagi sath koʻrsatkichlari
esa uni III guruhga qoʻyadi. Mashʼal massivi uchun tavsiyanomadagi
zovurlarni tozalash va suvni meʼyorlash 4.3-jadvaldagi xulosaga zid
emas: bu tarmoqni kengaytirish emas, mavjud tarmoqni saqlash. Istiqlol
massivida 2026-yil noyabrida ZKD-6-2 zovurini tozalash boshlandi; uning
samarasi №76, №67 va №32a kuzatuv quduqlarida 2027-yil vegetatsiya
davridagi sathni 2024--2026-yillar bilan solishtirib baholanadi. Bu T-13
tavsiyasining amaldagi namunasi. Ikki hujjatning solishtiruvi
55-ilovada, tavsiyanoma jadvallari va rasmlari 57--74-ilovalarda.

\textbf{Ustuvorlik va chora-tadbirlar xaritasi.} Ustuvorlik guruhlari va
birinchi navbatdagi chora-tadbirlar joylashuvi 4.5-rasmda
umumlashtirilgan. Chora-tadbirlar ham hududiy jihatdan ajraladi.
Minerallashuvni oʻlchash kerak boʻlgan shoʻrlanish xavfli 56 ta quduq
(T-8) asosan tumanning shimoliy va gʻarbiy qismida: Z.M.Bobur va Mashʼal
massivlarida 17 tadan, Istiqlolda 12 ta. Tubigacha zaxirasi 0,5 m dan
kam boʻlgan va chuqurlashtirilishi kerak 31 ta quduq (T-5) esa asosan
janubi-sharqda: M.Toʻjiboyevda 8 ta, T.Mirzayevda 6 ta, U.Qoraboyevda 5
ta. Talabni qoplamaydigan 19 ta kollektor (xaritada 27,7 km) texnik
holati va sarfi birinchi navbatda tekshirilishi kerak boʻlgan
obyektlardir (T-9, T-10). Ichki zovurlar inventarizatsiyasi (T-14)
taʼmirsiz zovurlar ulushi 70 foizdan yuqori boʻlgan beshta massivda
ustuvor: Z.M.Bobur, Istiqlol, Paxtakor, U.Qoraboyev va Y.Oxunboboyev.
Bunday xarita meliorativ ekspeditsiya va monitoring xizmatining yillik
ish rejasi uchun asos boʻla oladi; u har yili yangi maʼlumotlar bilan
qayta hisoblanadi.

\oakfig{figures/map8_priority.png}{4.5-rasm. Massivlarning ustuvorlik guruhlari (a) va birinchi navbatdagi
chora-tadbirlar joylashuvi (b); T-5, T-8, T-9, T-10, T-14 ---
52-ilovadagi tavsiyalar}

\textbf{Taʼmir samarasini aniqlash uchun zarur quduqlar soni.} 4.3-§
dagi bahoning standart xatosi 0,105 m, taʼmirlangan kollektorlar yonida
36 ta quduq. Bunday tarmoqda ishonchli (80 foiz ehtimol bilan)
aniqlanadigan eng kichik taʼsir standart xatoning 2,8 baravari, yaʼni
0,29 m. Sinov hisobi (2.5-§; 53-ilova) buni tasdiqlaydi: hozir 0,30 m
lik taʼsir 82 foiz, 0,20 m lik taʼsir 47 foiz ehtimol bilan aniqlanadi.
Oddiy formula zarur quduqlar sonini kam koʻrsatadi (0,20 m uchun 78 ta,
bunda ehtimol 75 foiz); 80 foizga taxminan 95 ta quduq bilan erishiladi.

Tumanda kollektorlardan 300 m gacha masofada 108 ta quduq bor, lekin
ulardan faqat 36 tasi nomi va taʼmir yili aniqlangan kollektor yonida.
0,20 m lik taʼsirni aniqlash uchun yangi quduq qurish shart emas: barcha
kollektorlar boʻyicha taʼmir daftarini yuritish (T-10) va xaritada
kollektor nomlarini toʻliq koʻrsatish yetarli boʻlishi mumkin. 0,10 m
lik taʼsirni aniqlash esa mavjud tarmoq bilan imkonsiz.

Bu hisob taxminiy: u quduqlar bir-biridan mustaqil va taʼsir mahalliy
deb oladi. Taʼsir hududiy boʻlsa (4.3-§), quduqlar sonini oshirish
yetarli emas; bunday holda taʼmirni bosqichma-bosqich, tumanning bir
qismida oldin, boshqa qismida keyin oʻtkazish va shu ikki qismni
solishtirish kerak boʻladi.

\textbf{Monitoring tizimining GAT asosidagi tuzilmasi.} Tavsiyalarni
amalga oshirish uchun dissertatsiyada yaratilgan GAT bazasi (2.3-§)
monitoring tizimining asosi sifatida taklif etiladi. Unda toʻrtta oʻzaro
bogʻlangan qatlam boʻlishi kerak: kuzatuv quduqlari (pasport
maʼlumotlari, oʻlchovlar va har bir oʻlchovning sifat belgisi bilan),
kollektor va zovurlar (kadastr koʻrsatkichlari, taʼmir va nosozliklar
daftari, oʻzan holati bilan), massivlar (maydon, ekinlar tarkibi,
meliorativ xavf koʻrsatkichlari bilan) va kanallar. Har yili yangi
oʻlchovlar kiritilganda tekshiruv (T-1) avtomatik bajariladi, doimiy
kuzatilgan quduqlar boʻyicha oʻzgarish, sayoz sath ehtimoli xaritasi va
4.3-jadvaldagi koʻrsatkichlar qayta hisoblanadi. Barcha hisoblar ochiq
dasturda (Python) bajarilgan va takrorlanadigan dasturlar shaklida
saqlangan, shuning uchun ularni boshqa tumanlarga ham qoʻllash mumkin;
3.6-§ dagi viloyat tahlili buning amaliy namunasidir.

\oakpar{4.6-§. Iqtisodiy samaradorlik}

Tavsiyalarning iqtisodiy maʼnosi sarflanayotgan mablagʻ bilan uning
natijasini oʻlchash xarajati orasidagi nisbatda koʻrinadi. Tumanda
birgina 2017--2020-yillarda 53 ta kollektorni taʼmirlashga 2026-yil
narxlarida 48,2 mlrd soʻm (yiliga oʻrtacha 12,0 mlrd soʻm) sarflangan.
Bu mablagʻning meliorativ natijasi oʻlchanmagan (4.3-§), uning massivlar
boʻyicha taqsimoti esa meliorativ ehtiyojga mos kelmagan (4.4-§).

\textbf{Tavsiyalarni joriy qilish xarajatlari.} 52-ilovadagi 15 ta
tavsiyaning toʻqqiztasi (T-1, T-2, T-3, T-4, T-7, T-10, T-11, T-12,
T-13) qoʻshimcha mablagʻ talab qilmaydi: ular hisobot yuritish,
maʼlumotlarni qayta ishlash va taʼmirni rejalashtirish tartibini
oʻzgartiradi. Tekshiruv tartibi dasturiy ravishda amalga oshirilgan va
bir tuman maʼlumotini bir necha daqiqada tekshiradi; 3.6-§ da u
viloyatning 14 ta tumani boʻyicha 854 676 ta oʻlchovga qoʻllanildi.
Mablagʻ talab qiladigan tavsiyalar: kuzatuv quduqlarini chuqurlashtirish
va maqsadli joylashtirish (T-5, T-6), minerallashuvni oʻlchash (T-8),
kollektorlarda sarfni oʻlchash (T-9), ichki zovurlarni inventarizatsiya
qilish (T-14) va sugʻorish samaradorligini oshirish (T-15). Ularning
qiymati quduq qurish va chuqurlashtirish, laboratoriya tahlili,
gidrometrik oʻlchov va inventarizatsiya ishlarining smetasi asosida
aniqlanadi; bu smetalar ushbu ishda tuzilmagan.

\textbf{Taʼmir xarajatlarini kamaytirish imkoniyatlari.} Mavjud
maʼlumotlar ikki yoʻnalishdagi tejash imkoniyatini miqdoriy baholashga
imkon beradi.

Birinchisi hajm samarasi bilan bogʻliq (4.4-§). 2017--2020-yillarda 1 km
dan qisqa 12 ta uchastka (jami 7,1 km) alohida taʼmirlangan va ularga
2026-yil narxlarida 3 593 mln soʻm sarflangan. Agar ular qoʻshni
kollektorlar bilan birlashtirilib, 1--3 km li uchastkalar narxida (354
mln soʻm/km) taʼmirlanganida, xarajat 2 515 mln soʻmni tashkil etar va 1
078 mln soʻm (30 foiz) tejalar edi.

Ikkinchisi taʼmir turini oʻzan holatiga qarab tanlashdir (4.2, 4.4-§).
Kapital taʼmir sifatida baholangan 45 ta kollektorda (127,8 km) oʻrtacha
narx 368 mln soʻm/km, yengil tozalashda (8 ta kollektor, 35,5 km) 34 mln
soʻm/km. 4.2-§ ga koʻra oʻzanning oʻsimlik bilan qoplanishi oʻtkazish
qobiliyatini 29 foizga kamaytiradi va bu yoʻqotish tozalash bilan toʻliq
tiklanadi; kapital taʼmir esa asosan loyqa qatlami qalin boʻlganda
zarur. Bu hisob taxminiy: kapital taʼmirlangan kollektorlarning
qanchasida oʻzan faqat oʻsimlik bilan qoplangani nomaʼlum, shuning uchun
ulush shartli ravishda 25 va 50 foiz olindi. Agar texnik holat daftari
(T-10) asosida kapital taʼmir uzunligining 25 foizi tozalash bilan
almashtirilsa, 10,7 mlrd soʻm, 50 foizi almashtirilsa 21,3 mlrd soʻm
tejaladi (2026-yil narxlarida); bu 2017--2020-yillardagi umumiy
xarajatning 22--44 foiziga teng. Bu koʻrsatkichlar taxminiy: qaysi
kollektorlarda tozalash yetarli ekanini faqat oʻzan holatini qayd etish
orqali aniqlash mumkin. Ammo ular texnik holat daftarini yuritish
xarajati tejaladigan mablagʻga nisbatan juda kichik ekanini koʻrsatadi.

\textbf{Notoʻgʻri qarorlarning oldini olish.} Iqtisodiy samara notoʻgʻri
maʼlumot asosidagi qarorlarning oldini olish hisobiga ham yuzaga keladi.
Tekshirilmagan maʼlumot sathning oʻzgarish tezligini 1,4 baravar oshirib
koʻrsatadi (3.2-§); kriging xaritasi sayoz sathli maydonni oʻn
baravardan ortiq kam koʻrsatadi (3.3-§); bir yildagi joylarni
solishtirish esa tarmogʻi zich hududlarni «yomon» deb koʻrsatadi
(3.4-§). Bunday maʼlumotlar asosida mablagʻ taqsimlansa, u Z.M.Bobur
massividagi kabi haqiqiy ehtiyoj boʻlgan joyga yetib bormaydi.

Samarani soʻmda toʻliq ifodalash uchun taʼmir turlarining haqiqiy
meliorativ natijasi maʼlum boʻlishi kerak; u esa aynan taklif etilgan
monitoring joriy qilingandan keyin oʻlchanadi. Shuning uchun bu
bosqichda faqat xarajatlarni kamaytirish imkoniyati baholandi.

\oakpar{IV bob boʻyicha xulosalar}

\begin{enumerate}
\def\labelenumi{\arabic{enumi}.}
\tightlist
\item
  Meliorativ kadastr maʼlumotlari ishonchlilikning odatdagi
  koʻrsatkichlarini (buzilmasdan ishlash ehtimoli, buzilishlar soni)
  hisoblashga yetmaydi, chunki buzilishlar va taʼmirlar tarixi qayd
  etilmaydi.
\item
  Tumandagi 90 ta kollektorning (265,8 km) yoshi medianasi 61 yil,
  uzunlikning 90 foizi 50 yildan eski; barcha kollektorlar
  2010--2020-yillarda taʼmirlangan, uzunlikning 71 foizi
  2016--2019-yillarda. Taʼmir yili kollektor yoshi, uzunligi va sarfiga
  bogʻliq emas: taʼmir holatga qarab emas, dasturlar boʻyicha
  oʻtkazilgan. Xoʻjalik ichki zovurlarining 64,7 foizida taʼmir qayd
  etilmagan.
\item
  Ochiq tarmoqning solishtirma uzunligi massivlar boʻyicha 6,8 m/ga dan
  48,6 m/ga gacha; eng zaif tarmoq Y.Oxunboboyev, M.Toʻjiboyev va
  Z.M.Bobur massivlarida, eng yaxshisi Mashʼal massivida.
\item
  Suv balansi boʻyicha iyuldagi talab etiladigan drenaj moduli
  0,16--0,66 l/(s·ga), oʻrta variantda 0,26--0,43 l/(s·ga); sugʻorish
  samaradorligi oshirilganda u 2,5 baravar kamayadi. Kadastrdagi sarf
  qiymatlari toʻgʻri deb qabul qilinganda 90 ta kollektordan 27--48
  tasining (variantga qarab) sarfi talabdan past.
\item
  Kollektor taʼmiridan keyin sath uzoqdagi quduqlarga nisbatan
  chuqurlashmagan (−0,12 m; 95 foizlik ishonch oraligʻi −0,33 \ldots{}
  +0,09 m; boshqa tekshiruvlar ham shunday), lekin 0,2 m atrofidagi
  taʼsir istisno qilinmaydi: uni aniqlash uchun taʼmirlangan
  kollektorlar yonida taxminan 95 ta quduq kerak (hozir 36 ta). Bu
  taʼmir samarasizligini emas, mavjud monitoring bilan uni oʻlchab
  boʻlmasligini koʻrsatadi.
\item
  Taʼmirning 2026-yil narxlaridagi haqiqiy narxi medianasi 348 mln
  soʻm/km; yengil tozalash va kapital taʼmir orasida taxminan oʻn
  baravar farq bor. 2017--2020-yillarda sarflangan 48,2 mlrd soʻmning
  massivlar boʻyicha taqsimoti meliorativ xavf bilan bogʻlanmagan: xavf
  eng yuqori Z.M.Bobur massiviga tuman oʻrtachasidan kam mablagʻ
  ajratilgan. Kapital taʼmirning bir qismini tozalash bilan almashtirib,
  22--44 foiz tejash mumkin.
\item
  Monitoring tizimini takomillashtirish boʻyicha 15 ta tavsiya ishlab
  chiqildi (toʻqqiztasi mablagʻ talab qilmaydi).
\end{enumerate}

\oakbob{UMUMIY XULOSALAR}

«GAT texnologiyalari asosida meliorativ tizimlar samaradorligi va
ishonchliligini taʼminlash (Andijon viloyati Marhamat tumani misolida)»
mavzusidagi dissertatsiya boʻyicha olib borilgan tadqiqotlar natijasida
quyidagi xulosalar taqdim etildi:

\begin{enumerate}
\def\labelenumi{\arabic{enumi}.}
\tightlist
\item
  Meliorativ tizimlarni baholashning mavjud usullari kirish
  maʼlumotlarini toʻgʻri deb oladi, drenaj tarmogʻining joylashuvini kam
  hisobga oladi va taʼmir samarasini bajarilgan ish hajmi bilan
  oʻlchaydi. Shuning uchun tadqiqotda asosiy qoida belgilandi: tizim
  ishonchliligini baholashdan oldin maʼlumotlarning ishonchliligini
  baholash kerak.
\item
  Mavjud xaritalar va kuzatuv quduqlarining koordinatalari asosida
  drenaj tarmogʻining GAT bazasini tiklash usuli ishlab chiqildi.
  Marhamat tumani xaritasi 181 ta quduq orqali koordinataga bogʻlandi
  (oʻrtacha xato 1,9 m). 214,6 km kollektor va 170,9 km kanal
  raqamlandi, sakkizta massiv chegarasi tiklandi, 60 ta kollektorga
  kadastr maʼlumotlari biriktirildi. Baza kollektor koʻrsatkichlari va
  taʼmir xarajatlarini massivlar boʻyicha taqsimlash imkonini berdi.
\item
  Monitoring maʼlumotlarini tekshirishning uchta qoidasi ishlab
  chiqildi: yil ichidagi oʻzgarmaslik, oʻlchov birligi va kuzatuv
  tarmogʻi tarkibi; quduq qurigan holatlar alohida belgilanadi.
  2009--2024-yillardagi 73 868 ta sath oʻlchovining 87,7 foizi yaroqli,
  ularning 8,2 foizi quduq qurigan holatni koʻrsatadi. Nuqsonlar
  tasodifiy emas: sathi sayoz massivlarda yaroqli oʻlchovlar 97--98
  foiz, sathi chuqur massivlarda 50--60 foiz. Bu qonuniyat Andijon
  viloyatining 14 ta tumanida ham (1 972 ta quduq, 854 676 ta oʻlchov)
  tasdiqlandi. Sababi --- kuzatuv quduqlari sath chuqur joylarda
  yetarlicha chuqur emas. 2025-yil bahorida qiymati oʻzgarmagan quduqlar
  ulushi 32,4 foizga yetdi, yaʼni tekshiruv yangi maʼlumotlar uchun ham
  zarur.
\item
  Doimiy kuzatilgan 84 ta quduq boʻyicha sizot suvi sathi 2009-yildagi
  1,60 m dan 2024-yilda 1,99 m gacha chuqurlashgan, yiliga taxminan 3
  sm. Oʻzgarish asosan 2016-yildan keyin bir martalik pogʻona
  koʻrinishida boʻlgan. Tekshiruvsiz hisob bu tezlikni 1,4 baravar
  oshirib koʻrsatadi. Sath viloyatning 14 ta tumanidan 12 tasida
  chuqurlashgan; Marhamatdagi tezlik viloyat oʻrtachasiga teng.
\item
  Tuman oʻrtachasi massivlardagi qarama-qarshi oʻzgarishlarni yashiradi:
  oltita massivda sath yiliga 3,1--5,3 sm chuqurlashgan, Z.M.Bobur
  massivida esa yiliga 3,1 sm koʻtarilib, 2019-yildan keyin 1,27--1,71 m
  ga yetgan. Sath yozda, sugʻorish mavsumida eng sayoz boʻladi; yil
  davomidagi tebranish tuman boʻyicha 0,48 m, massivlarda 0,32--0,85 m.
  Shuning uchun drenaj tizimini yillik oʻrtacha sath bilan emas, yozgi
  eng sayoz sath bilan baholash kerak.
\item
  Sizot suvi sathi qoʻshni quduqlar orasida ham keskin farq qiladi,
  shuning uchun oddiy sath xaritasining (kriging) xatosi katta --- 0,82
  m. Bu xarita sayoz sathli maydonni oʻn baravardan ortiq kam
  koʻrsatadi: yozda quduqlarning 23 foizida sath 1,5 m dan sayoz,
  xaritada esa maydonning atigi 1,1 foizida. Shuning uchun meliorativ
  holatni sathning kritik chuqurlikdan sayoz boʻlish ehtimoli xaritasi
  bilan baholash taklif etildi. Quduqlar sonini 94 tadan oshirish xarita
  aniqligini oshirmaydi.
\item
  Sizot suvi kollektorlarga yaqin va tarmoq zich joylarda sayozroq;
  bogʻlanish taxminan 850 m gacha kuzatiladi va massivlar ichida ham
  saqlanadi. Buning sababi, ehtimol, kollektorlar tabiatan sayoz sathli
  joylarga qurilgani. Shuning uchun bir yildagi turli joylarni
  solishtirib drenaj samarasini baholab boʻlmaydi. Kollektordan 800 m
  dan uzoqdagi quduqlarda sath kollektor yaqinidagiga qaraganda tezroq
  chuqurlashgan (yiliga 6,3 va 1,7 sm).
\item
  Sizot suvlari minerallashuvi boʻyicha faqat 2004--2010-yillar
  maʼlumotlari ishonchli. Bu davrda massivlar orasidagi farq uch
  baravarga yetadi: Mashʼalda 2,47 g/l, T.Mirzayev va U.Qoraboyevda
  0,80--0,84 g/l. Yozgi sath va minerallashuv birga olinganda 113 ta
  quduqdan 56 tasi shoʻrlanish xavfi ostida, ularning 82 foizi Mashʼal,
  Z.M.Bobur va Istiqlol massivlarida. Minerallashuv keyingi yillarda
  oʻzgargan boʻlishi mumkin, shuning uchun bu quduqlarda u qayta
  oʻlchanishi kerak.
\item
  Tumandagi 90 ta kollektorning (265,8 km) uzunlik boʻyicha 90 foizi 50
  yildan eski. Barcha kollektorlar 2010--2020-yillarda taʼmirlangan,
  lekin taʼmir yili kollektorning yoshi, uzunligi va sarfiga bogʻliq
  emas: taʼmir holatga qarab emas, dasturlar boʻyicha oʻtkazilgan. Eng
  zaif boʻgʻin xoʻjalik ichki zovurlari: ularning 64,7 foizida taʼmir
  qayd etilmagan. Kadastr maʼlumotlari ishonchlilikning odatdagi
  koʻrsatkichlarini hisoblashga yetarli emas.
\item
  Suv balansi boʻyicha iyulda talab etiladigan drenaj moduli 0,16--0,66
  l/(s·ga), oʻrta variantda 0,26--0,43 l/(s·ga); sugʻorish samaradorligi
  oshirilsa, u 2,5 baravar kamayadi. Kadastrdagi sarflar boʻyicha 90 ta
  kollektordan 27--48 tasining sarfi talabdan past. Oʻzandagi oʻsimlik
  va loyqa kollektorning oʻtkazish qobiliyatini 29--61 foizga
  kamaytiradi.
\item
  Mavjud monitoring maʼlumotlarida kollektor taʼmiridan keyin sizot
  sathi chuqurlashgani aniqlanmadi: taʼmirlangan kollektorlar yonidagi
  quduqlarda sath taʼmirlanmaganlar yonidagiga nisbatan chuqurlashmagan
  (baho −0,12 m, ishonch oraligʻi −0,33 dan +0,09 m gacha). Boshqa
  tekshiruvlar ham shu natijani berdi. Bu taʼmir samarasiz degani emas:
  mavjud kuzatuv tarmogʻi uni oʻlchashga yetarli emas. 0,20 m lik
  taʼsirni ishonchli aniqlash uchun taʼmirlangan kollektorlar yonida
  taxminan 95 ta quduq kerak, hozir 36 ta.
\item
  1 km kollektor taʼmirining 2026-yil narxlaridagi oʻrtacha narxi 348
  mln soʻm; yengil tozalash va kapital taʼmir narxi taxminan oʻn baravar
  farq qiladi, uzun uchastkalarda 1 km arzonroq. 2017--2020-yillarda
  sarflangan 48,2 mlrd soʻm (2026-yil narxlarida) massivlar boʻyicha
  meliorativ xavfga qarab taqsimlanmagan: xavf eng yuqori Z.M.Bobur
  massiviga tuman oʻrtachasidan kam mablagʻ ajratilgan. Kapital
  taʼmirning bir qismini tozalash bilan almashtirib, xarajatlarni 22--44
  foizga kamaytirish mumkin.
\item
  Massivlar meliorativ xavf, drenaj tarmogʻi holati va xarajatlar
  boʻyicha uch guruhga ajratildi. Birinchi navbatda Z.M.Bobur va
  Istiqlol massivlarida drenaj tarmogʻining texnik holatini tekshirish
  va taʼmirlash, Mashʼal massivida esa sugʻorish samaradorligini
  oshirish va minerallashuvni oʻlchashni tiklash tavsiya etiladi. Bu
  xulosalar tumanga topshirilgan tavsiyanomadagi massiv tadbirlari bilan
  mos keladi.
\item
  Monitoring tizimini takomillashtirish boʻyicha 15 ta tavsiya ishlab
  chiqildi: maʼlumotlarni majburiy tekshirish, quruq quduqlarni maxsus
  belgilash, sath chuqur joylarda quduqlarni chuqurlashtirish, yangi
  quduqlarni kerakli joylarga qurish, sayoz sath ehtimoli xaritalaridan
  foydalanish, kollektorlarda sarfni oʻlchash, taʼmir daftarini
  yuritish, taʼmir navbatini holat va xavf boʻyicha belgilash va
  boshqalar. Ulardan toʻqqiztasi qoʻshimcha mablagʻ talab qilmaydi.
  Ularni amalga oshirishga dissertatsiyada yaratilgan GAT bazasi va
  hisob dasturlari asos boʻladi.
\end{enumerate}

\textbf{Tadqiqotni davom ettirish yoʻnalishlari.} Ayrim kuzatishlarning
boshqa mumkin boʻlgan sabablarini mavjud maʼlumotlar bilan toʻliq
ajratib boʻlmaydi (54-ilova). Shuning uchun keyingi tadqiqotlar quyidagi
yoʻnalishlarda olib borilishi kerak:

\begin{itemize}
\tightlist
\item
  2016--2018-yillardagi pogʻonali chuqurlashishga sugʻorish suvi hajmi,
  iqlim yoki oʻlchash tartibi oʻzgarishidan qaysi biri sabab boʻlganini
  yillik suv olish va meteorologik maʼlumotlar bilan aniqlash;
\item
  taʼmirni bosqichma-bosqich oʻtkazib, juft quduqlar orqali uning
  mahalliy va hududiy taʼsirini oʻlchash;
\item
  Z.M.Bobur massividagi sath koʻtarilishining sababini kollektorlarda
  sarfni oʻlchash va sunʼiy yoʻldosh maʼlumotlari bilan aniqlash;
\item
  sayoz sath ehtimoli va shoʻrlanish xavfi xaritalarini Andijon
  viloyatining barcha tumanlari uchun tuzish;
\item
  kollektorlar kesimi va sarflarini oʻlchab, gidravlik hisobni
  aniqlashtirish;
\item
  xoʻjalik ichki zovurlarining raqamli xaritasini yaratish.
\end{itemize}

\oakbob{FOYDALANILGAN ADABIYOTLAR ROʻYXATI}

\textbf{I. Normativ-huquqiy hujjatlar va meʼyoriy hujjatlar}

\begin{enumerate}
\def\labelenumi{\arabic{enumi}.}
\item
  Oʻzbekiston Respublikasi Prezidentining «``Oʻzbekiston --- 2030''
  strategiyasi toʻgʻrisida»gi PF-158-son Farmoni. -- Toshkent, 2023-yil
  11-sentabr. {[}Elektron resurs{]}. -- URL: https://lex.uz
\item
  Oʻzbekiston Respublikasi Prezidentining «Oʻzbekiston Respublikasi suv
  xoʻjaligini rivojlantirishning 2020--2030-yillarga moʻljallangan
  konsepsiyasini tasdiqlash toʻgʻrisida»gi PF-6024-son Farmoni. --
  Toshkent, 2020-yil 10-iyul. {[}Elektron resurs{]}. -- URL:
  https://lex.uz/docs/4892953 (murojaat sanasi: 16.11.2026).
\item
  Oʻzbekiston Respublikasi Prezidentining «2020--2030-yillarga
  moʻljallangan Oʻzbekiston Respublikasi qishloq xoʻjaligini
  rivojlantirish strategiyasini tasdiqlash toʻgʻrisida»gi PF-5853-son
  Farmoni. -- Toshkent, 2019-yil 23-oktabr.
\item
  Oʻzbekiston Respublikasi Prezidentining «Oʻzbekiston Respublikasida
  suv resurslarini boshqarish va irrigatsiya sektorini
  rivojlantirishning 2025--2028-yillarga moʻljallangan dasturini
  tasdiqlash toʻgʻrisida»gi PQ-250-son qarori. -- Toshkent, 2025-yil
  15-avgust. {[}Elektron resurs{]}. -- URL:
  https://www.lex.uz/uz/docs/7686130
\item
  Oʻzbekiston Respublikasi Prezidentining «2018--2019-yillarda
  irrigatsiyani rivojlantirish va sugʻoriladigan yerlarning meliorativ
  holatini yaxshilash Davlat dasturi toʻgʻrisida»gi PQ-3405-son qarori.
  -- Toshkent, 2017-yil 27-noyabr. {[}Elektron resurs{]}. -- URL:
  https://lex.uz/docs/3426211 (murojaat sanasi: 16.11.2026).
\item
  ShNQ 2.06.03-12. Sugʻorish tizimlari. Loyihalash normalari (Суғориш
  тизимлари. Лойиҳалаш нормалари). -- Toshkent: Oʻzbekiston Respublikasi
  Davlat arxitektura va qurilish qoʻmitasi, 2012. -- Yakuniy tahriri.
  («Oʻzsuvloyiha» OAJ tomonidan QMQ 2.06.03-97 oʻrniga qayta ishlangan).
\item
  ISO 19157:2013. Geographic information -- Data quality. -- Geneva:
  International Organization for Standardization, 2013. -- 146 p.
\item
  ГОСТ 27.002-2015. Надёжность в технике. Термины и определения. -- М.:
  Стандартинформ, 2016.
\end{enumerate}

\textbf{II. Ilmiy adabiyotlar}

\begin{enumerate}
\def\labelenumi{\arabic{enumi}.}
\setcounter{enumi}{8}
\item
  Аверьянов С.Ф. Борьба с засолением орошаемых земель. -- М.: Колос,
  1978. -- 288 с.
\item
  Аверьянов С.Ф. Фильтрация из каналов и её влияние на режим грунтовых
  вод. -- М.: Колос, 1982. -- 237 с.
\item
  Айдаров И.П. Регулирование водно-солевого и питательного режимов
  орошаемых земель. -- М.: Агропромиздат, 1985. -- 304 с.
\item
  Аҳмедов Х.А. Зах қочириш мелиорацияси: дарслик. -- Тошкент: «Ўқитувчи»
  нашриёти, 1975. -- 224 б. (Русча нашри: Ахмедов Х.А. Осушительные
  мелиорации. -- Ташкент: Ўқитувчи, 1974).
\item
  Голованов А.И., Айдаров И.П., Григоров М.С. и др. Мелиорация земель /
  под ред. А.И. Голованова. -- М.: КолосС, 2011. -- 824 с.
\item
  Духовный В.А., Соколов В.И. Интегрированное управление водными
  ресурсами: опыт и уроки Центральной Азии. -- Ташкент: НИЦ МКВК, 2005.
\item
  Егоров В.В. Засоление и солонцеватость почв орошаемых территорий. --
  М.: Наука, 1977.
\item
  Икрамов Р.К., Бараев Ф.А., Юсупов Г.У., Каримова Н.М. Мелиоративный
  мониторинг и кадастр орошаемых земель: учебное пособие. -- Ташкент:
  Министерство высшего и среднего специального образования Республики
  Узбекистан, 2008. -- 224 с. -- УДК 631.6.
\item
  Кац Д.М. Влияние орошения на грунтовые воды. -- М.: Колос, 1976. --
  271 с.
\item
  Ковда В.А. Основы учения о почвах. Общая теория почвообразовательного
  процесса. Кн. 2. -- М.: Наука, 1973. -- 468 с.
\item
  Ковда В.А. Проблемы борьбы с опустыниванием и засолением орошаемых
  почв. -- М.: Колос, 1984. -- 304 с.
\item
  Костяков А.Н. Основы мелиораций. -- 6-е изд. -- М.: Сельхозиздат,
  1960. -- 622 с.
\item
  Маслов Б.С., Минаев И.В., Губер К.В. Справочник по мелиорации. -- М.:
  Росагропромиздат, 1989. -- 384 с.
\item
  Минашина Н.Г. Мелиорация засолённых почв. -- М.: Колос, 1978. -- 263
  с.
\item
  Мирцхулава Ц.Е. Надёжность гидромелиоративных сооружений. -- М.:
  Колос, 1974. -- 424 с.
\item
  Мирцхулава Ц.Е. О надёжности крупных каналов. -- М.: Колос, 1981. --
  318 с.
\item
  Панкова Е.И., Айдаров И.П., Ямнова И.А. и др. Природное и
  антропогенное засоление почв бассейна Аральского моря. -- М.:
  Почвенный институт им. В.В. Докучаева, 1996. -- 187 с.
\item
  Полубаринова-Кочина П.Я. Теория движения грунтовых вод. -- 2-е изд. --
  М.: Наука, 1977. -- 664 с.
\item
  Шестаков В.М. Гидрогеодинамика. -- 3-е изд. -- М.: Изд-во МГУ, 1995.
  -- 368 с.
\item
  Akramkhanov A., Vlek P.L.G. The assessment of spatial distribution of
  soil salinity risk using neural network // Environmental Monitoring
  and Assessment. -- 2012. -- Vol. 184, No.~4. -- P. 2475--2485. --
  https://doi.org/10.1007/s10661-011-2132-5
\item
  Alexandersson H. A homogeneity test applied to precipitation data //
  Journal of Climatology. -- 1986. -- Vol. 6, No.~6. -- P. 661--675. --
  https://doi.org/10.1002/joc.3370060607
\item
  Allbed A., Kumar L. Soil salinity mapping and monitoring in arid and
  semi-arid regions using remote sensing technology: a review //
  Advances in Remote Sensing. -- 2013. -- Vol. 2, No.~4. -- P. 373--385.
  -- https://doi.org/10.4236/ars.2013.24040
\item
  Allen R.G., Pereira L.S., Raes D., Smith M. Crop Evapotranspiration:
  Guidelines for Computing Crop Water Requirements. FAO Irrigation and
  Drainage Paper 56. -- Rome: FAO, 1998. -- 300 p.
\item
  Amanov B., Ghirardelli A., Sherov A., Gafarova A., Gulomov D., Tarolli
  P. An integrated framework for agroclimatic, salinity, and groundwater
  assessment in the Jizzakh Region, Uzbekistan // iScience. -- 2026. --
  Vol. 29. -- Art. 117582. -- https://doi.org/10.1016/j.isci.2026.117582
\item
  Amanov B.T., Sherov A.G., G'ulomov D., Karabayev A., Sirochov A.
  Marhamat tumanidagi sug'oriladigan yerlarning meliorativ holatini
  o'rganish va rejimini maqbullashtirish: tavsiyanoma. -- Toshkent:
  TIQXMMI MTU, Andijon qishloq xo'jaligi va agrotexnologiyalar
  instituti, 2026.
\item
  Angrist J.D., Pischke J.-S. Mostly Harmless Econometrics: An
  Empiricist's Companion. -- Princeton: Princeton University Press,
  2009. -- 373 p.
\item
  Ayers R.S., Westcot D.W. Water Quality for Agriculture. FAO Irrigation
  and Drainage Paper 29 (Rev.~1). -- Rome: FAO, 1985. -- 174 p.
\item
  Bear J. Dynamics of Fluids in Porous Media. -- New York: American
  Elsevier, 1972. -- 764 p.
\item
  Bear J. Hydraulics of Groundwater. -- New York: McGraw-Hill, 1979. --
  567 p.
\item
  Benjamini Y., Hochberg Y. Controlling the false discovery rate: a
  practical and powerful approach to multiple testing // Journal of the
  Royal Statistical Society: Series B (Methodological). -- 1995. -- Vol.
  57, No.~1. -- P. 289--300. --
  https://doi.org/10.1111/j.2517-6161.1995.tb02031.x
\item
  Borusyak K., Jaravel X., Spiess J. Revisiting event-study designs:
  robust and efficient estimation // Review of Economic Studies. --
  2024. -- Vol. 91, No.~6. -- P. 3253--3285. --
  https://doi.org/10.1093/restud/rdae007
\item
  Bos M.G., Burton M.A., Molden D.J. Irrigation and Drainage Performance
  Assessment: Practical Guidelines. -- Wallingford: CABI Publishing,
  2005. -- 176 p.
\item
  Burrough P.A., McDonnell R.A. Principles of Geographical Information
  Systems. -- Oxford: Oxford University Press, 1998. -- 333 p.
\item
  Burrough P.A., McDonnell R.A., Lloyd C.D. Principles of Geographical
  Information Systems. -- 3rd ed.~-- Oxford: Oxford University Press,
  2015. -- 352 p.
\item
  Callaway B., Sant'Anna P.H.C. Difference-in-differences with multiple
  time periods // Journal of Econometrics. -- 2021. -- Vol. 225, No.~2.
  -- P. 200--230. -- https://doi.org/10.1016/j.jeconom.2020.12.001
\item
  Cambardella C.A., Moorman T.B., Novak J.M. et al.~Field-scale
  variability of soil properties in central Iowa soils // Soil Science
  Society of America Journal. -- 1994. -- Vol. 58, No.~5. -- P.
  1501--1511. -- https://doi.org/10.2136/sssaj1994.03615995005800050033x
\item
  Chilès J.-P., Delfiner P. Geostatistics: Modeling Spatial Uncertainty.
  -- 2nd ed.~-- Hoboken: John Wiley \& Sons, 2012. -- 699 p.
\item
  Corwin D.L., Lesch S.M. Apparent soil electrical conductivity
  measurements in agriculture // Computers and Electronics in
  Agriculture. -- 2005. -- Vol. 46, No.~1--3. -- P. 11--43. --
  https://doi.org/10.1016/j.compag.2004.10.005
\item
  Corwin D.L., Scudiero E. Review of soil salinity assessment for
  agriculture across multiple scales using proximal and/or remote
  sensors // Advances in Agronomy. -- 2019. -- Vol. 158. -- P. 1--130.
  -- https://doi.org/10.1016/bs.agron.2019.07.001
\item
  Cressie N.A.C. Statistics for Spatial Data. -- Rev.~ed.~-- New York:
  John Wiley \& Sons, 1993. -- 900 p.
\item
  Cui T., Pagendam D., Gilfedder M. Gaussian process machine learning
  and Kriging for groundwater salinity interpolation // Environmental
  Modelling \& Software. -- 2021. -- Vol. 144. -- Art. 105170. --
  https://doi.org/10.1016/j.envsoft.2021.105170
\item
  Duan Y., Ma L., Abuduwaili J., Liu W., Saparov G., Smanov Z. Driving
  factor identification for the spatial distribution of soil salinity in
  the irrigation area of the Syr Darya River, Kazakhstan // Agronomy. --
  2022. -- Vol. 12, No.~8. -- Art. 1912. --
  https://doi.org/10.3390/agronomy12081912
\item
  Dukhovny V., Umarov P., Yakubov H., Madramootoo C.A. Drainage in the
  Aral Sea Basin // Irrigation and Drainage. -- 2007. -- Vol. 56,
  No.~S1. -- P. S91--S100. -- https://doi.org/10.1002/ird.367
\item
  Ernst L.F. Groundwater flow in the saturated zone and its calculation
  when horizontal parallel open conduits are present. -- Wageningen,
  1962.
\item
  FAO. Global status of salt-affected soils -- Main report. -- Rome:
  FAO, 2024. -- https://doi.org/10.4060/cd3044en
\item
  Farifteh J., Farshad A., George R.J. Assessing salt-affected soils
  using remote sensing, solute modelling, and geophysics // Geoderma. --
  2006. -- Vol. 130, No.~3--4. -- P. 191--206. --
  https://doi.org/10.1016/j.geoderma.2005.02.003
\item
  Forkutsa I., Sommer R., Shirokova Y.I. et al.~Modeling irrigated
  cotton with shallow groundwater in the Aral Sea Basin of Uzbekistan //
  Irrigation Science. -- 2009. -- Vol. 27, No.~4. -- P. 331--346. --
  https://doi.org/10.1007/s00271-009-0148-1
\item
  Freeze R.A., Cherry J.A. Groundwater. -- Englewood Cliffs:
  Prentice-Hall, 1979. -- 604 p.
\item
  Ghassemi F., Jakeman A.J., Nix H.A. Salinisation of Land and Water
  Resources: Human Causes, Extent, Management and Case Studies. --
  Sydney: University of New South Wales Press, 1995. -- 526 p.
\item
  Goodchild M.F. Geographical information science // International
  Journal of Geographical Information Systems. -- 1992. -- Vol. 6,
  No.~1. -- P. 31--45. -- https://doi.org/10.1080/02693799208901893
\item
  Goodman-Bacon A. Difference-in-differences with variation in treatment
  timing // Journal of Econometrics. -- 2021. -- Vol. 225, No.~2. -- P.
  254--277. -- https://doi.org/10.1016/j.jeconom.2021.03.014
\item
  Goovaerts P. Geostatistics for Natural Resources Evaluation. -- New
  York: Oxford University Press, 1997. -- 483 p.
\item
  Guptill S.C., Morrison J.L. (eds.). Elements of Spatial Data Quality.
  -- Oxford: Elsevier, 1995. -- 202 p.
\item
  Hamed K.H., Rao A.R. A modified Mann-Kendall trend test for
  autocorrelated data // Journal of Hydrology. -- 1998. -- Vol. 204,
  No.~1--4. -- P. 182--196. --
  https://doi.org/10.1016/s0022-1694(97)00125-x
\item
  Hamidov A., Kasymov U., Djumaboev K., Paul C. Rebound effects in
  irrigated agriculture in Uzbekistan: a stakeholder-based assessment //
  Sustainability. -- 2022. -- Vol. 14, No.~14. -- Art. 8375. --
  https://doi.org/10.3390/su14148375
\item
  Hamidov A., Khamidov M., Ishchanov J. Impact of climate change on
  groundwater management in the northwestern part of Uzbekistan //
  Agronomy. -- 2020. -- Vol. 10, No.~8. -- Art. 1173. --
  https://doi.org/10.3390/agronomy10081173
\item
  Helsel D.R. Nondetects and Data Analysis: Statistics for Censored
  Environmental Data. -- Hoboken: John Wiley \& Sons, 2005. -- 250 p.
\item
  Helsel D.R., Hirsch R.M. Statistical Methods in Water Resources. --
  Amsterdam: Elsevier, 1992. -- 522 p.
\item
  Hengl T. A Practical Guide to Geostatistical Mapping. -- Amsterdam,
  2009. -- 291 p.
\item
  Hirsch R.M., Slack J.R., Smith R.A. Techniques of trend analysis for
  monthly water quality data // Water Resources Research. -- 1982. --
  Vol. 18, No.~1. -- P. 107--121. --
  https://doi.org/10.1029/wr018i001p00107
\item
  Hooghoudt S.B. Bijdragen tot de kennis van eenige natuurkundige
  grootheden van den grond. -- Verslagen van Landbouwkundige
  Onderzoekingen, 1940. -- No.~46. -- P. 515--707.
\item
  Ibrakhimov M., Khamzina A., Forkutsa I. et al.~Groundwater table and
  salinity: Spatial and temporal distribution and influence on soil
  salinization in Khorezm region (Uzbekistan, Aral Sea Basin) //
  Irrigation and Drainage Systems. -- 2007. -- Vol. 21, No.~3--4. -- P.
  219--236. -- https://doi.org/10.1007/s10795-007-9033-3
\item
  Isaaks E.H., Srivastava R.M. An Introduction to Applied Geostatistics.
  -- New York: Oxford University Press, 1989. -- 561 p.
\item
  Jensen J.R. Remote Sensing of the Environment: An Earth Resource
  Perspective. -- 2nd ed.~-- Upper Saddle River: Prentice Hall, 2007. --
  592 p.
\item
  Journel A.G., Huijbregts C.J. Mining Geostatistics. -- London:
  Academic Press, 1978. -- 600 p.
\item
  Kendall M.G. Rank Correlation Methods. -- 4th ed.~-- London: Charles
  Griffin, 1975. -- 202 p.
\item
  Kenjabaev Sh., Sultanov M., Bauer Ch., Abduraimova D., Akhmedkhodjaeva
  I., Ibragimova Z. Assessment of the land reclamation condition using
  GIS techniques and environmental variables: case study in Kulavat
  canal irrigation system, Khorezm // InterCarto. InterGIS. GI support
  of sustainable development of territories: Proceedings of the
  International Conference. -- Moscow: Moscow University Press, 2020. --
  Vol. 26, Part 3. -- P. 264--270. --
  https://doi.org/10.35595/2414-9179-2020-3-26-264-270
\item
  Kenjabaev Sh., Sultonov M. Estimation of groundwater contribution to
  crop water use in Kulavat irrigation command area in Khorezm,
  Uzbekistan // International Journal of Agricultural Policy and
  Research. -- 2016. -- Vol. 4, No.~11. -- P. 249--255.
\item
  Kholdorov Sh., Lakshmi G., Jabbarov Z., Yamaguchi T., Yamashita M.,
  Samatov N., Katsura K. Analysis of irrigated salt-affected soils in
  the Central Fergana Valley, Uzbekistan, using Landsat 8 and Sentinel-2
  satellite images, laboratory studies, and spectral index-based
  approaches // Eurasian Soil Science. -- 2023. -- Vol. 56, No.~8. -- P.
  1178--1189. -- https://doi.org/10.1134/s1064229323600185
\item
  Kirkham D. Theory of land drainage // Drainage of Agricultural Lands.
  -- Madison: American Society of Agronomy, 1957. -- P. 139--181.
\item
  Krige D.G. A statistical approach to some basic mine valuation
  problems on the Witwatersrand // Journal of the Chemical,
  Metallurgical and Mining Society of South Africa. -- 1951. -- Vol. 52,
  No.~6. -- P. 119--139.
\item
  Kulmatov R. Problems of sustainable use and management of water and
  land resources in Uzbekistan // Journal of Water Resource and
  Protection. -- 2014. -- Vol. 6, No.~1. -- P. 35--42. --
  https://doi.org/10.4236/jwarp.2014.61006
\item
  Kulmatov R., Alimov A., Opp C. Determination and assessment of the
  groundwater table and mineralization in irrigated areas (Xojeli
  district, Republic of Karakalpakstan) // BIO Web of Conferences. --
  2024. -- Vol. 93. -- Art. 04014. --
  https://doi.org/10.1051/bioconf/20249304014
\item
  Li J., Heap A.D. A review of comparative studies of spatial
  interpolation methods in environmental sciences: Performance and
  impact factors // Ecological Informatics. -- 2011. -- Vol. 6,
  No.~3--4. -- P. 228--241. --
  https://doi.org/10.1016/j.ecoinf.2010.12.003
\item
  Longley P.A., Goodchild M.F., Maguire D.J., Rhind D.W. Geographic
  Information Systems and Science. -- 2nd ed.~-- Chichester: John Wiley
  \& Sons, 2005. -- 517 p.
\item
  Loáiciga H.A., Charbeneau R.J., Everett L.G., Fogg G.E., Hobbs B.F.,
  Rouhani S. Review of ground-water quality monitoring network design //
  Journal of Hydraulic Engineering. -- 1992. -- Vol. 118, No.~1. -- P.
  11--37. -- https://doi.org/10.1061/(ASCE)0733-9429(1992)118:1(11)
\item
  Luthin J.N. Drainage Engineering. -- New York: John Wiley \& Sons,
  1966. -- 250 p.
\item
  Maas E.V., Hoffman G.J. Crop salt tolerance -- current assessment //
  Journal of the Irrigation and Drainage Division, ASCE. -- 1977. --
  Vol. 103, No.~2. -- P. 115--134. --
  https://doi.org/10.1061/JRCEA4.0001137
\item
  Malczewski J. GIS and Multicriteria Decision Analysis. -- New York:
  John Wiley \& Sons, 1999. -- 392 p.
\item
  Malczewski J. GIS-based multicriteria decision analysis: a survey of
  the literature // International Journal of Geographical Information
  Science. -- 2006. -- Vol. 20, No.~7. -- P. 703--726. --
  https://doi.org/10.1080/13658810600661508
\item
  Malczewski J., Rinner C. Multicriteria Decision Analysis in Geographic
  Information Science. -- Berlin: Springer, 2015. -- 331 p.
\item
  Mann H.B. Nonparametric tests against trend // Econometrica. -- 1945.
  -- Vol. 13, No.~3. -- P. 245--259. -- https://doi.org/10.2307/1907187
\item
  Matheron G. Principles of geostatistics // Economic Geology. -- 1963.
  -- Vol. 58, No.~8. -- P. 1246--1266. --
  https://doi.org/10.2113/gsecongeo.58.8.1246
\item
  Measho S., Li F., Pellikka P. et al.~Soil salinity variations and
  associated implications for agriculture and land resources development
  using remote sensing datasets in Central Asia // Remote Sensing. --
  2022. -- Vol. 14, No.~10. -- Art. 2501. --
  https://doi.org/10.3390/rs14102501
\item
  Metternicht G.I., Zinck J.A. Remote sensing of soil salinity:
  potentials and constraints // Remote Sensing of Environment. -- 2003.
  -- Vol. 85, No.~1. -- P. 1--20. --
  https://doi.org/10.1016/s0034-4257(02)00188-8
\item
  Moran P.A.P. Notes on continuous stochastic phenomena // Biometrika.
  -- 1950. -- Vol. 37, No.~1--2. -- P. 17--23. --
  https://doi.org/10.1093/biomet/37.1-2.17
\item
  Nardo M., Saisana M., Saltelli A. et al.~Handbook on Constructing
  Composite Indicators: Methodology and User Guide. -- Paris: OECD
  Publishing, 2008. -- 158 p.
\item
  Okuda Y., Onishi J., Omori K., Oya T., Fukuo A., Kurvantaev R.,
  Shirokova Y., Nasonov V.G. Current status and problems of the drainage
  system in Uzbekistan // Journal of Arid Land Studies. -- 2015. -- Vol.
  25, No.~3. -- P. 81--84. -- https://doi.org/10.14976/jals.25.3\_81
\item
  Oliver M.A., Webster R. A tutorial guide to geostatistics: Computing
  and modelling variograms and kriging // Catena. -- 2014. -- Vol. 113.
  -- P. 56--69. -- https://doi.org/10.1016/j.catena.2013.09.006
\item
  Oliver M.A., Webster R. Kriging: a method of interpolation for
  geographical information systems // International Journal of
  Geographical Information Systems. -- 1990. -- Vol. 4, No.~3. -- P.
  313--332. -- https://doi.org/10.1080/02693799008941549
\item
  Oosterbaan R.J. SALTMOD: Description of Principles, User Manual and
  Examples of Application. -- Wageningen: ILRI, 2000. -- 68 p.
\item
  Pan X., Wang W., Liu T., Huang Y., De Maeyer P., Guo C., Ling Y.,
  Akmalov S. Quantitative detection and attribution of groundwater level
  variations in the Amu Darya Delta // Water. -- 2020. -- Vol. 12,
  No.~10. -- Art. 2869. -- https://doi.org/10.3390/w12102869
\item
  Pereira L.S., Oweis T., Zairi A. Irrigation management under water
  scarcity // Agricultural Water Management. -- 2002. -- Vol. 57, No.~3.
  -- P. 175--206. -- https://doi.org/10.1016/s0378-3774(02)00075-6
\item
  Pettitt A.N. A non-parametric approach to the change-point problem //
  Journal of the Royal Statistical Society. Series C (Applied
  Statistics). -- 1979. -- Vol. 28, No.~2. -- P. 126--135. --
  https://doi.org/10.2307/2346729
\item
  Pulatov A., Khamidov A., Akhmatov D., Pulatov B., Vasenev V. Soil
  salinity mapping by different interpolation methods in Mirzaabad
  district, Syrdarya Province // IOP Conference Series: Materials
  Science and Engineering. -- 2020. -- Vol. 883. -- Art. 012089. --
  https://doi.org/10.1088/1757-899X/883/1/012089
\item
  Qadir M., Ghafoor A., Murtaza G. Amelioration strategies for saline
  soils: a review // Land Degradation \& Development. -- 2000. -- Vol.
  11, No.~6. -- P. 501--521. --
  https://doi.org/10.1002/1099-145x(200011/12)11:6\textless501::aid-ldr405\textgreater3.0.co;2-s
\item
  Rengasamy P. World salinization with emphasis on Australia // Journal
  of Experimental Botany. -- 2006. -- Vol. 57, No.~5. -- P. 1017--1023.
  -- https://doi.org/10.1093/jxb/erj108
\item
  Rhoades J.D., Chanduvi F., Lesch S. Soil Salinity Assessment: Methods
  and Interpretation of Electrical Conductivity Measurements. FAO
  Irrigation and Drainage Paper 57. -- Rome: FAO, 1999. -- 150 p.
\item
  Richards L.A. (ed.). Diagnosis and Improvement of Saline and Alkali
  Soils. USDA Agriculture Handbook No.~60. -- Washington: U.S.
  Government Printing Office, 1954. -- 160 p.
\item
  Ritzema H.P. (ed.). Drainage Principles and Applications. ILRI
  Publication 16. -- Wageningen: International Institute for Land
  Reclamation and Improvement, 2006. -- 1125 p.
\item
  Roth J., Sant'Anna P.H.C., Bilinski A., Poe J. What's trending in
  difference-in-differences? A synthesis of the recent econometrics
  literature // Journal of Econometrics. -- 2023. -- Vol. 235, No.~2. --
  P. 2218--2244. -- https://doi.org/10.1016/j.jeconom.2023.03.008
\item
  Rouse J.W., Haas R.H., Schell J.A., Deering D.W. Monitoring vegetation
  systems in the Great Plains with ERTS // Third ERTS Symposium, NASA
  SP-351. -- Washington, 1974. -- Vol. 1. -- P. 309--317.
\item
  Runfola D., Anderson A., Baier H. et al.~geoBoundaries: A global
  database of political administrative boundaries // PLoS ONE. -- 2020.
  -- Vol. 15, No.~4. -- Art. e0231866. --
  https://doi.org/10.1371/journal.pone.0231866
\item
  Saaty T.L. Decision making with the analytic hierarchy process //
  International Journal of Services Sciences. -- 2008. -- Vol. 1, No.~1.
  -- P. 83--98. -- https://doi.org/10.1504/ijssci.2008.017590
\item
  Saaty T.L. How to make a decision: The Analytic Hierarchy Process //
  European Journal of Operational Research. -- 1990. -- Vol. 48, No.~1.
  -- P. 9--26. -- https://doi.org/10.1016/0377-2217(90)90057-i
\item
  Saaty T.L. The Analytic Hierarchy Process: Planning, Priority Setting,
  Resource Allocation. -- New York: McGraw-Hill, 1980. -- 287 p.
\item
  Saaty T.L., Vargas L.G. Models, Methods, Concepts and Applications of
  the Analytic Hierarchy Process. -- 2nd ed.~-- New York: Springer,
  2012. -- 346 p.
\item
  Saltelli A., Ratto M., Andres T. et al.~Global Sensitivity Analysis:
  The Primer. -- Chichester: John Wiley \& Sons, 2008. -- 292 p.
\item
  Sen P.K. Estimates of the regression coefficient based on Kendall's
  tau // Journal of the American Statistical Association. -- 1968. --
  Vol. 63, No.~324. -- P. 1379--1389. --
  https://doi.org/10.1080/01621459.1968.10480934
\item
  Shahid S.A., Zaman M., Heng L. Soil Salinity: Historical Perspectives
  and a World Overview of the Problem // Guideline for Salinity
  Assessment, Mitigation and Adaptation. -- Cham: Springer, 2018. -- P.
  43--53.
\item
  Shepard D. A two-dimensional interpolation function for
  irregularly-spaced data // Proceedings of the 23rd ACM National
  Conference. -- New York, 1968. -- P. 517--524.
\item
  Skaggs R.W. A water management model for shallow water table soils. --
  Raleigh: Water Resources Research Institute, 1978.
\item
  Skaggs R.W., Youssef M.A., Chescheir G.M. DRAINMOD: Model use,
  calibration and validation // Transactions of the ASABE. -- 2012. --
  Vol. 55, No.~4. -- P. 1509--1522. --
  https://doi.org/10.13031/2013.42259
\item
  Smedema L.K., Vlotman W.F., Rycroft D.W. Modern Land Drainage:
  Planning, Design and Management of Agricultural Drainage Systems. --
  Leiden: A.A. Balkema Publishers, 2004. -- 446 p.
\item
  Szabolcs I. Salt-Affected Soils. -- Boca Raton: CRC Press, 1989. --
  274 p.
\item
  Teimoori S., Olya M.H., Miller C.J. Groundwater level monitoring
  network design with machine learning methods // Journal of Hydrology.
  -- 2023. -- Vol. 625. -- Art. 130145. --
  https://doi.org/10.1016/j.jhydrol.2023.130145
\item
  Tilekeyev K., Kirbasheva D. Rehabilitating the Irrigation Canals in
  Southern Kyrgyzstan: An Impact Evaluation Study. Working Paper No.~69.
  -- Bishkek: University of Central Asia, 2022.
\item
  Tomlin C.D. Geographic Information Systems and Cartographic Modeling.
  -- Englewood Cliffs: Prentice-Hall, 1990. -- 249 p.
\item
  Van Schilfgaarde J. (ed.). Drainage for Agriculture. Agronomy
  Monograph No.~17. -- Madison: American Society of Agronomy, 1974. --
  700 p.
\item
  Vincent B., Vlotman W.F., Zimmer D., Hornbuckle J.W. Performance
  Assessment and Potential Indicators for Drainage Systems. -- ICID
  Working Group on Drainage, 2007. -- 98 p.
\item
  Vlotman W.F., Smedema L.K., Rycroft D.W. Modern Land Drainage:
  Planning, Design and Management of Agricultural Drainage Systems. --
  2nd ed.~-- Leiden: CRC Press, 2020.
\item
  Wang J., Khan S., O'Connell N. Quantifying rainfall and flooding
  impacts on groundwater levels in irrigation areas: GIS approach //
  Journal of Irrigation and Drainage Engineering. -- 2007. -- Vol. 133,
  No.~4. -- P. 359--367. --
  https://doi.org/10.1061/(asce)0733-9437(2007)133:4(359)
\item
  Webster R., Oliver M.A.~Geostatistics for Environmental Scientists. --
  2nd ed.~-- Chichester: John Wiley \& Sons, 2007. -- 315 p.
\item
  World Meteorological Organization. Guide to Hydrological Practices.
  Vol. I: Hydrology -- From Measurement to Hydrological Information.
  WMO-No.~168. -- 6th ed.~-- Geneva: WMO, 2008.
\item
  Yang Y., Li D., Huang W., Zhou X., Li Z., Dong X., Wang X. Effects of
  subsurface drainage on soil salinity and groundwater table in drip
  irrigated cotton fields in oasis regions of Tarim Basin //
  Agriculture. -- 2022. -- Vol. 12, No.~12. -- Art. 2167. --
  https://doi.org/10.3390/agriculture12122167
\item
  Zhang J., Goodchild M.F. Uncertainty in Geographical Information. --
  London: Taylor \& Francis, 2002. -- 266 p.
\end{enumerate}

\textbf{III. Internet resurslari}

\begin{enumerate}
\def\labelenumi{\arabic{enumi}.}
\setcounter{enumi}{134}
\item
  Marhamat District // Wikipedia. {[}Elektron resurs{]}. -- URL:
  https://en.wikipedia.org/wiki/Marhamat\_District (murojaat sanasi:
  16.11.2026); Мархаматский район // Википедия. -- URL:
  https://ru.wikipedia.org/wiki/Мархаматский\_район
\item
  Андижан: климат (данные 1961--1990 гг.) // Метео-ТВ. {[}Elektron
  resurs{]}. -- URL:
  https://meteo-tv.ru/uzbekistan/andizhan/weather/climate/ (murojaat
  sanasi: 16.11.2026)
\item
  Marhamat tumani meliorativ tizimi: interaktiv xarita // GIS Cloud.
  {[}Elektron resurs{]}. -- URL: https://editor.giscloud.com/map/3370958
  (yaratilgan sana: 16.11.2026)
\item
  Copernicus DEM -- Global and European Digital Elevation Model
  (COP-DEM), GLO-30 // European Space Agency. {[}Elektron resurs{]}. --
  URL: https://doi.org/10.5270/ESA-c5d3d65 (murojaat sanasi:
  17.11.2026).
\item
  Iste'mol narxlari indeksi: 2017--2026-yillar bo`yicha oylik va yillik
  press-relizlar // O`zbekiston Respublikasi Prezidenti huzuridagi
  Statistika agentligi. {[}Elektron resurs{]}. -- URL: https://stat.uz
  (murojaat sanasi: 17.11.2026).
\end{enumerate}

\oakbob{ILOVALAR}

\oakilova{1-ilova}

\oakilovanom{\textbf{Andijon viloyati kollektor-drenaj tarmogʻining tumanlar boʻyicha
tarkibi (kadastr, 2021-yil 1-yanvar), km}}

{\fontsize{9}{11}\selectfont\setlength{\tabcolsep}{2.5pt}

\noindent\begin{longtable}[]{|
  >{\raggedright\arraybackslash}p{(\columnwidth - 16\tabcolsep - 9\arrayrulewidth) * \real{0.1420}}|>{\raggedright\arraybackslash}p{(\columnwidth - 16\tabcolsep - 9\arrayrulewidth) * \real{0.0812}}|>{\raggedright\arraybackslash}p{(\columnwidth - 16\tabcolsep - 9\arrayrulewidth) * \real{0.1044}}|>{\raggedright\arraybackslash}p{(\columnwidth - 16\tabcolsep - 9\arrayrulewidth) * \real{0.1275}}|>{\raggedright\arraybackslash}p{(\columnwidth - 16\tabcolsep - 9\arrayrulewidth) * \real{0.1623}}|>{\raggedright\arraybackslash}p{(\columnwidth - 16\tabcolsep - 9\arrayrulewidth) * \real{0.1391}}|>{\raggedright\arraybackslash}p{(\columnwidth - 16\tabcolsep - 9\arrayrulewidth) * \real{0.1275}}|>{\raggedright\arraybackslash}p{(\columnwidth - 16\tabcolsep - 9\arrayrulewidth) * \real{0.1159}}|}
\hline
\begin{minipage}[b]{\linewidth}\raggedright
\textbf{Tuman}
\end{minipage} & \begin{minipage}[b]{\linewidth}\raggedright
\textbf{Ochiq tarmoq, jami}
\end{minipage} & \begin{minipage}[b]{\linewidth}\raggedright
\textbf{Magistral}
\end{minipage} & \begin{minipage}[b]{\linewidth}\raggedright
\textbf{Tumanlararo}
\end{minipage} & \begin{minipage}[b]{\linewidth}\raggedright
\textbf{Xoʻjaliklararo (davlat)}
\end{minipage} & \begin{minipage}[b]{\linewidth}\raggedright
\textbf{Xoʻjalik ichki (uyushmalar)}
\end{minipage} & \begin{minipage}[b]{\linewidth}\raggedright
\textbf{Yopiq-yotiq drenaj}
\end{minipage} & \begin{minipage}[b]{\linewidth}\raggedright
\textbf{Inshootlar soni}
\end{minipage} \\
\hline
\endhead
\hline
\endlastfoot
Ulugʻnor & 1 259,3 & 39,0 & 55,9 & 259,9 & 904,5 & 3,1 & 580 \\ \hline
Baliqchi & 873,9 & 33,0 & 71,0 & 355,4 & 414,6 & 62,8 & 844 \\ \hline
Shahrixon & 675,3 & 7,6 & 164,3 & 184,0 & 319,4 & 65,7 & 473 \\ \hline
Boʻston & 622,1 & --- & 60,7 & 234,6 & 326,8 & 10,4 & 376 \\ \hline
Marhamat & 621,1 & 41,8 & 21,5 & 202,5 & 355,3 & 91,7 & 455 \\ \hline
Oltinkoʻl & 606,1 & --- & 115,1 & 199,7 & 291,4 & 25,1 & 652 \\ \hline
Jalaquduq & 565,1 & 21,2 & 37,3 & 184,5 & 322,1 & 110,0 & 624 \\ \hline
Izboskan & 531,9 & 26,0 & 52,1 & 201,3 & 252,6 & 33,4 & 405 \\ \hline
Buloqboshi & 415,3 & 3,5 & 6,4 & 158,8 & 246,6 & 209,8 & 293 \\ \hline
Andijon & 401,5 & --- & 11,0 & 101,1 & 289,4 & 17,6 & 304 \\ \hline
Qoʻrgʻontepa & 337,0 & 7,4 & 23,0 & 59,1 & 247,5 & 66,7 & 342 \\ \hline
Asaka & 299,1 & 7,1 & 63,0 & 71,1 & 157,8 & 7,7 & 261 \\ \hline
Xoʻjaobod & 241,0 & 2,8 & 3,1 & 79,9 & 155,1 & 62,3 & 334 \\ \hline
Paxtaobod & 167,5 & 10,9 & 29,7 & 52,7 & 74,2 & 20,5 & 133 \\ \hline
Viloyat boʻyicha & 7 616,3 & 200,4 & 714,0 & 2 344,6 & 4 357,4 & 786,6 &
6 076 \\
\end{longtable}

}

\oakilova{2-ilova}

\oakilovanom{\textbf{Markaziy Osiyoda sizot suvlari rejimi va shoʻrlanishni GAT va
geostatistika yordamida oʻrganish boʻyicha tadqiqotlar}}

{\fontsize{9}{11}\selectfont\setlength{\tabcolsep}{2.5pt}

\noindent\begin{longtable}[]{|
  >{\raggedright\arraybackslash}p{(\columnwidth - 10\tabcolsep - 6\arrayrulewidth) * \real{0.1721}}|>{\raggedright\arraybackslash}p{(\columnwidth - 10\tabcolsep - 6\arrayrulewidth) * \real{0.2065}}|>{\raggedright\arraybackslash}p{(\columnwidth - 10\tabcolsep - 6\arrayrulewidth) * \real{0.2071}}|>{\raggedright\arraybackslash}p{(\columnwidth - 10\tabcolsep - 6\arrayrulewidth) * \real{0.1955}}|>{\raggedright\arraybackslash}p{(\columnwidth - 10\tabcolsep - 6\arrayrulewidth) * \real{0.2188}}|}
\hline
\begin{minipage}[b]{\linewidth}\raggedright
\textbf{Mualliflar, yil}
\end{minipage} & \begin{minipage}[b]{\linewidth}\raggedright
\textbf{Hudud}
\end{minipage} & \begin{minipage}[b]{\linewidth}\raggedright
\textbf{Maʼlumotlar}
\end{minipage} & \begin{minipage}[b]{\linewidth}\raggedright
\textbf{Usul}
\end{minipage} & \begin{minipage}[b]{\linewidth}\raggedright
\textbf{Asosiy natija}
\end{minipage} \\
\hline
\endhead
\hline
\endlastfoot
M.Ibrakhimov va boshq., 2007 {[}70{]} & Xorazm viloyati &
1990--2000-yillar, 1 972 ta kuzatuv qudugʻi & Kriging & Vegetatsiya
davrida hududning taxminan uchdan ikki qismida sizot sathi 1,5 m kritik
chuqurlikdan sayoz \\ \hline
I.Forkutsa va boshq., 2009 {[}55{]} & Xorazm viloyati & Dala tajribasi &
Modellashtirish & Sayoz sizot suvi sharoitida gʻoʻzaning suv isteʼmoli
baholangan \\ \hline
A.Akramkhanov, P.L.G.Vlek, 2012 {[}28{]} & Xorazm viloyati & Tuproq
oʻlchovlari va atrof-muhit koʻrsatkichlari & Neyron tarmogʻi & Tuproq
shoʻrlanishi xavfining hududiy taqsimoti \\ \hline
A.Hamidov va boshq., 2020 {[}64{]} & Xorazm viloyati & 1991--2017-yillar
& Mann--Kendall mezoni, regressiya & 2100-yilgacha sath 1,72--1,79 m,
minerallashuv 1,85--2,1 g/l boʻlishi prognoz qilingan \\ \hline
Sh.Kenjabayev va boshq. {[}75, 76{]} & Xorazm, Qulovot & Monitoring va
atrof-muhit koʻrsatkichlari & GAT tahlili & Yerlarning meliorativ holati
baholangan \\ \hline
X.Pan va boshq., 2020 {[}100{]} & Amudaryo deltasi & 1999--2017-yillar,
sath, iqlim, masofadan zondlash & Trend, korrelyatsiya, bosh
komponentlar & Sath 1999--2002-yillarda 1,89 m dan 3,50 m gacha
pasaygan, 2017-yilgacha 2,11 m ga qaytgan; sugʻorish bilan bogʻliq \\ \hline
Y.Duan va boshq., 2022 {[}50{]} & Sirdaryo boʻyi (Qozogʻiston) & 84 ta
tuproq namunasi & Interpolyatsiya, geodetektor & Shoʻrlanishni kanal va
drenajgacha masofa eng koʻp tushuntiradi \\ \hline
R.Kulmatov va boshq., 2024 {[}81{]} & Qoraqalpogʻiston, Xoʻjayli &
2015--2023-yillar, 142 ta kuzatuv qudugʻi & IDW, korrelyatsiya & Sath va
minerallashuvning mavsumiy xaritalari \\ \hline
A.Pulatov va boshq. {[}103{]} & Sirdaryo, Mirzaobod & Tuproq namunalari
& Interpolyatsiya usullari qiyosi & Shoʻrlanish xaritalari, usullar
aniqligi solishtirilgan \\ \hline
Sh.Kholdorov va boshq., 2023 {[}77{]} & Markaziy Fargʻona & Landsat 8,
Sentinel-2, 2019-yil tuproq namunalari & Spektral indekslar & Landsat 8
boʻyicha NDSI va SI12 indekslari eng yaxshi natija bergan \\
\end{longtable}

}

\oakilova{3-ilova}

\oakilovanom{\textbf{Tadqiqot obyektining massivlar boʻyicha tavsifi}}

{\fontsize{12}{14}\selectfont\setlength{\tabcolsep}{4pt}

\noindent\begin{longtable}[]{|
  >{\raggedright\arraybackslash}p{(\columnwidth - 12\tabcolsep - 7\arrayrulewidth) * \real{0.2163}}|>{\raggedright\arraybackslash}p{(\columnwidth - 12\tabcolsep - 7\arrayrulewidth) * \real{0.1477}}|>{\raggedright\arraybackslash}p{(\columnwidth - 12\tabcolsep - 7\arrayrulewidth) * \real{0.2122}}|>{\raggedright\arraybackslash}p{(\columnwidth - 12\tabcolsep - 7\arrayrulewidth) * \real{0.1630}}|>{\raggedright\arraybackslash}p{(\columnwidth - 12\tabcolsep - 7\arrayrulewidth) * \real{0.1468}}|>{\raggedright\arraybackslash}p{(\columnwidth - 12\tabcolsep - 7\arrayrulewidth) * \real{0.1140}}|}
\hline
\begin{minipage}[b]{\linewidth}\raggedright
\textbf{Massiv}
\end{minipage} & \begin{minipage}[b]{\linewidth}\raggedright
\textbf{Kuzatuv quduqlari (pasport)}
\end{minipage} & \begin{minipage}[b]{\linewidth}\raggedright
\textbf{Quduqlarga biriktirilgan maydon, ga}
\end{minipage} & \begin{minipage}[b]{\linewidth}\raggedright
\textbf{Median balandlik, m}
\end{minipage} & \begin{minipage}[b]{\linewidth}\raggedright
\textbf{Ichki zovurlar, km}
\end{minipage} & \begin{minipage}[b]{\linewidth}\raggedright
\textbf{Bonitet balli}
\end{minipage} \\
\hline
\endhead
\hline
\endlastfoot
Z.M.Bobur & 27 & 2 376 & 565 & 35,2 & 67,9 \\ \hline
Paxtakor & 19 & 1 588 & 587 & 41,9 & 64,1 \\ \hline
Mashʼal & 29 & 3 134 & 596 & 92,0 & 59,7 \\ \hline
U.Qoraboyev & 26 & 1 944 & 601 & 21,9 & 60,8 \\ \hline
M.Toʻjiboyev & 33 & 2 027 & 603 & ---* & 58,8 \\ \hline
T.Mirzayev & 20 & 1 982 & 610 & 30,4 & 63,3 \\ \hline
Istiqlol & 27 & 2 400 & 614 & 70,2 & 62,9 \\ \hline
Y.Oxunboboyev & 7 & 534 & 641 & 20,1 & 41,3 \\ \hline
Jami & 188 & 15 985 & 597 & & 59,7 \\
\end{longtable}

}

Izoh: * --- kadastrda M.Toʻjiboyev massivi boʻyicha ichki zovurlar
alohida koʻrsatilmagan.

\oakilova{4-ilova}

\oakilovanom{\textbf{Andijon meteostansiyasining iqlim koʻrsatkichlari (1961--1990)
va hisoblangan etalon bugʻlanuvchanlik}}

{\fontsize{12}{14}\selectfont\setlength{\tabcolsep}{4pt}

\noindent\begin{longtable}[]{|
  >{\raggedright\arraybackslash}p{(\columnwidth - 14\tabcolsep - 8\arrayrulewidth) * \real{0.1319}}|>{\raggedright\arraybackslash}p{(\columnwidth - 14\tabcolsep - 8\arrayrulewidth) * \real{0.1605}}|>{\raggedright\arraybackslash}p{(\columnwidth - 14\tabcolsep - 8\arrayrulewidth) * \real{0.1605}}|>{\raggedright\arraybackslash}p{(\columnwidth - 14\tabcolsep - 8\arrayrulewidth) * \real{0.1522}}|>{\raggedright\arraybackslash}p{(\columnwidth - 14\tabcolsep - 8\arrayrulewidth) * \real{0.1254}}|>{\raggedright\arraybackslash}p{(\columnwidth - 14\tabcolsep - 8\arrayrulewidth) * \real{0.0761}}|>{\raggedright\arraybackslash}p{(\columnwidth - 14\tabcolsep - 8\arrayrulewidth) * \real{0.1934}}|}
\hline
\begin{minipage}[b]{\linewidth}\raggedright
\textbf{Oy}
\end{minipage} & \begin{minipage}[b]{\linewidth}\raggedright
\textbf{Oʻrtacha maksimal harorat, °C}
\end{minipage} & \begin{minipage}[b]{\linewidth}\raggedright
\textbf{Oʻrtacha minimal harorat, °C}
\end{minipage} & \begin{minipage}[b]{\linewidth}\raggedright
\textbf{Oʻrtacha harorat, °C}
\end{minipage} & \begin{minipage}[b]{\linewidth}\raggedright
\textbf{Yogʻin, mm}
\end{minipage} & \begin{minipage}[b]{\linewidth}\raggedright
\textbf{ET₀, mm}
\end{minipage} & \begin{minipage}[b]{\linewidth}\raggedright
\textbf{Namlik tanqisligi (ET₀ − yogʻin), mm}
\end{minipage} \\
\hline
\endhead
\hline
\endlastfoot
Yanvar & 2 & −7 & −2,5 & 29 & 20 & −9 \\ \hline
Fevral & 5 & −5 & 0,0 & 36 & 30 & −6 \\ \hline
Mart & 12 & 2 & 7,0 & 45 & 62 & 17 \\ \hline
Aprel & 21 & 9 & 15,0 & 46 & 110 & 64 \\ \hline
May & 27 & 14 & 20,5 & 38 & 159 & 121 \\ \hline
Iyun & 32 & 17 & 24,5 & 17 & 193 & 176 \\ \hline
Iyul & 33 & 19 & 26,0 & 12 & 194 & 182 \\ \hline
Avgust & 31 & 16 & 23,5 & 6 & 170 & 164 \\ \hline
Sentabr & 27 & 11 & 19,0 & 7 & 123 & 116 \\ \hline
Oktabr & 19 & 6 & 12,5 & 32 & 70 & 38 \\ \hline
Noyabr & 11 & 0 & 5,5 & 27 & 35 & 8 \\ \hline
Dekabr & 4 & −4 & 0,0 & 28 & 19 & −9 \\ \hline
Yil & --- & --- & 12,6 & 323 & 1 185 & 862 \\
\end{longtable}

}

\oakilova{5-ilova}

\oakilovanom{\textbf{Andijon meteostansiyasi boʻyicha iqlimning yillik yoʻli
(1961--1990): a --- oylik yogʻin va hisoblangan etalon bugʻlanuvchanlik;
b --- oʻrtacha oylik harorat}}

\oakimg{figures/ch2_climate.png}

\oakilova{6-ilova}

\oakilovanom{\textbf{Sizot suvlari sathining kuzatuv quduqlari boʻyicha taqsimoti
(2024-yil 1-aprel, 182 ta quduq)}}

{\fontsize{12}{14}\selectfont\setlength{\tabcolsep}{4pt}

\noindent\begin{longtable}[]{|
  >{\raggedright\arraybackslash}p{(\columnwidth - 4\tabcolsep - 3\arrayrulewidth) * \real{0.6174}}|>{\raggedright\arraybackslash}p{(\columnwidth - 4\tabcolsep - 3\arrayrulewidth) * \real{0.3826}}|}
\hline
\begin{minipage}[b]{\linewidth}\raggedright
\textbf{Sath chuqurligi}
\end{minipage} & \begin{minipage}[b]{\linewidth}\raggedright
\textbf{Quduqlar ulushi, \%}
\end{minipage} \\
\hline
\endhead
\hline
\endlastfoot
1,0 m gacha & 2,7 \\ \hline
1,0--1,5 m & 16,5 \\ \hline
1,5--2,0 m & 18,7 \\ \hline
2,0--3,0 m & 34,1 \\ \hline
3,0--5,0 m & 22,5 \\ \hline
5,0 m dan chuqur & 5,5 \\
\end{longtable}

}

\oakilova{7-ilova}

\oakilovanom{\textbf{Davlat hisobidagi kollektorlarning uzunlik va suv sarfi boʻyicha
taqsimoti}}

{\fontsize{12}{14}\selectfont\setlength{\tabcolsep}{4pt}

\noindent\begin{longtable}[]{|
  >{\raggedright\arraybackslash}p{(\columnwidth - 10\tabcolsep - 6\arrayrulewidth) * \real{0.2248}}|>{\raggedright\arraybackslash}p{(\columnwidth - 10\tabcolsep - 6\arrayrulewidth) * \real{0.2476}}|>{\raggedright\arraybackslash}p{(\columnwidth - 10\tabcolsep - 6\arrayrulewidth) * \real{0.0459}}|>{\raggedright\arraybackslash}p{(\columnwidth - 10\tabcolsep - 6\arrayrulewidth) * \real{0.2341}}|>{\raggedright\arraybackslash}p{(\columnwidth - 10\tabcolsep - 6\arrayrulewidth) * \real{0.2476}}|}
\hline
\begin{minipage}[b]{\linewidth}\raggedright
\textbf{Uzunlik, km}
\end{minipage} & \begin{minipage}[b]{\linewidth}\raggedright
\textbf{Kollektorlar soni}
\end{minipage} & \begin{minipage}[b]{\linewidth}\raggedright
\end{minipage} & \begin{minipage}[b]{\linewidth}\raggedright
\textbf{Oxiridagi sarf, m³/s}
\end{minipage} & \begin{minipage}[b]{\linewidth}\raggedright
\textbf{Kollektorlar soni}
\end{minipage} \\
\hline
\endhead
\hline
\endlastfoot
1,0 gacha & 21 & & 0,01 gacha & 31 \\ \hline
1,0--2,0 & 24 & & 0,01--0,05 & 24 \\ \hline
2,0--5,0 & 30 & & 0,05--0,2 & 25 \\ \hline
5,0--10,0 & 14 & & 0,2--1,0 & 5 \\ \hline
10,0 dan ortiq & 1 & & 1,0 dan ortiq & 5 \\ \hline
Jami & 90 & & Jami & 90 \\
\end{longtable}

}

\oakilova{8-ilova}

\oakilovanom{\textbf{Uzunligi boʻyicha eng yirik oʻnta kollektor}}

{\fontsize{12}{14}\selectfont\setlength{\tabcolsep}{4pt}

\noindent\begin{longtable}[]{|
  >{\raggedright\arraybackslash}p{(\columnwidth - 12\tabcolsep - 7\arrayrulewidth) * \real{0.2202}}|>{\raggedright\arraybackslash}p{(\columnwidth - 12\tabcolsep - 7\arrayrulewidth) * \real{0.1518}}|>{\raggedright\arraybackslash}p{(\columnwidth - 12\tabcolsep - 7\arrayrulewidth) * \real{0.1564}}|>{\raggedright\arraybackslash}p{(\columnwidth - 12\tabcolsep - 7\arrayrulewidth) * \real{0.1392}}|>{\raggedright\arraybackslash}p{(\columnwidth - 12\tabcolsep - 7\arrayrulewidth) * \real{0.2254}}|>{\raggedright\arraybackslash}p{(\columnwidth - 12\tabcolsep - 7\arrayrulewidth) * \real{0.1070}}|}
\hline
\begin{minipage}[b]{\linewidth}\raggedright
\textbf{Kollektor}
\end{minipage} & \begin{minipage}[b]{\linewidth}\raggedright
\textbf{Uzunligi, km}
\end{minipage} & \begin{minipage}[b]{\linewidth}\raggedright
\textbf{Oxiridagi sarf, m³/s}
\end{minipage} & \begin{minipage}[b]{\linewidth}\raggedright
\textbf{Xizmat maydoni, ga}
\end{minipage} & \begin{minipage}[b]{\linewidth}\raggedright
\textbf{Foydalanishga topshirilgan yil}
\end{minipage} & \begin{minipage}[b]{\linewidth}\raggedright
\textbf{Oxirgi taʼmir yili}
\end{minipage} \\
\hline
\endhead
\hline
\endlastfoot
Gʻarbiy Qoraqoʻrgʻon & 12,50 & 4,60 & 580 & 1918 & 2020 \\ \hline
Nayman & 9,80 & 1,20 & 480 & 1948 & 2019 \\ \hline
Gʻarbiy & 9,70 & 2,00 & 545 & 1944 & 2019 \\ \hline
Markaziy & 8,50 & 1,50 & 375 & 1944 & 2012 \\ \hline
Aravon & 8,40 & 0,20 & 304 & 1943 & 2013 \\ \hline
Shukur mergan & 8,36 & 0,20 & 440 & 1948 & 2017 \\ \hline
JFK chap & 8,30 & 0,06 & 512 & 1969 & 2019 \\ \hline
Sharqiy Qoraqoʻrgʻon & 8,11 & 0,05 & 450 & 1950 & 2019 \\ \hline
OD-2 & 6,70 & 0,30 & 151 & 1957 & 2016 \\ \hline
Shoʻr & 6,70 & 0,20 & 619 & 1951 & 2019 \\
\end{longtable}

}

\oakilova{9-ilova}

\oakilovanom{\textbf{Kollektor-drenaj tarmogʻidagi gidrotexnik inshootlar (kadastr,
2021)}}

{\fontsize{12}{14}\selectfont\setlength{\tabcolsep}{4pt}

\noindent\begin{longtable}[]{|
  >{\raggedright\arraybackslash}p{(\columnwidth - 8\tabcolsep - 5\arrayrulewidth) * \real{0.3916}}|>{\raggedright\arraybackslash}p{(\columnwidth - 8\tabcolsep - 5\arrayrulewidth) * \real{0.1821}}|>{\raggedright\arraybackslash}p{(\columnwidth - 8\tabcolsep - 5\arrayrulewidth) * \real{0.3410}}|>{\raggedright\arraybackslash}p{(\columnwidth - 8\tabcolsep - 5\arrayrulewidth) * \real{0.0852}}|}
\hline
\begin{minipage}[b]{\linewidth}\raggedright
\textbf{Inshoot turi}
\end{minipage} & \begin{minipage}[b]{\linewidth}\raggedright
\textbf{Davlat hisobida}
\end{minipage} & \begin{minipage}[b]{\linewidth}\raggedright
\textbf{Suv isteʼmolchilari uyushmalari hisobida}
\end{minipage} & \begin{minipage}[b]{\linewidth}\raggedright
\textbf{Jami}
\end{minipage} \\
\hline
\endhead
\hline
\endlastfoot
Nazorat tarmogʻi inshootlari & 238 & --- & 238 \\ \hline
Ochiq koʻpriklar & 26 & 117 & 143 \\ \hline
Quvurli oʻtish joylari & 7 & 35 & 42 \\ \hline
Dyukerlar & 10 & --- & 10 \\ \hline
Suv oʻlchash inshootlari & 6 & --- & 6 \\ \hline
Akveduklar & 2 & --- & 2 \\ \hline
Suv toʻsish inshootlari & 1 & --- & 1 \\ \hline
Boshqalar & --- & 13 & 13 \\ \hline
Jami & 290 & 165 & 455 \\
\end{longtable}

}

\oakilova{10-ilova}

\oakilovanom{\textbf{Kuzatuv quduqlarining qurilish davri, materiali va chuqurligi}}

{\fontsize{12}{14}\selectfont\setlength{\tabcolsep}{4pt}

\noindent\begin{longtable}[]{|
  >{\raggedright\arraybackslash}p{(\columnwidth - 8\tabcolsep - 5\arrayrulewidth) * \real{0.2460}}|>{\raggedright\arraybackslash}p{(\columnwidth - 8\tabcolsep - 5\arrayrulewidth) * \real{0.1988}}|>{\raggedright\arraybackslash}p{(\columnwidth - 8\tabcolsep - 5\arrayrulewidth) * \real{0.2734}}|>{\raggedright\arraybackslash}p{(\columnwidth - 8\tabcolsep - 5\arrayrulewidth) * \real{0.2818}}|}
\hline
\begin{minipage}[b]{\linewidth}\raggedright
\textbf{Qurilish davri}
\end{minipage} & \begin{minipage}[b]{\linewidth}\raggedright
\textbf{Quduqlar soni}
\end{minipage} & \begin{minipage}[b]{\linewidth}\raggedright
\textbf{Material}
\end{minipage} & \begin{minipage}[b]{\linewidth}\raggedright
\textbf{Chuqurlik medianasi, m}
\end{minipage} \\
\hline
\endhead
\hline
\endlastfoot
1971--1990 & 65 & Metall & 3,1 \\ \hline
1991--2012 & 3 & Metall & 3,2 \\ \hline
2013--2016 & 120 & Polietilen & 3,2 \\ \hline
Jami & 188 & & 3,2 \\
\end{longtable}

}

\oakilova{11-ilova}

\oakilovanom{\textbf{GAT bazasi uchun boshlangʻich manbalar}}

{\fontsize{12}{14}\selectfont\setlength{\tabcolsep}{4pt}

\noindent\begin{longtable}[]{|
  >{\raggedright\arraybackslash}p{(\columnwidth - 8\tabcolsep - 5\arrayrulewidth) * \real{0.2779}}|>{\raggedright\arraybackslash}p{(\columnwidth - 8\tabcolsep - 5\arrayrulewidth) * \real{0.1819}}|>{\raggedright\arraybackslash}p{(\columnwidth - 8\tabcolsep - 5\arrayrulewidth) * \real{0.2779}}|>{\raggedright\arraybackslash}p{(\columnwidth - 8\tabcolsep - 5\arrayrulewidth) * \real{0.2622}}|}
\hline
\begin{minipage}[b]{\linewidth}\raggedright
\textbf{Manba}
\end{minipage} & \begin{minipage}[b]{\linewidth}\raggedright
\textbf{Shakli}
\end{minipage} & \begin{minipage}[b]{\linewidth}\raggedright
\textbf{Joylashuv maʼlumoti}
\end{minipage} & \begin{minipage}[b]{\linewidth}\raggedright
\textbf{Qoʻshimcha maʼlumot}
\end{minipage} \\
\hline
\endhead
\hline
\endlastfoot
Kuzatuv quduqlari koordinatalari & Jadval (GAT eksporti) & 205 ta
nuqtaning geografik koordinatalari & Quduq raqami, tuman \\ \hline
Meliorativ ekspeditsiya xaritalari, 1:50 000 & PDF (2 varaq) & Quduqlar,
kollektorlar, kanallar, massiv chegaralari tasviri & 2024-yil 1-aprel
holatiga sath va minerallashuv yozuvlari, kollektor nomlari \\ \hline
Meliorativ kadastr, 2021 & Jadvallar & Yoʻq & Kollektorlar, zovurlar,
quduqlar pasporti \\ \hline
Monitoring jadvallari, 2009--2024 & Jadvallar & Yoʻq & Oʻn kunlik sath
va yillik minerallashuv \\
\end{longtable}

}

\oakilova{12-ilova}

\oakilovanom{\textbf{GAT bazasini yaratish bosqichlari}}

{\fontsize{12}{14}\selectfont\setlength{\tabcolsep}{4pt}

\noindent\begin{longtable}[]{|
  >{\raggedright\arraybackslash}p{(\columnwidth - 6\tabcolsep - 4\arrayrulewidth) * \real{0.3385}}|>{\raggedright\arraybackslash}p{(\columnwidth - 6\tabcolsep - 4\arrayrulewidth) * \real{0.3385}}|>{\raggedright\arraybackslash}p{(\columnwidth - 6\tabcolsep - 4\arrayrulewidth) * \real{0.3230}}|}
\hline
\begin{minipage}[b]{\linewidth}\raggedright
\textbf{Bosqich}
\end{minipage} & \begin{minipage}[b]{\linewidth}\raggedright
\textbf{Mazmuni}
\end{minipage} & \begin{minipage}[b]{\linewidth}\raggedright
\textbf{Natija}
\end{minipage} \\
\hline
\endhead
\hline
\endlastfoot
1. Rastrlash & Xaritaning yozuvsiz tasviri 300 nuqta/dyuym aniqlikda
olindi; yozuvlar (quduq raqamlari, qiymatlar, kollektor nomlari) oʻrni
bilan alohida ajratildi & Piksel oʻlchami joyda 4,23 m \\ \hline
2. Quduq belgilarini aniqlash & Quduq belgilari rang boʻyicha ajratildi,
markazi hisoblandi va eng yaqin raqam hamda qiymat yozuvi bilan
juftlandi & 202 ta quduq belgisi \\ \hline
3. Koordinataga bogʻlash & Raqami koordinatalar jadvalida topilgan
quduqlar tayanch nuqta sifatida olinib, affin almashtirish eng kichik
kvadratlar usulida aniqlandi & 181 ta tayanch nuqta \\ \hline
4. Tarmoqni raqamlash & Kollektor va kanal chiziqlari rang boʻyicha
ajratildi, oʻq chizigʻiga keltirildi va siniq chiziqlarga aylantirildi &
Kollektorlar 214,6 km, kanallar 170,9 km \\ \hline
5. Atributlarni biriktirish & Kollektor nomlari yozuvlari chiziq
boʻlaklariga biriktirildi va kadastr yozuvlari bilan bogʻlandi & 60 ta
kollektor, 114,8 km \\ \hline
6. Massiv chegaralarini tiklash & Chegara chizigʻi boʻlaklari va
quduqlarning massivga mansubligi asosida chegaralar chizildi & 8 ta
massiv chegarasi \\
\end{longtable}

}

\oakilova{13-ilova}

\oakilovanom{\textbf{Xaritani koordinataga bogʻlash aniqligi}}

{\fontsize{12}{14}\selectfont\setlength{\tabcolsep}{4pt}

\noindent\begin{longtable}[]{|
  >{\raggedright\arraybackslash}p{(\columnwidth - 4\tabcolsep - 3\arrayrulewidth) * \real{0.7700}}|>{\raggedright\arraybackslash}p{(\columnwidth - 4\tabcolsep - 3\arrayrulewidth) * \real{0.2300}}|}
\hline
\begin{minipage}[b]{\linewidth}\raggedright
\textbf{Koʻrsatkich}
\end{minipage} & \begin{minipage}[b]{\linewidth}\raggedright
\textbf{Qiymat}
\end{minipage} \\
\hline
\endhead
\hline
\endlastfoot
Juftlangan tayanch nuqtalar & 190 \\ \hline
Hisobda qoldirilgan tayanch nuqtalar & 181 \\ \hline
Xato medianasi & 1,9 m \\ \hline
Oʻrtacha kvadratik xato (181 nuqta) & 11,9 m \\ \hline
Qoldigʻi 5 m dan kichik nuqtalar ulushi (190 nuqta) & 91 \% \\ \hline
Aniqlangan piksel oʻlchami & 4,232 m \\ \hline
1:50 000 masshtab va 300 nuqta/dyuym uchun nazariy piksel oʻlchami &
4,233 m \\
\end{longtable}

}

\oakilova{14-ilova}

\oakilovanom{\textbf{Xaritani koordinataga bogʻlash xatolari (190 ta tayanch nuqta):
a --- 10 m gacha boʻlgan xatolar taqsimoti; b --- barcha xatolar,
logarifmik shkala}}

\oakimg{figures/ch2_georef_residuals.png}

\oakilova{15-ilova}

\oakilovanom{\textbf{Xarita boʻyicha massivlar maydoni va kollektorlar zichligi}}

{\fontsize{12}{14}\selectfont\setlength{\tabcolsep}{4pt}

\noindent\begin{longtable}[]{|
  >{\raggedright\arraybackslash}p{(\columnwidth - 10\tabcolsep - 6\arrayrulewidth) * \real{0.2373}}|>{\raggedright\arraybackslash}p{(\columnwidth - 10\tabcolsep - 6\arrayrulewidth) * \real{0.1287}}|>{\raggedright\arraybackslash}p{(\columnwidth - 10\tabcolsep - 6\arrayrulewidth) * \real{0.2478}}|>{\raggedright\arraybackslash}p{(\columnwidth - 10\tabcolsep - 6\arrayrulewidth) * \real{0.2184}}|>{\raggedright\arraybackslash}p{(\columnwidth - 10\tabcolsep - 6\arrayrulewidth) * \real{0.1678}}|}
\hline
\begin{minipage}[b]{\linewidth}\raggedright
\textbf{Massiv}
\end{minipage} & \begin{minipage}[b]{\linewidth}\raggedright
\textbf{Umumiy maydon, ga}
\end{minipage} & \begin{minipage}[b]{\linewidth}\raggedright
\textbf{Sugʻoriladigan maydon, ga}
\end{minipage} & \begin{minipage}[b]{\linewidth}\raggedright
\textbf{Kollektorlar (xaritada), km}
\end{minipage} & \begin{minipage}[b]{\linewidth}\raggedright
\textbf{Kollektor zichligi, m/ga}
\end{minipage} \\
\hline
\endhead
\hline
\endlastfoot
Mashʼal & 5 100 & 3 790 & 54,8 & 14,5 \\ \hline
Z.M.Bobur & 3 928 & 3 610 & 39,2 & 10,9 \\ \hline
T.Mirzayev & 4 553 & 3 371 & 24,0 & 7,1 \\ \hline
M.Toʻjiboyev & 4 603 & 2 532 & 17,3 & 6,8 \\ \hline
U.Qoraboyev & 2 698 & 2 420 & 32,2 & 13,3 \\ \hline
Istiqlol & 4 458 & 2 146 & 34,1 & 15,9 \\ \hline
Paxtakor & 1 854 & 1 601 & 12,9 & 8,0 \\ \hline
Y.Oxunboboyev & 5 677 & 1 191 & 0,0 & 0,0 \\ \hline
Jami & 32 872 & 20 662 & 214,6 & 10,4 \\
\end{longtable}

}

\oakilova{16-ilova}

\oakilovanom{\textbf{Qoʻllanilgan tahlil usullari}}

{\fontsize{12}{14}\selectfont\setlength{\tabcolsep}{4pt}

\noindent\begin{longtable}[]{|
  >{\raggedright\arraybackslash}p{(\columnwidth - 6\tabcolsep - 4\arrayrulewidth) * \real{0.3282}}|>{\raggedright\arraybackslash}p{(\columnwidth - 6\tabcolsep - 4\arrayrulewidth) * \real{0.5203}}|>{\raggedright\arraybackslash}p{(\columnwidth - 6\tabcolsep - 4\arrayrulewidth) * \real{0.1515}}|}
\hline
\begin{minipage}[b]{\linewidth}\raggedright
\textbf{Masala}
\end{minipage} & \begin{minipage}[b]{\linewidth}\raggedright
\textbf{Usul}
\end{minipage} & \begin{minipage}[b]{\linewidth}\raggedright
\textbf{Paragraf}
\end{minipage} \\
\hline
\endhead
\hline
\endlastfoot
Koʻp yillik oʻzgarish & Mann--Kendall mezoni, Sen qiyaligi, chiziqli
regressiya & 3.2 \\ \hline
Mavsumiy oʻzgarish & Oʻn kunliklar boʻyicha koʻp yillik oʻrtacha &
3.2 \\ \hline
Sath xaritasi & Variogramma, oddiy kriging, xarita aniqligini tekshirish
(kross-validatsiya) & 3.3 \\ \hline
Tarmoq bilan bogʻliqlik & Spearman korrelyatsiyasi, koʻp omilli
regressiya & 3.4 \\ \hline
Taʼmir taʼsiri & Taʼmirlangan va taʼmirlanmagan kollektorlar yonidagi
quduqlarni taʼmirdan oldin va keyin solishtirish («farqlar farqi») &
4.3 \\ \hline
Zarur quduqlar soni & Taʼsirni aniqlash ehtimoli hisobi (Monte Karlo) &
4.5 \\
\end{longtable}

}

\oakilova{17-ilova}

\oakilovanom{\textbf{Monitoring maʼlumotlarining yillar boʻyicha tekshiruvi}}

{\fontsize{11}{13}\selectfont\setlength{\tabcolsep}{2.5pt}

\noindent\begin{longtable}[]{|
  >{\raggedright\arraybackslash}p{(\columnwidth - 14\tabcolsep - 8\arrayrulewidth) * \real{0.0638}}|>{\raggedright\arraybackslash}p{(\columnwidth - 14\tabcolsep - 8\arrayrulewidth) * \real{0.1199}}|>{\raggedright\arraybackslash}p{(\columnwidth - 14\tabcolsep - 8\arrayrulewidth) * \real{0.1512}}|>{\raggedright\arraybackslash}p{(\columnwidth - 14\tabcolsep - 8\arrayrulewidth) * \real{0.1618}}|>{\raggedright\arraybackslash}p{(\columnwidth - 14\tabcolsep - 8\arrayrulewidth) * \real{0.1715}}|>{\raggedright\arraybackslash}p{(\columnwidth - 14\tabcolsep - 8\arrayrulewidth) * \real{0.1558}}|>{\raggedright\arraybackslash}p{(\columnwidth - 14\tabcolsep - 8\arrayrulewidth) * \real{0.1761}}|}
\hline
\begin{minipage}[b]{\linewidth}\raggedright
\textbf{Yil}
\end{minipage} & \begin{minipage}[b]{\linewidth}\raggedright
\textbf{Quduqlar}
\end{minipage} & \begin{minipage}[b]{\linewidth}\raggedright
\textbf{Oʻlchovlar}
\end{minipage} & \begin{minipage}[b]{\linewidth}\raggedright
\textbf{Oʻzgarmas qatordagi oʻlchovlar}
\end{minipage} & \begin{minipage}[b]{\linewidth}\raggedright
\textbf{Yaroqli oʻlchovlar, \%}
\end{minipage} & \begin{minipage}[b]{\linewidth}\raggedright
\textbf{Kesilgan oʻlchovlar}
\end{minipage} & \begin{minipage}[b]{\linewidth}\raggedright
\textbf{5 m dan chuqur oʻlchovlar, \%}
\end{minipage} \\
\hline
\endhead
\hline
\endlastfoot
2009 & 132 & 4 410 & 392 & 91,1 & 155 & 0,0 \\ \hline
2011 & 138 & 4 907 & 1 044 & 78,7 & 29 & 0,0 \\ \hline
2012 & 133 & 4 733 & 1 440 & 69,6 & 35 & 0,0 \\ \hline
2014 & 137 & 4 788 & 756 & 84,2 & 95 & 0,1 \\ \hline
2015 & 134 & 4 755 & 711 & 85,0 & 84 & 0,0 \\ \hline
2016 & 128 & 4 508 & 612 & 86,4 & 137 & 0,0 \\ \hline
2018 & 188 & 6 763 & 468 & 93,1 & 662 & 3,7 \\ \hline
2019 & 188 & 6 764 & 576 & 91,5 & 643 & 4,0 \\ \hline
2020 & 188 & 5 736 & 327 & 94,3 & 500 & 3,4 \\ \hline
2021 & 188 & 6 764 & 504 & 92,5 & 1 023 & 4,5 \\ \hline
2022 & 188 & 6 768 & 792 & 88,3 & 691 & 4,2 \\ \hline
2023 & 188 & 6 204 & 693 & 88,8 & 615 & 4,0 \\ \hline
2024 & 188 & 6 768 & 792 & 88,3 & 657 & 4,1 \\ \hline
Jami & 188 & 73 868 & 9 107 & 87,7 & 5 326 & --- \\
\end{longtable}

}

\oakilova{18-ilova}

\oakilovanom{\textbf{Oʻzgarmas qatorlarda eng koʻp uchraydigan qiymatlar}}

{\fontsize{12}{14}\selectfont\setlength{\tabcolsep}{4pt}

\noindent\begin{longtable}[]{|
  >{\raggedright\arraybackslash}p{(\columnwidth - 4\tabcolsep - 3\arrayrulewidth) * \real{0.4803}}|>{\raggedright\arraybackslash}p{(\columnwidth - 4\tabcolsep - 3\arrayrulewidth) * \real{0.5197}}|}
\hline
\begin{minipage}[b]{\linewidth}\raggedright
\textbf{Yozilgan qiymat, m}
\end{minipage} & \begin{minipage}[b]{\linewidth}\raggedright
\textbf{Quduq-yillar soni}
\end{minipage} \\
\hline
\endhead
\hline
\endlastfoot
3,0 & 65 \\ \hline
2,8 & 17 \\ \hline
3,2 & 12 \\ \hline
2,6 & 12 \\ \hline
5,4 & 10 \\ \hline
2,5 & 9 \\
\end{longtable}

}

\oakilova{19-ilova}

\oakilovanom{\textbf{Monitoring tarmogʻi yillar boʻyicha: a --- kuzatilgan quduqlar
soni; b --- yaroqsiz oʻlchovlar ulushi; c --- 5 m dan chuqur oʻlchovlar
ulushi (uzuq chiziq --- tarmoq tarkibi oʻzgargan yil)}}

\oakimg{figures/ch3_audit.png}

\emph{Kuzatuv quduqlarining tahlillar boʻyicha oqimi (manba → filtr →
tahlil)}

{\fontsize{12}{14}\selectfont\setlength{\tabcolsep}{4pt}

\noindent\begin{longtable}[]{|
  >{\raggedright\arraybackslash}p{(\columnwidth - 8\tabcolsep - 5\arrayrulewidth) * \real{0.3096}}|>{\raggedright\arraybackslash}p{(\columnwidth - 8\tabcolsep - 5\arrayrulewidth) * \real{0.1409}}|>{\raggedright\arraybackslash}p{(\columnwidth - 8\tabcolsep - 5\arrayrulewidth) * \real{0.3411}}|>{\raggedright\arraybackslash}p{(\columnwidth - 8\tabcolsep - 5\arrayrulewidth) * \real{0.2084}}|}
\hline
\begin{minipage}[b]{\linewidth}\raggedright
\textbf{Bosqich}
\end{minipage} & \begin{minipage}[b]{\linewidth}\raggedright
\textbf{Quduqlar soni}
\end{minipage} & \begin{minipage}[b]{\linewidth}\raggedright
\textbf{Tanlash sharti}
\end{minipage} & \begin{minipage}[b]{\linewidth}\raggedright
\textbf{Qoʻllanilgan joyi}
\end{minipage} \\
\hline
\endhead
\hline
\endlastfoot
Quduqlar pasporti & 188 & pasportda qayd etilgan & 2.2-§ \\ \hline
Koordinatalar jadvali (GAT eksporti) & 205 & koordinatasi bor & 2.2-§ \\ \hline
2024-yilgi xarita & 202 & xaritada belgilangan & 2.2-§ \\ \hline
Kuzatuv jadvallari, 2009--2024 & 188 & kamida bitta sath oʻlchovi (jami
qatorlari chiqarilgan) & 3.1-§ \\ \hline
Xaritani koordinataga bogʻlash & 190 → 181 & raqami koordinatalar
jadvali bilan mos; xatosi kattalar chiqarilgan & 2.3-§ \\ \hline
Pasport chuqurligi bilan solishtirish & 182 & pasport chuqurligi va
2019--2024-yillar sathi bor & 3.1-§ \\ \hline
2019--2024-yillar tozalangan sathi & 168 & tekshiruvdan oʻtgan
oʻlchovlar, koordinata bor & 3.3--3.4-§ \\ \hline
Yozgi (iyul--avgust) sath & 167 & 2019--2024-yillarda iyul--avgust
oʻlchovi bor & 3.3--3.5-§ \\ \hline
Quduqlar boʻyicha oʻzgarish & 122 & kamida 8 yil yaroqli maʼlumot &
3.2-§ \\ \hline
Sath oʻzgarishi va tarmoq & 112 & 122 tadan koordinatasi va
2019--2024-yillar sathi borlari & 3.4-§ \\ \hline
Doimiy kuzatilgan quduqlar (barqaror panel) & 84 & barcha 13 yilda
yaroqli oʻlchov, hech qachon oʻzgarmas qiymat bermagan & 3.2-§ \\ \hline
Minerallashuv, 2004--2010 & 122 & ishonchli davrda oʻlchov bor &
3.5-§ \\ \hline
Shoʻrlanish xavfi & 113 & minerallashuv va yozgi sath bor & 3.5-§,
4.5-§ \\ \hline
Taʼmir taʼsiri (taʼsir / nazorat) & 36 / 49 & kollektordan 300 m gacha /
600 m dan uzoq & 4.3-§ \\
\end{longtable}

}

\oakilova{20-ilova}

\oakilovanom{\textbf{Monitoring maʼlumotlari ishonchliligining massivlar boʻyicha
koʻrsatkichlari}}

{\fontsize{11}{13}\selectfont\setlength{\tabcolsep}{2.5pt}

\noindent\begin{longtable}[]{|
  >{\raggedright\arraybackslash}p{(\columnwidth - 12\tabcolsep - 7\arrayrulewidth) * \real{0.2070}}|>{\raggedright\arraybackslash}p{(\columnwidth - 12\tabcolsep - 7\arrayrulewidth) * \real{0.1191}}|>{\raggedright\arraybackslash}p{(\columnwidth - 12\tabcolsep - 7\arrayrulewidth) * \real{0.1504}}|>{\raggedright\arraybackslash}p{(\columnwidth - 12\tabcolsep - 7\arrayrulewidth) * \real{0.1713}}|>{\raggedright\arraybackslash}p{(\columnwidth - 12\tabcolsep - 7\arrayrulewidth) * \real{0.1906}}|>{\raggedright\arraybackslash}p{(\columnwidth - 12\tabcolsep - 7\arrayrulewidth) * \real{0.1616}}|}
\hline
\begin{minipage}[b]{\linewidth}\raggedright
\textbf{Massiv}
\end{minipage} & \begin{minipage}[b]{\linewidth}\raggedright
\textbf{Quduqlar}
\end{minipage} & \begin{minipage}[b]{\linewidth}\raggedright
\textbf{Oʻlchovlar}
\end{minipage} & \begin{minipage}[b]{\linewidth}\raggedright
\textbf{Yaroqli oʻlchovlar, \%}
\end{minipage} & \begin{minipage}[b]{\linewidth}\raggedright
\textbf{Yaroqli oʻlchovlar ichida kesilganlar, \%}
\end{minipage} & \begin{minipage}[b]{\linewidth}\raggedright
\textbf{Oʻrtacha sath, 2019--2024, m}
\end{minipage} \\
\hline
\endhead
\hline
\endlastfoot
Mashʼal & 27 & 11 075 & 97,9 & 2,9 & 1,91 \\ \hline
Istiqlol & 25 & 9 539 & 97,2 & 5,8 & 2,08 \\ \hline
Paxtakor & 19 & 7 806 & 93,9 & 5,6 & 2,39 \\ \hline
Z.M.Bobur & 26 & 10 384 & 92,4 & 5,8 & 2,08 \\ \hline
M.Toʻjiboyev & 28 & 10 154 & 87,8 & 10,5 & 2,36 \\ \hline
U.Qoraboyev & 22 & 8 545 & 84,1 & 13,3 & 2,94 \\ \hline
T.Mirzayev & 18 & 7 414 & 60,0 & 20,8 & 3,16 \\ \hline
Y.Oxunboboyev & 7 & 2 844 & 50,3 & 29,2 & 3,36 \\
\end{longtable}

}

\oakilova{21-ilova}

\oakilovanom{\textbf{Doimiy kuzatilgan 84 ta quduq boʻyicha sizot suvining oʻrtacha
sathi, m}}

{\fontsize{10}{12}\selectfont\setlength{\tabcolsep}{2.5pt}

\noindent\begin{longtable}[]{|
  >{\raggedright\arraybackslash}p{(\columnwidth - 28\tabcolsep - 15\arrayrulewidth) * \real{0.1787}}|>{\raggedright\arraybackslash}p{(\columnwidth - 28\tabcolsep - 15\arrayrulewidth) * \real{0.0632}}|>{\raggedright\arraybackslash}p{(\columnwidth - 28\tabcolsep - 15\arrayrulewidth) * \real{0.0632}}|>{\raggedright\arraybackslash}p{(\columnwidth - 28\tabcolsep - 15\arrayrulewidth) * \real{0.0632}}|>{\raggedright\arraybackslash}p{(\columnwidth - 28\tabcolsep - 15\arrayrulewidth) * \real{0.0632}}|>{\raggedright\arraybackslash}p{(\columnwidth - 28\tabcolsep - 15\arrayrulewidth) * \real{0.0632}}|>{\raggedright\arraybackslash}p{(\columnwidth - 28\tabcolsep - 15\arrayrulewidth) * \real{0.0632}}|>{\raggedright\arraybackslash}p{(\columnwidth - 28\tabcolsep - 15\arrayrulewidth) * \real{0.0632}}|>{\raggedright\arraybackslash}p{(\columnwidth - 28\tabcolsep - 15\arrayrulewidth) * \real{0.0632}}|>{\raggedright\arraybackslash}p{(\columnwidth - 28\tabcolsep - 15\arrayrulewidth) * \real{0.0632}}|>{\raggedright\arraybackslash}p{(\columnwidth - 28\tabcolsep - 15\arrayrulewidth) * \real{0.0632}}|>{\raggedright\arraybackslash}p{(\columnwidth - 28\tabcolsep - 15\arrayrulewidth) * \real{0.0632}}|>{\raggedright\arraybackslash}p{(\columnwidth - 28\tabcolsep - 15\arrayrulewidth) * \real{0.0632}}|>{\raggedright\arraybackslash}p{(\columnwidth - 28\tabcolsep - 15\arrayrulewidth) * \real{0.0632}}|}
\hline
\begin{minipage}[b]{\linewidth}\raggedright
\textbf{Yil}
\end{minipage} & \begin{minipage}[b]{\linewidth}\raggedright
\textbf{2009}
\end{minipage} & \begin{minipage}[b]{\linewidth}\raggedright
\textbf{2011}
\end{minipage} & \begin{minipage}[b]{\linewidth}\raggedright
\textbf{2012}
\end{minipage} & \begin{minipage}[b]{\linewidth}\raggedright
\textbf{2014}
\end{minipage} & \begin{minipage}[b]{\linewidth}\raggedright
\textbf{2015}
\end{minipage} & \begin{minipage}[b]{\linewidth}\raggedright
\textbf{2016}
\end{minipage} & \begin{minipage}[b]{\linewidth}\raggedright
\textbf{2018}
\end{minipage} & \begin{minipage}[b]{\linewidth}\raggedright
\textbf{2019}
\end{minipage} & \begin{minipage}[b]{\linewidth}\raggedright
\textbf{2020}
\end{minipage} & \begin{minipage}[b]{\linewidth}\raggedright
\textbf{2021}
\end{minipage} & \begin{minipage}[b]{\linewidth}\raggedright
\textbf{2022}
\end{minipage} & \begin{minipage}[b]{\linewidth}\raggedright
\textbf{2023}
\end{minipage} & \begin{minipage}[b]{\linewidth}\raggedright
\textbf{2024}
\end{minipage} \\
\hline
\endhead
\hline
\endlastfoot
Yillik oʻrtacha & 1,60 & 1,53 & 1,69 & 1,60 & 1,62 & 1,73 & 1,85 & 1,82
& 1,95 & 2,03 & 1,88 & 1,93 & 1,99 \\ \hline
Iyul--avgust & 1,44 & 1,32 & 1,47 & 1,37 & 1,38 & 1,55 & 1,69 & 1,60 &
1,67 & 1,92 & 1,69 & 1,71 & 1,80 \\
\end{longtable}

}

\oakilova{22-ilova}

\oakilovanom{\textbf{Quduqlar boʻyicha sizot suvi sathi oʻzgarishining massivlardagi
taqsimoti (kamida 8 yil maʼlumotli quduqlar)}}

{\fontsize{12}{14}\selectfont\setlength{\tabcolsep}{4pt}

\noindent\begin{longtable}[]{|
  >{\raggedright\arraybackslash}p{(\columnwidth - 10\tabcolsep - 6\arrayrulewidth) * \real{0.2163}}|>{\raggedright\arraybackslash}p{(\columnwidth - 10\tabcolsep - 6\arrayrulewidth) * \real{0.1297}}|>{\raggedright\arraybackslash}p{(\columnwidth - 10\tabcolsep - 6\arrayrulewidth) * \real{0.1654}}|>{\raggedright\arraybackslash}p{(\columnwidth - 10\tabcolsep - 6\arrayrulewidth) * \real{0.2758}}|>{\raggedright\arraybackslash}p{(\columnwidth - 10\tabcolsep - 6\arrayrulewidth) * \real{0.2127}}|}
\hline
\begin{minipage}[b]{\linewidth}\raggedright
\textbf{Massiv}
\end{minipage} & \begin{minipage}[b]{\linewidth}\raggedright
\textbf{Quduqlar soni}
\end{minipage} & \begin{minipage}[b]{\linewidth}\raggedright
\textbf{Sen qiyaligi medianasi, m/yil}
\end{minipage} & \begin{minipage}[b]{\linewidth}\raggedright
\textbf{Chuqurlashayotgan quduqlar ulushi, \%}
\end{minipage} & \begin{minipage}[b]{\linewidth}\raggedright
\textbf{Ahamiyatli chuqurlashish (p \textless{} 0,05), quduq}
\end{minipage} \\
\hline
\endhead
\hline
\endlastfoot
Istiqlol & 19 & +0,014 & 74 & 7 \\ \hline
M.Toʻjiboyev & 13 & +0,041 & 77 & 6 \\ \hline
Mashʼal & 21 & +0,034 & 86 & 13 \\ \hline
Paxtakor & 14 & +0,022 & 79 & 7 \\ \hline
T.Mirzayev & 7 & +0,037 & 86 & 5 \\ \hline
U.Qoraboyev & 13 & +0,034 & 92 & 9 \\ \hline
Y.Oxunboboyev & 2 & +0,068 & 100 & 1 \\ \hline
Z.M.Bobur & 23 & −0,019 & 22 & 1 \\ \hline
Jami & 122 & +0,023 & 70 & 54 \\
\end{longtable}

}

\oakilova{23-ilova}

\oakilovanom{\textbf{Sizot suvi sathining massivlar boʻyicha mavsumiy koʻrsatkichlari
(doimiy kuzatilgan quduqlar, 2009--2024-yillar oʻrtachasi)}}

{\fontsize{12}{14}\selectfont\setlength{\tabcolsep}{4pt}

\noindent\begin{longtable}[]{|
  >{\raggedright\arraybackslash}p{(\columnwidth - 12\tabcolsep - 7\arrayrulewidth) * \real{0.2341}}|>{\raggedright\arraybackslash}p{(\columnwidth - 12\tabcolsep - 7\arrayrulewidth) * \real{0.1067}}|>{\raggedright\arraybackslash}p{(\columnwidth - 12\tabcolsep - 7\arrayrulewidth) * \real{0.1746}}|>{\raggedright\arraybackslash}p{(\columnwidth - 12\tabcolsep - 7\arrayrulewidth) * \real{0.1243}}|>{\raggedright\arraybackslash}p{(\columnwidth - 12\tabcolsep - 7\arrayrulewidth) * \real{0.1808}}|>{\raggedright\arraybackslash}p{(\columnwidth - 12\tabcolsep - 7\arrayrulewidth) * \real{0.1794}}|}
\hline
\begin{minipage}[b]{\linewidth}\raggedright
\textbf{Massiv}
\end{minipage} & \begin{minipage}[b]{\linewidth}\raggedright
\textbf{Eng sayoz sath, m}
\end{minipage} & \begin{minipage}[b]{\linewidth}\raggedright
\textbf{Eng sayoz davr}
\end{minipage} & \begin{minipage}[b]{\linewidth}\raggedright
\textbf{Eng chuqur sath, m}
\end{minipage} & \begin{minipage}[b]{\linewidth}\raggedright
\textbf{Eng chuqur davr}
\end{minipage} & \begin{minipage}[b]{\linewidth}\raggedright
\textbf{Amplituda, m}
\end{minipage} \\
\hline
\endhead
\hline
\endlastfoot
Istiqlol & 1,72 & iyul, 3-dekada & 2,29 & fevral, 2-dekada & 0,57 \\ \hline
M.Toʻjiboyev & 1,55 & iyul, 3-dekada & 1,98 & fevral, 1-dekada & 0,43 \\ \hline
Mashʼal & 1,48 & may, 2-dekada & 1,84 & yanvar, 3-dekada & 0,36 \\ \hline
Paxtakor & 1,82 & iyun, 2-dekada & 2,47 & yanvar, 3-dekada & 0,65 \\ \hline
T.Mirzayev & 1,53 & may, 1-dekada & 1,86 & dekabr, 3-dekada & 0,32 \\ \hline
U.Qoraboyev & 1,45 & iyul, 3-dekada & 2,29 & fevral, 2-dekada & 0,85 \\ \hline
Z.M.Bobur & 1,46 & may, 2-dekada & 1,97 & yanvar, 3-dekada & 0,51 \\
\end{longtable}

}

\oakilova{24-ilova}

\oakilovanom{\textbf{Sizot suvi sathining massivlar boʻyicha yil ichidagi oʻrtacha
yoʻli (doimiy kuzatilgan quduqlar, 36 dekada)}}

\oakimg{figures/ch3_season_massif.png}

\oakilova{25-ilova}

\oakilovanom{\textbf{Kriging uchun boshlangʻich maʼlumotlarning statistik tavsifi}}

{\fontsize{9}{11}\selectfont\setlength{\tabcolsep}{2.5pt}

\noindent\begin{longtable}[]{|
  >{\raggedright\arraybackslash}p{(\columnwidth - 18\tabcolsep - 10\arrayrulewidth) * \real{0.1385}}|>{\raggedright\arraybackslash}p{(\columnwidth - 18\tabcolsep - 10\arrayrulewidth) * \real{0.0969}}|>{\raggedright\arraybackslash}p{(\columnwidth - 18\tabcolsep - 10\arrayrulewidth) * \real{0.1088}}|>{\raggedright\arraybackslash}p{(\columnwidth - 18\tabcolsep - 10\arrayrulewidth) * \real{0.0967}}|>{\raggedright\arraybackslash}p{(\columnwidth - 18\tabcolsep - 10\arrayrulewidth) * \real{0.0972}}|>{\raggedright\arraybackslash}p{(\columnwidth - 18\tabcolsep - 10\arrayrulewidth) * \real{0.1328}}|>{\raggedright\arraybackslash}p{(\columnwidth - 18\tabcolsep - 10\arrayrulewidth) * \real{0.1335}}|>{\raggedright\arraybackslash}p{(\columnwidth - 18\tabcolsep - 10\arrayrulewidth) * \real{0.0978}}|>{\raggedright\arraybackslash}p{(\columnwidth - 18\tabcolsep - 10\arrayrulewidth) * \real{0.0978}}|}
\hline
\begin{minipage}[b]{\linewidth}\raggedright
\textbf{Koʻrsatkich}
\end{minipage} & \begin{minipage}[b]{\linewidth}\raggedright
\textbf{Quduqlar soni}
\end{minipage} & \begin{minipage}[b]{\linewidth}\raggedright
\textbf{Oʻrtacha, m}
\end{minipage} & \begin{minipage}[b]{\linewidth}\raggedright
\textbf{Mediana, m}
\end{minipage} & \begin{minipage}[b]{\linewidth}\raggedright
\textbf{Standart ogʻish, m}
\end{minipage} & \begin{minipage}[b]{\linewidth}\raggedright
\textbf{Asimmetriya}
\end{minipage} & \begin{minipage}[b]{\linewidth}\raggedright
\textbf{Asimmetriya (logarifm)}
\end{minipage} & \begin{minipage}[b]{\linewidth}\raggedright
\textbf{1,5 m dan sayoz quduqlar soni}
\end{minipage} & \begin{minipage}[b]{\linewidth}\raggedright
\textbf{2,0 m dan sayoz quduqlar ulushi, \%}
\end{minipage} \\
\hline
\endhead
\hline
\endlastfoot
Oʻrtacha yillik sath, 2019--2024 & 168 & 2,39 & 2,25 & 0,99 & +1,31 &
+0,36 & 28 & 40 \\ \hline
Iyul--avgust sathi, 2019--2024 & 167 & 2,14 & 1,93 & 0,92 & +1,59 &
+0,46 & 39 & 55 \\
\end{longtable}

}

\oakilova{26-ilova}

\oakilovanom{\textbf{Sizot suvi sathining (2019--2024-yillar oʻrtachasi)
variogrammalari: a --- oʻlchovlar boʻyicha variogramma va nazariy egri
chiziqlar; b --- toʻrt yoʻnalish boʻyicha variogrammalar}}

\oakimg{figures/ch3_variogram.png}

\oakilova{27-ilova}

\oakilovanom{\textbf{Variogramma egri chiziqlarini solishtirish va xarita aniqligini
tekshirish natijalari (oʻrtacha yillik sath, n = 168)}}

{\fontsize{10}{12}\selectfont\setlength{\tabcolsep}{2.5pt}

\noindent\begin{longtable}[]{|
  >{\raggedright\arraybackslash}p{(\columnwidth - 14\tabcolsep - 8\arrayrulewidth) * \real{0.1923}}|>{\raggedright\arraybackslash}p{(\columnwidth - 14\tabcolsep - 8\arrayrulewidth) * \real{0.0922}}|>{\raggedright\arraybackslash}p{(\columnwidth - 14\tabcolsep - 8\arrayrulewidth) * \real{0.1198}}|>{\raggedright\arraybackslash}p{(\columnwidth - 14\tabcolsep - 8\arrayrulewidth) * \real{0.1060}}|>{\raggedright\arraybackslash}p{(\columnwidth - 14\tabcolsep - 8\arrayrulewidth) * \real{0.1063}}|>{\raggedright\arraybackslash}p{(\columnwidth - 14\tabcolsep - 8\arrayrulewidth) * \real{0.1201}}|>{\raggedright\arraybackslash}p{(\columnwidth - 14\tabcolsep - 8\arrayrulewidth) * \real{0.2633}}|}
\hline
\begin{minipage}[b]{\linewidth}\raggedright
\textbf{Egri chiziq}
\end{minipage} & \begin{minipage}[b]{\linewidth}\raggedright
\textbf{Nagget, m²}
\end{minipage} & \begin{minipage}[b]{\linewidth}\raggedright
\textbf{Nagget / umumiy tarqoqlik}
\end{minipage} & \begin{minipage}[b]{\linewidth}\raggedright
\textbf{Taʼsir radiusi, km}
\end{minipage} & \begin{minipage}[b]{\linewidth}\raggedright
\textbf{Oʻrtacha xato (ME), m}
\end{minipage} & \begin{minipage}[b]{\linewidth}\raggedright
\textbf{Oʻrtacha kvadratik xato (RMSE), m}
\end{minipage} & \begin{minipage}[b]{\linewidth}\raggedright
\textbf{Standartlashtirilgan xato (RMSSE)}
\end{minipage} \\
\hline
\endhead
\hline
\endlastfoot
Sferik & 0,512 & 0,53 & \textgreater{} 30 & −0,001 & 0,843 & 1,06 \\ \hline
Eksponensial & 0,496 & 0,51 & \textgreater{} 30 & +0,002 & 0,819 &
0,93 \\ \hline
Gauss & 0,601 & 0,62 & 7,5 & −0,001 & 0,860 & 1,04 \\ \hline
Eksponensial (avtomatik moslashtirish) & 0,267 & 0,28 & \textgreater{}
19 & +0,005 & 0,813 & 1,19 \\
\end{longtable}

}

\oakilova{28-ilova}

\oakilovanom{\textbf{Xaritalash usullarining aniqligini tekshirish natijalari}}

{\fontsize{12}{14}\selectfont\setlength{\tabcolsep}{4pt}

\noindent\begin{longtable}[]{|
  >{\raggedright\arraybackslash}p{(\columnwidth - 10\tabcolsep - 6\arrayrulewidth) * \real{0.4064}}|>{\raggedright\arraybackslash}p{(\columnwidth - 10\tabcolsep - 6\arrayrulewidth) * \real{0.1807}}|>{\raggedright\arraybackslash}p{(\columnwidth - 10\tabcolsep - 6\arrayrulewidth) * \real{0.0944}}|>{\raggedright\arraybackslash}p{(\columnwidth - 10\tabcolsep - 6\arrayrulewidth) * \real{0.2241}}|>{\raggedright\arraybackslash}p{(\columnwidth - 10\tabcolsep - 6\arrayrulewidth) * \real{0.0944}}|}
\hline
\begin{minipage}[b]{\linewidth}\raggedright
\textbf{Usul}
\end{minipage} & \begin{minipage}[b]{\linewidth}\raggedright
\textbf{Oʻrtacha yillik sath: RMSE, m}
\end{minipage} & \begin{minipage}[b]{\linewidth}\raggedright
\textbf{RMSSE}
\end{minipage} & \begin{minipage}[b]{\linewidth}\raggedright
\textbf{Iyul--avgust sathi: RMSE, m}
\end{minipage} & \begin{minipage}[b]{\linewidth}\raggedright
\textbf{RMSSE}
\end{minipage} \\
\hline
\endhead
\hline
\endlastfoot
Barcha quduqlar oʻrtachasi & 0,991 & --- & 0,927 & --- \\ \hline
Eng yaqin quduq qiymati & 1,052 & --- & 0,948 & --- \\ \hline
Oddiy kriging & 0,819 & 0,93 & 0,765 & 0,93 \\ \hline
Universal kriging (chiziqli trend) & 0,828 & 0,96 & 0,782 & 0,97 \\ \hline
Regressiya-kriging (tarmoq koʻrsatkichlari) & 0,817 & 0,95 & 0,780 &
0,97 \\
\end{longtable}

}

\oakilova{29-ilova}

\oakilovanom{\textbf{Yozgi sizot suvi sathining massivlar boʻyicha xarita
koʻrsatkichlari}}

{\fontsize{11}{13}\selectfont\setlength{\tabcolsep}{2.5pt}

\noindent\begin{longtable}[]{|
  >{\raggedright\arraybackslash}p{(\columnwidth - 14\tabcolsep - 8\arrayrulewidth) * \real{0.2014}}|>{\raggedright\arraybackslash}p{(\columnwidth - 14\tabcolsep - 8\arrayrulewidth) * \real{0.1412}}|>{\raggedright\arraybackslash}p{(\columnwidth - 14\tabcolsep - 8\arrayrulewidth) * \real{0.1551}}|>{\raggedright\arraybackslash}p{(\columnwidth - 14\tabcolsep - 8\arrayrulewidth) * \real{0.1126}}|>{\raggedright\arraybackslash}p{(\columnwidth - 14\tabcolsep - 8\arrayrulewidth) * \real{0.1269}}|>{\raggedright\arraybackslash}p{(\columnwidth - 14\tabcolsep - 8\arrayrulewidth) * \real{0.1412}}|>{\raggedright\arraybackslash}p{(\columnwidth - 14\tabcolsep - 8\arrayrulewidth) * \real{0.1216}}|}
\hline
\begin{minipage}[b]{\linewidth}\raggedright
\textbf{Massiv}
\end{minipage} & \begin{minipage}[b]{\linewidth}\raggedright
\textbf{Xarita bilan qoplangan maydon, \%}
\end{minipage} & \begin{minipage}[b]{\linewidth}\raggedright
\textbf{Baholangan oʻrtacha sath, m}
\end{minipage} & \begin{minipage}[b]{\linewidth}\raggedright
\textbf{Sathi 2,0 m dan sayoz maydon, \%}
\end{minipage} & \begin{minipage}[b]{\linewidth}\raggedright
\textbf{Sayoz sath ehtimoli 0,3 dan katta maydon, \%}
\end{minipage} & \begin{minipage}[b]{\linewidth}\raggedright
\textbf{Quduqlar: jami / 1,5 m dan sayoz}
\end{minipage} & \begin{minipage}[b]{\linewidth}\raggedright
\textbf{Standart xato, m}
\end{minipage} \\
\hline
\endhead
\hline
\endlastfoot
Istiqlol & 74 & 1,96 & 61 & 55 & 25 / 10 & 0,82 \\ \hline
M.Toʻjiboyev & 81 & 2,43 & 20 & 17 & 27 / 4 & 0,81 \\ \hline
Mashʼal & 86 & 1,84 & 80 & 71 & 27 / 9 & 0,82 \\ \hline
Paxtakor & 100 & 2,03 & 44 & 21 & 19 / 1 & 0,78 \\ \hline
T.Mirzayev & 81 & 2,66 & 30 & 28 & 14 / 1 & 0,84 \\ \hline
U.Qoraboyev & 96 & 2,43 & 45 & 19 & 22 / 2 & 0,80 \\ \hline
Y.Oxunboboyev & 57 & 3,15 & 0 & 0 & 7 / 0 & 0,85 \\ \hline
Z.M.Bobur & 85 & 2,00 & 61 & 58 & 26 / 12 & 0,82 \\ \hline
Jami & 82 & --- & 49 & 40 & 167 / 39 & --- \\
\end{longtable}

}

\emph{Kritik chuqurlik va ehtimollik usuliga sezgirlik.} Ehtimoli 0,3
dan katta maydon, \%; shoʻrlanish xavfli quduqlar --- yozgi sath
chegaradan sayoz va minerallashuv 1 g/l dan yuqori (113 ta quduqdan).

{\fontsize{12}{14}\selectfont\setlength{\tabcolsep}{4pt}

\noindent\begin{longtable}[]{|
  >{\raggedright\arraybackslash}p{(\columnwidth - 12\tabcolsep - 7\arrayrulewidth) * \real{0.2906}}|>{\raggedright\arraybackslash}p{(\columnwidth - 12\tabcolsep - 7\arrayrulewidth) * \real{0.0913}}|>{\raggedright\arraybackslash}p{(\columnwidth - 12\tabcolsep - 7\arrayrulewidth) * \real{0.0913}}|>{\raggedright\arraybackslash}p{(\columnwidth - 12\tabcolsep - 7\arrayrulewidth) * \real{0.1279}}|>{\raggedright\arraybackslash}p{(\columnwidth - 12\tabcolsep - 7\arrayrulewidth) * \real{0.2100}}|>{\raggedright\arraybackslash}p{(\columnwidth - 12\tabcolsep - 7\arrayrulewidth) * \real{0.1888}}|}
\hline
\begin{minipage}[b]{\linewidth}\raggedright
\textbf{Koʻrsatkich}
\end{minipage} & \begin{minipage}[b]{\linewidth}\raggedright
\textbf{1,5 m}
\end{minipage} & \begin{minipage}[b]{\linewidth}\raggedright
\textbf{2,0 m}
\end{minipage} & \begin{minipage}[b]{\linewidth}\raggedright
\textbf{2,5 m}
\end{minipage} & \begin{minipage}[b]{\linewidth}\raggedright
\textbf{1,5 m, log-kriging}
\end{minipage} & \begin{minipage}[b]{\linewidth}\raggedright
\textbf{1,5 m, indikator kriging}
\end{minipage} \\
\hline
\endhead
\hline
\endlastfoot
Yozda sathi chegaradan sayoz quduqlar, \% (soni) & 23 (39) & 55 (92) &
78 (130) & --- & --- \\ \hline
Kriging bahosi chegaradan sayoz maydon, \% & 1,1 & 48,6 & 74,4 & --- &
--- \\ \hline
Ehtimoli 0,3 dan katta maydon: tuman & 40 & 71 & 86 & 36 & 19 \\ \hline
Mashʼal & 71 & 98 & 100 & 60 & 27 \\ \hline
Z.M.Bobur & 58 & 82 & 96 & 59 & 42 \\ \hline
Istiqlol & 55 & 89 & 100 & 55 & 31 \\ \hline
Boshqa massivlar (oraliq) & 0--28 & 0--95 & 29--100 & 0--25 & 0--13 \\ \hline
Shoʻrlanish xavfli quduqlar & 29 & 56 & 73 & --- & --- \\
\end{longtable}

}

\oakilova{30-ilova}

\oakilovanom{\textbf{Modelga kiritilgan quduqlar soni va nazorat quduqlaridagi
baholash xatosi (40 ta takrorlash oʻrtachasi va standart ogʻishi).
Raqamlar --- qoʻshni quduqlar orasidagi masofaning medianasi}}

\oakimg{figures/ch3_density.png}

\oakilova{31-ilova}

\oakilovanom{\textbf{Sizot suvi sathining kollektorgacha masofaga bogʻliqligi: a ---
oʻrtacha sath va boʻlinma regressiya chizigʻi; b --- qishki va yozgi
sath (A --- mavsumiy koʻtarilish); c --- quduqlar boʻyicha Sen
qiyaligining medianasi}}

\oakimg{figures/ch3_network.png}

\oakilova{32-ilova}

\oakilovanom{\textbf{Sizot suvi sathining koʻp omilli regressiya modellari (n = 168)}}

{\fontsize{12}{14}\selectfont\setlength{\tabcolsep}{4pt}

\noindent\begin{longtable}[]{|
  >{\raggedright\arraybackslash}p{(\columnwidth - 10\tabcolsep - 6\arrayrulewidth) * \real{0.3179}}|>{\raggedright\arraybackslash}p{(\columnwidth - 10\tabcolsep - 6\arrayrulewidth) * \real{0.2390}}|>{\raggedright\arraybackslash}p{(\columnwidth - 10\tabcolsep - 6\arrayrulewidth) * \real{0.1021}}|>{\raggedright\arraybackslash}p{(\columnwidth - 10\tabcolsep - 6\arrayrulewidth) * \real{0.2390}}|>{\raggedright\arraybackslash}p{(\columnwidth - 10\tabcolsep - 6\arrayrulewidth) * \real{0.1021}}|}
\hline
\begin{minipage}[b]{\linewidth}\raggedright
\textbf{Omil}
\end{minipage} & \begin{minipage}[b]{\linewidth}\raggedright
\textbf{1-model: koeffitsiyent}
\end{minipage} & \begin{minipage}[b]{\linewidth}\raggedright
\textbf{p}
\end{minipage} & \begin{minipage}[b]{\linewidth}\raggedright
\textbf{2-model: koeffitsiyent}
\end{minipage} & \begin{minipage}[b]{\linewidth}\raggedright
\textbf{p}
\end{minipage} \\
\hline
\endhead
\hline
\endlastfoot
Kollektorgacha masofa, lg(m) & +0,73 & \textless{} 0,001 & & \\ \hline
1 km radiusdagi kollektor uzunligi, km & & & −0,18 & \textless{}
0,001 \\ \hline
Kanalgacha masofa, lg(m) & −0,15 & 0,28 & −0,06 & 0,71 \\ \hline
Mutlaq balandlik, m & +0,001 & 0,76 & +0,002 & 0,66 \\ \hline
R² & 0,15 & & 0,20 & \\
\end{longtable}

}

Izoh: standart xatolar HC3 usulida. Qoldiqlar oʻngga qiyshiq
(asimmetriya 1,1--1,3; Shapiro--Wilk p \textless{} 0,001), 2-modelda
xatolar tarqoqligi notekis (Breusch--Pagan p = 0,009); VIF ≤ 1,12. Sath
logarifmi boʻyicha modelda qoldiqlar normal taqsimotdan farq qilmaydi (p
= 0,22) va kollektorgacha masofa koeffitsiyenti ahamiyatli qoladi
(+0,29; 95 \% ishonch oraligʻi 0,19--0,40). Kuk masofasi 4/n dan katta
6--10 ta quduq chiqarilganda koeffitsiyentlar 10--18 foizga kamayadi va
ahamiyatli qoladi.

\oakilova{33-ilova}

\oakilovanom{\textbf{Sath va tarmoq koʻrsatkichlari bogʻlanishining massiv omilini
hisobga olgan holda bahosi (n = 168)}}

{\fontsize{12}{14}\selectfont\setlength{\tabcolsep}{4pt}

\noindent\begin{longtable}[]{|
  >{\raggedright\arraybackslash}p{(\columnwidth - 8\tabcolsep - 5\arrayrulewidth) * \real{0.4963}}|>{\raggedright\arraybackslash}p{(\columnwidth - 8\tabcolsep - 5\arrayrulewidth) * \real{0.2607}}|>{\raggedright\arraybackslash}p{(\columnwidth - 8\tabcolsep - 5\arrayrulewidth) * \real{0.1432}}|>{\raggedright\arraybackslash}p{(\columnwidth - 8\tabcolsep - 5\arrayrulewidth) * \real{0.0998}}|}
\hline
\begin{minipage}[b]{\linewidth}\raggedright
\textbf{Model}
\end{minipage} & \begin{minipage}[b]{\linewidth}\raggedright
\textbf{Koeffitsiyent}
\end{minipage} & \begin{minipage}[b]{\linewidth}\raggedright
\textbf{p}
\end{minipage} & \begin{minipage}[b]{\linewidth}\raggedright
\textbf{R²}
\end{minipage} \\
\hline
\endhead
\hline
\endlastfoot
Faqat massiv omili & --- & --- & 0,20 \\ \hline
Kollektor uzunligi (1 km radiusda), km & −0,185 & \textless{} 0,001 &
0,20 \\ \hline
Kollektor uzunligi, massiv ichida solishtirish & −0,169 & \textless{}
0,001 & 0,33 \\ \hline
Kollektorgacha masofa, lg(m) & +0,78 & \textless{} 0,001 & 0,15 \\ \hline
Kollektorgacha masofa, massiv ichida solishtirish & +0,71 & \textless{}
0,001 & 0,29 \\
\end{longtable}

}

\emph{Qoʻshni quduqlar oʻxshashligini hisobga olgan tekshiruv.} Moran I
--- eng yaqin 8 ta qoʻshni boʻyicha, ahamiyatlilik 999 ta tasodifiy
almashtirish bilan; ishonch oraligʻi --- 2×2 km kataklar boʻyicha blokli
bootstrap (2 000 takror); klasterli p --- kataklar boʻyicha
klasterlangan standart xato bilan.

{\fontsize{11}{13}\selectfont\setlength{\tabcolsep}{2.5pt}

\noindent\begin{longtable}[]{|
  >{\raggedright\arraybackslash}p{(\columnwidth - 14\tabcolsep - 8\arrayrulewidth) * \real{0.2294}}|>{\raggedright\arraybackslash}p{(\columnwidth - 14\tabcolsep - 8\arrayrulewidth) * \real{0.0463}}|>{\raggedright\arraybackslash}p{(\columnwidth - 14\tabcolsep - 8\arrayrulewidth) * \real{0.1187}}|>{\raggedright\arraybackslash}p{(\columnwidth - 14\tabcolsep - 8\arrayrulewidth) * \real{0.1261}}|>{\raggedright\arraybackslash}p{(\columnwidth - 14\tabcolsep - 8\arrayrulewidth) * \real{0.1674}}|>{\raggedright\arraybackslash}p{(\columnwidth - 14\tabcolsep - 8\arrayrulewidth) * \real{0.1691}}|>{\raggedright\arraybackslash}p{(\columnwidth - 14\tabcolsep - 8\arrayrulewidth) * \real{0.1429}}|}
\hline
\begin{minipage}[b]{\linewidth}\raggedright
\textbf{Bogʻlanish}
\end{minipage} & \begin{minipage}[b]{\linewidth}\raggedright
\textbf{n}
\end{minipage} & \begin{minipage}[b]{\linewidth}\raggedright
\textbf{Spearman ρ}
\end{minipage} & \begin{minipage}[b]{\linewidth}\raggedright
\textbf{ρ ning 95 \% blokli ishonch oraligʻi}
\end{minipage} & \begin{minipage}[b]{\linewidth}\raggedright
\textbf{Qoldiqlarda Moran I (p)}
\end{minipage} & \begin{minipage}[b]{\linewidth}\raggedright
\textbf{Massiv omili bilan qoldiqlarda Moran I (p)}
\end{minipage} & \begin{minipage}[b]{\linewidth}\raggedright
\textbf{Qiyalik p: oddiy / klasterli}
\end{minipage} \\
\hline
\endhead
\hline
\endlastfoot
Sath --- kollektorgacha masofa & 168 & +0,40 & +0,26 \ldots{} +0,52 &
0,15 (0,002) & 0,10 (0,003) & \textless{} 0,001 / \textless{} 0,001 \\ \hline
Sath --- tarmoq zichligi & 168 & −0,47 & −0,58 \ldots{} −0,33 & 0,15
(0,001) & 0,08 (0,017) & \textless{} 0,001 / \textless{} 0,001 \\ \hline
Yozgi sath --- kollektorgacha masofa & 167 & +0,37 & +0,23 \ldots{}
+0,49 & 0,14 (0,001) & 0,11 (0,005) & \textless{} 0,001 / \textless{}
0,001 \\ \hline
Yozgi sath --- tarmoq zichligi & 167 & −0,45 & −0,56 \ldots{} −0,31 &
0,13 (0,003) & 0,08 (0,007) & \textless{} 0,001 / \textless{} 0,001 \\ \hline
Sath oʻzgarishi --- kollektorgacha masofa & 112 & +0,25 & +0,01 \ldots{}
+0,45 & 0,16 (0,002) & −0,01 (0,53) & 0,002 / 0,006 \\ \hline
Sath oʻzgarishi --- tarmoq zichligi & 112 & −0,34 & −0,51 \ldots{} −0,13
& 0,14 (0,002) & −0,04 (0,74) & \textless{} 0,001 / \textless{} 0,001 \\
\end{longtable}

}

Massivlar ichidagi korrelyatsiyalar (yettita test) uchun
Benjamini--Hochberg tuzatishidan keyin (q \textless{} 0,05) U.Qoraboyev
(q = 0,021) va M.Toʻjiboyev (q = 0,025) massivlaridagi bogʻlanish
ahamiyatli qoladi; Istiqlol va Z.M.Bobur massivlarida q = 0,083.

\oakilova{34-ilova}

\oakilovanom{\textbf{Kollektordan 500 m gacha masofadagi quduqlarda sizot suvi sathi
kollektorning kadastr koʻrsatkichlari boʻyicha}}

{\fontsize{12}{14}\selectfont\setlength{\tabcolsep}{4pt}

\noindent\begin{longtable}[]{|
  >{\raggedright\arraybackslash}p{(\columnwidth - 10\tabcolsep - 6\arrayrulewidth) * \real{0.2620}}|>{\raggedright\arraybackslash}p{(\columnwidth - 10\tabcolsep - 6\arrayrulewidth) * \real{0.1481}}|>{\raggedright\arraybackslash}p{(\columnwidth - 10\tabcolsep - 6\arrayrulewidth) * \real{0.2179}}|>{\raggedright\arraybackslash}p{(\columnwidth - 10\tabcolsep - 6\arrayrulewidth) * \real{0.1649}}|>{\raggedright\arraybackslash}p{(\columnwidth - 10\tabcolsep - 6\arrayrulewidth) * \real{0.2072}}|}
\hline
\begin{minipage}[b]{\linewidth}\raggedright
\textbf{Guruh}
\end{minipage} & \begin{minipage}[b]{\linewidth}\raggedright
\textbf{Quduqlar soni}
\end{minipage} & \begin{minipage}[b]{\linewidth}\raggedright
\textbf{Kollektorlar soni}
\end{minipage} & \begin{minipage}[b]{\linewidth}\raggedright
\textbf{Yillik oʻrtacha sath, m}
\end{minipage} & \begin{minipage}[b]{\linewidth}\raggedright
\textbf{Iyul--avgust sathi, m}
\end{minipage} \\
\hline
\endhead
\hline
\endlastfoot
Kollektor yoshi 55 yildan kam & 12 & 8 & 2,05 & 1,86 \\ \hline
Kollektor yoshi 55--65 yil & 16 & 7 & 1,94 & 1,80 \\ \hline
Kollektor yoshi 65 yildan ortiq & 59 & 26 & 2,31 & 2,10 \\ \hline
Drenaj moduli past (0,04--0,69 l/(s·ga)) & 30 & --- & 1,94 & 1,69 \\ \hline
Drenaj moduli oʻrta (0,71--1,43 l/(s·ga)) & 28 & --- & 2,42 & 2,26 \\ \hline
Drenaj moduli yuqori (1,48--10,2 l/(s·ga)) & 29 & --- & 2,26 & 2,10 \\
\end{longtable}

}

\oakilova{35-ilova}

\oakilovanom{\textbf{Minerallashuv qatorlarining tekshiruvi}}

{\fontsize{12}{14}\selectfont\setlength{\tabcolsep}{4pt}

\noindent\begin{longtable}[]{|
  >{\raggedright\arraybackslash}p{(\columnwidth - 6\tabcolsep - 4\arrayrulewidth) * \real{0.2486}}|>{\raggedright\arraybackslash}p{(\columnwidth - 6\tabcolsep - 4\arrayrulewidth) * \real{0.3679}}|>{\raggedright\arraybackslash}p{(\columnwidth - 6\tabcolsep - 4\arrayrulewidth) * \real{0.3834}}|}
\hline
\begin{minipage}[b]{\linewidth}\raggedright
\textbf{Davr}
\end{minipage} & \begin{minipage}[b]{\linewidth}\raggedright
\textbf{Holat}
\end{minipage} & \begin{minipage}[b]{\linewidth}\raggedright
\textbf{Baho}
\end{minipage} \\
\hline
\endhead
\hline
\endlastfoot
1999--2003 & Mustaqil oʻlchov maʼlumotlari mavjud emas & Tahlilga
kiritilmadi \\ \hline
2004--2010 & Yillar boʻyicha qiymatlar xilma-xil, nuqson aniqlanmadi &
Tahlil uchun yaroqli \\ \hline
2011--2015 & Mustaqil oʻlchov maʼlumotlari mavjud emas & Tahlilga
kiritilmadi \\ \hline
2016--2020 (yillik) & Mediana 1,25 dan 2,49 g/l gacha oʻsgan &
Ehtiyotkorlik bilan \\ \hline
2021--2022 & Mediana 2,49 dan 0,80 g/l ga keskin tushgan & Usul yoki
birlik oʻzgargan boʻlishi mumkin \\ \hline
2023 & Barcha 188 ta quduqda bitta qiymat: 0,20 g/l & Yaroqsiz \\ \hline
2024 & Qiymatlar kam xilma-xil, mediana 0,20 g/l & Shubhali \\
\end{longtable}

}

\oakilova{36-ilova}

\oakilovanom{\textbf{Sizot suvlari minerallashuvi sinflarining yillar boʻyicha
ulushi, \%}}

{\fontsize{12}{14}\selectfont\setlength{\tabcolsep}{4pt}

\noindent\begin{longtable}[]{|
  >{\raggedright\arraybackslash}p{(\columnwidth - 12\tabcolsep - 7\arrayrulewidth) * \real{0.2577}}|>{\raggedright\arraybackslash}p{(\columnwidth - 12\tabcolsep - 7\arrayrulewidth) * \real{0.1451}}|>{\raggedright\arraybackslash}p{(\columnwidth - 12\tabcolsep - 7\arrayrulewidth) * \real{0.1146}}|>{\raggedright\arraybackslash}p{(\columnwidth - 12\tabcolsep - 7\arrayrulewidth) * \real{0.1146}}|>{\raggedright\arraybackslash}p{(\columnwidth - 12\tabcolsep - 7\arrayrulewidth) * \real{0.1660}}|>{\raggedright\arraybackslash}p{(\columnwidth - 12\tabcolsep - 7\arrayrulewidth) * \real{0.2020}}|}
\hline
\begin{minipage}[b]{\linewidth}\raggedright
\textbf{Yil}
\end{minipage} & \begin{minipage}[b]{\linewidth}\raggedright
\textbf{1 g/l dan kam}
\end{minipage} & \begin{minipage}[b]{\linewidth}\raggedright
\textbf{1--3 g/l}
\end{minipage} & \begin{minipage}[b]{\linewidth}\raggedright
\textbf{3--5 g/l}
\end{minipage} & \begin{minipage}[b]{\linewidth}\raggedright
\textbf{5 g/l va ortiq}
\end{minipage} & \begin{minipage}[b]{\linewidth}\raggedright
\textbf{Mediana, g/l}
\end{minipage} \\
\hline
\endhead
\hline
\endlastfoot
2004 & 12,2 & 75,7 & 12,2 & 0,0 & 1,38 \\ \hline
2005 & 27,3 & 56,9 & 15,4 & 0,3 & 1,45 \\ \hline
2006 & 32,5 & 59,6 & 7,6 & 0,3 & 1,26 \\ \hline
2007 & 33,9 & 59,3 & 6,2 & 0,7 & 1,24 \\ \hline
2008 & 36,2 & 56,4 & 7,4 & 0,0 & 1,22 \\ \hline
2009 & 46,4 & 51,0 & 2,6 & 0,0 & 1,02 \\ \hline
2010 & 26,3 & 64,0 & 9,8 & 0,0 & 1,35 \\ \hline
2004--2010 & 30,7 & 60,4 & 8,8 & 0,2 & --- \\
\end{longtable}

}

\oakilova{37-ilova}

\oakilovanom{\textbf{Sizot suvlari minerallashuvi va shoʻrlanish xavfining massivlar
boʻyicha koʻrsatkichlari}}

{\fontsize{9}{11}\selectfont\setlength{\tabcolsep}{2.5pt}

\noindent\begin{longtable}[]{|
  >{\raggedright\arraybackslash}p{(\columnwidth - 18\tabcolsep - 10\arrayrulewidth) * \real{0.1520}}|>{\raggedright\arraybackslash}p{(\columnwidth - 18\tabcolsep - 10\arrayrulewidth) * \real{0.0988}}|>{\raggedright\arraybackslash}p{(\columnwidth - 18\tabcolsep - 10\arrayrulewidth) * \real{0.1619}}|>{\raggedright\arraybackslash}p{(\columnwidth - 18\tabcolsep - 10\arrayrulewidth) * \real{0.1138}}|>{\raggedright\arraybackslash}p{(\columnwidth - 18\tabcolsep - 10\arrayrulewidth) * \real{0.0741}}|>{\raggedright\arraybackslash}p{(\columnwidth - 18\tabcolsep - 10\arrayrulewidth) * \real{0.0619}}|>{\raggedright\arraybackslash}p{(\columnwidth - 18\tabcolsep - 10\arrayrulewidth) * \real{0.0863}}|>{\raggedright\arraybackslash}p{(\columnwidth - 18\tabcolsep - 10\arrayrulewidth) * \real{0.1619}}|>{\raggedright\arraybackslash}p{(\columnwidth - 18\tabcolsep - 10\arrayrulewidth) * \real{0.0894}}|}
\hline
\begin{minipage}[b]{\linewidth}\raggedright
\textbf{Massiv}
\end{minipage} & \begin{minipage}[b]{\linewidth}\raggedright
\textbf{Quduqlar soni}
\end{minipage} & \begin{minipage}[b]{\linewidth}\raggedright
\textbf{Minerallashuv medianasi, g/l}
\end{minipage} & \begin{minipage}[b]{\linewidth}\raggedright
\textbf{3 g/l dan yuqori quduqlar, \%}
\end{minipage} & \begin{minipage}[b]{\linewidth}\raggedright
\textbf{Aprel, g/l}
\end{minipage} & \begin{minipage}[b]{\linewidth}\raggedright
\textbf{Iyul, g/l}
\end{minipage} & \begin{minipage}[b]{\linewidth}\raggedright
\textbf{Oktabr, g/l}
\end{minipage} & \begin{minipage}[b]{\linewidth}\raggedright
\textbf{Xavfli quduqlar (yozgi sath \textless{} 2,0 m, minerallashuv ≥ 1
g/l)}
\end{minipage} & \begin{minipage}[b]{\linewidth}\raggedright
\textbf{Shundan yozgi sath \textless{} 1,5 m}
\end{minipage} \\
\hline
\endhead
\hline
\endlastfoot
Mashʼal & 21 & 2,47 & 10 & 2,36 & 2,15 & 2,45 & 17 & 8 \\ \hline
Z.M.Bobur & 23 & 1,41 & 9 & 1,51 & 1,45 & 1,48 & 17 & 10 \\ \hline
Paxtakor & 15 & 1,40 & 0 & 1,24 & 1,33 & 1,42 & 4 & 1 \\ \hline
Istiqlol & 19 & 1,22 & 0 & 1,15 & 1,22 & 1,25 & 12 & 8 \\ \hline
M.Toʻjiboyev & 14 & 1,03 & 0 & 1,10 & 1,07 & 1,00 & 4 & 2 \\ \hline
U.Qoraboyev & 12 & 0,84 & 0 & 0,84 & 0,90 & 0,86 & 1 & 0 \\ \hline
T.Mirzayev & 9 & 0,80 & 0 & 1,02 & 0,74 & 0,80 & 1 & 0 \\ \hline
Jami & 113 & --- & --- & 1,27 & 1,25 & 1,32 & 56 & 29 \\
\end{longtable}

}

Izoh: minerallashuv --- 2004--2010-yillar, yozgi sath ---
2019--2024-yillar (iyul--avgust) boʻyicha.

\oakilova{38-ilova}

\oakilovanom{\textbf{Quduqlarning yozgi sath va minerallashuv boʻyicha taqsimoti
(quduqlar soni, n = 113)}}

{\fontsize{12}{14}\selectfont\setlength{\tabcolsep}{4pt}

\noindent\begin{longtable}[]{|
  >{\raggedright\arraybackslash}p{(\columnwidth - 10\tabcolsep - 6\arrayrulewidth) * \real{0.3690}}|>{\raggedright\arraybackslash}p{(\columnwidth - 10\tabcolsep - 6\arrayrulewidth) * \real{0.1755}}|>{\raggedright\arraybackslash}p{(\columnwidth - 10\tabcolsep - 6\arrayrulewidth) * \real{0.1163}}|>{\raggedright\arraybackslash}p{(\columnwidth - 10\tabcolsep - 6\arrayrulewidth) * \real{0.2169}}|>{\raggedright\arraybackslash}p{(\columnwidth - 10\tabcolsep - 6\arrayrulewidth) * \real{0.1222}}|}
\hline
\begin{minipage}[b]{\linewidth}\raggedright
\textbf{Yozgi sath, m}
\end{minipage} & \begin{minipage}[b]{\linewidth}\raggedright
\textbf{1 g/l dan kam}
\end{minipage} & \begin{minipage}[b]{\linewidth}\raggedright
\textbf{1--3 g/l}
\end{minipage} & \begin{minipage}[b]{\linewidth}\raggedright
\textbf{3 g/l va ortiq}
\end{minipage} & \begin{minipage}[b]{\linewidth}\raggedright
\textbf{Jami}
\end{minipage} \\
\hline
\endhead
\hline
\endlastfoot
1,5 dan sayoz & 6 & 28 & 1 & 35 \\ \hline
1,5--2,0 & 12 & 24 & 3 & 39 \\ \hline
2,0--3,0 & 8 & 22 & 0 & 30 \\ \hline
3,0 va chuqur & 1 & 8 & 0 & 9 \\ \hline
Jami & 27 & 82 & 4 & 113 \\
\end{longtable}

}

\oakilova{39-ilova}

\oakilovanom{\textbf{Andijon viloyati tumanlarida monitoring maʼlumotlarining
ishonchlilik koʻrsatkichlari, 2009--2023-yillar}}

{\fontsize{11}{13}\selectfont\setlength{\tabcolsep}{2.5pt}

\noindent\begin{longtable}[]{|
  >{\raggedright\arraybackslash}p{(\columnwidth - 12\tabcolsep - 7\arrayrulewidth) * \real{0.1846}}|>{\raggedright\arraybackslash}p{(\columnwidth - 12\tabcolsep - 7\arrayrulewidth) * \real{0.1193}}|>{\raggedright\arraybackslash}p{(\columnwidth - 12\tabcolsep - 7\arrayrulewidth) * \real{0.1406}}|>{\raggedright\arraybackslash}p{(\columnwidth - 12\tabcolsep - 7\arrayrulewidth) * \real{0.2396}}|>{\raggedright\arraybackslash}p{(\columnwidth - 12\tabcolsep - 7\arrayrulewidth) * \real{0.1655}}|>{\raggedright\arraybackslash}p{(\columnwidth - 12\tabcolsep - 7\arrayrulewidth) * \real{0.1504}}|}
\hline
\begin{minipage}[b]{\linewidth}\raggedright
\textbf{Tuman}
\end{minipage} & \begin{minipage}[b]{\linewidth}\raggedright
\textbf{Quduqlar soni}
\end{minipage} & \begin{minipage}[b]{\linewidth}\raggedright
\textbf{Mustaqil maʼlumoti yoʻq yillar}
\end{minipage} & \begin{minipage}[b]{\linewidth}\raggedright
\textbf{Oʻzgarmas qatorlar, quduq-yillarning \%i}
\end{minipage} & \begin{minipage}[b]{\linewidth}\raggedright
\textbf{Yaroqli oʻlchovlar, \%}
\end{minipage} & \begin{minipage}[b]{\linewidth}\raggedright
\textbf{Yaxlit qiymatlar, \%}
\end{minipage} \\
\hline
\endhead
\hline
\endlastfoot
Oltinkoʻl & 130 & 2 & 0,0 & 100,0 & 18,1 \\ \hline
Boʻston & 148 & 1 & 0,2 & 99,8 & 11,9 \\ \hline
Ulugʻnor & 186 & 2 & 2,5 & 97,3 & 19,9 \\ \hline
Shahrixon & 165 & 2 & 2,7 & 97,5 & 17,1 \\ \hline
Baliqchi & 211 & 2 & 3,8 & 96,4 & 19,4 \\ \hline
Buloqboshi & 154 & 2 & 4,3 & 96,5 & 21,0 \\ \hline
Xoʻjaobod & 60 & 2 & 9,6 & 90,8 & 36,1 \\ \hline
Izboskan & 150 & 2 & 11,5 & 89,5 & 26,6 \\ \hline
Marhamat & 188 & 3 & 12,5 & 88,2 & 29,1 \\ \hline
Andijon & 119 & 2 & 19,0 & 82,1 & 28,7 \\ \hline
Jalaquduq & 134 & 3 & 22,7 & 79,1 & 25,9 \\ \hline
Qoʻrgʻontepa & 108 & 2 & 26,3 & 74,4 & 42,7 \\ \hline
Paxtaobod & 88 & 2 & 28,0 & 71,9 & 19,2 \\ \hline
Asaka & 131 & 2 & 29,7 & 70,6 & 33,1 \\ \hline
Mediana & --- & 2 & 10,6 & 90,2 & --- \\
\end{longtable}

}

Izoh: Boʻston va Paxtaobod tumanlari boʻyicha 14 yillik, Buloqboshi
tumani boʻyicha 13 yillik jadvallar mavjud.

\oakilova{40-ilova}

\oakilovanom{\textbf{Andijon viloyati tumanlarida sizot suvi sathining asosiy
koʻrsatkichlari}}

{\fontsize{10}{12}\selectfont\setlength{\tabcolsep}{2.5pt}

\noindent\begin{longtable}[]{|
  >{\raggedright\arraybackslash}p{(\columnwidth - 16\tabcolsep - 9\arrayrulewidth) * \real{0.1623}}|>{\raggedright\arraybackslash}p{(\columnwidth - 16\tabcolsep - 9\arrayrulewidth) * \real{0.1327}}|>{\raggedright\arraybackslash}p{(\columnwidth - 16\tabcolsep - 9\arrayrulewidth) * \real{0.1193}}|>{\raggedright\arraybackslash}p{(\columnwidth - 16\tabcolsep - 9\arrayrulewidth) * \real{0.0677}}|>{\raggedright\arraybackslash}p{(\columnwidth - 16\tabcolsep - 9\arrayrulewidth) * \real{0.1460}}|>{\raggedright\arraybackslash}p{(\columnwidth - 16\tabcolsep - 9\arrayrulewidth) * \real{0.1196}}|>{\raggedright\arraybackslash}p{(\columnwidth - 16\tabcolsep - 9\arrayrulewidth) * \real{0.1196}}|>{\raggedright\arraybackslash}p{(\columnwidth - 16\tabcolsep - 9\arrayrulewidth) * \real{0.1327}}|}
\hline
\begin{minipage}[b]{\linewidth}\raggedright
\textbf{Tuman}
\end{minipage} & \begin{minipage}[b]{\linewidth}\raggedright
\textbf{Doimiy kuzatilgan quduqlar}
\end{minipage} & \begin{minipage}[b]{\linewidth}\raggedright
\textbf{Sen qiyaligi, sm/yil}
\end{minipage} & \begin{minipage}[b]{\linewidth}\raggedright
\textbf{p}
\end{minipage} & \begin{minipage}[b]{\linewidth}\raggedright
\textbf{Iyul--avgust sathi medianasi, m}
\end{minipage} & \begin{minipage}[b]{\linewidth}\raggedright
\textbf{Yozda 1,5 m dan sayoz quduqlar, \%}
\end{minipage} & \begin{minipage}[b]{\linewidth}\raggedright
\textbf{Yozda 2,0 m dan sayoz quduqlar, \%}
\end{minipage} & \begin{minipage}[b]{\linewidth}\raggedright
\textbf{Yil ichidagi tebranish, m}
\end{minipage} \\
\hline
\endhead
\hline
\endlastfoot
Oltinkoʻl & 128 & +1,2 & 0,10 & 1,68 & 11,5 & 93,8 & 0,11 \\ \hline
Boʻston & 141 & +0,8 & 0,08 & 1,57 & 28,4 & 98,6 & 0,21 \\ \hline
Ulugʻnor & 54 & +0,7 & 0,024 & 1,82 & 5,6 & 87,0 & 0,07 \\ \hline
Shahrixon & 111 & +1,8 & 0,001 & 1,80 & 7,3 & 72,7 & 0,43 \\ \hline
Baliqchi & 100 & +3,5 & \textless{} 0,001 & 1,73 & 13,0 & 79,7 & 0,18 \\ \hline
Buloqboshi & 34 & +3,5 & 0,003 & 2,05 & 15,6 & 46,9 & 0,24 \\ \hline
Xoʻjaobod & 25 & +4,5 & 0,009 & 2,10 & 11,1 & 42,6 & 0,60 \\ \hline
Izboskan & 41 & +2,7 & 0,006 & 2,06 & 12,4 & 46,9 & 0,35 \\ \hline
Marhamat & 87 & +3,3 & 0,001 & 1,92 & 22,9 & 54,7 & 0,48 \\ \hline
Andijon & 39 & +3,3 & 0,012 & 1,86 & 15,5 & 62,9 & 0,10 \\ \hline
Jalaquduq & 57 & +5,6 & \textless{} 0,001 & 2,06 & 9,6 & 48,2 & 0,37 \\ \hline
Qoʻrgʻontepa & 24 & +3,2 & 0,012 & 2,14 & 13,3 & 40,0 & 0,21 \\ \hline
Paxtaobod & 32 & +10,3 & \textless{} 0,001 & 2,65 & 4,5 & 18,2 & 0,44 \\ \hline
Asaka & 31 & +6,4 & \textless{} 0,001 & 2,65 & 0,0 & 15,8 & 0,33 \\
\end{longtable}

}

\oakilova{41-ilova}

\oakilovanom{\textbf{Drenaj tarmogʻi ishonchliligi koʻrsatkichlari va ularni
hisoblash uchun maʼlumotlarning mavjudligi}}

{\fontsize{12}{14}\selectfont\setlength{\tabcolsep}{4pt}

\noindent\begin{longtable}[]{|
  >{\raggedright\arraybackslash}p{(\columnwidth - 8\tabcolsep - 5\arrayrulewidth) * \real{0.2658}}|>{\raggedright\arraybackslash}p{(\columnwidth - 8\tabcolsep - 5\arrayrulewidth) * \real{0.2290}}|>{\raggedright\arraybackslash}p{(\columnwidth - 8\tabcolsep - 5\arrayrulewidth) * \real{0.2447}}|>{\raggedright\arraybackslash}p{(\columnwidth - 8\tabcolsep - 5\arrayrulewidth) * \real{0.2605}}|}
\hline
\begin{minipage}[b]{\linewidth}\raggedright
\textbf{Ishonchlilik xossasi}
\end{minipage} & \begin{minipage}[b]{\linewidth}\raggedright
\textbf{Koʻrsatkich}
\end{minipage} & \begin{minipage}[b]{\linewidth}\raggedright
\textbf{Talab etiladigan maʼlumot}
\end{minipage} & \begin{minipage}[b]{\linewidth}\raggedright
\textbf{Meliorativ kadastrda mavjudligi}
\end{minipage} \\
\hline
\endhead
\hline
\endlastfoot
Buzilmaslik & Buzilmasdan ishlash ehtimoli, nosozliklar oqimi parametri
& Nosozliklar (loyqa bosish, qirgʻoq oʻpirilishi, inshoot buzilishi)
sanasi va turi & Mavjud emas \\ \hline
Uzoq muddatlilik & Xizmat muddati, chegaraviy holatgacha resurs &
Foydalanishga topshirilgan yil, chegaraviy holat mezoni & Qisman: faqat
foydalanishga topshirilgan yil \\ \hline
Taʼmirboplik & Taʼmirlararo muddat, taʼmir narxi & Barcha taʼmirlar
sanasi, hajmi va narxi & Qisman: faqat oxirgi taʼmir yili, 53 ta
kollektor uchun narx \\ \hline
Saqlanuvchanlik & --- & --- & Ochiq kollektorlar uchun qoʻllanilmaydi \\ \hline
Vazifaviy samaradorlik & Sizot suvi sathining meʼyoriy chuqurlikda
ushlanishi & Sath monitoringi & Mavjud (III bob) \\
\end{longtable}

}

\oakilova{42-ilova}

\oakilovanom{\textbf{Kollektorlarning foydalanishga topshirilgan davri boʻyicha
taqsimoti}}

{\fontsize{12}{14}\selectfont\setlength{\tabcolsep}{4pt}

\noindent\begin{longtable}[]{|
  >{\raggedright\arraybackslash}p{(\columnwidth - 6\tabcolsep - 4\arrayrulewidth) * \real{0.3908}}|>{\raggedright\arraybackslash}p{(\columnwidth - 6\tabcolsep - 4\arrayrulewidth) * \real{0.3524}}|>{\raggedright\arraybackslash}p{(\columnwidth - 6\tabcolsep - 4\arrayrulewidth) * \real{0.2568}}|}
\hline
\begin{minipage}[b]{\linewidth}\raggedright
\textbf{Davr}
\end{minipage} & \begin{minipage}[b]{\linewidth}\raggedright
\textbf{Kollektorlar soni}
\end{minipage} & \begin{minipage}[b]{\linewidth}\raggedright
\textbf{Uzunligi, km}
\end{minipage} \\
\hline
\endhead
\hline
\endlastfoot
1910-yillar & 1 & 12,5 \\ \hline
1940-yillar & 10 & 65,6 \\ \hline
1950-yillar & 31 & 90,8 \\ \hline
1960-yillar & 30 & 68,5 \\ \hline
1970-yillar & 15 & 22,5 \\ \hline
1980-yillar & 3 & 6,0 \\
\end{longtable}

}

\oakilova{43-ilova}

\oakilovanom{\textbf{Kollektorlar uzunligining taqsimoti: a --- foydalanishga
topshirilgan oʻn yillik boʻyicha; b --- oxirgi taʼmir yili boʻyicha; c
--- 1 km taʼmir narxining 2026-yil narxlaridagi taqsimoti (uzuq chiziq
--- mediana)}}

\oakimg{figures/ch4_reliability_costs.png}

\oakilova{44-ilova}

\oakilovanom{\textbf{Ichki zovurlarning tizimlar boʻyicha holati}}

{\fontsize{12}{14}\selectfont\setlength{\tabcolsep}{4pt}

\noindent\begin{longtable}[]{|
  >{\raggedright\arraybackslash}p{(\columnwidth - 8\tabcolsep - 5\arrayrulewidth) * \real{0.4211}}|>{\raggedright\arraybackslash}p{(\columnwidth - 8\tabcolsep - 5\arrayrulewidth) * \real{0.1671}}|>{\raggedright\arraybackslash}p{(\columnwidth - 8\tabcolsep - 5\arrayrulewidth) * \real{0.1792}}|>{\raggedright\arraybackslash}p{(\columnwidth - 8\tabcolsep - 5\arrayrulewidth) * \real{0.2326}}|}
\hline
\begin{minipage}[b]{\linewidth}\raggedright
\textbf{Tizim}
\end{minipage} & \begin{minipage}[b]{\linewidth}\raggedright
\textbf{Zovurlar soni}
\end{minipage} & \begin{minipage}[b]{\linewidth}\raggedright
\textbf{Uzunligi, km}
\end{minipage} & \begin{minipage}[b]{\linewidth}\raggedright
\textbf{Taʼmir qayd etilmagan uzunlik ulushi}
\end{minipage} \\
\hline
\endhead
\hline
\endlastfoot
Mashʼal massivi & 72 & 92,0 & 21 \% \\ \hline
T.Mirzayev massivi & 36 & 30,4 & 47 \% \\ \hline
U.Qoraboyev massivi & 29 & 21,9 & 72 \% \\ \hline
Istiqlol massivi & 57 & 70,2 & 74 \% \\ \hline
Paxtakor massivi & 33 & 41,9 & 74 \% \\ \hline
Z.M.Bobur massivi & 32 & 35,2 & 82 \% \\ \hline
Markaziy zovur tizimi & 54 & 52,8 & 100 \% \\ \hline
Y.Oxunboboyev massivi & 12 & 20,1 & 100 \% \\ \hline
Xoʻjaliklararo boʻrdoqichilik bazasi & 7 & 6,1 & 100 \% \\
\end{longtable}

}

\oakilova{45-ilova}

\oakilovanom{\textbf{Massivlar boʻyicha drenaj tarmogʻining ishonchlilik
koʻrsatkichlari}}

{\fontsize{9}{11}\selectfont\setlength{\tabcolsep}{2.5pt}

\noindent\begin{longtable}[]{|
  >{\raggedright\arraybackslash}p{(\columnwidth - 18\tabcolsep - 10\arrayrulewidth) * \real{0.1260}}|>{\raggedright\arraybackslash}p{(\columnwidth - 18\tabcolsep - 10\arrayrulewidth) * \real{0.1330}}|>{\raggedright\arraybackslash}p{(\columnwidth - 18\tabcolsep - 10\arrayrulewidth) * \real{0.1235}}|>{\raggedright\arraybackslash}p{(\columnwidth - 18\tabcolsep - 10\arrayrulewidth) * \real{0.1235}}|>{\raggedright\arraybackslash}p{(\columnwidth - 18\tabcolsep - 10\arrayrulewidth) * \real{0.1520}}|>{\raggedright\arraybackslash}p{(\columnwidth - 18\tabcolsep - 10\arrayrulewidth) * \real{0.0855}}|>{\raggedright\arraybackslash}p{(\columnwidth - 18\tabcolsep - 10\arrayrulewidth) * \real{0.0855}}|>{\raggedright\arraybackslash}p{(\columnwidth - 18\tabcolsep - 10\arrayrulewidth) * \real{0.0855}}|>{\raggedright\arraybackslash}p{(\columnwidth - 18\tabcolsep - 10\arrayrulewidth) * \real{0.0855}}|}
\hline
\begin{minipage}[b]{\linewidth}\raggedright
\textbf{Massiv}
\end{minipage} & \begin{minipage}[b]{\linewidth}\raggedright
\textbf{Sugʻoriladigan maydon, ga}
\end{minipage} & \begin{minipage}[b]{\linewidth}\raggedright
\textbf{Kollektorlar, km}
\end{minipage} & \begin{minipage}[b]{\linewidth}\raggedright
\textbf{Kollektorlar, m/ga}
\end{minipage} & \begin{minipage}[b]{\linewidth}\raggedright
\textbf{Kollektorlarning oʻrtacha yoshi, yil*}
\end{minipage} & \begin{minipage}[b]{\linewidth}\raggedright
\textbf{Oxirgi taʼmirdan beri, yil*}
\end{minipage} & \begin{minipage}[b]{\linewidth}\raggedright
\textbf{Ichki zovurlar, km}
\end{minipage} & \begin{minipage}[b]{\linewidth}\raggedright
\textbf{Ichki zovurlar, m/ga}
\end{minipage} & \begin{minipage}[b]{\linewidth}\raggedright
\textbf{Taʼmirsiz ichki zovurlar, \%}
\end{minipage} \\
\hline
\endhead
\hline
\endlastfoot
Mashʼal & 3 790 & 54,8 & 14,5 & 63 & 3,9 & 92,0 & 24,3 & 21 \\ \hline
Istiqlol & 2 146 & 34,1 & 15,9 & 75 & 5,8 & 70,2 & 32,7 & 74 \\ \hline
U.Qoraboyev & 2 420 & 32,2 & 13,3 & 64 & 4,2 & 21,9 & 9,1 & 72 \\ \hline
Z.M.Bobur & 3 610 & 39,3 & 10,9 & 67 & 4,9 & 35,2 & 9,8 & 82 \\ \hline
Paxtakor & 1 601 & 12,9 & 8,0 & 72 & 5,0 & 41,9 & 26,2 & 74 \\ \hline
T.Mirzayev & 3 371 & 24,0 & 7,1 & 62 & 4,6 & 30,4 & 9,0 & 47 \\ \hline
M.Toʻjiboyev & 2 532 & 17,3 & 6,8 & 66 & 5,8 & --- & --- & --- \\ \hline
Y.Oxunboboyev & 1 191 & 0,0 & 0,0 & --- & --- & 20,1 & 16,9 & 100 \\
\end{longtable}

}

Izoh: * --- kadastrda nomi aniqlangan kollektor boʻlaklari boʻyicha
uzunlik bilan vaznlangan oʻrtacha (massivdagi kollektorlar uzunligining
35--90 foizi). Markaziy zovur tizimi (52,8 km) va boʻrdoqichilik bazasi
zovurlari (6,1 km) kadastrda massivlarga bogʻlanmagan.

\oakilova{46-ilova}

\oakilovanom{\textbf{Massivlar boʻyicha: a --- kollektorlar va ichki zovurlarning
solishtirma uzunligi; b --- taʼmir qayd etilmagan ichki zovurlar ulushi
(toʻq sariq, xoch bilan --- 70 foizdan ortiq)}}

\oakimg{figures/ch4_massif_reliability.png}

\oakilova{47-ilova}

\oakilovanom{\textbf{Marhamat tumani boʻyicha oylik suv balansi tarkibiy qismlari,
mm}}

{\fontsize{10}{12}\selectfont\setlength{\tabcolsep}{2.5pt}

\noindent\begin{longtable}[]{|
  >{\raggedright\arraybackslash}p{(\columnwidth - 14\tabcolsep - 8\arrayrulewidth) * \real{0.1164}}|>{\raggedright\arraybackslash}p{(\columnwidth - 14\tabcolsep - 8\arrayrulewidth) * \real{0.2197}}|>{\raggedright\arraybackslash}p{(\columnwidth - 14\tabcolsep - 8\arrayrulewidth) * \real{0.1941}}|>{\raggedright\arraybackslash}p{(\columnwidth - 14\tabcolsep - 8\arrayrulewidth) * \real{0.1529}}|>{\raggedright\arraybackslash}p{(\columnwidth - 14\tabcolsep - 8\arrayrulewidth) * \real{0.0806}}|>{\raggedright\arraybackslash}p{(\columnwidth - 14\tabcolsep - 8\arrayrulewidth) * \real{0.1105}}|>{\raggedright\arraybackslash}p{(\columnwidth - 14\tabcolsep - 8\arrayrulewidth) * \real{0.1256}}|}
\hline
\begin{minipage}[b]{\linewidth}\raggedright
\textbf{Oy}
\end{minipage} & \begin{minipage}[b]{\linewidth}\raggedright
\textbf{Etalon bugʻlanuvchanlik ET₀}
\end{minipage} & \begin{minipage}[b]{\linewidth}\raggedright
\textbf{Oʻrtacha ekin koeffitsiyenti}
\end{minipage} & \begin{minipage}[b]{\linewidth}\raggedright
\textbf{Ekinlar bugʻlanishi ETc}
\end{minipage} & \begin{minipage}[b]{\linewidth}\raggedright
\textbf{Yogʻin}
\end{minipage} & \begin{minipage}[b]{\linewidth}\raggedright
\textbf{Samarali yogʻin}
\end{minipage} & \begin{minipage}[b]{\linewidth}\raggedright
\textbf{Sof sugʻorish talabi}
\end{minipage} \\
\hline
\endhead
\hline
\endlastfoot
Mart & 62 & 0,42 & 26 & 45 & 42 & 0 \\ \hline
Aprel & 110 & 0,74 & 82 & 46 & 43 & 39 \\ \hline
May & 159 & 0,83 & 132 & 38 & 36 & 96 \\ \hline
Iyun & 193 & 0,72 & 139 & 17 & 17 & 123 \\ \hline
Iyul & 194 & 0,62 & 120 & 12 & 12 & 108 \\ \hline
Avgust & 170 & 0,62 & 105 & 6 & 6 & 99 \\ \hline
Sentabr & 123 & 0,47 & 58 & 7 & 7 & 51 \\ \hline
Oktabr & 70 & 0,54 & 38 & 32 & 30 & 7 \\ \hline
Yil boʻyicha & 1 185 & --- & 719 & 323 & 306 & 524 \\
\end{longtable}

}

\oakilova{48-ilova}

\oakilovanom{\textbf{a --- oylik ekinlar bugʻlanishi, samarali yogʻin va sof
sugʻorish talabi; b --- kollektorlarning kadastr boʻyicha solishtirma
sarfi va iyuldagi talab etiladigan drenaj moduli (takroriy ekinlarsiz
variant)}}

\oakimg{figures/ch4_hydraulic.png}

\oakilova{49-ilova}

\oakilovanom{\textbf{Namunaviy kollektor kesimining oʻzan holatiga qarab oʻtkazish
qobiliyati}}

{\fontsize{12}{14}\selectfont\setlength{\tabcolsep}{4pt}

\noindent\begin{longtable}[]{|
  >{\raggedright\arraybackslash}p{(\columnwidth - 12\tabcolsep - 7\arrayrulewidth) * \real{0.1798}}|>{\raggedright\arraybackslash}p{(\columnwidth - 12\tabcolsep - 7\arrayrulewidth) * \real{0.2527}}|>{\raggedright\arraybackslash}p{(\columnwidth - 12\tabcolsep - 7\arrayrulewidth) * \real{0.1394}}|>{\raggedright\arraybackslash}p{(\columnwidth - 12\tabcolsep - 7\arrayrulewidth) * \real{0.1421}}|>{\raggedright\arraybackslash}p{(\columnwidth - 12\tabcolsep - 7\arrayrulewidth) * \real{0.1253}}|>{\raggedright\arraybackslash}p{(\columnwidth - 12\tabcolsep - 7\arrayrulewidth) * \real{0.1608}}|}
\hline
\begin{minipage}[b]{\linewidth}\raggedright
\textbf{Oʻzan holati}
\end{minipage} & \begin{minipage}[b]{\linewidth}\raggedright
\textbf{gʻadir-budurlik n}
\end{minipage} & \begin{minipage}[b]{\linewidth}\raggedright
\textbf{Loyqa qatlami, m}
\end{minipage} & \begin{minipage}[b]{\linewidth}\raggedright
\textbf{Q (i = 0,0005), m³/s}
\end{minipage} & \begin{minipage}[b]{\linewidth}\raggedright
\textbf{Q (i = 0,001), m³/s}
\end{minipage} & \begin{minipage}[b]{\linewidth}\raggedright
\textbf{Toza oʻzanga nisbatan, \%}
\end{minipage} \\
\hline
\endhead
\hline
\endlastfoot
Toza oʻzan & 0,025 & 0 & 2,21 & 3,13 & 100 \\ \hline
Oʻsimlik bilan qoplangan & 0,035 & 0 & 1,58 & 2,23 & 71 \\ \hline
Oʻsimlik va loyqa & 0,035 & 0,3 & 1,28 & 1,81 & 58 \\ \hline
Kuchli oʻsimlik va loyqa & 0,040 & 0,5 & 0,87 & 1,23 & 39 \\
\end{longtable}

}

\oakilova{50-ilova}

\oakilovanom{\textbf{Kollektorlarni taʼmirlash xarajatlari yillar boʻyicha
(kadastrdagi narxlarda va 2026-yil oktabr narxlarida)}}

{\fontsize{9}{11}\selectfont\setlength{\tabcolsep}{2.5pt}

\noindent\begin{longtable}[]{|
  >{\raggedright\arraybackslash}p{(\columnwidth - 16\tabcolsep - 9\arrayrulewidth) * \real{0.0495}}|>{\raggedright\arraybackslash}p{(\columnwidth - 16\tabcolsep - 9\arrayrulewidth) * \real{0.1447}}|>{\raggedright\arraybackslash}p{(\columnwidth - 16\tabcolsep - 9\arrayrulewidth) * \real{0.1457}}|>{\raggedright\arraybackslash}p{(\columnwidth - 16\tabcolsep - 9\arrayrulewidth) * \real{0.1105}}|>{\raggedright\arraybackslash}p{(\columnwidth - 16\tabcolsep - 9\arrayrulewidth) * \real{0.1107}}|>{\raggedright\arraybackslash}p{(\columnwidth - 16\tabcolsep - 9\arrayrulewidth) * \real{0.1700}}|>{\raggedright\arraybackslash}p{(\columnwidth - 16\tabcolsep - 9\arrayrulewidth) * \real{0.1344}}|>{\raggedright\arraybackslash}p{(\columnwidth - 16\tabcolsep - 9\arrayrulewidth) * \real{0.1344}}|}
\hline
\begin{minipage}[b]{\linewidth}\raggedright
\textbf{Yil}
\end{minipage} & \begin{minipage}[b]{\linewidth}\raggedright
\textbf{Kollektorlar soni}
\end{minipage} & \begin{minipage}[b]{\linewidth}\raggedright
\textbf{Taʼmirlangan uzunlik, km}
\end{minipage} & \begin{minipage}[b]{\linewidth}\raggedright
\textbf{Xarajat, mln soʻm (kadastr)}
\end{minipage} & \begin{minipage}[b]{\linewidth}\raggedright
\textbf{Oʻrtacha narx, mln soʻm/km (kadastr)}
\end{minipage} & \begin{minipage}[b]{\linewidth}\raggedright
\textbf{Qayta hisoblash koeffitsiyenti}
\end{minipage} & \begin{minipage}[b]{\linewidth}\raggedright
\textbf{Xarajat, mln soʻm (2026-yil narxlarida)}
\end{minipage} & \begin{minipage}[b]{\linewidth}\raggedright
\textbf{Oʻrtacha narx, mln soʻm/km (2026-yil narxlarida)}
\end{minipage} \\
\hline
\endhead
\hline
\endlastfoot
2017 & 28 & 69,0 & 9 809 & 142,3 & 2,60 & 25 454 & 369,2 \\ \hline
2018 & 9 & 15,6 & 2 742 & 175,8 & 2,27 & 6 224 & 399,0 \\ \hline
2019 & 13 & 65,0 & 5 127 & 78,9 & 1,98 & 10 141 & 156,1 \\ \hline
2020 & 3 & 13,8 & 3 636 & 263,5 & 1,75 & 6 356 & 460,6 \\ \hline
Jami & 53 & 163,3 & 21 314 & 130,5 & --- & 48 176 & 295,0 \\
\end{longtable}

}

Izoh: qayta hisoblash isteʼmol narxlari indeksi boʻyicha (4.4-§); taʼmir
yili xarajati shu yilning oʻrtacha narxlarida deb olingan.

\oakilova{51-ilova}

\oakilovanom{\textbf{Kollektorlar taʼmir xarajatlarining massivlar boʻyicha taqsimoti
(2017--2020-yillar) va meliorativ xavf koʻrsatkichlari}}

{\fontsize{9}{11}\selectfont\setlength{\tabcolsep}{2.5pt}

\noindent\begin{longtable}[]{|
  >{\raggedright\arraybackslash}p{(\columnwidth - 14\tabcolsep - 8\arrayrulewidth) * \real{0.1620}}|>{\raggedright\arraybackslash}p{(\columnwidth - 14\tabcolsep - 8\arrayrulewidth) * \real{0.1704}}|>{\raggedright\arraybackslash}p{(\columnwidth - 14\tabcolsep - 8\arrayrulewidth) * \real{0.1589}}|>{\raggedright\arraybackslash}p{(\columnwidth - 14\tabcolsep - 8\arrayrulewidth) * \real{0.1360}}|>{\raggedright\arraybackslash}p{(\columnwidth - 14\tabcolsep - 8\arrayrulewidth) * \real{0.1360}}|>{\raggedright\arraybackslash}p{(\columnwidth - 14\tabcolsep - 8\arrayrulewidth) * \real{0.1008}}|>{\raggedright\arraybackslash}p{(\columnwidth - 14\tabcolsep - 8\arrayrulewidth) * \real{0.1360}}|}
\hline
\begin{minipage}[b]{\linewidth}\raggedright
\textbf{Massiv}
\end{minipage} & \begin{minipage}[b]{\linewidth}\raggedright
\textbf{Sugʻoriladigan maydon, ga}
\end{minipage} & \begin{minipage}[b]{\linewidth}\raggedright
\textbf{Taʼmirlangan kollektorlar*}
\end{minipage} & \begin{minipage}[b]{\linewidth}\raggedright
\textbf{Xarajat, mln soʻm (2026-yil narxlarida)}
\end{minipage} & \begin{minipage}[b]{\linewidth}\raggedright
\textbf{Xarajat, ming soʻm/ga (2026-yil narxlarida)}
\end{minipage} & \begin{minipage}[b]{\linewidth}\raggedright
\textbf{Yozgi sayoz sath ehtimoli 0,3 dan katta maydon, \% (3.3-§)}
\end{minipage} & \begin{minipage}[b]{\linewidth}\raggedright
\textbf{Shoʻrlanish xavfli quduqlar (3.5-§)}
\end{minipage} \\
\hline
\endhead
\hline
\endlastfoot
Istiqlol & 2 146 & 6 & 7 828 & 3 648 & 55 & 12 \\ \hline
U.Qoraboyev & 2 420 & 9 & 6 570 & 2 715 & 19 & 1 \\ \hline
Mashʼal & 3 790 & 13 & 9 187 & 2 424 & 71 & 17 \\ \hline
Z.M.Bobur & 3 610 & 6 & 6 472 & 1 793 & 58 & 17 \\ \hline
T.Mirzayev & 3 371 & 5 & 5 525 & 1 639 & 28 & 1 \\ \hline
Paxtakor & 1 601 & 5 & 2 338 & 1 461 & 21 & 4 \\ \hline
M.Toʻjiboyev & 2 532 & 3 & 662 & 261 & 17 & 4 \\ \hline
Y.Oxunboboyev & 1 191 & 0 & 0 & 0 & 0 & --- \\ \hline
Jami & 20 661 & 36 & 38 582 & 1 867 & --- & 56 \\
\end{longtable}

}

Izoh: * --- bir kollektor bir necha massivdan oʻtishi mumkin, shuning
uchun massivlar boʻyicha sonlar yigʻindisi jami sondan katta.

\oakilova{52-ilova}

\oakilovanom{\textbf{Amaliy tavsiyalar va ularning asosi}}

{\fontsize{12}{14}\selectfont\setlength{\tabcolsep}{4pt}

\noindent\begin{longtable}[]{|
  >{\raggedright\arraybackslash}p{(\columnwidth - 6\tabcolsep - 4\arrayrulewidth) * \real{0.0867}}|>{\raggedright\arraybackslash}p{(\columnwidth - 6\tabcolsep - 4\arrayrulewidth) * \real{0.4644}}|>{\raggedright\arraybackslash}p{(\columnwidth - 6\tabcolsep - 4\arrayrulewidth) * \real{0.4489}}|}
\hline
\begin{minipage}[b]{\linewidth}\raggedright
\textbf{№}
\end{minipage} & \begin{minipage}[b]{\linewidth}\raggedright
\textbf{Tavsiya}
\end{minipage} & \begin{minipage}[b]{\linewidth}\raggedright
\textbf{Asos}
\end{minipage} \\
\hline
\endhead
\hline
\endlastfoot
& \textbf{Maʼlumotlar sifati} & \\ \hline
T-1 & Monitoring jadvallarini qabul qilishda dissertatsiyada ishlab
chiqilgan tekshiruv tartibini majburiy qoʻllash: yil ichidagi
oʻzgarmaslik, birlik xatosi, tarmoq tarkibi va kesilgan qiymatlar &
Oʻlchovlarning 12,1 foizi yaroqsiz, yaroqlilarining 8,2 foizi kesilgan
(3.1-§); viloyatning 14 ta tumanida nuqsonlar 0--29 foiz (3.6-§) \\ \hline
T-2 & Quduq qurigan yoki oʻlchov oʻtkazilmagan hollarda qiymat oʻrniga
maxsus belgi qoʻyish (masalan, «quruq, quduq chuqurligi 3,0 m») & 270 ta
quduq-yilda yil boʻyi bitta qiymat; oʻzgarmas qatorlarning yarmidan
koʻpi quduq chuqurligiga teng (3.1-§) \\ \hline
T-3 & Kuzatuv tarmogʻidagi har bir oʻzgarishni (yangi quduq,
chuqurlashtirish, koʻchirish) sanasi bilan qayd etish & 2018-yildagi
oʻzgarish tekshirilmagan maʼlumotda oʻzgarish tezligini 1,4 baravar
oshirgan (3.2-§) \\ \hline
T-4 & Tuman yoki viloyat boʻyicha oʻrtacha sathni faqat doimiy
kuzatilgan quduqlar boʻyicha hisoblash va hisobotlarda yozgi eng sayoz
sathni alohida koʻrsatish & Yillik oʻrtacha massivlardagi yozgi xavfli
holatni yashiradi (3.2-§) \\ \hline
& \textbf{Kuzatuv tarmogʻi} & \\ \hline
T-5 & Sath chuqur hududlarda kuzatuv quduqlarini 5--6 m gacha
chuqurlashtirish yoki qayta qurish & Ishonchlilik sath chuqurligi bilan
bogʻliq: massivlarda 50--98 foiz, tumanlarda ρ = +0,85 (3.1, 3.6-§) \\ \hline
T-6 & Tarmoqni bir tekis zichlashtirish oʻrniga yangi quduqlarni
maqsadli joylashtirish: ustuvor massivlarga, taʼmirlanadigan
kollektorlar yaqiniga va standart xato eng katta chekka hududlarga & 94
ta quduqdan keyin xarita aniqligi oshmaydi (3.3-§) \\ \hline
T-7 & Meliorativ holatni sath xaritasi bilan emas, sathning kritik
chuqurlikdan sayoz boʻlish ehtimoli xaritasi bilan baholash & Kriging
xaritasi sayoz maydonni oʻn baravardan ortiq kam koʻrsatadi (3.3-§) \\ \hline
T-8 & Minerallashuv oʻlchovlarini tiklash: yiliga bir marta, birinchi
navbatda shoʻrlanish xavfi ostidagi 56 ta quduqda; dala elektr
oʻtkazuvchanlik asboblaridan foydalanish & Ishonchli qator faqat
2004--2010-yillar uchun; mavsumiy farq ahamiyatsiz (3.5-§) \\ \hline
T-9 & Asosiy kollektorlar oxirida suv sarfini muntazam oʻlchash &
Kadastrdagi sarflar 0,03--50 l/(s·ga) oraligʻida, ishonchliligi nomaʼlum
(4.2-§) \\ \hline
& \textbf{Tizimni boshqarish} & \\ \hline
T-10 & Kollektor va zovurlar boʻyicha taʼmir va nosozliklar daftarini
yuritish: sana, piketlar, ish turi, hajmi, xarajati, oʻzan holati
(oʻsimlik, loyqa qatlami) & Kadastrda faqat oxirgi taʼmir yili;
ishonchlilik koʻrsatkichlarini hisoblab boʻlmaydi (4.1-§) \\ \hline
T-11 & Taʼmir navbatini kollektor holati va massivning meliorativ xavfi
boʻyicha aniqlash & Taʼmir yili yoshga, xarajatlar xavfga bogʻliq emas
(4.1, 4.4-§) \\ \hline
T-12 & Qoʻshni kichik kollektorlarni bitta shartnoma doirasida
taʼmirlash; yengil tozalash va kapital taʼmirni oʻzan holatiga qarab
tanlash & Uchastka uzunligi ikki baravar oshganda narx 25 foiz kamayadi;
oʻsimlik va loyqa oʻtkazish qobiliyatini 29--61 foiz kamaytiradi (4.2,
4.4-§) \\ \hline
T-13 & Taʼmir samarasini oʻlchash uchun taʼmir rejasini monitoring bilan
muvofiqlashtirish: har bir taʼmirlanadigan kollektor yonida va undan
uzoqda juft quduqlar & Taʼsir guruhida atigi 36 ta quduq; aniqlanadigan
eng kichik taʼsir 0,29 m (4.3-§) \\ \hline
T-14 & Xoʻjalik ichki zovurlarini inventarizatsiya qilish va taʼmir
rejasiga kiritish & 64,7 foizida taʼmir qayd etilmagan (4.1-§) \\ \hline
T-15 & Drenaj yuklamasini kamaytirish uchun sugʻorish samaradorligini
oshirish tadbirlarini ustuvor massivlarda qoʻllash & Samaradorlik
oshirilganda talab etiladigan drenaj moduli 2,5 baravar kamayadi
(4.2-§) \\
\end{longtable}

}

\oakilova{53-ilova}

\oakilovanom{\textbf{Taʼmir taʼsirini aniqlash ehtimoli (Monte Karlo sinov hisobi,
har bir variant 1 000 takror, ikki tomonlama α = 0,05)}}

{\fontsize{12}{14}\selectfont\setlength{\tabcolsep}{4pt}

\noindent\begin{longtable}[]{|
  >{\raggedright\arraybackslash}p{(\columnwidth - 12\tabcolsep - 7\arrayrulewidth) * \real{0.3219}}|>{\raggedright\arraybackslash}p{(\columnwidth - 12\tabcolsep - 7\arrayrulewidth) * \real{0.2106}}|>{\raggedright\arraybackslash}p{(\columnwidth - 12\tabcolsep - 7\arrayrulewidth) * \real{0.1169}}|>{\raggedright\arraybackslash}p{(\columnwidth - 12\tabcolsep - 7\arrayrulewidth) * \real{0.1169}}|>{\raggedright\arraybackslash}p{(\columnwidth - 12\tabcolsep - 7\arrayrulewidth) * \real{0.1169}}|>{\raggedright\arraybackslash}p{(\columnwidth - 12\tabcolsep - 7\arrayrulewidth) * \real{0.1169}}|}
\hline
\begin{minipage}[b]{\linewidth}\raggedright
\textbf{Taʼsir / nazorat quduqlari}
\end{minipage} & \begin{minipage}[b]{\linewidth}\raggedright
\textbf{Standart xato, m}
\end{minipage} & \begin{minipage}[b]{\linewidth}\raggedright
\textbf{Δ = 0,10 m}
\end{minipage} & \begin{minipage}[b]{\linewidth}\raggedright
\textbf{Δ = 0,15 m}
\end{minipage} & \begin{minipage}[b]{\linewidth}\raggedright
\textbf{Δ = 0,20 m}
\end{minipage} & \begin{minipage}[b]{\linewidth}\raggedright
\textbf{Δ = 0,30 m}
\end{minipage} \\
\hline
\endhead
\hline
\endlastfoot
36 / 49 (hozirgi) & 0,104 & 0,16 & 0,31 & 0,47 & 0,82 \\ \hline
60 / 49 & 0,085 & 0,22 & 0,42 & 0,66 & 0,94 \\ \hline
78 / 49 & 0,076 & 0,25 & 0,52 & 0,75 & 0,97 \\ \hline
78 / 78 & 0,073 & 0,26 & 0,52 & 0,78 & 0,98 \\ \hline
100 / 49 & 0,068 & 0,33 & 0,59 & 0,82 & 0,99 \\ \hline
100 / 100 & 0,065 & 0,33 & 0,67 & 0,87 & 0,99 \\ \hline
120 / 60 & 0,062 & 0,35 & 0,68 & 0,89 & 0,99 \\ \hline
150 / 100 & 0,055 & 0,43 & 0,78 & 0,96 & 1,00 \\
\end{longtable}

}

Izoh: jadvalda taʼsirni aniqlash ehtimoli (statistik quvvat) berilgan;
Monte Karlo xatosi ±0,03 dan oshmaydi. Δ = 0 da (36 / 49) birinchi tur
xato 0,050 (95 \% oraliq 0,041--0,060). Simulyatsiya haqiqiy quduq
qatorlarini qayta tanlashga asoslangan, shuning uchun taʼsir mahalliy va
quduqlar mustaqil degan faraz saqlanadi.

\oakilova{54-ilova}

\oakilovanom{\textbf{Asosiy kuzatishlarning boshqa mumkin boʻlgan sabablari, ularni
ajratadigan belgilar va dalillar holati}}

{\fontsize{9}{11}\selectfont\setlength{\tabcolsep}{2.5pt}

\noindent\begin{longtable}[]{|
  >{\raggedright\arraybackslash}p{(\columnwidth - 10\tabcolsep - 6\arrayrulewidth) * \real{0.2490}}|>{\raggedright\arraybackslash}p{(\columnwidth - 10\tabcolsep - 6\arrayrulewidth) * \real{0.1790}}|>{\raggedright\arraybackslash}p{(\columnwidth - 10\tabcolsep - 6\arrayrulewidth) * \real{0.1674}}|>{\raggedright\arraybackslash}p{(\columnwidth - 10\tabcolsep - 6\arrayrulewidth) * \real{0.1907}}|>{\raggedright\arraybackslash}p{(\columnwidth - 10\tabcolsep - 6\arrayrulewidth) * \real{0.2140}}|}
\hline
\begin{minipage}[b]{\linewidth}\raggedright
\textbf{Kuzatish}
\end{minipage} & \begin{minipage}[b]{\linewidth}\raggedright
\textbf{Mumkin boʻlgan sabab}
\end{minipage} & \begin{minipage}[b]{\linewidth}\raggedright
\textbf{Ajratadigan belgi}
\end{minipage} & \begin{minipage}[b]{\linewidth}\raggedright
\textbf{Mavjud dalil (holat)}
\end{minipage} & \begin{minipage}[b]{\linewidth}\raggedright
\textbf{Keyingi tekshiruv}
\end{minipage} \\
\hline
\endhead
\hline
\endlastfoot
Sathning 2016--2018-yillardagi pogʻonali chuqurlashishi (0,30 m) & A.
Hududiy suv taʼminoti va sugʻorish rejimining oʻzgarishi & Pogʻona suv
olishning kamayishi bilan bir vaqtda va koʻp tumanlarda & 14 ta
tumanning 12 tasida chuqurlashish (mos) & 2009--2024-yillar suv olish
hajmlari bilan yillik regressiya \\ \hline
& B. Kollektorlar taʼmiri & Kollektor yaqinida kuchliroq chuqurlashish &
Uzoqdagi quduqlarda tezroq; taʼmirdan keyin chuqurlashish aniqlanmadi
(qarshi, lekin quduqlar kam) & Bosqichma-bosqich taʼmir va juft
quduqlar \\ \hline
& C. 2018-yilda oʻlchash tartibining oʻzgarishi & Pogʻona barcha
quduqlarda bir yilda, mavsumiy yoʻl shakli oʻzgaradi & Doimiy kuzatilgan
quduqlarda ham pogʻona bor (A va C ni ajratmaydi) & Kuzatuvchi
jurnallari, nazorat oʻlchovlari \\ \hline
& D. Yogʻin va bugʻlanish anomaliyasi & Pogʻona iqlim anomaliyasi bilan
bir vaqtda & Tekshirilmagan: faqat koʻp yillik normalar mavjud &
2009--2024-yillar meteorologik qatorlari \\ \hline
Kollektorga yaqin joyda sath sayozroq & A. Teskari bogʻliqlik:
kollektorlar zax joylarga qurilgan & Bogʻlanish qurilishdan oldin ham
boʻlgan; massiv ichida saqlanadi & Massiv ichida saqlanadi; relyef
kuchsiz (mos) & Qurilishgacha boʻlgan xaritalar va arxiv oʻlchovlari \\ \hline
& B. Loyqa bosgan kollektor sathni koʻtaradi & Taʼmirdan keyin kollektor
yaqinida sath chuqurlashadi & Chuqurlashish aniqlanmadi (zaif qarshi
dalil) & Taʼmir oldidan va keyin oʻzan holati va sarf oʻlchovi \\ \hline
Z.M.Bobur massivida sathning koʻtarilishi & A. Drenaj quvvatining
yetishmasligi & Koʻtarilish sugʻorish mavsumida kuchliroq & Qishda −0,16
m (−0,49 \ldots{} +0,17), yozda −0,07 m (−0,29 \ldots{} +0,15): noaniq &
Massiv kollektorlarining oxiridagi sarf, oʻzan koʻrigi \\ \hline
& B. Sugʻoriladigan maydon yoki ekin tarkibining oʻzgarishi &
Koʻtarilish yer foydalanish oʻzgarishi bilan bir vaqtda & Tekshirilmagan
& Masofadan zondlash (oʻsimlik indekslari), kontur jadvallari \\ \hline
& C. Kuzatuv tarmogʻi tarkibining oʻzgarishi & Bir xil quduqlarda
koʻtarilish yoʻq & Doimiy kuzatilgan quduqlarda ham koʻtarilish (p =
0,024): qisman istisno qilingan & --- \\ \hline
Yil boʻyi oʻzgarmas qatorlar & A. Quruq quduqda shartli qiymat yozilishi
& Oʻzgarmas qator quduq tubigacha zaxira kichik boʻlganda & Zaxira 0,5 m
dan kam: 44 \%, 1,5--2,0 m: 2 \%; OR = 0,31 / m (mos) & Quduqlarni
chuqurlashtirish va natijani kuzatish \\ \hline
& B. Kuzatuvchining hisobot amaliyoti & Oʻzgarmas qatorlar zaxiradan
qatʼi nazar toʻplanadi & Zaxira 2 m dan katta quduqlarda ham 9 \%
(qisman) & Kuzatuvchilar boʻyicha tahlil \\
\end{longtable}

}

Izoh: «mos» --- dalil izohga zid emas, lekin uni isbotlamaydi; bir
nechta sabab bir vaqtda taʼsir qilishi mumkin. Tekshiruvlar: 57-skript
(results/hypothesis\_checks.json).

\oakilova{55-ilova}

\oakilovanom{\textbf{Dissertatsiya natijalari va tumanga topshirilgan tavsiyanoma
{[}33{]} muvofiqligi}}

{\fontsize{12}{14}\selectfont\setlength{\tabcolsep}{4pt}

\noindent\begin{longtable}[]{|
  >{\raggedright\arraybackslash}p{(\columnwidth - 8\tabcolsep - 5\arrayrulewidth) * \real{0.2573}}|>{\raggedright\arraybackslash}p{(\columnwidth - 8\tabcolsep - 5\arrayrulewidth) * \real{0.2581}}|>{\raggedright\arraybackslash}p{(\columnwidth - 8\tabcolsep - 5\arrayrulewidth) * \real{0.2581}}|>{\raggedright\arraybackslash}p{(\columnwidth - 8\tabcolsep - 5\arrayrulewidth) * \real{0.2265}}|}
\hline
\begin{minipage}[b]{\linewidth}\raggedright
\textbf{Masala}
\end{minipage} & \begin{minipage}[b]{\linewidth}\raggedright
\textbf{Dissertatsiya}
\end{minipage} & \begin{minipage}[b]{\linewidth}\raggedright
\textbf{Tavsiyanoma}
\end{minipage} & \begin{minipage}[b]{\linewidth}\raggedright
\textbf{Izoh}
\end{minipage} \\
\hline
\endhead
\hline
\endlastfoot
Kuzatuv quduqlari & 188 ta (kuzatuv jadvallari va pasport) & 188 ta
(hisobotlar), GAT bazasida 213 ta & Mos keladi \\ \hline
Sugʻoriladigan maydon & 20 662 ga (xarita boʻyicha) & 17 965 ga
(meliorativ ekspeditsiya hisoboti) & Maydon chegarasi turlicha
olingan \\ \hline
Kollektor-drenaj tarmogʻi & Kadastr: 72 ta xoʻjaliklararo kollektor,
202,5 km; xaritadan 214,6 km raqamlandi & GAT bazasi: 119 ta zovur,
240,3 km & Tarmoqning turli qismi hisobga olingan \\ \hline
Yillik yogʻin & 323 mm (bitta manba, izchil davr) & 196 mm (Andijon
meteostansiyasi, 2015--2023) & Hisoblash davri va manba turlicha \\ \hline
Xavfli massivlar & Z.M.Bobur, Istiqlol (I guruh), Mashʼal (II guruh) &
Mashʼal, Paxtakor, Z.M.Bobur (1-navbat), Istiqlol (2-navbat) & Mashʼal,
Z.M.Bobur va Istiqlol ikkala hujjatda ham xavfli \\ \hline
Paxtakor & III guruh: sath chuqur, shoʻrlanish xavfli quduqlar 4 ta &
1-navbat: 2024-yilda minerallashgan sizot suvli yerlar 49,5 \% & Mezon
turlicha: quduqlardagi sath va hisobotdagi minerallashuv maydoni \\ \hline
Mashʼal & Tarmoq yaxshi holatda; sugʻorish samaradorligi (T-15),
minerallashuv kuzatuvi (T-8) & Zovurlarni tozalash, suvni meʼyorlash,
har oy oʻlchash & Zid emas: tozalash mavjud tarmoqni saqlash \\ \hline
Monitoring & Minerallashuvni oʻlchashni tiklash (T-8), quduqlarni
chuqurlashtirish (T-5), maʼlumotlarni tekshirish (T-1--T-4) &
Minerallashuvni yiliga kamida 2 marta oʻlchash, Diver datchiklari,
hisobotlarda avtomatik tekshiruv & Mos keladi \\ \hline
Taʼmir samarasini oʻlchash & Taʼmirni kuzatuv bilan bogʻlash (T-13),
taʼmir daftari (T-10) & Istiqlol, ZKD-6-2: №76, №67, №32a quduqlarda
2027-yil sathini 2024--2026-yillar bilan solishtirish & Mos keladi \\ \hline
Sugʻorish & Ustuvor massivlarda sugʻorish samaradorligini oshirish
(T-15) & Sath chuqurligiga qarab sugʻorish meʼyori; qiyaliklarda
tomchilatib sugʻorish & Tavsiyanoma bu tavsiyani meʼyorlar bilan
aniqlashtiradi \\
\end{longtable}

}

\oakilova{56-ilova}

\oakilovanom{\textbf{2025-yil bahori: kuzatuv maʼlumotlarini 2009--2024-yillarning
shu oylari (mart--aprel) bilan solishtirish}}

{\fontsize{12}{14}\selectfont\setlength{\tabcolsep}{4pt}

\noindent\begin{longtable}[]{|
  >{\raggedright\arraybackslash}p{(\columnwidth - 12\tabcolsep - 7\arrayrulewidth) * \real{0.0802}}|>{\raggedright\arraybackslash}p{(\columnwidth - 12\tabcolsep - 7\arrayrulewidth) * \real{0.1433}}|>{\raggedright\arraybackslash}p{(\columnwidth - 12\tabcolsep - 7\arrayrulewidth) * \real{0.1957}}|>{\raggedright\arraybackslash}p{(\columnwidth - 12\tabcolsep - 7\arrayrulewidth) * \real{0.2761}}|>{\raggedright\arraybackslash}p{(\columnwidth - 12\tabcolsep - 7\arrayrulewidth) * \real{0.1578}}|>{\raggedright\arraybackslash}p{(\columnwidth - 12\tabcolsep - 7\arrayrulewidth) * \real{0.1469}}|}
\hline
\begin{minipage}[b]{\linewidth}\raggedright
\textbf{Yil}
\end{minipage} & \begin{minipage}[b]{\linewidth}\raggedright
\textbf{Quduqlar}
\end{minipage} & \begin{minipage}[b]{\linewidth}\raggedright
\textbf{Oltita oʻn kunlikda qiymati bir xil quduqlar, \%}
\end{minipage} & \begin{minipage}[b]{\linewidth}\raggedright
\textbf{Yaroqli quduqlar (doimiy kuzatilganlar)}
\end{minipage} & \begin{minipage}[b]{\linewidth}\raggedright
\textbf{Oʻrtacha sath, m}
\end{minipage} & \begin{minipage}[b]{\linewidth}\raggedright
\textbf{Mediana, m}
\end{minipage} \\
\hline
\endhead
\hline
\endlastfoot
2009 & 121 & 22,3 & 90 & 1,63 & 1,53 \\ \hline
2011 & 137 & 23,4 & 97 & 1,57 & 1,45 \\ \hline
2012 & 133 & 31,6 & 89 & 1,75 & 1,62 \\ \hline
2014 & 135 & 27,4 & 92 & 1,69 & 1,57 \\ \hline
2015 & 134 & 19,4 & 98 & 1,75 & 1,72 \\ \hline
2016 & 125 & 19,2 & 91 & 1,75 & 1,66 \\ \hline
2018 & 188 & 16,5 & 95 & 1,96 & 1,95 \\ \hline
2019 & 188 & 16,5 & 98 & 1,87 & 1,84 \\ \hline
2020 & 141 & 0,0 & 92 & 2,04 & 1,90 \\ \hline
2021 & 188 & 12,8 & 100 & 2,20 & 2,06 \\ \hline
2022 & 188 & 12,8 & 99 & 1,99 & 1,88 \\ \hline
2023 & 188 & 14,9 & 98 & 2,04 & 1,98 \\ \hline
2024 & 188 & 13,8 & 97 & 2,09 & 2,00 \\ \hline
2025 & 188 & 32,4 & 79 & 2,03 & 1,89 \\
\end{longtable}

}

Massivlar boʻyicha bahorgi sath medianasi, m (2023 / 2024 / 2025):
Z.M.Bobur 1,30 / 1,62 / 2,00; Istiqlol 1,58 / 1,74 / 1,78; M.Toʻjiboyev
2,05 / 2,09 / 1,84; Mashʼal 1,97 / 1,86 / 1,88; Paxtakor 2,57 / 2,62 /
2,27; T.Mirzayev 1,89 / 1,95 / 1,88; U.Qoraboyev 2,82 / 2,92 / 2,24.

Izoh: 2010, 2013 va 2017-yillar boʻyicha mustaqil kuzatuv maʼlumotlari
mavjud emas. Doimiy kuzatilgan quduqlar --- 14 yilning kamida 10 tasida
yaroqli oʻlchovi bor 101 ta quduq. Manba: meliorativ ekspeditsiyaning
2025-yil mart va aprel holatidagi jadvallari; hisob: 61-skript
(results/spring\_2025.json).

\oakilova{57-ilova}

\oakilovanom{\textbf{Tavsiyanoma {[}33{]}: qisqartmalar va massivlarning asosiy
tavsifi}}

\begin{enumerate}
\def\labelenumi{\alph{enumi})}
\tightlist
\item
  Qisqartmalar
\end{enumerate}

{\fontsize{12}{14}\selectfont\setlength{\tabcolsep}{4pt}

\noindent\begin{longtable}[]{|
  >{\raggedright\arraybackslash}p{(\columnwidth - 4\tabcolsep - 3\arrayrulewidth) * \real{0.2153}}|>{\raggedright\arraybackslash}p{(\columnwidth - 4\tabcolsep - 3\arrayrulewidth) * \real{0.7847}}|}
\hline
\begin{minipage}[b]{\linewidth}\raggedright
\textbf{Qisqartma}
\end{minipage} & \begin{minipage}[b]{\linewidth}\raggedright
\textbf{Toʻliq nomi}
\end{minipage} \\
\hline
\endhead
\hline
\endlastfoot
BKP & Bioklimatik potensial \\ \hline
GAT & Geoaxborot tizimi \\ \hline
GTK & Gidrotermik koeffitsiyent (Selyaninov) \\ \hline
ITHB & Irrigatsiya tizimlari havza boshqarmasi \\ \hline
KDT & Kollektor-drenaj tarmogʻi \\ \hline
MSOZ & Sugʻoriladigan yerlarning meliorativ holati \\ \hline
ME & Meliorativ ekspeditsiya \\ \hline
PQ & Prezident qarori \\ \hline
SES & Sizot (yer osti) suvlari \\ \hline
FIK & Foydali ish koeffitsiyenti \\ \hline
Cm & Namlanish koeffitsiyenti (Ivanov) \\ \hline
ETc & Ekinning suv isteʼmoli (jami bugʻlanish) \\ \hline
ETo & Etalon jami bugʻlanish (FAO-56) \\ \hline
Im & Namlik indeksi (Shashko) \\ \hline
Kc & Ekin koeffitsiyenti \\ \hline
Pe & Foydali yogʻin \\ \hline
Rn & Radiatsiya balansi \\ \hline
VPD & Havoning toʻyinish tanqisligi \\
\end{longtable}

}

\begin{enumerate}
\def\labelenumi{\alph{enumi})}
\setcounter{enumi}{1}
\tightlist
\item
  Massivlarning asosiy tavsifi (tavsiyanoma, 1.1-jadval)
\end{enumerate}

{\fontsize{10}{12}\selectfont\setlength{\tabcolsep}{2.5pt}

\noindent\begin{longtable}[]{|
  >{\raggedright\arraybackslash}p{(\columnwidth - 16\tabcolsep - 9\arrayrulewidth) * \real{0.1832}}|>{\raggedright\arraybackslash}p{(\columnwidth - 16\tabcolsep - 9\arrayrulewidth) * \real{0.1926}}|>{\raggedright\arraybackslash}p{(\columnwidth - 16\tabcolsep - 9\arrayrulewidth) * \real{0.1363}}|>{\raggedright\arraybackslash}p{(\columnwidth - 16\tabcolsep - 9\arrayrulewidth) * \real{0.1112}}|>{\raggedright\arraybackslash}p{(\columnwidth - 16\tabcolsep - 9\arrayrulewidth) * \real{0.0990}}|>{\raggedright\arraybackslash}p{(\columnwidth - 16\tabcolsep - 9\arrayrulewidth) * \real{0.0964}}|>{\raggedright\arraybackslash}p{(\columnwidth - 16\tabcolsep - 9\arrayrulewidth) * \real{0.0839}}|>{\raggedright\arraybackslash}p{(\columnwidth - 16\tabcolsep - 9\arrayrulewidth) * \real{0.0974}}|}
\hline
\begin{minipage}[b]{\linewidth}\raggedright
\textbf{Massiv}
\end{minipage} & \begin{minipage}[b]{\linewidth}\raggedright
\textbf{Sugʻoriladigan maydon, ga}
\end{minipage} & \begin{minipage}[b]{\linewidth}\raggedright
\textbf{Balandlik, m}
\end{minipage} & \begin{minipage}[b]{\linewidth}\raggedright
\textbf{Oʻrtacha nishab, °}
\end{minipage} & \begin{minipage}[b]{\linewidth}\raggedright
\textbf{Nishabi \textgreater3° yerlar, \%}
\end{minipage} & \begin{minipage}[b]{\linewidth}\raggedright
\textbf{Bonitet balli}
\end{minipage} & \begin{minipage}[b]{\linewidth}\raggedright
\textbf{Gʻoʻza 2024, ga}
\end{minipage} & \begin{minipage}[b]{\linewidth}\raggedright
\textbf{Bugʻdoy 2024, ga}
\end{minipage} \\
\hline
\endhead
\hline
\endlastfoot
M.Toʻjiboyev & 2 167 & 559--719 & 2,5 & 25 & 58,9 & 456 & 503 \\ \hline
T.Mirzayev & 2 946 & 592--794 & 2,3 & 20 & 62,6 & 804 & 1 062 \\ \hline
Mashʼal & 3 198 & 568--685 & 2,1 & 16 & 59,5 & 772 & 1 100 \\ \hline
Paxtakor & 1 679 & 552--684 & 2,3 & 19 & 67,5 & 320 & 418 \\ \hline
U.Qoraboyev & 1 843 & 558--658 & 1,8 & 12 & 62,2 & 491 & 622 \\ \hline
Z.M.Bobur & 2 474 & 521--675 & 2,1 & 17 & 65,9 & 583 & 743 \\ \hline
Istiqlol & 2 478 & 505--795 & 5,8 & 62 & 57,5 & 516 & 572 \\ \hline
Y.Oxunboboyev & 1 189 & 651--895 & 5,9 & 79 & 43,5 & 91 & 158 \\ \hline
Jami & 17 974 & 505--895 & --- & --- & 59,4 & 4 033 & 5 548 \\
\end{longtable}

}

Manba: tavsiyanoma {[}33{]}; matni kirill yozuvidan lotin yozuviga
oʻgirilgan, sonlar oʻzgartirilmagan.

\oakilova{58-ilova}

\oakilovanom{\textbf{Marhamat tumani qishloq xoʻjaligining sxematik xaritasi
(massivlar chegaralari)}}

\oakimg{figures/tv0_sxema.jpg}

Manba: Marhamat tumani qishloq xoʻjaligi sxematik xaritasi (meliorativ
ekspeditsiya); xaritadagi yozuvlar asl holida (kirill yozuvida)
berilgan.

\oakilova{59-ilova}

\oakilovanom{\textbf{Tavsiyanoma {[}33{]}: Marhamat tumani relyefi, irrigatsiya va
drenaj tarmogʻi}}

\oakimg{figures/tv1_tarmoq.png}

Manba: tavsiyanoma {[}33{]}, 1-rasm; yozuvlar asl holida (kirill
yozuvida) berilgan.

\oakilova{60-ilova}

\oakilovanom{\textbf{Tavsiyanoma {[}33{]}: tuproqlarning shoʻrlanishi (a), eroziyasi
va toshloqligi (b)}}

\oakimg{figures/tv2_tuproq.png}

Manba: tavsiyanoma {[}33{]}, 2-rasm (tuproq-bonitet xaritasi asosida);
yozuvlar asl holida (kirill yozuvida) berilgan.

\oakilova{61-ilova}

\oakilovanom{\textbf{Tavsiyanoma {[}33{]}: sugʻoriladigan yerlarning meliorativ
holati, 2014--2025-yillar}}

\oakimg{figures/tv3_holat.png}

\begin{enumerate}
\def\labelenumi{\alph{enumi})}
\item
  Meliorativ holat toifalari, \%; b) minerallashuvi 3 g/l dan yuqori
  sizot suvli yerlar, ming ga (tavsiyanoma, 3-rasm; yozuvlar asl
  holida).
\item
  Massivlar kesimida meliorativ holat, 2023 → 2024 (tavsiyanoma,
  3.1-jadval)
\end{enumerate}

{\fontsize{12}{14}\selectfont\setlength{\tabcolsep}{4pt}

\noindent\begin{longtable}[]{|
  >{\raggedright\arraybackslash}p{(\columnwidth - 12\tabcolsep - 7\arrayrulewidth) * \real{0.2329}}|>{\raggedright\arraybackslash}p{(\columnwidth - 12\tabcolsep - 7\arrayrulewidth) * \real{0.1299}}|>{\raggedright\arraybackslash}p{(\columnwidth - 12\tabcolsep - 7\arrayrulewidth) * \real{0.1920}}|>{\raggedright\arraybackslash}p{(\columnwidth - 12\tabcolsep - 7\arrayrulewidth) * \real{0.0992}}|>{\raggedright\arraybackslash}p{(\columnwidth - 12\tabcolsep - 7\arrayrulewidth) * \real{0.1115}}|>{\raggedright\arraybackslash}p{(\columnwidth - 12\tabcolsep - 7\arrayrulewidth) * \real{0.2345}}|}
\hline
\begin{minipage}[b]{\linewidth}\raggedright
\textbf{Massiv}
\end{minipage} & \begin{minipage}[b]{\linewidth}\raggedright
\textbf{Yaxshi, \% (2023 → 2024)}
\end{minipage} & \begin{minipage}[b]{\linewidth}\raggedright
\textbf{Qoniqarsiz, \% (2024)}
\end{minipage} & \begin{minipage}[b]{\linewidth}\raggedright
\textbf{SES \textless2 m, \% (2023 → 2024)}
\end{minipage} & \begin{minipage}[b]{\linewidth}\raggedright
\textbf{SES \textgreater3 g/l, \% (2024)}
\end{minipage} & \begin{minipage}[b]{\linewidth}\raggedright
\textbf{Dinamika}
\end{minipage} \\
\hline
\endhead
\hline
\endlastfoot
Mashʼal & 31,8 → 2,2 & 4,7 & 31,5 → 78,3 & 61,7 & yomonlashish \\ \hline
T.Mirzayev & 85,6 → 53,1 & 3,1 & 1,0 → 24,9 & 0 & yomonlashish \\ \hline
Paxtakor & 13,0 → 11,2 & 4,1 & 19,2 → 19,6 & 49,5 & yomonlashish (suv
sifati) \\ \hline
Istiqlol & 37,4 → 31,2 & 5,9 & 35,4 → 42,1 & 4,6 & yomonlashish \\ \hline
Z.M.Bobur & 6,4 → 9,7 & 4,3 & 71,5 → 62,6 & 0 & barqaror qoniqarli \\ \hline
M.Toʻjiboyev & 50,3 → 54,2 & 0,9 & 15,1 → 16,9 & 0 & barqaror \\ \hline
U.Qoraboyev & 41,9 → 45,1 & 0,3 & 11,6 → 7,4 & 0 & barqaror \\ \hline
Y.Oxunboboyev & 94,5 → 95,0 & 0 & 0 → 0 & 0 & barqaror yaxshi \\
\end{longtable}

}

Manba: tavsiyanoma {[}33{]}; matni kirill yozuvidan lotin yozuviga
oʻgirilgan, sonlar oʻzgartirilmagan.

\oakilova{62-ilova}

\oakilovanom{\textbf{Tavsiyanoma {[}33{]}: sizot suvi sathi 2 m dan sayoz (a) va
minerallashuvi 3 g/l dan yuqori (b) yerlar ulushi, vegetatsiya davri,
2024-yil}}

\oakimg{figures/tv4_ses.png}

Manba: tavsiyanoma {[}33{]}, 4-rasm; yozuvlar asl holida (kirill
yozuvida) berilgan.

\oakilova{63-ilova}

\oakilovanom{\textbf{Tavsiyanoma {[}33{]}: agroiqlim koʻrsatkichlari (NASA POWER,
2001--2023-yillar oʻrtachasi)}}

{\fontsize{12}{14}\selectfont\setlength{\tabcolsep}{4pt}

\noindent\begin{longtable}[]{|
  >{\raggedright\arraybackslash}p{(\columnwidth - 6\tabcolsep - 4\arrayrulewidth) * \real{0.4178}}|>{\raggedright\arraybackslash}p{(\columnwidth - 6\tabcolsep - 4\arrayrulewidth) * \real{0.1799}}|>{\raggedright\arraybackslash}p{(\columnwidth - 6\tabcolsep - 4\arrayrulewidth) * \real{0.4023}}|}
\hline
\begin{minipage}[b]{\linewidth}\raggedright
\textbf{Koʻrsatkich}
\end{minipage} & \begin{minipage}[b]{\linewidth}\raggedright
\textbf{Qiymati}
\end{minipage} & \begin{minipage}[b]{\linewidth}\raggedright
\textbf{Bahosi}
\end{minipage} \\
\hline
\endhead
\hline
\endlastfoot
Havoning oʻrtacha yillik harorati, °C & 15,2 & issiq moʻʼtadil
mintaqa \\ \hline
Samarali haroratlar yigʻindisi ΣT\textgreater10 °C & 5 051 & goʻza va
bogʻdorchilik uchun yetarli \\ \hline
Havoning toʻyinish tanqisligi VPD (IV--X), kPa & 2,20 & yuqori
bugʻlanish talabi \\ \hline
Namlanish koeffitsiyenti Cm & 0,11 & juda quruq \\ \hline
Gidrotermik koeffitsiyent GTK & 0,21 & juda quruq (\textless0,4) \\ \hline
Namlik indeksi Im & 0,05 & juda quruq \\ \hline
Qurgʻoqchilik indeksi Rn/(L·P) & 4,7 & choʻl tipidagi iqlim
(\textgreater3) \\ \hline
Bioklimatik potensial BKP & 0,54 & sugʻorishsiz past \\ \hline
Yillik yogʻin / ETo, mm & 247 / 1 348 & tanqislik 1 100 mm \\
\end{longtable}

}

Manba: tavsiyanoma {[}33{]}; matni kirill yozuvidan lotin yozuviga
oʻgirilgan, sonlar oʻzgartirilmagan.

\oakilova{64-ilova}

\oakilovanom{\textbf{Tavsiyanoma {[}33{]}: massivlarni baholash koʻrsatkichlari va
integral bahosi (2024)}}

\begin{enumerate}
\def\labelenumi{\alph{enumi})}
\tightlist
\item
  Koʻrsatkichlar va ularning chegaralari (tavsiyanoma, 2.1-jadval)
\end{enumerate}

{\fontsize{12}{14}\selectfont\setlength{\tabcolsep}{4pt}

\noindent\begin{longtable}[]{|
  >{\raggedright\arraybackslash}p{(\columnwidth - 12\tabcolsep - 7\arrayrulewidth) * \real{0.0397}}|>{\raggedright\arraybackslash}p{(\columnwidth - 12\tabcolsep - 7\arrayrulewidth) * \real{0.3816}}|>{\raggedright\arraybackslash}p{(\columnwidth - 12\tabcolsep - 7\arrayrulewidth) * \real{0.0746}}|>{\raggedright\arraybackslash}p{(\columnwidth - 12\tabcolsep - 7\arrayrulewidth) * \real{0.0992}}|>{\raggedright\arraybackslash}p{(\columnwidth - 12\tabcolsep - 7\arrayrulewidth) * \real{0.0793}}|>{\raggedright\arraybackslash}p{(\columnwidth - 12\tabcolsep - 7\arrayrulewidth) * \real{0.3255}}|}
\hline
\begin{minipage}[b]{\linewidth}\raggedright
\textbf{№}
\end{minipage} & \begin{minipage}[b]{\linewidth}\raggedright
\textbf{Koʻrsatkich}
\end{minipage} & \begin{minipage}[b]{\linewidth}\raggedright
\textbf{2 ball}
\end{minipage} & \begin{minipage}[b]{\linewidth}\raggedright
\textbf{1 ball}
\end{minipage} & \begin{minipage}[b]{\linewidth}\raggedright
\textbf{0 ball}
\end{minipage} & \begin{minipage}[b]{\linewidth}\raggedright
\textbf{Manba}
\end{minipage} \\
\hline
\endhead
\hline
\endlastfoot
1 & SES chuqurligi 2 m dan kam yerlar ulushi, \% & \textless20 & 20--40
& \textgreater40 & tuman ME hisoboti, 12-jadval \\ \hline
2 & SES mineralizatsiyasi 3 g/l dan yuqori yerlar ulushi, \% &
\textless5 & 5--25 & \textgreater25 & tuman ME hisoboti, 12-jadval \\ \hline
3 & MSOZ «yaxshi» yerlar ulushi, \% & ≥50 & 25--50 & \textless25 & tuman
ME hisoboti, 15-jadval \\ \hline
4 & «Yaxshi» yerlar oʻzgarishi 2023→2024, foiz punkti & ≥0 & −10\ldots0
& \textless−10 & tuman ME hisoboti, 15-jadval \\ \hline
5 & Shoʻrlangan (kuchsiz va undan yuqori) tuproqlar ulushi, \% &
\textless15 & 15--35 & \textgreater35 & tuproq-bonitet xaritasi \\ \hline
6 & Yuvilgan (eroziyaga uchragan) tuproqlar ulushi, \% & \textless10 &
10--30 & \textgreater30 & tuproq-bonitet xaritasi \\ \hline
7 & Oʻrtacha nishab, ° & \textless2 & 2--4 & \textgreater4 & SRTM relyef
modeli \\ \hline
8 & Bonitet balli & ≥60 & 50--60 & \textless50 & tuproq-bonitet
xaritasi \\
\end{longtable}

}

\begin{enumerate}
\def\labelenumi{\alph{enumi})}
\setcounter{enumi}{1}
\tightlist
\item
  Massivlarning integral bahosi; qavs ichida --- koʻrsatkich olgan ball
  (tavsiyanoma, 4.2-jadval)
\end{enumerate}

{\fontsize{9}{11}\selectfont\setlength{\tabcolsep}{2.5pt}

\noindent\begin{longtable}[]{|
  >{\raggedright\arraybackslash}p{(\columnwidth - 22\tabcolsep - 12\arrayrulewidth) * \real{0.1586}}|>{\raggedright\arraybackslash}p{(\columnwidth - 22\tabcolsep - 12\arrayrulewidth) * \real{0.0488}}|>{\raggedright\arraybackslash}p{(\columnwidth - 22\tabcolsep - 12\arrayrulewidth) * \real{0.0488}}|>{\raggedright\arraybackslash}p{(\columnwidth - 22\tabcolsep - 12\arrayrulewidth) * \real{0.0837}}|>{\raggedright\arraybackslash}p{(\columnwidth - 22\tabcolsep - 12\arrayrulewidth) * \real{0.0837}}|>{\raggedright\arraybackslash}p{(\columnwidth - 22\tabcolsep - 12\arrayrulewidth) * \real{0.1435}}|>{\raggedright\arraybackslash}p{(\columnwidth - 22\tabcolsep - 12\arrayrulewidth) * \real{0.0957}}|>{\raggedright\arraybackslash}p{(\columnwidth - 22\tabcolsep - 12\arrayrulewidth) * \real{0.0837}}|>{\raggedright\arraybackslash}p{(\columnwidth - 22\tabcolsep - 12\arrayrulewidth) * \real{0.0837}}|>{\raggedright\arraybackslash}p{(\columnwidth - 22\tabcolsep - 12\arrayrulewidth) * \real{0.0478}}|>{\raggedright\arraybackslash}p{(\columnwidth - 22\tabcolsep - 12\arrayrulewidth) * \real{0.1220}}|}
\hline
\begin{minipage}[b]{\linewidth}\raggedright
\textbf{Massiv}
\end{minipage} & \begin{minipage}[b]{\linewidth}\raggedright
\textbf{SES \textless2 m, \%}
\end{minipage} & \begin{minipage}[b]{\linewidth}\raggedright
\textbf{SES \textgreater3 g/l, \%}
\end{minipage} & \begin{minipage}[b]{\linewidth}\raggedright
\textbf{Yaxshi, \%}
\end{minipage} & \begin{minipage}[b]{\linewidth}\raggedright
\textbf{Δ yaxshi, f.p.}
\end{minipage} & \begin{minipage}[b]{\linewidth}\raggedright
\textbf{Shoʻrlangan, \%}
\end{minipage} & \begin{minipage}[b]{\linewidth}\raggedright
\textbf{Eroziya, \%}
\end{minipage} & \begin{minipage}[b]{\linewidth}\raggedright
\textbf{Nishab, °}
\end{minipage} & \begin{minipage}[b]{\linewidth}\raggedright
\textbf{Bonitet}
\end{minipage} & \begin{minipage}[b]{\linewidth}\raggedright
\textbf{Ball}
\end{minipage} & \begin{minipage}[b]{\linewidth}\raggedright
\textbf{Sinf}
\end{minipage} \\
\hline
\endhead
\hline
\endlastfoot
Mashʼal & 78,3 (0) & 61,7 (0) & 2,2 (0) & −29,7 (0) & 27,5 (1) & 6,7 (2)
& 2,0 (1) & 59,5 (1) & 5 & qoniqarsiz \\ \hline
Istiqlol & 42,1 (0) & 4,6 (2) & 31,2 (1) & −6,2 (1) & 15,6 (1) & 19,1
(1) & 5,8 (0) & 57,5 (1) & 7 & qoniqarsiz \\ \hline
Y.Oxunboboyev & 0 (2) & 0 (2) & 95,0 (2) & +0,5 (2) & 41,7 (0) & 50,3
(0) & 5,9 (0) & 43,5 (0) & 8 & qoniqarli \\ \hline
Paxtakor & 19,6 (2) & 49,5 (0) & 11,2 (0) & −1,8 (1) & 51,9 (0) & 0 (2)
& 2,3 (1) & 67,5 (2) & 8 & qoniqarli \\ \hline
Z.M.Bobur & 62,6 (0) & 0 (2) & 9,7 (0) & +3,3 (2) & 43,3 (0) & 0,2 (2) &
2,1 (1) & 65,9 (2) & 9 & qoniqarli \\ \hline
T.Mirzayev & 24,9 (1) & 0 (2) & 53,1 (2) & −32,5 (0) & 31,7 (1) & 6,4
(2) & 2,3 (1) & 62,6 (2) & 11 & qoniqarli \\ \hline
M.Toʻjiboyev & 16,9 (2) & 0 (2) & 54,2 (2) & +3,9 (2) & 6,5 (2) & 30,7
(0) & 2,5 (1) & 58,9 (1) & 12 & qoniqarli \\ \hline
U.Qoraboyev & 7,4 (2) & 0 (2) & 45,1 (1) & +3,2 (2) & 12,5 (2) & 7,9 (2)
& 1,8 (2) & 62,2 (2) & 15 & yaxshi \\
\end{longtable}

}

Jami ball boʻyicha sinflar: 14--16 --- yaxshi, 8--13 --- qoniqarli, 0--7
--- qoniqarsiz.

Manba: tavsiyanoma {[}33{]}; matni kirill yozuvidan lotin yozuviga
oʻgirilgan, sonlar oʻzgartirilmagan.

\oakilova{65-ilova}

\oakilovanom{\textbf{Tavsiyanoma {[}33{]}: massivlarning integral bahosi (a) va
maqbul meliorativ rejimi (b)}}

\oakimg{figures/tv5_integral.png}

Manba: tavsiyanoma {[}33{]}, 5-rasm; yozuvlar asl holida (kirill
yozuvida) berilgan.

\oakilova{66-ilova}

\oakilovanom{\textbf{Tavsiyanoma {[}33{]}: Marhamat tumaniga olingan sugʻorish suvi}}

\begin{enumerate}
\def\labelenumi{\alph{enumi})}
\tightlist
\item
  Yillar boʻyicha olingan suv (tavsiyanoma, 5.1-jadval)
\end{enumerate}

{\fontsize{9}{11}\selectfont\setlength{\tabcolsep}{2.5pt}

\noindent\begin{longtable}[]{|
  >{\raggedright\arraybackslash}p{(\columnwidth - 18\tabcolsep - 10\arrayrulewidth) * \real{0.0480}}|>{\raggedright\arraybackslash}p{(\columnwidth - 18\tabcolsep - 10\arrayrulewidth) * \real{0.0705}}|>{\raggedright\arraybackslash}p{(\columnwidth - 18\tabcolsep - 10\arrayrulewidth) * \real{0.0940}}|>{\raggedright\arraybackslash}p{(\columnwidth - 18\tabcolsep - 10\arrayrulewidth) * \real{0.1410}}|>{\raggedright\arraybackslash}p{(\columnwidth - 18\tabcolsep - 10\arrayrulewidth) * \real{0.1410}}|>{\raggedright\arraybackslash}p{(\columnwidth - 18\tabcolsep - 10\arrayrulewidth) * \real{0.1293}}|>{\raggedright\arraybackslash}p{(\columnwidth - 18\tabcolsep - 10\arrayrulewidth) * \real{0.1645}}|>{\raggedright\arraybackslash}p{(\columnwidth - 18\tabcolsep - 10\arrayrulewidth) * \real{0.0823}}|>{\raggedright\arraybackslash}p{(\columnwidth - 18\tabcolsep - 10\arrayrulewidth) * \real{0.1293}}|}
\hline
\begin{minipage}[b]{\linewidth}\raggedright
\textbf{Yil}
\end{minipage} & \begin{minipage}[b]{\linewidth}\raggedright
\textbf{Yillik limit, mln m³}
\end{minipage} & \begin{minipage}[b]{\linewidth}\raggedright
\textbf{Amalda olingan, mln m³}
\end{minipage} & \begin{minipage}[b]{\linewidth}\raggedright
\textbf{Limitdan foydalanish, \%}
\end{minipage} & \begin{minipage}[b]{\linewidth}\raggedright
\textbf{shu jumladan tik quduqlardan, mln m³}
\end{minipage} & \begin{minipage}[b]{\linewidth}\raggedright
\textbf{Vegetatsiya davri: limit / amalda, mln m³}
\end{minipage} & \begin{minipage}[b]{\linewidth}\raggedright
\textbf{Vegetatsiyadan tashqari, mln m³}
\end{minipage} & \begin{minipage}[b]{\linewidth}\raggedright
\textbf{Yillik, m³/ga}
\end{minipage} & \begin{minipage}[b]{\linewidth}\raggedright
\textbf{Vegetatsiya davrida, m³/ga}
\end{minipage} \\
\hline
\endhead
\hline
\endlastfoot
2014 & 151,1 & 157,6 & 104,3 & --- & 135,3 / 142,4 & 15,1 & 8 568 & 7
745 \\ \hline
2015 & 183,5 & 203,6 & 110,9 & --- & 143,5 / 159,0 & 44,6 & 11 065 & 8
639 \\ \hline
2016 & 213,5 & 203,4 & 95,3 & --- & --- & --- & 11 056 & --- \\ \hline
2019 & 206,7 & 191,5 & 92,7 & 5,18 & 156,2 / 138,2 & 53,3 & 10 945 & 7
900 \\ \hline
2020 & 174,1 & 159,8 & 91,8 & 5,69 & 136,5 / 121,6 & 38,2 & 9 132 & 6
950 \\ \hline
2022 & 148,1 & 138,5 & 93,5 & --- & 106,5 / 113,6 & 24,9 & 7 704 & 6
319 \\ \hline
2023 & 167,9 & 153,4 & 91,3 & --- & 120,5 / 120,5 & 32,9 & 8 532 & 6
704 \\ \hline
2024 & 159,0 & 147,3 & 92,6 & --- & 124,8 / 124,8 & 22,5 & 8 195 & 6
943 \\ \hline
2025 & 148,0 & 147,3 & 99,5 & 0,28 & 109,3 / 109,3 & 38,0 & 8 200 & 6
084 \\
\end{longtable}

}

\begin{enumerate}
\def\labelenumi{\alph{enumi})}
\setcounter{enumi}{1}
\tightlist
\item
  Kuzgi-qishki davrda olingan suv, shoʻr yuvish va nam toʻplash maydoni
  (tavsiyanoma, 5.2-jadval)
\end{enumerate}

{\fontsize{11}{13}\selectfont\setlength{\tabcolsep}{2.5pt}

\noindent\begin{longtable}[]{|
  >{\raggedright\arraybackslash}p{(\columnwidth - 14\tabcolsep - 8\arrayrulewidth) * \real{0.1419}}|>{\raggedright\arraybackslash}p{(\columnwidth - 14\tabcolsep - 8\arrayrulewidth) * \real{0.1106}}|>{\raggedright\arraybackslash}p{(\columnwidth - 14\tabcolsep - 8\arrayrulewidth) * \real{0.1273}}|>{\raggedright\arraybackslash}p{(\columnwidth - 14\tabcolsep - 8\arrayrulewidth) * \real{0.1851}}|>{\raggedright\arraybackslash}p{(\columnwidth - 14\tabcolsep - 8\arrayrulewidth) * \real{0.1307}}|>{\raggedright\arraybackslash}p{(\columnwidth - 14\tabcolsep - 8\arrayrulewidth) * \real{0.1594}}|>{\raggedright\arraybackslash}p{(\columnwidth - 14\tabcolsep - 8\arrayrulewidth) * \real{0.1450}}|}
\hline
\begin{minipage}[b]{\linewidth}\raggedright
\textbf{Mavsum}
\end{minipage} & \begin{minipage}[b]{\linewidth}\raggedright
\textbf{Suv limiti, mln m³}
\end{minipage} & \begin{minipage}[b]{\linewidth}\raggedright
\textbf{Amalda olingan, mln m³}
\end{minipage} & \begin{minipage}[b]{\linewidth}\raggedright
\textbf{Limitdan foydalanish, \%}
\end{minipage} & \begin{minipage}[b]{\linewidth}\raggedright
\textbf{Shoʻr yuvish va nam toʻplash maydoni: reja / amalda, ming ga}
\end{minipage} & \begin{minipage}[b]{\linewidth}\raggedright
\textbf{Nam toʻplashga talab etilgan suv, mln m³}
\end{minipage} & \begin{minipage}[b]{\linewidth}\raggedright
\textbf{Gektariga olingan suv, m³/ga}
\end{minipage} \\
\hline
\endhead
\hline
\endlastfoot
2023--2024 & 23,70 & 23,70 & 100 & 4,60 / 4,60 & 16,8 & 5 150 \\ \hline
2024--2025 & 27,52 & 19,76 & 72 & 3,75 / 3,75 & 21,1 & 5 270 \\
\end{longtable}

}

\begin{enumerate}
\def\labelenumi{\alph{enumi})}
\setcounter{enumi}{21}
\tightlist
\item
  Gektariga berilgan suv va ekinlarning sugʻorish meʼyorlari, m³/ga
  (tavsiyanoma, 5.3-jadval)
\end{enumerate}

{\fontsize{12}{14}\selectfont\setlength{\tabcolsep}{4pt}

\noindent\begin{longtable}[]{|
  >{\raggedright\arraybackslash}p{(\columnwidth - 4\tabcolsep - 3\arrayrulewidth) * \real{0.6110}}|>{\raggedright\arraybackslash}p{(\columnwidth - 4\tabcolsep - 3\arrayrulewidth) * \real{0.3890}}|}
\hline
\begin{minipage}[b]{\linewidth}\raggedright
\textbf{Koʻrsatkich}
\end{minipage} & \begin{minipage}[b]{\linewidth}\raggedright
\textbf{Qiymati, m³/ga}
\end{minipage} \\
\hline
\endhead
\hline
\endlastfoot
Tuman chegarasida olingan suv, yillik (2025) & 8 200 \\ \hline
shu jumladan vegetatsiya davrida (2025) & 6 080 \\ \hline
Isteʼmolchilarga yetkazilgan suv (Marhamat) & 7 600 \\ \hline
Isteʼmolchilarga yetkazilgan suv (viloyat boʻyicha oʻrtacha) & 8 100 \\ \hline
Oʻrtacha sugʻorish meʼyori (gidromodul rayonlashtirish boʻyicha) & 6
600 \\ \hline
Gʻoʻza, egatlab sugʻorish: SES 1--2 / 2--3 / \textgreater3 m & 7 430 / 9
220 / 10 420 \\ \hline
Gʻoʻza, tomchilatib sugʻorish: SES 1--2 / 2--3 / \textgreater3 m & 5 360
/ 6 660 / 7 520 \\ \hline
Kuzgi bugʻdoy, egatlab sugʻorish: SES 1--2 / 2--3 / \textgreater3 m & 2
350 / 3 250 / 3 850 \\
\end{longtable}

}

Manba: tavsiyanoma {[}33{]}; matni kirill yozuvidan lotin yozuviga
oʻgirilgan, sonlar oʻzgartirilmagan.

\oakilova{67-ilova}

\oakilovanom{\textbf{Marhamat tumaniga olingan sugʻorish suvi (a) va gektariga
olingan suv (b), 2014--2025-yillar}}

\oakimg{figures/tv_water.png}

Manba: tavsiyanoma {[}33{]}, 6-rasm; 66-ilovadagi sonlar boʻyicha lotin
yozuvida qayta chizilgan.

\oakilova{68-ilova}

\oakilovanom{\textbf{Tavsiyanoma {[}33{]}: kollektor-drenaj oqimi, suv minerallashuvi
va suv-tuz balansi}}

\begin{enumerate}
\def\labelenumi{\alph{enumi})}
\tightlist
\item
  Kollektor-drenaj oqimi va suv minerallashuvi (tavsiyanoma, 5.4-jadval)
\end{enumerate}

{\fontsize{9}{11}\selectfont\setlength{\tabcolsep}{2.5pt}

\noindent\begin{longtable}[]{|
  >{\raggedright\arraybackslash}p{(\columnwidth - 16\tabcolsep - 9\arrayrulewidth) * \real{0.0474}}|>{\raggedright\arraybackslash}p{(\columnwidth - 16\tabcolsep - 9\arrayrulewidth) * \real{0.0697}}|>{\raggedright\arraybackslash}p{(\columnwidth - 16\tabcolsep - 9\arrayrulewidth) * \real{0.1162}}|>{\raggedright\arraybackslash}p{(\columnwidth - 16\tabcolsep - 9\arrayrulewidth) * \real{0.1278}}|>{\raggedright\arraybackslash}p{(\columnwidth - 16\tabcolsep - 9\arrayrulewidth) * \real{0.1394}}|>{\raggedright\arraybackslash}p{(\columnwidth - 16\tabcolsep - 9\arrayrulewidth) * \real{0.0813}}|>{\raggedright\arraybackslash}p{(\columnwidth - 16\tabcolsep - 9\arrayrulewidth) * \real{0.2091}}|>{\raggedright\arraybackslash}p{(\columnwidth - 16\tabcolsep - 9\arrayrulewidth) * \real{0.2091}}|}
\hline
\begin{minipage}[b]{\linewidth}\raggedright
\textbf{Yil}
\end{minipage} & \begin{minipage}[b]{\linewidth}\raggedright
\textbf{Zovur oqimi, mln m³}
\end{minipage} & \begin{minipage}[b]{\linewidth}\raggedright
\textbf{shu jumladan tik drenajdan, mln m³}
\end{minipage} & \begin{minipage}[b]{\linewidth}\raggedright
\textbf{Daryoga tashlangan, mln m³}
\end{minipage} & \begin{minipage}[b]{\linewidth}\raggedright
\textbf{Sugʻorishga qaytarilgan, mln m³}
\end{minipage} & \begin{minipage}[b]{\linewidth}\raggedright
\textbf{Zovur oqimi / olingan suv, \%}
\end{minipage} & \begin{minipage}[b]{\linewidth}\raggedright
\textbf{Sugʻorish suvi mineralizatsiyasi, g/l}
\end{minipage} & \begin{minipage}[b]{\linewidth}\raggedright
\textbf{Zovur suvi mineralizatsiyasi, g/l}
\end{minipage} \\
\hline
\endhead
\hline
\endlastfoot
2014 & 9,3 & 0,80 & --- & --- & 5,9 & 1,18 & 1,09 \\ \hline
2015 & 10,6 & --- & --- & --- & 5,2 & 0,88 & 0,64 \\ \hline
2016 & 125,0 & --- & --- & --- & 61,5 & 0,85 & 0,78 \\ \hline
2019 & 200,1 & --- & --- & --- & 104,5 & 0,60 & 0,81 \\ \hline
2020 & 151,9 & --- & --- & --- & 95,1 & 0,60 & 0,76 \\ \hline
2022 & 90,1 & 0,81 & 90,1 & 0 & 65,1 & 0,38 & 0,66 \\ \hline
2023 & 98,3 & 0,57 & 98,3 & 0 & 64,1 & 0,38 & 0,66 \\ \hline
2024 & 98,2 & 0,63 & 97,6 & 0 & 66,7 & 0,40 & 0,75 \\ \hline
2025 & 118,5 & 0,74 & 117,5 & 0,28 & 80,4 & 0,52 & 0,99 \\
\end{longtable}

}

\begin{enumerate}
\def\labelenumi{\alph{enumi})}
\setcounter{enumi}{1}
\tightlist
\item
  Suv-tuz balansi, quruq qoldiq boʻyicha (tavsiyanoma, 5.5-jadval)
\end{enumerate}

{\fontsize{12}{14}\selectfont\setlength{\tabcolsep}{4pt}

\noindent\begin{longtable}[]{|
  >{\raggedright\arraybackslash}p{(\columnwidth - 10\tabcolsep - 6\arrayrulewidth) * \real{0.1218}}|>{\raggedright\arraybackslash}p{(\columnwidth - 10\tabcolsep - 6\arrayrulewidth) * \real{0.2239}}|>{\raggedright\arraybackslash}p{(\columnwidth - 10\tabcolsep - 6\arrayrulewidth) * \real{0.2468}}|>{\raggedright\arraybackslash}p{(\columnwidth - 10\tabcolsep - 6\arrayrulewidth) * \real{0.2109}}|>{\raggedright\arraybackslash}p{(\columnwidth - 10\tabcolsep - 6\arrayrulewidth) * \real{0.1966}}|}
\hline
\begin{minipage}[b]{\linewidth}\raggedright
\textbf{Yil}
\end{minipage} & \begin{minipage}[b]{\linewidth}\raggedright
\textbf{Kirgan tuz, ming t}
\end{minipage} & \begin{minipage}[b]{\linewidth}\raggedright
\textbf{Chiqqan tuz, ming t}
\end{minipage} & \begin{minipage}[b]{\linewidth}\raggedright
\textbf{Balans, ming t}
\end{minipage} & \begin{minipage}[b]{\linewidth}\raggedright
\textbf{Balans, t/ga}
\end{minipage} \\
\hline
\endhead
\hline
\endlastfoot
2014 & 185,9 & 10,1 & +175,8 & +9,56 \\ \hline
2015 & 179,2 & 6,8 & +172,4 & +9,37 \\ \hline
2016 & 172,9 & 97,5 & +75,4 & +4,10 \\ \hline
2019 & 114,9 & 162,1 & −47,2 & −2,70 \\ \hline
2020 & 95,9 & 115,4 & −19,6 & −1,12 \\ \hline
2022 & 52,6 & 59,5 & −6,9 & −0,38 \\ \hline
2023 & 58,3 & 64,9 & −6,6 & −0,37 \\ \hline
2024 & 57,3 & 73,6 & −16,3 & −0,91 \\ \hline
2025 & 76,6 & 87,4 & −10,8 & −0,60 \\
\end{longtable}

}

Manba: tavsiyanoma {[}33{]}; matni kirill yozuvidan lotin yozuviga
oʻgirilgan, sonlar oʻzgartirilmagan.

\oakilova{69-ilova}

\oakilovanom{\textbf{Zovur oqimining olingan suvga nisbati (a), sugʻorish va zovur
suvining minerallashuvi hamda tuz balansi (b)}}

\oakimg{figures/tv_drainage.png}

Manba: tavsiyanoma {[}33{]}, 7-rasm; 68-ilovadagi sonlar boʻyicha lotin
yozuvida qayta chizilgan.

\oakilova{70-ilova}

\oakilovanom{\textbf{Tavsiyanoma {[}33{]}: maqbul meliorativ rejim va sugʻorish
meʼyorlari}}

\begin{enumerate}
\def\labelenumi{\alph{enumi})}
\tightlist
\item
  Vegetatsiya davrida sizot suvi sathining maqbul va kritik chuqurligi,
  oʻrta qumoqli tuproqlar (tavsiyanoma, 6.1-jadval)
\end{enumerate}

{\fontsize{12}{14}\selectfont\setlength{\tabcolsep}{4pt}

\noindent\begin{longtable}[]{|
  >{\raggedright\arraybackslash}p{(\columnwidth - 8\tabcolsep - 5\arrayrulewidth) * \real{0.3086}}|>{\raggedright\arraybackslash}p{(\columnwidth - 8\tabcolsep - 5\arrayrulewidth) * \real{0.1767}}|>{\raggedright\arraybackslash}p{(\columnwidth - 8\tabcolsep - 5\arrayrulewidth) * \real{0.1767}}|>{\raggedright\arraybackslash}p{(\columnwidth - 8\tabcolsep - 5\arrayrulewidth) * \real{0.3380}}|}
\hline
\begin{minipage}[b]{\linewidth}\raggedright
\textbf{SES mineralizatsiyasi, g/l}
\end{minipage} & \begin{minipage}[b]{\linewidth}\raggedright
\textbf{Maqbul chuqurlik, m}
\end{minipage} & \begin{minipage}[b]{\linewidth}\raggedright
\textbf{Kritik chuqurlik, m}
\end{minipage} & \begin{minipage}[b]{\linewidth}\raggedright
\textbf{Rejim turi}
\end{minipage} \\
\hline
\endhead
\hline
\endlastfoot
\textless1 & 1,5--2,0 & 1,2 & gidromorf (SES suvidan foydalanish) \\ \hline
1--3 & 2,0--2,5 & 1,8 & yarim gidromorf \\ \hline
3--5 & 2,5--3,0 & 2,2 & yarim avtomorf, drenaj jadallashtiriladi \\ \hline
\textgreater5 & \textgreater3,0 & 2,5 & avtomorf, shoʻr yuvish rejimi \\
\end{longtable}

}

\begin{enumerate}
\def\labelenumi{\alph{enumi})}
\setcounter{enumi}{1}
\tightlist
\item
  Sugʻorish meʼyorlari, m³/ga (tavsiyanoma, 6.2-jadval)
\end{enumerate}

{\fontsize{12}{14}\selectfont\setlength{\tabcolsep}{4pt}

\noindent\begin{longtable}[]{|
  >{\raggedright\arraybackslash}p{(\columnwidth - 14\tabcolsep - 8\arrayrulewidth) * \real{0.1463}}|>{\raggedright\arraybackslash}p{(\columnwidth - 14\tabcolsep - 8\arrayrulewidth) * \real{0.1798}}|>{\raggedright\arraybackslash}p{(\columnwidth - 14\tabcolsep - 8\arrayrulewidth) * \real{0.0735}}|>{\raggedright\arraybackslash}p{(\columnwidth - 14\tabcolsep - 8\arrayrulewidth) * \real{0.1481}}|>{\raggedright\arraybackslash}p{(\columnwidth - 14\tabcolsep - 8\arrayrulewidth) * \real{0.1119}}|>{\raggedright\arraybackslash}p{(\columnwidth - 14\tabcolsep - 8\arrayrulewidth) * \real{0.1351}}|>{\raggedright\arraybackslash}p{(\columnwidth - 14\tabcolsep - 8\arrayrulewidth) * \real{0.2052}}|}
\hline
\begin{minipage}[b]{\linewidth}\raggedright
\textbf{Ekin}
\end{minipage} & \begin{minipage}[b]{\linewidth}\raggedright
\textbf{SES chuqurligi}
\end{minipage} & \begin{minipage}[b]{\linewidth}\raggedright
\textbf{ETc, mm}
\end{minipage} & \begin{minipage}[b]{\linewidth}\raggedright
\textbf{SES hissasi, mm}
\end{minipage} & \begin{minipage}[b]{\linewidth}\raggedright
\textbf{Netto meʼyor}
\end{minipage} & \begin{minipage}[b]{\linewidth}\raggedright
\textbf{Egatlab (η = 0,65)}
\end{minipage} & \begin{minipage}[b]{\linewidth}\raggedright
\textbf{Tomchilatib (η = 0,90)}
\end{minipage} \\
\hline
\endhead
\hline
\endlastfoot
Gʻoʻza & 1--2 m & 778 & 233 & 4 830 & 7 430 & 5 360 \\ \hline
Gʻoʻza & 2--3 m & 778 & 117 & 5 990 & 9 220 & 6 660 \\ \hline
Gʻoʻza & \textgreater3 m & 778 & 39 & 6 770 & 10 420 & 7 520 \\ \hline
Kuzgi bugʻdoy & 1--2 m & 392 & 118 & 1 530 & 2 350 & 1 700 \\ \hline
Kuzgi bugʻdoy & 2--3 m & 392 & 59 & 2 110 & 3 250 & 2 350 \\ \hline
Kuzgi bugʻdoy & \textgreater3 m & 392 & 20 & 2 510 & 3 850 & 2 780 \\
\end{longtable}

}

\begin{enumerate}
\def\labelenumi{\alph{enumi})}
\setcounter{enumi}{21}
\tightlist
\item
  Massivlar boʻyicha maqbul meliorativ rejim (tavsiyanoma, 6.3-jadval)
\end{enumerate}

{\fontsize{11}{13}\selectfont\setlength{\tabcolsep}{2.5pt}

\noindent\begin{longtable}[]{|
  >{\raggedright\arraybackslash}p{(\columnwidth - 12\tabcolsep - 7\arrayrulewidth) * \real{0.1963}}|>{\raggedright\arraybackslash}p{(\columnwidth - 12\tabcolsep - 7\arrayrulewidth) * \real{0.1367}}|>{\raggedright\arraybackslash}p{(\columnwidth - 12\tabcolsep - 7\arrayrulewidth) * \real{0.1644}}|>{\raggedright\arraybackslash}p{(\columnwidth - 12\tabcolsep - 7\arrayrulewidth) * \real{0.1643}}|>{\raggedright\arraybackslash}p{(\columnwidth - 12\tabcolsep - 7\arrayrulewidth) * \real{0.1306}}|>{\raggedright\arraybackslash}p{(\columnwidth - 12\tabcolsep - 7\arrayrulewidth) * \real{0.2077}}|}
\hline
\begin{minipage}[b]{\linewidth}\raggedright
\textbf{Massiv}
\end{minipage} & \begin{minipage}[b]{\linewidth}\raggedright
\textbf{SES 2024, m (oʻrtacha / eng sayoz)}
\end{minipage} & \begin{minipage}[b]{\linewidth}\raggedright
\textbf{Hozirgi rejim}
\end{minipage} & \begin{minipage}[b]{\linewidth}\raggedright
\textbf{Maqsadli SES chuqurligi, m}
\end{minipage} & \begin{minipage}[b]{\linewidth}\raggedright
\textbf{Gʻoʻza meʼyori, m³/ga (egat / tomchi)}
\end{minipage} & \begin{minipage}[b]{\linewidth}\raggedright
\textbf{Asosiy talab}
\end{minipage} \\
\hline
\endhead
\hline
\endlastfoot
Mashʼal & 1,93 / 1,75 & gidromorf, SES 3--5 g/l & ≥2,5 & 9 220 / 6 660 &
drenajni tiklash, SES ni pasaytirish \\ \hline
Paxtakor & 2,43 / 2,04 & yarim gidromorf, SES 3--5 g/l & ≥2,5 & 9 220 /
6 660 & SES ni pasaytirish, shoʻr yuvish \\ \hline
Z.M.Bobur & 1,95 / 1,70 & gidromorf, chuchuk SES & 1,8--2,2 & 7 430 / 5
360 & suvni kamaytirish, SES dan foydalanish \\ \hline
Istiqlol & 2,07 / 1,71 & aralash: tekislikda gidromorf, adirda avtomorf
& 2,0--2,5 & 7 430--10 420 / 5 360--7 520 & tekislik va adirni alohida
boshqarish \\ \hline
M.Toʻjiboyev & 2,45 / 2,27 & yarim gidromorf & 2,0--2,5 & 9 220 / 6 660
& eroziyaga qarshi sugʻorish \\ \hline
U.Qoraboyev & 2,82 / 2,37 & yarim gidromorf & 2,0--2,5 & 9 220 / 6 660 &
hozirgi holatni saqlash \\ \hline
T.Mirzayev & 3,26 / 3,02 & avtomorf (25\% yerda \textless2 m) & ≥2,0 &
10 420 / 7 520 & sayoz SES li uchastkalarni aniqlash \\ \hline
Y.Oxunboboyev & 3,32 / 2,63 & avtomorf & \textgreater3,0 & 10 420 / 7
520 & tomchilatib sugʻorish, eroziyani toʻxtatish \\
\end{longtable}

}

Manba: tavsiyanoma {[}33{]}; matni kirill yozuvidan lotin yozuviga
oʻgirilgan, sonlar oʻzgartirilmagan.

\oakilova{71-ilova}

\oakilovanom{\textbf{Tavsiyanoma {[}33{]}: massivlar boʻyicha tadbirlar rejasi}}

{\fontsize{12}{14}\selectfont\setlength{\tabcolsep}{4pt}

\noindent\begin{longtable}[]{|
  >{\raggedright\arraybackslash}p{(\columnwidth - 10\tabcolsep - 6\arrayrulewidth) * \real{0.0958}}|>{\raggedright\arraybackslash}p{(\columnwidth - 10\tabcolsep - 6\arrayrulewidth) * \real{0.2128}}|>{\raggedright\arraybackslash}p{(\columnwidth - 10\tabcolsep - 6\arrayrulewidth) * \real{0.3125}}|>{\raggedright\arraybackslash}p{(\columnwidth - 10\tabcolsep - 6\arrayrulewidth) * \real{0.2320}}|>{\raggedright\arraybackslash}p{(\columnwidth - 10\tabcolsep - 6\arrayrulewidth) * \real{0.1469}}|}
\hline
\begin{minipage}[b]{\linewidth}\raggedright
\textbf{Navbat}
\end{minipage} & \begin{minipage}[b]{\linewidth}\raggedright
\textbf{Massiv}
\end{minipage} & \begin{minipage}[b]{\linewidth}\raggedright
\textbf{Tadbir}
\end{minipage} & \begin{minipage}[b]{\linewidth}\raggedright
\textbf{Kutiladigan natija}
\end{minipage} & \begin{minipage}[b]{\linewidth}\raggedright
\textbf{Muddat}
\end{minipage} \\
\hline
\endhead
\hline
\endlastfoot
1 & Mashʼal & KDT va xoʻjaliklararo zovurlarni tozalash, tik drenaj
quduqlarini tiklash; zovur oqimi va mineralizatsiyasini har oy oʻlchash
& vegetatsiyada SES ≥2,5 m, \textgreater3 g/l maydon 1,97 ming ga dan
kamayadi & 2026--2027 \\ \hline
1 & Mashʼal & Suv taqsimotini meʼyorlash (gʻoʻza uchun 9 220 m³/ga dan
oshmasligi), oqova suvlarni kamaytirish & filtratsiya isrofi kamayadi &
har yil \\ \hline
1 & Paxtakor & Kuzgi tuproq suratmasi; shoʻrlangani tasdiqlangan
yerlarda shoʻr yuvish (≈ 2 200 m³/ga); zovurlarni tozalash & SES ≥2,5 m,
«yaxshi» yerlar ulushi 11\% dan oshadi & 2026--2027 \\ \hline
1 & Z.M.Bobur & Sugʻorish meʼyorini «1--2 m» guruhi darajasiga tushirish
(7 430 m³/ga); sugʻorishlar oraligʻini uzaytirish & SES suvidan
foydalaniladi, \textless2 m yerlar 63\% dan kamayadi & har yil \\ \hline
2 & Istiqlol & Adir qismida (nishab \textgreater3°) tomchilatib
sugʻorish; tekislikda zovurlarni tozalash & eroziya va pastki yerlarga
suv tushishi kamayadi & 2026--2028 \\ \hline
2 & Y.Oxunboboyev & Tomchilatib yoki yomgʻirlatib sugʻorish; nishab
yerlarda terrasalash, bogʻ qator oralarini chimzorlash & suv sarfi
gektarga ≈ 2 900 m³ ga kamayadi, eroziya toʻxtaydi & 2026--2028 \\ \hline
2 & M.Toʻjiboyev & Egat uzunligini qisqartirish va kichik oqimli
sugʻorish; nishab uchastkalarda tomchilatib sugʻorish & yuvilgan yerlar
(31\%) kamayadi & 2026--2028 \\ \hline
2 & T.Mirzayev & SES sathi 2 m dan sayoz uchastkalarni (25\%) joyida
aniqlash; zovursiz 1 120 ga da mahalliy drenaj zaruratini baholash &
2024 yildagi salbiy siljish sababi aniqlanadi & 2026 \\ \hline
3 & U.Qoraboyev & Hozirgi sugʻorish va drenaj rejimini saqlash,
monitoring & «yaxshi» baho saqlanadi & har yil \\ \hline
Tuman & Barcha massivlar & 188 ta kuzatuv qudugʻida SES
mineralizatsiyasini yiliga kamida 2 marta oʻlchash; natijalarni GAT da
1×1 km toʻr boʻyicha tahlil qilish & yomonlashish ilgari aniqlanadi &
har yil \\ \hline
Tuman & Barcha massivlar & ME bahosi va tuproq-bonitet xaritasidagi
shoʻrlanish maʼlumotlarini muvofiqlashtirish (5 097 ga) & shoʻr yuvish
aniq yerlarda oʻtkaziladi & 2026 \\ \hline
Tuman & Barcha massivlar & Integral ballni har yili ME hisoboti asosida
yangilash & tadbirlar samarasi baholanadi & har yil \\
\end{longtable}

}

Manba: tavsiyanoma {[}33{]}; matni kirill yozuvidan lotin yozuviga
oʻgirilgan, sonlar oʻzgartirilmagan.

\oakilova{72-ilova}

\oakilovanom{\textbf{Istiqlol massividagi ZKD-6-2 zovurini tozalash (a) va
yopiq-yotiq zovur qudugʻidagi beton ishlari (b), 11.11.2026}}

\oakimg{figures/tv9_zkd.jpg}

Manba: tavsiyanoma {[}33{]}, 9-rasm (dala kuzatuvi). Tozalash samarasi
№76, №67 va №32a kuzatuv quduqlarida 2027-yil vegetatsiya davridagi
sathni 2024--2026-yillar bilan solishtirib baholanadi (4.5-§).

\oakilova{73-ilova}

\oakilovanom{\textbf{Tavsiyanoma {[}33{]}: Andijon viloyatining asosiy meliorativ
koʻrsatkichlari va suv yetkazib berish}}

\begin{enumerate}
\def\labelenumi{\alph{enumi})}
\tightlist
\item
  Asosiy meliorativ koʻrsatkichlar, 2025-yil (tavsiyanoma, 8.1-jadval)
\end{enumerate}

{\fontsize{12}{14}\selectfont\setlength{\tabcolsep}{4pt}

\noindent\begin{longtable}[]{|
  >{\raggedright\arraybackslash}p{(\columnwidth - 6\tabcolsep - 4\arrayrulewidth) * \real{0.4849}}|>{\raggedright\arraybackslash}p{(\columnwidth - 6\tabcolsep - 4\arrayrulewidth) * \real{0.3219}}|>{\raggedright\arraybackslash}p{(\columnwidth - 6\tabcolsep - 4\arrayrulewidth) * \real{0.1932}}|}
\hline
\begin{minipage}[b]{\linewidth}\raggedright
\textbf{Koʻrsatkich}
\end{minipage} & \begin{minipage}[b]{\linewidth}\raggedright
\textbf{Qiymati}
\end{minipage} & \begin{minipage}[b]{\linewidth}\raggedright
\textbf{Viloyat ME jadvali}
\end{minipage} \\
\hline
\endhead
\hline
\endlastfoot
Sugʻoriladigan yer & 270,9 ming ga & J-2 \\ \hline
MSOZ: yaxshi / qoniqarli / qoniqarsiz & 41,2\% / 57,5\% / 1,3\% &
J-15 \\ \hline
SES sathi 2 m dan sayoz yerlar (2024) & 89,2 ming ga (32,8\%) & J-12 \\ \hline
Kollektor-drenaj tarmogʻi bilan taʼminlangan yer & 186,4 ming ga
(68,8\%) & J-2 \\ \hline
Kollektor-drenaj tarmogʻi / shundan taʼmir talab & 8 251,6 km / 807,6 km
& J-4b \\ \hline
Tik drenaj quduqlari / ishlagani / suv oʻlchagichlisi & 485 / 136 (28\%)
/ 32 & J-10, J-10a \\ \hline
Kuzatuv quduqlari / onlayn (``Diver'') / taʼmir talab & 2 210 / 725 /
241 & J-11 \\ \hline
Tizimli taʼmirlash: reja / bajarildi & 1 008,7 km / 807,9 km (80\%) &
J-5a \\ \hline
Kollektor-drenaj oqimi & 1 673,7 mln m³ & J-16a \\
\end{longtable}

}

\begin{enumerate}
\def\labelenumi{\alph{enumi})}
\setcounter{enumi}{1}
\tightlist
\item
  Sugʻorishga olingan va isteʼmolchiga yetkazilgan suv, mln m³
  (tavsiyanoma, 8.2-jadval)
\end{enumerate}

{\fontsize{12}{14}\selectfont\setlength{\tabcolsep}{4pt}

\noindent\begin{longtable}[]{|
  >{\raggedright\arraybackslash}p{(\columnwidth - 10\tabcolsep - 6\arrayrulewidth) * \real{0.0847}}|>{\raggedright\arraybackslash}p{(\columnwidth - 10\tabcolsep - 6\arrayrulewidth) * \real{0.2081}}|>{\raggedright\arraybackslash}p{(\columnwidth - 10\tabcolsep - 6\arrayrulewidth) * \real{0.2686}}|>{\raggedright\arraybackslash}p{(\columnwidth - 10\tabcolsep - 6\arrayrulewidth) * \real{0.1674}}|>{\raggedright\arraybackslash}p{(\columnwidth - 10\tabcolsep - 6\arrayrulewidth) * \real{0.2712}}|}
\hline
\begin{minipage}[b]{\linewidth}\raggedright
\textbf{Yil}
\end{minipage} & \begin{minipage}[b]{\linewidth}\raggedright
\textbf{Viloyatga olingan (amalda)}
\end{minipage} & \begin{minipage}[b]{\linewidth}\raggedright
\textbf{Isteʼmolchiga yetkazilgan}
\end{minipage} & \begin{minipage}[b]{\linewidth}\raggedright
\textbf{Yoʻqotish}
\end{minipage} & \begin{minipage}[b]{\linewidth}\raggedright
\textbf{Isteʼmolchiga yetgan ulush}
\end{minipage} \\
\hline
\endhead
\hline
\endlastfoot
2025 & 2 207,0 & 1 468,7 & 738,3 & 66,5\% \\ \hline
2024 & 2 161,9 & 1 478,0 & 683,9 & 68,4\% \\ \hline
2023 & 2 175,0 & 1 448,0 & 727,0 & 66,6\% \\
\end{longtable}

}

Manba: tavsiyanoma {[}33{]}; matni kirill yozuvidan lotin yozuviga
oʻgirilgan, sonlar oʻzgartirilmagan.

\oakilova{74-ilova}

\oakilovanom{\textbf{Andijon viloyati tumanlarida berilgan sugʻorish suvining
meʼyorga nisbati}}

\oakimg{figures/tv_region_ratio.png}

Manba: tavsiyanoma {[}33{]}, 8-rasm; lotin yozuvida qayta chizilgan.

\end{document}
